%% tutorial.tex -- Open-SERDES Design Tutorial (BRICD)
%% Compile with pdfLaTeX + BibTeX (Overleaf: set this file as the Main document).
%% All sources live in the tutorial/ folder.
\documentclass[11pt,a4paper,oneside,openany]{book}
\usepackage{tutorial/serdestut}

% Faster compiles while editing (Overleaf free plan has a 10 s limit): uncomment and list the
% chapters you are working on; cross-references to other chapters keep working.
%\includeonly{tutorial/chapters/03_channel,tutorial/chapters/04_linkanalysis}
% Partial builds (local): if <jobname>.only exists it may contain \includeonly{...}
\IfFileExists{\jobname.only}{\input{\jobname.only}}{}

\hypersetup{pdftitle={Designing a Serial Link from Scratch: The Open-SERDES Tutorial},
            pdfauthor={Open-SERDES Project, BRICD}}

\begin{document}
\frontmatter
%% ------------------------------------------------------------------ title page
\begin{titlepage}
\thispagestyle{empty}
\tikzexternaldisable
\begin{tikzpicture}[remember picture,overlay]
  \fill[navy] (current page.south west) rectangle (current page.north east);
  \foreach \x in {0,1,...,21}{\draw[slate,line width=0.3pt] ($(current page.south west)+(\x*1cm,0)$) -- ++(0,30cm);}
  \foreach \y in {0,1,...,30}{\draw[slate,line width=0.3pt] ($(current page.south west)+(0,\y*1cm)$) -- ++(21cm,0);}
  % stylised eye diagram
  \begin{scope}[shift={($(current page.south west)+(3.0cm,4.6cm)$)},x=1.15cm,y=1cm]
    \foreach \dx/\op in {-0.18/0.45,0/1,0.18/0.45}{
      \foreach \s in {0,4,8}{
        \draw[goldlt,line width=2.2pt,opacity=\op,line cap=round]
          (\s+\dx-0.6,0) .. controls (\s+\dx+1.0,0) and (\s+\dx+1.0,3.6) .. (\s+\dx+2.6,3.6);
        \draw[goldlt,line width=2.2pt,opacity=\op,line cap=round]
          (\s+\dx-0.6,3.6) .. controls (\s+\dx+1.0,3.6) and (\s+\dx+1.0,0) .. (\s+\dx+2.6,0);
      }}
    \draw[goldlt,line width=2.2pt,opacity=0.85] (-0.6,0) -- (12.6,0) (-0.6,3.6) -- (12.6,3.6);
  \end{scope}
  \node[anchor=north west,text=goldlt,font=\sffamily\bfseries\large] at ($(current page.north west)+(2.4cm,-3.2cm)$)
     {OPEN-SERDES \textperiodcentered{} IHP SG13G2 \textperiodcentered{} 5\,Gb/s NRZ};
  \node[anchor=north west,text=white,font=\sffamily\bfseries\fontsize{40}{46}\selectfont,text width=16.5cm,align=left]
     at ($(current page.north west)+(2.4cm,-4.4cm)$) {Designing a Serial Link from Scratch};
  \draw[gold,line width=2.4pt] ($(current page.north west)+(2.45cm,-8.75cm)$) -- ++(4.5cm,0);
  \node[anchor=north west,text=pale,font=\sffamily\Large,text width=15.5cm,align=left]
     at ($(current page.north west)+(2.4cm,-9.3cm)$)
     {A design tutorial for undergraduate engineers building an open-source high-speed SerDes in 130\,nm SiGe BiCMOS};
  \node[anchor=north west,text=white,font=\sffamily\normalsize,text width=15.5cm,align=left]
     at ($(current page.north west)+(2.4cm,-12.6cm)$)
     {Channels \textperiodcentered{} Link budgets \textperiodcentered{} Equalization \textperiodcentered{} TX and RX circuits \textperiodcentered{} Jitter, PLLs and CDR \textperiodcentered{} Open-source design flow \textperiodcentered{} Hands-on labs};
  \node[anchor=south west,text=white,font=\sffamily\bfseries\large] at ($(current page.south west)+(2.4cm,2.6cm)$) {BRICD};
  \node[anchor=south west,text=pale,font=\sffamily\small] at ($(current page.south west)+(2.4cm,1.9cm)$)
     {Biomedical Research in Integrated Circuit Design \textperiodcentered{} thebricd.org};
  \node[anchor=south east,text=pale,font=\sffamily\small,align=right] at ($(current page.south east)+(-2.4cm,1.9cm)$)
     {Open-SERDES Project\\Draft v0.1 \textperiodcentered{} October 2026};
\end{tikzpicture}
\tikzexternalenable
\end{titlepage}

\tikzsetfigurename{00_preface-}
%% 00_preface.tex -- Preface (front matter, unnumbered). Author: p1
\chapter{Preface}\label{ch:preface}

Every fast interface, from a PCIe slot to the backplane of an Ethernet switch, moves its data over a \textbf{serial link}: a transmitter, a lossy wire and a receiver that recovers both the bits and the clock. This book teaches you to design one by following a single, concrete chip, the \textbf{\OS{}} transceiver: a \qty{5}{\giga\bit\per\second} NRZ link built in the open-source IHP \sgx{} \qty{130}{\nano\meter} SiGe BiCMOS process~\cite{ihppdk} with open-source tools~\cite{iicosic}.

\paragraph{Who this book is for} Third- and fourth-year electrical engineering undergraduates who will help design the \OS{} chip, and anyone who wants a practical first course in high-speed links. After reading it you should be able to:
\begin{itemize}
  \item predict the eye a channel will produce from its loss and pulse response;
  \item choose and size the TX FFE, CTLE and DFE for a channel, and defend the split;
  \item turn a link budget into specifications for the driver, front end, samplers, PLL and CDR;
  \item model the link in Python, then carry it through schematic, layout and sign-off with open tools.
\end{itemize}

\paragraph{Prerequisites} Basic circuits (KCL, small-signal MOS and BJT models, the differential pair), signals and systems (Fourier, Laplace, convolution) and some Python. Transmission lines, equalization, jitter, PLLs and CDRs are built up from first principles.

\section*{How the book is organized}
The book has five parts, shown in \cref{fig:pre-roadmap}. \textbf{Part~I, Foundations} (Chapters~1--4), gives you the vocabulary of links (pulse response, eye diagram, bit-error rate), the physics of the channel and the tools for a link budget. \textbf{Part~II, Equalization} (Chapters~5--8), shows how the TX FFE, the CTLE and the DFE undo the channel, and how to split the work between them. \textbf{Part~III, Transceiver Circuits} (Chapters~9--10), turns the system model into HBT and CMOS circuits. \textbf{Part~IV, Clocking} (Chapters~11--13), covers jitter, the PLL and clock and data recovery. \textbf{Part~V, Putting It Together} (Chapters~14--16), assembles the full architecture, walks through the open-source design flow and collects the hands-on labs.

\begin{figure}[htbp]
  \centering
  \includegraphics[width=\linewidth]{canva_roadmap}
  \caption{\textbf{Your path through this tutorial.} The five parts are milestones: Foundations, Equalization, Transceiver circuits, Clocking, and Putting it together. Each part builds on the numbers and models of the one before it.}
  \label{fig:pre-roadmap}
\end{figure}

\paragraph{Two reading paths} Everyone should read Part~I first. After that, \cref{fig:pre-paths} suggests two routes. The \textbf{system and modeling path} suits readers who will own the link budget and the serdespy model. The \textbf{circuit path} suits readers who will design transistors and layout. Both paths meet again at the architecture chapter and the labs.

\begin{figure}[htbp]
\centering
\tikzset{prech/.style={circle,draw=#1,fill=white,line width=0.8pt,minimum size=6.4mm,inner sep=0pt,font=\sffamily\footnotesize\bfseries,text=#1},
  prelk/.style={line width=1.1pt,draw=#1}}
\begin{tikzpicture}[x=1cm,y=1cm]
% trunk
\foreach \n [count=\i] in {1,2,3,4}{\node[prech=navy] (t\n) at ({0.95*(\i-1)},0) {\n};}
\foreach \a/\b in {1/2,2/3,3/4}{\draw[prelk=navy] (t\a) -- (t\b);}
% system path (top)
\foreach \n [count=\i] in {5,6,8,11,13}{\node[prech=bblue] (s\n) at ({3.65+0.95*\i},0.75) {\n};}
\draw[prelk=bblue] (t4) -- (s5); \foreach \a/\b in {5/6,6/8,8/11,11/13}{\draw[prelk=bblue] (s\a) -- (s\b);}
% circuit path (bottom)
\foreach \n [count=\i] in {7,9,10,12,15}{\node[prech=golddk] (c\n) at ({3.65+0.95*\i},-0.75) {\n};}
\draw[prelk=golddk] (t4) -- (c7); \foreach \a/\b in {7/9,9/10,10/12,12/15}{\draw[prelk=golddk] (c\a) -- (c\b);}
% merge
\node[prech=navy] (m14) at (10.0,0) {14};
\node[prech=navy] (m16) at (10.95,0) {16};
\draw[prelk=bblue] (s13) -- (m14); \draw[prelk=golddk] (c15) -- (m14); \draw[prelk=navy] (m14) -- (m16);
\node[font=\sffamily\scriptsize,text=navy,anchor=north] at (1.42,-0.42) {Part I: everyone};
\node[font=\sffamily\scriptsize\bfseries,text=bblue,anchor=south] at (6.5,1.1) {system and modeling path};
\node[font=\sffamily\scriptsize\bfseries,text=golddk,anchor=north] at (6.5,-1.1) {circuit path};
\node[font=\sffamily\scriptsize,text=navy,anchor=north,align=center] at (10.48,-0.42) {architecture\\[-1pt]and labs};
\end{tikzpicture}
\caption{\textbf{Two ways to read the book after Part~I.} Circles are chapter numbers. The system path (blue) follows the equalization theory, jitter and the CDR; the circuit path (gold) follows the CTLE circuits, the TX and RX, the PLL and the design flow. Both meet at the architecture (\cref{ch:arch}) and the labs (\cref{ch:labs}).}
\label{fig:pre-paths}
\end{figure}

\section*{How to read the boxes}
Chapters open with what you will learn and close with key takeaways. In between, six kinds of boxes mark the role of each piece of text:

{\renewcommand{\thewexample}{P.1}%
\begin{keyeq}[Bit period]
\[ \Tb = \frac{1}{\fb} = \frac{1}{\qty{5}{\giga\bit\per\second}} = \qty{200}{\pico\second} \]
\end{keyeq}
\begin{example}{A number you will use on every page}
At \qty{5}{\giga\bit\per\second}, $\fN = \fb/2 = \qty{2.5}{\giga\hertz}$; channel loss is quoted there.
\end{example}}
\setcounter{wexample}{0}
\begin{insight}
An \emph{insight} box says what a result means for the \OS{} design.
\end{insight}
\begin{pitfall}
A \emph{pitfall} box warns about a common mistake, such as quoting loss at \qty{5}{\giga\hertz}.
\end{pitfall}
\begin{tryit}[Lab 0: set-up]
A \emph{try-it} box points to a lab script (\cref{sec:lab0}).
\end{tryit}
\begin{specnote}
A \emph{specnote} box states an \OS{} target: \qty{5}{\giga\bit\per\second} NRZ at $\BER \le 10^{-12}$.
\end{specnote}
Longer algebra sits in grey \emph{derivation} boxes.

\section*{How the labs work}
The labs of \cref{ch:labs} reproduce the key plots of the book with \tool{serdespy}~\cite{serdespy}, a Python library for system-level SerDes modeling, plus NumPy and SciPy. The scripts live in \code{tutorial/code/} and share one file, \code{common.py}, that defines the link, the reference channel and the helpers; the figure scripts in \code{tutorial/figsrc/} import the same model, so labs and figures use the same numbers. Run a lab, change a parameter and compare your plot with the text.

\section*{Sources and acknowledgments}
The link-modeling chapters follow the structure of Prof.\ Sam Palermo's ECEN~720 \emph{High-Speed Links} course at Texas A\&M University~\cite{palermo720}. The equalizer and CDR chapters draw on Benjamin Liu's lecture notes~\cite{liuEq,liuCDR}. The labs are built on serdespy by Richard Barrie and colleagues at the University of Toronto~\cite{serdespy}. The chip itself is possible because IHP opened its \sgx{} process design kit~\cite{ihppdk} and the Institute for Integrated Circuits at JKU Linz maintains the IIC-OSIC-TOOLS container and the analog chip template~\cite{iicosic,jkutemplate}.

This is a \textbf{living draft (v0.1)}; numbers will be refined as the design matures. If you find an error, an unclear derivation or a figure that does not explain what it should, please email \href{mailto:bricd.org@gmail.com}{bricd.org@gmail.com} or open an issue or a pull request on the project's GitHub repository, linked from thebricd.org.

\section*{How this book was made}
This book, its figures and lab code, and the accompanying slides were developed with the help of AI tools. The authors reviewed the text, figures and code, but errors may remain; please report any you find as described above, and we will correct them. The topics, structure and many of the results follow the sources acknowledged above and listed in the References.

{\hypersetup{linkcolor=navy}\tableofcontents}

\mainmatter
\tutpart{Foundations}{What a serial link is, the language we use to describe it, and the channel that makes it hard. Everything later in this tutorial is built on the pulse response, the eye diagram and the bit-error rate introduced here.}
\tikzsetfigurename{01_intro-}
%% 01_intro.tex -- Chapter 1: Why Serial Links? Author: p1
\chapter{Why Serial Links?}\label{ch:intro}
\chapterlead{Chips have far more bandwidth inside than pins to move it out, so we send bits one after another, very fast, over a few carefully engineered wires.}

\begin{objectives}
\begin{itemize}
  \item explain why high-speed I/O moved from wide parallel buses to serial links with an embedded clock;
  \item name every block of a SerDes link and the data rate at each point of the \OS{} design;
  \item place \qty{5}{\giga\bit\per\second} NRZ among real standards such as PCIe, USB and backplane Ethernet;
  \item state the \OS{} targets and why \qty{5}{\giga\bit\per\second} NRZ in \sgx{} is the right first chip;
  \item compare the speed of \sgx{} HBTs with \qty{130}{\nano\meter} CMOS, and describe the back-end metal stack.
\end{itemize}
\end{objectives}

\section{From parallel buses to serial links}\label{sec:ch1-parallel}
The obvious way to move data between chips is a \textbf{parallel bus}: one wire per bit plus a clock wire. A 16-bit bus at \qty{312.5}{\mega\bit\per\second} per wire moves \qty{5}{\giga\bit\per\second}. As rates grew, four problems stopped it from scaling (\cref{fig:ch1-parser}a).

\begin{figure}[htbp]
\centering
% ch01_parser.tex -- parallel bus with forwarded clock vs serial link with embedded clock
\begin{subfigure}[t]{0.5\linewidth}\centering
\begin{tikzpicture}[x=1cm,y=1cm,
  b/.style={blk,minimum width=9mm,minimum height=26mm,font=\sffamily\footnotesize}]
\node[b,fill=gtint,draw=golddk] (tx) at (0,-0.95) {TX};
\node[b,fill=btint] (rx) at (3.0,-0.95) {RX};
\foreach \i/\y in {0/0.15,1/-0.3,2/-0.75}{\draw[sig] (0.45,\y) -- node[lbl,above,pos=0.5,yshift=-2pt] {D\i} (2.55,\y);}
\foreach \yy in {-0.98,-1.1,-1.22}{\fill[muted] (1.5,\yy) circle (0.8pt);}
\draw[sig] (0.45,-1.6) -- node[lbl,above,pos=0.5,yshift=-2pt] {D15} (2.55,-1.6);
\draw[clk] (0.45,-2.05) -- node[lbl,above,pos=0.5,text=golddk,yshift=-2pt] {CLK} (2.55,-2.05);
% timing at the RX: data-valid windows with skew
\begin{scope}[shift={(3.75,0)}]
  \foreach \i/\y/\s in {0/0.15/0.0,1/-0.3/0.28,2/-0.75/-0.18,3/-1.6/0.4}{
    \draw[navy,line width=0.6pt,fill=btint] (\s,\y) -- (\s+0.15,\y+0.14) -- (\s+1.5,\y+0.14) -- (\s+1.65,\y) -- (\s+1.5,\y-0.14) -- (\s+0.15,\y-0.14) -- cycle;}
  \fill[gold,opacity=0.28] (0.55,-1.85) rectangle (1.32,0.4);
  \draw[gold,dashed,line width=0.7pt] (0.93,-2.25) -- (0.93,0.45);
  \draw[golddk,line width=0.8pt] (-0.15,-2.15) -- (0.9,-2.15) -- (0.96,-1.95) -- (1.35,-1.95);
  \node[lbl,text=golddk,anchor=north] at (0.93,-2.25) {sampling clock edge};
  \draw[<->,c4,line width=0.6pt] (-0.18,0.5) -- node[lbl,above,text=c4] {skew} (0.4,0.5);
  \node[lbl,text=golddk,anchor=west] at (1.0,0.5) {common valid window};
\end{scope}
\end{tikzpicture}
\caption{Parallel bus, forwarded clock}
\end{subfigure}%
\begin{subfigure}[t]{0.5\linewidth}\centering
\begin{tikzpicture}[x=1cm,y=1cm,
  b/.style={blk,minimum width=13mm,minimum height=11mm,font=\sffamily\footnotesize}]
\node[b,fill=gtint,draw=golddk] (tx) at (0,-0.95) {SER\\16:1};
\node[b,fill=btint] (rx) at (4.6,-0.95) {CDR\\+ DES};
\draw[sig] (tx.east|-0,-0.88) -- (rx.west|-0,-0.88);
\draw[sig] (tx.east|-0,-1.02) -- (rx.west|-0,-1.02);
\node[lbl,anchor=north] at (2.3,-1.1) {one differential pair};
\draw[sig,line width=1.5pt] (-1.25,-0.95) -- (tx.west); \node[lbl,anchor=south] at (-0.95,-0.92) {16};
\draw[sig,line width=1.5pt] (rx.east) -- (5.85,-0.95);
% NRZ waveform above the pair
\begin{scope}[shift={(0.95,-0.25)},x=0.34cm,y=0.5cm]
  \draw[c1,line width=0.9pt] (0,0) -- (1,0) -- (1.15,1) -- (3,1) -- (3.15,0) -- (4,0) -- (4.15,1) -- (5,1) -- (5.15,0) -- (7,0) -- (7.15,1) -- (8,1) -- (8.15,0) -- (10,0);
  \foreach \x in {1.07,3.07,4.07,5.07,7.07,8.07}{\fill[goldlt,draw=golddk,line width=0.3pt] (\x,0.5) circle (1.2pt);}
\end{scope}
\node[lbl,text=golddk,anchor=south] at (2.65,0.32) {clock information lives in the transitions};
\draw[clk] (4.6,-1.5) -- (4.6,-2.05) node[lbl,text=golddk,anchor=north] {recovered clock};
\end{tikzpicture}
\caption{Serial link, embedded clock}
\end{subfigure}

\caption{\textbf{A parallel bus is limited by skew; a serial link carries its own clock.} (a)~Each lane of a bus arrives with a slightly different delay. The receiver can only sample inside the window where every lane is valid, and that window shrinks as the bit time shrinks. (b)~A serial link sends one fast stream on one differential pair. The receiver extracts the clock from the data transitions (gold dots), so there is no clock wire to match.}
\label{fig:ch1-parser}
\end{figure}

\begin{itemize}
  \item \textbf{Skew.} Lanes differ in length, load and driver strength, so bits launched together arrive at different times; the usable window is the bit time minus the worst skew.
  \item \textbf{Pin count.} Pins and board area are expensive and grow far more slowly than on-chip bandwidth, and a wide bus needs many return pins.
  \item \textbf{Power and noise.} Many single-ended drivers switching together cause supply bounce and crosstalk, and the energy per bit stays high.
  \item \textbf{The clock wire.} The forwarded clock must be matched to every lane across process, voltage and temperature.
\end{itemize}

A \textbf{serial link} \textbf{serializes} the word into one full-rate stream on a single \textbf{differential pair}, which carries its own return path and rejects common-mode noise. No clock is sent: a \textbf{clock and data recovery} (CDR) circuit in the receiver locks a local clock to the data transitions, so the clock is \textbf{embedded} in the data and lane-to-lane skew no longer matters.

\begin{example}{Skew budget: parallel bus versus serial link}\label{we:ch1-skew}
A signal on FR-4 ($\varepsilon_r \approx 4$) travels at $v = c/\sqrt{\varepsilon_r} \approx \qty{0.15}{\milli\meter\per\pico\second}$, or about \qty{6.7}{\pico\second} per millimeter. Suppose two lanes differ in length by \qty{10}{\milli\meter}, a modest mismatch on a crowded board. The skew is $\qty{10}{\milli\meter}\times\qty{6.7}{\pico\second\per\milli\meter} \approx \qty{67}{\pico\second}$.
\begin{itemize}
  \item 16-bit bus at \qty{312.5}{\mega\bit\per\second}: the bit time is \qty{3.2}{\nano\second}, so the skew costs $67/3200 \approx 2\%$ of it.
  \item The same bus pushed to \qty{5}{\giga\bit\per\second} per lane: the bit time is \qty{200}{\pico\second} and the same skew eats \emph{one third} of it, before any jitter or noise.
\end{itemize}
The pins tell a similar story: 16 data wires plus a clock and their returns, against one differential pair (2 pins) for the serial link.
\end{example}

The price is speed: each bit lasts only $\Tb = 1/\fb$, and the wire that looked ideal at \qty{312.5}{\mega\bit\per\second} now smears each bit into its neighbors. Paying that price efficiently is measured by the energy per bit.

\begin{keyeq}[Energy efficiency of a link]
\begin{equation}
  E_b = \frac{P_\mathrm{link}}{\fb}, \qquad \qty{1}{\milli\watt} \text{ per } \unit{\giga\bit\per\second} = \qty{1}{\pico\joule\per\bit}
  \label{eq:ch1-energy}
\end{equation}
\end{keyeq}
A transceiver burning \qty{50}{\milli\watt} at \qty{5}{\giga\bit\per\second} spends \qty{10}{\pico\joule} per bit.

\begin{insight}
Serial links win because they trade many slow, skew-limited wires for one fast, self-clocked wire. The cost moves into the transceiver: it must drive, equalize and re-time a \qty{200}{\pico\second} bit after a lossy channel. Every chapter that follows is about doing that within a power budget.
\end{insight}

\section{Anatomy of a SerDes link}\label{sec:anatomy}
\Cref{fig:link-anatomy} shows the complete \OS{} link. Follow a 16-bit word from left to right.

\begin{figure}[htbp]
\centering
% ch01_anatomy.tex -- fig:link-anatomy : the Open-SERDES link, TX -> channel -> RX, with clocking
\begin{tikzpicture}[x=1cm,y=1cm,
  b/.style={blk,minimum width=12mm,minimum height=10mm,font=\sffamily\footnotesize,inner sep=2pt},
  bb/.style={b,fill=btint},
  bg/.style={b,fill=gtint,draw=golddk},
  bn/.style={b,fill=navy,text=white},
  seg/.style={draw=navy,line width=0.5pt,minimum height=6mm,inner sep=0pt,font=\sffamily\scriptsize},
  rate/.style={font=\sffamily\footnotesize\bfseries,text=golddk,align=center},
  ann/.style={font=\sffamily\scriptsize,text=muted,align=center}]
% ---------------- regions
\draw[region,fill=bgpale!60] (-0.05,-2.85) rectangle (5.0,1.05);
\draw[region,fill=gtint!35]  (5.15,-2.85) rectangle (9.05,1.05);
\draw[region,fill=bgpale!60] (9.2,-2.85) rectangle (16.0,1.05);
\node[lblb,anchor=north west] at (0.0,1.0) {TRANSMITTER (TX)};
\node[lblb,anchor=north west] at (5.2,1.0) {CHANNEL};
\node[lblb,anchor=north west] at (9.25,1.0) {RECEIVER (RX)};
% ---------------- rate annotations (above the regions)
\node[rate,anchor=west] at (-0.05,1.55) {16 $\times$ 312.5\,Mb/s};
\node[rate] at (7.1,1.55) {5\,Gb/s serial NRZ \;(1\,UI = 200\,ps)};
\node[rate,anchor=east] at (16.0,1.55) {16 $\times$ 312.5\,Mb/s};
% ---------------- TX data path
\node[bg] (ser) at (1.75,0) {Serializer\\16:1};
\node[bb,minimum width=15mm] (drv) at (3.75,0) {3-tap FFE\\CML driver};
\draw[sig,line width=1.6pt] (0.05,0) -- (ser.west);
\draw[navy] (0.45,-0.14) -- (0.62,0.14); \node[lbl,anchor=south] at (0.53,0.12) {16};
\draw[sig] (ser.east) -- (drv.west);
\node[draw=navy,fill=navy,minimum size=1.6mm,inner sep=0pt] (padt) at (4.85,0) {};
\draw[navy,line width=0.75pt] (drv.east) -- (padt.west);
% ---------------- channel: package | PCB | connector | PCB | package
\begin{scope}
  \draw[navy,line width=0.6pt] (4.93,0.07) -- (9.12,0.07);
  \draw[navy,line width=0.6pt] (4.93,-0.07) -- (9.12,-0.07);
  \node[seg,fill=gtint,minimum width=4.5mm]  (pk1) at (5.525,0) {};
  \node[seg,fill=btint,minimum width=8.5mm] (pc1) at (6.275,0) {};
  \node[seg,fill=pale,minimum width=4.5mm]   (cn)  at (7.025,0) {};
  \node[seg,fill=btint,minimum width=8.5mm] (pc2) at (7.775,0) {};
  \node[seg,fill=gtint,minimum width=4.5mm]  (pk2) at (8.525,0) {};
  \foreach \n/\t in {pk1/pkg,pc1/PCB,cn/conn.,pc2/PCB,pk2/pkg}{\node[ann,anchor=base] at (\n.center |- 0,-0.58) {\t};}
  \node[ann] at (7.05,-1.3) {REF channel: 1\,m FR-4 + pads\\[-1pt]\textbf{\textcolor{c4}{$-$17.4\,dB at 2.5\,GHz}}};
\end{scope}
\node[draw=navy,fill=navy,minimum size=1.6mm,inner sep=0pt] (padr) at (9.3,0) {};
% ---------------- RX data path
\draw[navy,line width=0.75pt] (padr.east) -- (9.75,0);
\node[dot] at (9.75,0) {};
\draw[navy,line width=0.7pt] (9.75,0) -- (9.75,-0.25);
\draw[navy,line width=0.7pt] (9.75,-0.25) -- ++(0.09,-0.05) -- ++(-0.18,-0.1) -- ++(0.18,-0.1) -- ++(-0.18,-0.1) -- ++(0.18,-0.1) -- ++(-0.09,-0.05) -- ++(0,-0.15);
\draw[navy,line width=0.7pt] (9.6,-1.05) -- (9.9,-1.05) (9.66,-1.12) -- (9.84,-1.12) (9.72,-1.19) -- (9.78,-1.19);
\draw[navy,line width=0.7pt] (9.75,-0.9) -- (9.75,-1.05);
\node[ann,anchor=north] at (9.75,-1.22) {100\,$\Omega$\\[-1pt]+ ESD};
\node[bb] (ctle) at (10.85,0) {CTLE\\0--12\,dB};
\node[bb,minimum width=16mm] (smp) at (12.65,0) {Samplers\\+ 2-tap DFE};
\node[bg] (des) at (14.5,0) {Deser.\\1:16};
\draw[sig] (9.75,0) -- (ctle.west);
\draw[sig] (ctle.east) -- (smp.west);
\draw[sig] (smp.east) -- (des.west);
\draw[sig,line width=1.6pt] (des.east) -- (15.95,0);
\draw[navy] (15.45,-0.14) -- (15.62,0.14); \node[lbl,anchor=south] at (15.53,0.12) {16};
% ---------------- clocking
\node[bn,minimum width=15mm] (pll) at (1.75,-2.0) {PLL\\LC VCO};
\draw[sig] (0.05,-2.0) -- node[lbl,above,pos=0.42] {ref} node[lbl,below,pos=0.42] {50\,MHz} (pll.west);
\draw[clk] (pll.north) -- node[lbl,right,text=golddk,align=left] {2.5\,GHz\\(+ divided\\clocks)} (ser.south);
\node[lbl,align=left,anchor=west] at (2.65,-2.0) {$f_\mathrm{ref}\times N$\\$= 50\,\mathrm{MHz}\times 50$};
\node[bn,minimum width=26mm] (cdr) at (12.65,-2.0) {CDR: BB-PD +\\loop filter + PI};
\draw[sig] ([xshift=-4mm]smp.south) -- node[lbl,left,align=right] {data +\\edge} ([xshift=-4mm]cdr.north);
\draw[clk] ([xshift=4mm]cdr.north) -- node[lbl,right,text=golddk,align=left] {4-phase\\2.5\,GHz} ([xshift=4mm]smp.south);
\draw[clk] (cdr.east) -| node[lbl,text=golddk,pos=0.75,right] {$\div$8} (des.south);
\draw[clk] (15.95,-2.45) -- node[lbl,text=golddk,below,pos=0.45] {I/Q from PLL} (cdr.east |- 0,-2.45) ;
\end{tikzpicture}

\caption{\textbf{Anatomy of the \OS{} link, with the data rate at every point.} Sixteen bits at \qty{312.5}{\mega\bit\per\second} are serialized to one \qty{5}{\giga\bit\per\second} NRZ stream, driven through a 3-tap FFE driver, attenuated by the package, board and connector, then equalized (CTLE and DFE), sampled with a recovered clock and deserialized back to 16 bits. Gold dashed lines are clocks: the TX PLL multiplies a \qty{50}{\mega\hertz} reference to \qty{2.5}{\giga\hertz}, and the RX CDR steers four sampling phases using I/Q clocks from a local PLL. All high-speed paths are differential.}
\label{fig:link-anatomy}
\end{figure}

\paragraph{Transmitter (TX)} The \textbf{serializer} works in stages: CMOS logic does 16:4, current-mode logic (CML) does 4:2, and a \textbf{half-rate} 2:1 CML multiplexer clocked at \qty{2.5}{\giga\hertz} emits one bit per clock edge (\cref{sec:serializer}). The \textbf{driver} launches a \qty{1}{\Vppd} signal into the \qty{50}{\ohm}-per-side channel through a 3-tap \textbf{feed-forward equalizer} (FFE) that pre-distorts it against the channel loss (\cref{ch:txffe,sec:cml-driver}). A \textbf{phase-locked loop} (PLL) with an LC oscillator makes the \qty{2.5}{\giga\hertz} clock (\cref{ch:pll}).

\paragraph{Channel} Packages, board (PCB) traces and vias and a connector together act as a low-pass filter with reflections; the \OS{} reference channel loses \qty{17.4}{\deci\bel} at \qty{2.5}{\giga\hertz} (\cref{ch:channel}).

\paragraph{Receiver (RX)} A \qty{100}{\ohm} \textbf{termination} with ESD protection absorbs the wave (\cref{sec:termination}). A \textbf{continuous-time linear equalizer} (CTLE) restores the lost high frequencies (\cref{ch:ctle}). Four \textbf{samplers} on the \ang{0}/\ang{90}/\ang{180}/\ang{270} phases of a \qty{2.5}{\giga\hertz} clock take data and edge samples, while a 2-tap \textbf{decision-feedback equalizer} (DFE) subtracts the echo of previous bits (\cref{ch:dfe,sec:sampler}). The \textbf{CDR} steers a phase interpolator to keep the samples centered (\cref{ch:cdr}), and the \textbf{deserializer} rebuilds the 16-bit word (\cref{sec:deserializer}).

\begin{example}{Following the rates through the link}\label{we:ch1-rates}
\begin{itemize}
  \item Parallel side: $16 \times \qty{312.5}{\mega\bit\per\second} = \qty{5000}{\mega\bit\per\second} = \qty{5}{\giga\bit\per\second}$.
  \item Serial side: $\Tb = 1/\qty{5}{\giga\bit\per\second} = \qty{200}{\pico\second}$, the \textbf{unit interval} (UI).
  \item Half-rate clock: \qty{2.5}{\giga\hertz} has a \qty{400}{\pico\second} = 2~UI period; both edges give one bit per \qty{200}{\pico\second}. The four CDR phases are 0.5~UI apart: two data centers and two edges.
  \item PLL: $f_\mathrm{ref} \times N = \qty{50}{\mega\hertz}\times 50 = \qty{2.5}{\giga\hertz}$. The deserializer output clock is $\qty{2.5}{\giga\hertz}/8 = \qty{312.5}{\mega\hertz}$.
\end{itemize}
The half-rate clock frequency equals the Nyquist frequency $\fN = \fb/2$ of the data (\cref{sec:nrz}).
\end{example}

\Cref{fig:ch1-glance} summarizes the link with the four numbers you will meet on almost every page.

\begin{figure}[htbp]
  \centering
  \includegraphics[width=\linewidth]{canva_glance}
  \caption{\textbf{A SerDes link at a glance.} The transmitter (serializer, FFE driver), the channel (package, PCB, connector) and the receiver (CTLE, DFE, samplers, deserializer), held together by the PLL and the CDR. Keep the four key numbers in mind: \qty{5}{\giga\bit\per\second}, a \qty{200}{\pico\second} UI, a \qty{2.5}{\giga\hertz} Nyquist frequency and a target BER of $10^{-12}$.}
  \label{fig:ch1-glance}
\end{figure}

\section{Where SerDes are used}\label{sec:ch1-standards}
Serial links connect a CPU to its graphics card (PCIe), a computer to its drives (USB, SATA) and the line cards of a network switch (Ethernet). \Cref{tab:ch1-standards} lists common standards. \textbf{GT/s} counts symbols on the wire and Gb/s counts bits; they are equal for NRZ, while PAM4 carries two bits per symbol (\cref{sec:pam4}). Line rates include coding overhead: 8b/10b sends 10 line bits per 8 data bits, 64b/66b sends 66 per 64.

\begin{table}[tbp]
\centering\small
\caption{\textbf{Per-lane line rates of common serial standards.} The Nyquist frequency is half the symbol rate. \OS{}, at \qty{5}{\giga\bit\per\second} NRZ with $\fN = \qty{2.5}{\giga\hertz}$, sits next to PCIe~2.0, USB~3.2 Gen~1 and 5GBASE-KR.}
\label{tab:ch1-standards}
\begin{tabularx}{\linewidth}{@{}l l L r@{}}
\toprule
Standard (year) & Rate per lane & Modulation and coding & $\fN$ (GHz) \\
\midrule
PCIe 1.0a / 2.0 (2003 / 2007) & 2.5 / 5\,GT/s & NRZ, 8b/10b & 1.25 / 2.5 \\
PCIe 3.0 / 4.0 / 5.0 (2010--19) & 8 / 16 / 32\,GT/s & NRZ, 128b/130b & 4 / 8 / 16 \\
PCIe 6.0 (2022) & 64\,GT/s & PAM4, 32\,GBd, FEC & 16 \\
USB 3.2 Gen 1 / Gen 2 (2008 / 2013) & 5 / 10\,Gb/s & NRZ, 8b/10b / 128b/132b & 2.5 / 5 \\
USB4 v2 (2022) & 40\,Gb/s & PAM3, 25.6\,GBd & 12.8 \\
SATA I / II / III (2003--09) & 1.5 / 3 / 6\,Gb/s & NRZ, 8b/10b & 0.75 / 1.5 / 3 \\
2.5GBASE-KX (802.3cb, 2018) & 3.125\,Gb/s & NRZ, 8b/10b & 1.56 \\
5GBASE-KR (802.3cb, 2018) & 5.156\,Gb/s & NRZ, 64b/66b & 2.58 \\
10GBASE-KR (802.3ap, 2007) & 10.3125\,Gb/s & NRZ, 64b/66b & 5.16 \\
100GBASE-KR4 (802.3bj, 2014) & 25.78\,Gb/s & NRZ, RS-FEC & 12.9 \\
400GBASE-KR4 (802.3ck, 2022) & 106.25\,Gb/s & PAM4, 53.1\,GBd, FEC & 26.6 \\
UCIe 1.0--3.0 (2022--25) & 4--64\,GT/s & NRZ, single-ended, fwd.\ clock & $\le$\,32 \\
\bottomrule
\end{tabularx}
\end{table}

The last row is a useful counterexample: UCIe links chiplets millimeters apart in one package, where thousands of short, matched wires are cheap, so the parallel bus with a forwarded clock wins again. Serial links win when the wire is long, lossy and expensive, which is the situation \OS{} targets.

\Cref{fig:ch1-rate-year} shows per-lane rates doubling roughly every three to four years. Around \qty{50}{\giga\bit\per\second} the industry moved to PAM4, because the loss at the NRZ Nyquist frequency became too high.

\begin{figure}[htbp]
  \centering
  \includegraphics[width=\linewidth]{ch01_rate_vs_year}
  \caption{\textbf{Per-lane data rates double every three to four years.} PCIe generations, IEEE~802.3 backplane Ethernet lanes and USB, on a log axis. Squares mark multi-level signaling (PAM4, or PAM3 for USB4~v2); the open square is 802.3dj (\qty{200}{\giga\bit\per\second} per lane), whose completion is planned for 2026. The gold band marks the \OS{} target, close to 5GBASE-KR (802.3cb) and PCIe~2.0. Sources: \cite{p1:pcie,ieee8023,p1:usb4v2}.}
  \label{fig:ch1-rate-year}
\end{figure}

\section{The Open-SERDES project}\label{sec:os-targets}
\OS{} is an open project to design a complete \qty{5}{\giga\bit\per\second} serial-link IP for Ethernet-class backplane links in IHP \sgx{}, using only open-source PDKs and tools, with every schematic, layout, model and test bench public. The scope runs from specifications and architecture, through a Python link-simulation framework, to the circuits, custom analog layout, synthesized and routed digital blocks, and post-layout verification with DRC and LVS. The \textbf{deliverables} are a tape-out-ready GDSII, Verilog RTL with test benches, a demonstration of the modeling framework, a documented GitHub repository with a thesis, and a paper on the whole flow. The specnote collects the targets; \cref{ch:arch} expands them.

\begin{specnote}[Open-SERDES target (reference design)]
\small
\begin{tabularx}{\linewidth}{@{}l L@{}}
Data rate & \qty{5}{\giga\bit\per\second} NRZ, $\Tb = 1\,\UI = \qty{200}{\pico\second}$, $\fN = \qty{2.5}{\giga\hertz}$\\
Target BER & $\le 10^{-12}$ ($Q = 7.03$); PRBS7 for bring-up, PRBS31 for compliance\\
Channel & REF: \qty{1}{\meter} FR-4 stripline + \qty{300}{\femto\farad} pads, \qty{-17.4}{\deci\bel} at \qty{2.5}{\giga\hertz}\\
TX & 16:1 serializer (half-rate 2:1 CML mux), HBT CML driver, 3-tap FFE, 1\,V$_\mathrm{ppd}$, \qty{100}{\ohm} diff.\\
RX & \qty{100}{\ohm} termination + ESD, HBT CTLE (0--12\,dB), half-rate 2-tap DFE, 4 samplers, 1:16 deserializer\\
Clocking & integer-$N$ CP-PLL, LC VCO at \qty{2.5}{\giga\hertz} ($f_\mathrm{ref}$ = \qty{50}{\mega\hertz}, $N$ = 50); PI-based bang-bang CDR, 64 PI steps per 2\,UI\\
Test & on-chip PRBS7/31 generator and checker, near-end loopback\\
\end{tabularx}
\end{specnote}

\paragraph{Why \qty{5}{\giga\bit\per\second} NRZ is the right first target}
\begin{itemize}
  \item \textbf{Fast enough to be real.} It matches PCIe~2.0, USB~3.2 Gen~1 and 5GBASE-KR, and a \qty{1}{\meter} board loses \qty{17}{\deci\bel} at \qty{2.5}{\giga\hertz}: the raw eye is closed, so every equalization idea in this book is needed.
  \item \textbf{Slow enough to be safe.} A \qty{200}{\pico\second} UI is easy for \sgx{} HBTs but at the edge of \qty{130}{\nano\meter} CMOS (\cref{sec:sg13g2}), which shows clearly why BiCMOS helps.
  \item \textbf{NRZ is simple and robust.} PAM4 costs about \qty{9.5}{\deci\bel} of eye, three times the samplers and better linearity, which pays off only on much lossier channels (\cref{sec:pam4}).
  \item \textbf{Proven blocks.} The team's \qty{2.4}{\giga\hertz} LC-VCO PLL in \qty{130}{\nano\meter} BiCMOS~\cite{smacd2026pll} moves to \qty{2.5}{\giga\hertz} with little change.
\end{itemize}

\paragraph{Prior art} A previous undergraduate project modeled a \qty{5}{\giga\bit\per\second} SerDes in a commercial \qty{28}{\nano\meter} library~\cite{p1:serasinghe2016}. The Purdue \textbf{OpenSerDes}~\cite{openserdes2021} was the first open-source SerDes: an all-digital, synthesizable SkyWater \qty{130}{\nano\meter} CMOS design with inverter-based drivers and an oversampling CDR. In post-layout simulation it ran at \qty{2}{\giga\bit\per\second} and drew about \qty{438}{\milli\watt}, mostly in the serializer, deserializer and CDR logic; standard \qty{130}{\nano\meter} CMOS cells capped its speed. Neither project produced a silicon-proven, open, analog \qty{5}{\giga\bit\per\second} SerDes. \OS{} keeps OpenSerDes's open flow but moves the fastest circuits to SiGe HBTs, which breaks the CMOS speed limit at the same \qty{130}{\nano\meter} node.

\paragraph{The plan} We organize the scope into five phases (\cref{fig:ch1-phases}): specifications, channel and supply plan (Phase~0); the serdespy system model that checks every later decision (Phase~1); circuits and digital RTL (Phase~2); layout, place-and-route, extraction and EM simulation (Phase~3); and sign-off, tape-out and lab test (Phase~4).

\begin{figure}[htbp]
\centering
% ch01_phases.tex -- the five project phases as chevrons
\begin{tikzpicture}[x=1cm,y=1cm]
\foreach \i/\t/\d/\c/\col in {
  0/{Phase 0}/{Specs and\\feasibility}/{Ch.\,1, 4, 14}/gold,
  1/{Phase 1}/{System model\\(serdespy)}/{Ch.\,2--8, 16}/bblue,
  2/{Phase 2}/{Circuit design\\+ digital RTL}/{Ch.\,9--13}/navy,
  3/{Phase 3}/{Layout, APR,\\extraction, EM}/{Ch.\,15}/slate,
  4/{Phase 4}/{Sign-off, tape-out,\\lab test}/{Ch.\,15, 16}/golddk}{
  \pgfmathsetmacro{\x}{\i*3.22}
  \fill[\col] (\x,0) -- (\x+2.85,0) -- (\x+3.2,0.45) -- (\x+2.85,0.9) -- (\x,0.9) -- (\x+0.35,0.45) -- cycle;
  \node[font=\sffamily\footnotesize\bfseries,text=white] at (\x+1.6,0.45) {\t};
  \node[font=\sffamily\scriptsize,text=navy,align=center,anchor=north] at (\x+1.6,-0.08) {\d};
  \node[font=\sffamily\scriptsize\itshape,text=muted,anchor=north] at (\x+1.6,-0.78) {\c};
}
\end{tikzpicture}

\caption{\textbf{The \OS{} project in five phases.} Each phase ends in a reviewable result: a specification, a system model, schematics, layouts and a signed-off GDSII. The italic line names the chapters that support each phase.}
\label{fig:ch1-phases}
\end{figure}

\section{IHP SG13G2 at a glance}\label{sec:sg13g2}
\sgx{} is IHP's \qty{130}{\nano\meter} SiGe:C BiCMOS process, and its PDK is open source~\cite{ihppdk}. It offers two kinds of transistor on one die: standard \qty{130}{\nano\meter} CMOS (\qty{1.2}{\volt} thin-oxide core devices and \qty{3.3}{\volt} thick-oxide I/O devices), and the \textbf{npn13G2} SiGe heterojunction bipolar transistor (HBT), with $f_T \approx \qty{300}{\giga\hertz}$ and $f_{max} \approx$ 450--\qty{500}{\giga\hertz}~\cite{p1:heinemann2010}.

\paragraph{HBT versus CMOS speed} The transit frequency $f_T$, where the current gain falls to one, sets the device's intrinsic time scale $1/(2\pi f_T)$: about \qty{0.5}{\pico\second} for the HBT. A \qty{130}{\nano\meter} NMOS peaks on the order of \qty{100}{\giga\hertz}, and only at its best bias, and logic gates never operate there. A better yardstick for logic is the \textbf{FO4 delay}, the delay of an inverter driving four copies of itself: about \qty{45}{\pico\second} in \qty{130}{\nano\meter} CMOS. \Cref{fig:ch1-speed} puts these times next to the \qty{200}{\pico\second} UI.

\begin{figure}[htbp]
  \centering
  \includegraphics[width=\linewidth]{ch01_speed_ladder}
  \caption{\textbf{A \qty{200}{\pico\second} UI is about four CMOS gate delays, but about a hundred HBT gate delays.} Half-rate CMOS circuits need a bit time of roughly 3--4 FO4 delays, so \qty{130}{\nano\meter} CMOS has almost no margin at \qty{5}{\giga\bit\per\second}. The \sgx{} CML ring-oscillator gate delay of about \qty{2}{\pico\second}~\cite{p1:heinemann2010} is a best case; a loaded CML gate is several times slower and still fits comfortably. FO4 is a typical value for the node.}
  \label{fig:ch1-speed}
\end{figure}

\paragraph{Why CML with HBTs} The fast blocks (final mux, driver, CTLE, samplers, clock buffers) use \textbf{current-mode logic}: a differential pair steers a constant tail current into one of two load resistors. The swing is small (a few hundred millivolts), the supply current is constant, so little switching noise is injected, and an HBT pair switches fully with only about $4V_T \approx \qty{100}{\milli\volt}$ of input, because $\gm = I_C/V_T$ is several times a MOSFET's per unit current (\cref{sec:cml-driver,sec:ctle-hbt}).

\paragraph{The back end} \sgx{} has seven aluminum metal layers (\cref{fig:ch1-beol}): thin Metal1--Metal5 for dense wiring, and thick \textbf{TopMetal1} (\qty{2}{\micro\meter}) and \textbf{TopMetal2} (\qty{3}{\micro\meter}), whose low resistance suits spiral inductors (LC VCO, inductive peaking), on-chip transmission lines and supply rails. A MIM capacitor of roughly \qty{1.5}{\femto\farad\per\square\micro\meter} (check the PDK) sits below TopMetal1.

\begin{figure}[htbp]
\centering
% ch01_beol.tex -- stylized SG13G2 cross-section (not to scale) + HBT symbol and spiral inductor
\begin{tikzpicture}[x=0.95cm,y=1cm,
  met/.style={fill=c1!55,draw=navy,line width=0.4pt},
  tmet/.style={fill=c1!75,draw=navy,line width=0.5pt},
  via/.style={fill=muted!60,draw=none},
  ann/.style={font=\sffamily\scriptsize,text=navy,anchor=west,align=left}]
% oxide background
\fill[bgpale] (0,0.9) rectangle (8.6,6.15);
% substrate and device layer
\fill[muted!18] (0,0) rectangle (8.6,0.9);
\node[font=\sffamily\scriptsize,text=muted] at (6.9,0.3) {p-substrate};
% NMOS (left)
\fill[c3!30] (0.5,0.62) rectangle (1.05,0.9); \fill[c3!30] (1.55,0.62) rectangle (2.1,0.9);
\fill[muted!70] (1.05,0.95) rectangle (1.55,1.15);
\node[font=\sffamily\scriptsize,text=muted] at (1.3,0.4) {NMOS};
% HBT (right of NMOS): buried collector, SiGe base, T-shaped emitter
\fill[c2!35] (2.8,0.25) rectangle (5.2,0.6);
\fill[c2!55] (3.1,0.6) rectangle (4.9,0.9);
\fill[goldlt] (3.2,0.9) rectangle (4.8,1.05);
\fill[c4!60] (3.85,1.05) rectangle (4.15,1.25); \fill[c4!60] (3.65,1.25) rectangle (4.35,1.38);
\node[font=\sffamily\scriptsize,text=brown] at (4.0,0.08) {SiGe:C HBT};
% contacts
\foreach \x in {0.7,1.85,3.3,4.0,4.65}{\fill[via] (\x,1.05) rectangle (\x+0.12,1.6);}
% thin metals M1..M5
\foreach \y/\n in {1.6/1,1.95/2,2.3/3,2.65/4,3.0/5}{
  \fill[met] (0.4,\y) rectangle (2.3,\y+0.13);
  \fill[met] (3.0,\y) rectangle (5.0,\y+0.13);
  \fill[via] (1.2,\y+0.13) rectangle (1.32,\y+0.35);
  \fill[via] (3.9,\y+0.13) rectangle (4.02,\y+0.35);}
% MIM capacitor between M5 and TM1
\fill[gold] (5.7,3.28) rectangle (7.6,3.36);
\fill[met] (5.6,3.0) rectangle (7.8,3.13);
\fill[via] (6.5,3.36) rectangle (6.62,3.6);
\node[font=\sffamily\scriptsize,text=golddk] at (6.7,2.75) {MIM cap};
% TopMetal1 and TopMetal2
\fill[via] (1.2,3.13) rectangle (1.32,3.6); \fill[via] (3.9,3.13) rectangle (4.02,3.6);
\fill[tmet] (0.4,3.6) rectangle (8.0,4.05);
\fill[via] (2.0,4.05) rectangle (2.2,4.6); \fill[via] (6.0,4.05) rectangle (6.2,4.6);
\foreach \x in {0.6,2.6,4.6,6.6}{\fill[tmet] (\x,4.6) rectangle (\x+1.4,5.45);}
\draw[muted,line width=0.6pt] (0,5.75) -- (8.6,5.75);
\node[font=\sffamily\scriptsize,text=muted,anchor=south east] at (8.55,5.77) {passivation};
% labels on the right of the stack
\draw[decorate,decoration={brace,amplitude=4pt},navy] (8.75,3.13) -- (8.75,1.6);
\node[ann] at (8.95,2.37) {Metal1--Metal5: thin Al,\\[-1pt]dense local routing};
\node[ann] at (8.95,3.82) {TopMetal1, \textbf{2\,\textmu m}};
\node[ann] at (8.95,5.03) {TopMetal2, \textbf{3\,\textmu m}: inductors,\\[-1pt]T-lines, supply rails};
\node[ann] at (8.95,3.25) {\textcolor{golddk}{MIM $\approx$1.5\,fF/\textmu m$^2$}};
\node[font=\sffamily\scriptsize\itshape,text=muted,anchor=west] at (0.0,6.35) {stylized, not to scale};
% ---- HBT symbol
\begin{scope}[shift={(13.45,4.4)}]
  \draw[color=navy] (0,0) node[npn,scale=1.0](Q){};
  \draw[navy] (Q.B) -- ++(-0.35,0) node[left,font=\sffamily\scriptsize]{B};
  \draw[navy] (Q.C) -- ++(0,0.25) node[above,font=\sffamily\scriptsize]{C};
  \draw[navy] (Q.E) -- ++(0,-0.25) node[below,font=\sffamily\scriptsize]{E};
  \node[font=\sffamily\footnotesize\bfseries,text=navy,anchor=west] at (0.5,0.45) {npn13G2};
  \node[font=\sffamily\scriptsize,text=navy,anchor=west,align=left] at (0.5,-0.25)
     {$f_T\approx300$\,GHz\\[-1pt]$f_{max}\approx450$--$500$\,GHz\\[-1pt]$BV_{CEO}\approx1.7$\,V};
\end{scope}
% ---- spiral inductor (top view)
\begin{scope}[shift={(13.25,1.25)}]
  \draw[c1!80!navy,line width=1.6pt,line cap=round]
    (-1.0,-0.8) -- (-0.75,-0.8) -- (-0.75,0.75) -- (0.75,0.75) -- (0.75,-0.75) -- (-0.5,-0.75)
    -- (-0.5,0.5) -- (0.5,0.5) -- (0.5,-0.5) -- (-0.25,-0.5) -- (-0.25,0.25) -- (0.25,0.25) -- (0.25,-0.1);
  \node[font=\sffamily\scriptsize,text=navy,anchor=west,align=left] at (1.0,0.0) {spiral inductor\\[-1pt](TopMetal2,\\[-1pt]top view)};
\end{scope}
\end{tikzpicture}

\caption{\textbf{The \sgx{} stack: fast HBTs at the bottom, thick aluminum on top.} Stylized cross-section (not to scale) with an NMOS, a SiGe:C HBT, five thin metals, the MIM capacitor and the two thick top metals; right, the HBT symbol with its headline figures and a spiral inductor drawn in TopMetal2. Values are approximate; use the PDK documentation~\cite{ihppdk} for design.}
\label{fig:ch1-beol}
\end{figure}

\paragraph{Supplies} CMOS logic (slow SER/DES stages, PRBS, digital CDR loop) runs from \qty{1.2}{\volt}. HBT CML stacks a load, a switching pair and a tail source, so it uses a separate analog supply, nominally \qty{2.5}{\volt}, to be confirmed in Phase~0.

\begin{pitfall}[Respect $BV_{CEO}$]
The npn13G2 breakdown voltage $BV_{CEO}$ is only about \qty{1.7}{\volt}. On a \qty{2.5}{\volt} supply, an HBT that sees the whole supply (at power-up, or when a tail source saturates) can break down. Check every $V_{CE}$ in every corner and transient, not just at the operating point.
\end{pitfall}

\begin{insight}
Use HBTs where the UI is short (the last mux stage, the driver, the CTLE, the samplers, the VCO and the clock buffers), and CMOS where the rate is low and density matters (the slow serializer and deserializer stages, the PRBS logic and the CDR filter). Put inductors and long high-speed routes on the thick top metals.
\end{insight}

\section{How to use the rest of the book}\label{sec:ch1-roadmap}
\Cref{ch:signals} builds the vocabulary of links (pulse response, eye, BER), \cref{ch:channel} explains the \qty{17.4}{\deci\bel} channel and \cref{ch:linkanalysis} turns pulse responses into link budgets. Part~II (\cref{ch:eqfund,ch:txffe,ch:ctle,ch:dfe}) opens the closed eye with equalization. Part~III (\cref{ch:tx,ch:rx}) designs the transmitter and receiver circuits. Part~IV (\cref{ch:jitter,ch:pll,ch:cdr}) handles jitter and clocking. Part~V (\cref{ch:arch,ch:flow,ch:labs}) assembles the architecture, runs the open-source flow and provides the labs.

\begin{takeaways}
\begin{itemize}
  \item Parallel buses stop scaling because of skew, pin count and switching noise; a serial link sends one fast differential stream and recovers the clock from the data.
  \item A SerDes link is a serializer, an FFE driver, a lossy channel, a CTLE, samplers with a DFE, a deserializer, a PLL and a CDR.
  \item In \OS{}, 16 bits at \qty{312.5}{\mega\bit\per\second} become \qty{5}{\giga\bit\per\second} NRZ with a \qty{200}{\pico\second} UI and a \qty{2.5}{\giga\hertz} Nyquist frequency, using a half-rate \qty{2.5}{\giga\hertz} clock.
  \item \qty{5}{\giga\bit\per\second} NRZ matches PCIe~2.0, USB~3.2 Gen~1 and 5GBASE-KR, and the \qty{17.4}{\deci\bel} reference channel forces real equalization.
  \item The targets are BER $\le 10^{-12}$ over the REF channel with a 3-tap FFE, a 0--12\,dB CTLE, a 2-tap DFE, an LC PLL and a PI-based CDR.
  \item \sgx{} HBTs ($f_T \approx \qty{300}{\giga\hertz}$) far outrun \qty{130}{\nano\meter} CMOS (FO4 $\approx$ \qty{45}{\pico\second}), so the fast blocks use HBT CML; $BV_{CEO} \approx \qty{1.7}{\volt}$ limits the stack, and the thick top metals carry inductors and lines.
\end{itemize}
\end{takeaways}

\tikzsetfigurename{02_signals-}
%% 02_signals.tex -- Chapter 2: Signals, Systems and the Language of Links. Author: p1
\chapter{Signals, Systems and the Language of Links}\label{ch:signals}
\chapterlead{Five ideas carry the whole book: the unit interval, the pulse response, the cursors, the eye diagram and the bit-error rate.}

\begin{objectives}
\begin{itemize}
  \item describe an NRZ stream by its unit interval, bit rate and Nyquist frequency, and explain why most of its energy lies below $\fN$;
  \item generate PRBS test patterns with an LFSR and know what PRBS7 and PRBS31 stress;
  \item predict how a simple low-pass channel slows edges, using impulse, step and frequency responses;
  \item build a received waveform from shifted pulse responses and read the main, pre- and post-cursors;
  \item read an eye diagram (height, width, crossing, slope) and turn an eye opening and an rms noise into a BER.
\end{itemize}
\end{objectives}

Here the channel is deliberately simple, a one-pole RC filter or a short line; the realistic REF channel waits for \cref{sec:refchannel}, and the tools carry over unchanged.

\section{NRZ signaling: bits, symbols and the unit interval}\label{sec:nrz}
A link carries \textbf{bits} $b_k \in \{0,1\}$. The simplest way to put them on a wire is \textbf{non-return-to-zero} (NRZ) signaling, also called PAM2: each bit becomes a \textbf{symbol} $a_k = 2b_k - 1 \in \{-1,+1\}$, and the transmitter holds the voltage $A\,a_k$ for one bit period (\cref{fig:ch2-nrz}). In \OS{}, $A = \qty{0.5}{\volt}$ (\qty{1}{\Vppd}). One symbol lasts one \textbf{unit interval} (UI). NRZ sends one bit per symbol, so the \textbf{bit rate} $\fb$ equals the \textbf{baud rate}; PAM4 sends two (\cref{sec:pam4}).

\begin{figure}[htbp]
  \centering
  \includegraphics[width=\linewidth]{ch02_nrz_wave}
  \caption{\textbf{An NRZ waveform is a sequence of $\pm A$ levels, one per UI.} Each bit $b_k$ maps to a symbol $a_k = \pm1$ held for $\Tb = \qty{200}{\pico\second}$. Gold dots mark the transitions that clock recovery feeds on; a run of identical bits carries no timing information.}
  \label{fig:ch2-nrz}
\end{figure}

The fastest NRZ pattern, \code{1010\ldots}, is a square wave with a period of $2\Tb$. Its fundamental frequency is the \textbf{Nyquist frequency}, the single most important frequency in link design.

\begin{keyeq}[Unit interval and Nyquist frequency]
\begin{equation}
  \Tb = 1\,\UI = \frac{1}{\fb}, \qquad \fN = \frac{1}{2\Tb} = \frac{\fb}{2}
  \label{eq:ch2-nyquist}
\end{equation}
For \OS{}: $\fb = \qty{5}{\giga\bit\per\second}$, $\Tb = \qty{200}{\pico\second}$ and $\fN = \qty{2.5}{\giga\hertz}$.
\end{keyeq}

\paragraph{The spectrum of random NRZ} An NRZ waveform is a sum of rectangular pulses,
\[ x(t) = A\sum_n a_n\, g(t - n\Tb), \]
where $g(t)$ is a 1-UI rectangle with Fourier transform $G(f) = \Tb\,\mathrm{sinc}(f\Tb)\,e^{-j\pi f\Tb}$ and $\mathrm{sinc}(u) = \sin(\pi u)/(\pi u)$. If the symbols are independent and equally likely to be $\pm1$, then $\mathrm{E}[a_n a_m]$ is 1 for $n = m$ and 0 otherwise. The cross terms average away, and each pulse contributes its own energy spectrum $A^2\lvert G(f)\rvert^2$ once every $\Tb$ seconds. The power spectral density (PSD) is therefore $A^2\lvert G(f)\rvert^2/\Tb$:

\begin{keyeq}[PSD of random NRZ]
\begin{equation}
  S_x(f) = A^2\,\Tb\,\mathrm{sinc}^2(f\Tb), \qquad \text{nulls at } f = \fb, 2\fb, \dots
  \label{eq:ch2-psd}
\end{equation}
\end{keyeq}

\Cref{fig:ch2-psd}a plots \cref{eq:ch2-psd}. The main lobe runs from DC to the first null at $\fb = \qty{5}{\giga\hertz}$, and at $\fN$ the PSD is $\mathrm{sinc}^2(0.5) = (2/\pi)^2$, or \qty{-3.9}{\deci\bel}. Integrating the PSD (\cref{fig:ch2-psd}b) shows that 77\% of the power lies below $\fN$ and 90\% below $\fb$.

\begin{figure}[htbp]
  \centering
  \includegraphics[width=\linewidth]{ch02_nrz_psd}
  \caption{\textbf{Random NRZ puts most of its power below the Nyquist frequency.} (a)~The PSD follows $\mathrm{sinc}^2(f\Tb)$, with nulls at multiples of $\fb = \qty{5}{\giga\hertz}$ and only \qty{3.9}{\deci\bel} of roll-off at $\fN = \qty{2.5}{\giga\hertz}$. (b)~Cumulative power: 77\% lies below $\fN$ and 90\% below $\fb$. This is why channel loss is always quoted at $\fN$.}
  \label{fig:ch2-psd}
\end{figure}

\begin{example}{Where the power of a \qty{5}{\giga\bit\per\second} stream sits}\label{we:ch2-nyquist}
The fraction of power below $\fN$ is $2\int_0^{1/2}\mathrm{sinc}^2(u)\,du = 0.774$, so an ideal low-pass filter at \qty{2.5}{\giga\hertz} would still pass three quarters of the power. The REF channel's \qty{17.4}{\deci\bel} at \qty{2.5}{\giga\hertz}, a factor of 7.4 in amplitude, attacks exactly that band.
\end{example}

\begin{insight}
Treat $\fN = \qty{2.5}{\giga\hertz}$ as the frequency your link lives at: quote channel loss there, and give every block enough bandwidth around it (\cref{sec:ch2-isi}). The energy above $\fb$ is under 10\% and mostly edge sharpness.
\end{insight}

\section{PRBS test patterns}\label{sec:prbs}
Testing needs data that looks random but is reproducible, so the receiver can check every bit. A \textbf{pseudo-random bit sequence} (PRBS) does this. It comes from a \textbf{linear-feedback shift register} (LFSR): a chain of $n$ flip-flops whose input is the XOR of two of its outputs. \Cref{fig:ch2-lfsr} shows PRBS7, defined by the polynomial $x^7 + x^6 + 1$: the new bit is stage~6 XOR stage~7. With the right taps, the register steps through all $2^n - 1$ non-zero states before it repeats, producing a \textbf{maximal-length sequence}.

\begin{figure}[htbp]
\centering
% ch02_lfsr.tex -- 7-stage Fibonacci LFSR for PRBS7, polynomial x^7 + x^6 + 1
\begin{tikzpicture}[x=1cm,y=1cm,
  ff/.style={blk,minimum width=11mm,minimum height=9mm,font=\sffamily\footnotesize,fill=btint}]
\foreach \i in {1,...,7}{
  \node[ff] (f\i) at ({1.65*\i},0) {D\quad Q};
  \node[lbl,anchor=south] at ($(f\i.north)+(0,0.0)$) {stage \i};
}
\foreach \i [evaluate=\i as \j using int(\i+1)] in {1,...,6}{\draw[sig] (f\i.east) -- (f\j.west);}
% clock
\draw[clk,-] (1.0,-0.85) -- (11.75,-0.85);
\foreach \i in {1,...,7}{\draw[clk] ({1.65*\i},-0.85) -- (f\i.south);}
\node[lbl,text=golddk,anchor=east] at (1.0,-0.85) {clock $f_b$};
% taps from stage 6 and 7 outputs into XOR, back to stage 1
\node[sumnode,font=\small] (x) at (10.725,1.55) {$\oplus$};
\draw[sig] ($(f6.east)!0.5!(f7.west)$) node[dot]{} -- (x.south);
\draw[sig] (f7.east) -- ++(0.35,0) node[dot](o){} |- (x.east);
\draw[sig] (x.north) -- ++(0,0.35) -| ($(f1.west)+(-0.45,0)$) -- (f1.west);
\draw[sig] (o) -- ++(0.9,0) node[right,font=\sffamily\footnotesize,text=navy] {PRBS7 out};
\node[lbl,text=navy,anchor=south] at (6.0,2.3) {feedback $=$ stage 6 $\oplus$ stage 7 \;\; $\Leftrightarrow$\;\; $x^7+x^6+1$};
\end{tikzpicture}

\caption{\textbf{A 7-stage LFSR generates PRBS7.} Each clock edge shifts the bits right and feeds stage~6 XOR stage~7 into stage~1; from any non-zero state the output repeats after $2^7 - 1 = 127$ bits. A 31-stage register with taps 28 and 31 gives PRBS31.}
\label{fig:ch2-lfsr}
\end{figure}

A PRBS-$n$ sequence has a few properties worth remembering (\cref{fig:ch2-prbs}):
\begin{itemize}
  \item \textbf{Period} $2^n - 1$ bits, during which every $n$-bit word except all zeros appears exactly once.
  \item \textbf{Nearly balanced:} $2^{n-1}$ ones and $2^{n-1} - 1$ zeros (64 and 63 for PRBS7).
  \item \textbf{Run lengths:} the longest run is $n$ ones and $n-1$ zeros. Half of all runs have length 1, a quarter length 2, and so on, like a fair coin.
  \item \textbf{Line spectrum:} because the pattern repeats, its spectrum is a set of lines spaced $\fb/(2^n-1)$ apart under the same $\mathrm{sinc}^2$ envelope as random NRZ.
\end{itemize}

\begin{figure}[htbp]
  \centering
  \includegraphics[width=\linewidth]{ch02_prbs7}
  \caption{\textbf{PRBS7 behaves like random data with short memory.} (a)~One full period of 127 bits, with the longest runs (seven 1s, six 0s) shaded. (b)~Run-length histogram: 32 runs of length 1, 16 of length 2, halving at each step, as for coin flips. (c)~The line spectrum follows the $\mathrm{sinc}^2$ envelope of \cref{fig:ch2-psd}. (d)~Zoom near DC: lines every $\fb/127 = \qty{39.4}{\mega\hertz}$, with an almost empty DC line because the pattern is balanced.}
  \label{fig:ch2-prbs}
\end{figure}

\paragraph{Choosing a PRBS} PRBS7 repeats every \qty{25.4}{\nano\second} and never holds a level longer than 7~UI: ideal for bring-up. PRBS31 ($x^{31}+x^{28}+1$) has runs of up to 31 bits and a period of \qty{0.43}{\second}; its low-frequency content stresses AC coupling (baseline wander), long ISI tails and the CDR's ability to hold phase without transitions, so it is the compliance pattern. Model work uses PRBS13 (8191 bits), far longer than the channel's ISI memory. Real protocols also \emph{encode} the data: \textbf{8b/10b} (PCIe~1/2, USB~3.2 Gen~1, SATA) maps each byte to 10 bits with at most five equal bits in a row and long-term DC balance, at 25\% overhead.

\begin{pitfall}[A short pattern can hide a long problem]
If the link remembers more bits than the pattern varies over, the eye looks too good: PRBS7 can pass where PRBS31 fails on baseline wander or a long reflection. Confirm designs with a pattern longer than every memory in the link.
\end{pitfall}

\section{Linear systems in one section}\label{sec:lti}
Between the TX output and the RX sampler, the link is close to a \textbf{linear time-invariant} (LTI) system, fully described by its \textbf{impulse response} $h(t)$, the output for a very short unit-area pulse. For any input $x(t)$ the output is the \textbf{convolution}
\begin{equation}
  y(t) = (x * h)(t) = \int_{-\infty}^{\infty} x(\tau)\,h(t-\tau)\,d\tau ,
  \qquad Y(f) = X(f)\,H(f),
  \label{eq:ch2-conv}
\end{equation}
where the \textbf{frequency response} $H(f)$ is the Fourier transform of $h(t)$. Convolution becomes multiplication, so cascaded blocks multiply. Gains are quoted in decibels, $\lvert H\rvert_\mathrm{dB} = 20\log_{10}\lvert H(f)\rvert$: \qty{-3}{\deci\bel} is a factor 0.707, \qty{-6}{\deci\bel} one half, \qty{-20}{\deci\bel} one tenth, and in a cascade the dB values add. A \textbf{Bode plot} shows $\lvert H\rvert_\mathrm{dB}$ against log frequency.

\paragraph{The one-pole system} A resistor driving a capacitor is the simplest low-pass filter and a good first model of a slow channel or a CML stage with load $\RL\CL$ (\cref{fig:ch2-onepole}):

\begin{keyeq}[One-pole low-pass filter]
\begin{equation}
  H(f) = \frac{1}{1 + jf/f_\mathrm{3dB}},\quad
  h(t) = \frac{1}{\tau}e^{-t/\tau},\quad
  s(t) = 1 - e^{-t/\tau},\quad
  t_r^{10\text{--}90} = \tau\ln 9 \approx \frac{0.35}{f_\mathrm{3dB}}
  \label{eq:ch2-onepole}
\end{equation}
\end{keyeq}
Here $\tau = RC = 1/(2\pi f_\mathrm{3dB})$, $s(t)$ is the unit step response (for $t \ge 0$) and $t_r$ is the 10--90\% rise time. Above $f_\mathrm{3dB}$ the gain falls by \qty{20}{\deci\bel} per decade.

\begin{figure}[htbp]
  \centering
  \includegraphics[width=\linewidth]{ch02_onepole}
  \caption{\textbf{Bandwidth and rise time are two views of the same pole.} (a)~A one-pole response with $f_\mathrm{3dB} = \qty{2.5}{\giga\hertz}$: \qty{-3}{\deci\bel} at the corner, then \qty{-20}{\deci\bel} per decade. (b)~Its step response rises from 10\% to 90\% in \qty{140}{\pico\second} $= 0.35/f_\mathrm{3dB}$, which is 0.7~UI at \qty{5}{\giga\bit\per\second}.}
  \label{fig:ch2-onepole}
\end{figure}

\begin{example}{What bandwidth keeps ISI small at \qty{5}{\giga\bit\per\second}?}\label{we:ch2-bw}
For a bit to barely affect the next, the output must almost settle within one UI. One UI after a step, a one-pole system still has an error of $e^{-\Tb/\tau}$. Asking for at most 5\%:
\begin{itemize}
  \item $e^{-\Tb/\tau} \le 0.05 \;\Rightarrow\; \Tb/\tau \ge \ln 20 = 3.0 \;\Rightarrow\; \tau \le \qty{66.8}{\pico\second}$;
  \item $f_\mathrm{3dB} = 1/(2\pi\tau) \ge 3.0/(2\pi\times\qty{200}{\pico\second}) = \qty{2.38}{\giga\hertz}$;
  \item the rise time is then $0.35/\qty{2.38}{\giga\hertz} = \qty{147}{\pico\second} \approx 0.73\,\UI$.
\end{itemize}
A one-pole bandwidth close to the Nyquist frequency, $f_\mathrm{3dB} \approx \fN = \qty{2.5}{\giga\hertz}$, keeps the ISI to a few percent. \Cref{sec:ch2-isi} turns this into an exact eye opening.
\end{example}

\section{The pulse response}\label{sec:pulse-response}
A link never sends impulses; it sends 1-UI rectangles. The response to one isolated rectangle of \qty{1}{\volt} and width $\Tb$ is the \textbf{pulse response} $p(t)$, the most useful single picture of a link:
\begin{equation}
  p(t) = h(t) * \Pi_{\Tb}(t) = \int_{t-\Tb}^{t} h(\tau)\,d\tau = s(t) - s(t - \Tb),
  \label{eq:ch2-pulse}
\end{equation}
where $\Pi_{\Tb}(t)$ is 1 for $0 \le t < \Tb$: a step up, then a step down one UI later. An NRZ waveform is a sum of shifted, signed rectangles and the channel is linear, so the received signal is a sum of shifted, signed pulse responses (\cref{fig:ch2-superposition}). Sample it once per UI at $\ts + n\Tb$, with $\ts$ usually at the peak of $p(t)$. The samples of the pulse response at integer UI offsets are the \textbf{cursors} $\pk{k} = p(\ts + k\Tb)$ (\cref{fig:ch2-pulse}). The \textbf{main cursor} $\pk{0}$ carries the wanted bit. \textbf{Post-cursors} $\pk{k}$, $k > 0$, are the tail of a pulse sent $k$ bits \emph{earlier}. \textbf{Pre-cursors} $\pk{k}$, $k < 0$, are the leading edge of a pulse sent $\lvert k\rvert$ bits \emph{later}.

\begin{keyeq}[Superposition of pulse responses]
\begin{equation}
  y(t) = A\sum_n a_n\, p(t - n\Tb)
  \qquad\Longrightarrow\qquad
  y_n = y(\ts + n\Tb) = \underbrace{A\,a_n\,\pk{0}}_{\text{wanted}} \;+\; \underbrace{A\sum_{k\neq0} a_{n-k}\,\pk{k}}_{\text{intersymbol interference (ISI)}}
  \label{eq:ch2-isi}
\end{equation}
\end{keyeq}

The channel acts like an FIR filter whose taps are the cursors; everything except the main cursor is \textbf{intersymbol interference} (ISI), the echo of neighboring bits at the sampling instant.

\begin{figure}[htbp]
  \centering
  \includegraphics[width=\linewidth]{ch02_pulse_cursors}
  \caption{\textbf{The pulse response and its cursors.} (a)~A \qty{1}{\volt}, 1-UI pulse through a mild \qty{10}{\deci\bel}-at-Nyquist line ($0.57$\,m of the loss family) comes out lower, wider and with a long tail. Sampling it once per UI around the peak gives the cursors. (b)~The main cursor $\pk{0} = 0.556$ is the useful signal; $\pk{-1} = 0.020$ is a pre-cursor and $\pk{1} = 0.160$, $\pk{2} = 0.063$, \ldots{} are post-cursors.}
  \label{fig:ch2-pulse}
\end{figure}

\begin{figure}[htbp]
  \centering
  \includegraphics[width=\linewidth]{ch02_superposition}
  \caption{\textbf{A received waveform is a sum of shifted pulse responses.} (a)~Each bit launches its own pulse $A\,a_n\,p(t-n\Tb)$ (blue for $+1$, amber for $-1$). (b)~Their sum is the received signal; gold dots are the samples at $\ts + n\Tb$. A 1 that follows another 1 gets help from its neighbor's tail (\qty{334}{\milli\volt}), while a lone 1 after three 0s is pulled down to \qty{156}{\milli\volt}. The line is idle before bit~0.}
  \label{fig:ch2-superposition}
\end{figure}

\begin{example}{Reading a sample from the cursors}\label{we:ch2-cursors}
Bit~7 in \cref{fig:ch2-superposition} is a lone 1 after three 0s, followed by a 0. With $A = \qty{0.5}{\volt}$ and the cursors of \cref{fig:ch2-pulse}:
\begin{itemize}
  \item wanted signal: $A\pk{0} = 0.5\times0.5561 = \qty{278}{\milli\volt}$;
  \item the three previous 0s and the next 0 subtract: $-A(\pk{1}+\pk{2}+\pk{3}+\pk{-1}) = -0.5\,(0.1605+0.0629+0.0340+0.0205) = \qty{-139}{\milli\volt}$;
  \item the older bits ($a_3 = a_2 = a_0 = +1$, $a_1 = -1$) add $0.5\,(0.0216+0.0151-0.0112+0.0087) = \qty{+17}{\milli\volt}$.
\end{itemize}
The sample is $278 - 139 + 17 = \qty{156}{\milli\volt}$, as simulated: ISI took 44\% of the wanted signal.
\end{example}

\begin{insight}
The cursors are the channel's fingerprint, and every equalizer in Part~II acts on them: the TX FFE reshapes the pulse (\cref{ch:txffe}), the CTLE sharpens it (\cref{ch:ctle}) and the DFE subtracts post-cursors (\cref{ch:dfe}). Faced with a new channel, plot its pulse response first.
\end{insight}

\begin{tryit}[Lab 1: pulse response and cursors]
\code{lab1\_channel.py} turns the REF channel's $S_{21}$ into a pulse response and its cursors. Lengthen the line and watch $\pk{1}/\pk{0}$ grow (\cref{sec:lab1}).
\end{tryit}

\section{ISI: how pulse spreading closes the eye}\label{sec:ch2-isi}
The one-pole channel shows in closed form how bandwidth turns into ISI. Its pulse response peaks at the end of the bit, $t = \Tb$. Sampling there and writing $r = e^{-\Tb/\tau}$ gives $\pk{0} = 1 - r$ and $\pk{k} = (1-r)\,r^{k}$ for $k \ge 1$, with no pre-cursors. The post-cursors form a geometric series that sums to $r$. The worst case for a 1 is that every earlier bit is a 0, so each tail pulls it down. The smallest 1 is then $A(\pk{0} - \sum_{k\ge1}\pk{k})$, and by symmetry the eye opening at the sampling instant is

\begin{keyeq}[Worst-case eye of a one-pole channel]
\begin{equation}
  \mathrm{EH} = 2A\Bigl(\pk{0} - \sum_{k\ge1}\pk{k}\Bigr) = 2A\,\bigl(1 - 2e^{-\Tb/\tau}\bigr), \qquad \frac{\Tb}{\tau} = 2\pi f_\mathrm{3dB}\Tb = \pi\,\frac{f_\mathrm{3dB}}{\fN}
  \label{eq:ch2-onepole-eye}
\end{equation}
\end{keyeq}

\Cref{fig:ch2-isi-tau} shows the pulse and the eye for three time constants. With $\tau = 0.2\,\UI$ ($f_\mathrm{3dB} = \qty{4.0}{\giga\hertz}$) the eye is 98.7\% open. At $\tau = 0.5\,\UI$ (\qty{1.6}{\giga\hertz}) it is 73\% open, and at $\tau = 1\,\UI$ (\qty{0.8}{\giga\hertz}) only 26\%. The eye closes completely when $r = 1/2$, that is $\tau = \Tb/\ln 2 = 1.44\,\UI$, or $f_\mathrm{3dB} = \qty{552}{\mega\hertz}$. For the bandwidth of \cref{we:ch2-bw}, $f_\mathrm{3dB} = \fN$, the exponent is $\pi$ and $\mathrm{EH} = 2A(1 - 2e^{-\pi}) = 0.91 \times 2A$.

\begin{figure}[htbp]
  \centering
  \includegraphics[width=\linewidth]{ch02_isi_tau}
  \caption{\textbf{Slower channels spread each pulse over more UIs and close the eye.} Top: one-pole pulse responses for $\tau = 0.2$, 0.5 and 1~UI; filled dots are the main cursor, open dots the post-cursors. Bottom: PRBS7 eyes through the same channels ($2A = \qty{1}{\volt}$). The measured openings, 987, 729 and \qty{265}{\milli\volt}, match $2A(1-2e^{-\Tb/\tau})$ from \cref{eq:ch2-onepole-eye}. The arrow sits at the best sampling instant, which for a one-pole channel is the pulse peak at the end of the bit, not the visual center of the eye.}
  \label{fig:ch2-isi-tau}
\end{figure}

Subtracting the magnitudes of all other cursors from the main cursor works for any channel; this is \textbf{peak-distortion analysis} (\cref{sec:pda}). For the \qty{10}{\deci\bel} line of \cref{fig:ch2-pulse} it gives \qty{230}{\milli\volt} out of \qty{1}{\volt}; for the REF channel it is negative: the eye is closed.

\begin{insight}[Design rule of thumb]
Each one-pole block in the data path needs $f_\mathrm{3dB} \gtrsim \fN = \qty{2.5}{\giga\hertz}$ to cost less than about 10\% of eye height. A cascade of several poles has less bandwidth than any single one, so in practice each CML stage of \OS{} should reach 0.7--1\,$\fb$ (3.5--\qty{5}{\giga\hertz}).
\end{insight}

\section{Eye diagrams}\label{sec:eye}
An \textbf{eye diagram} cuts the received waveform into 2-UI windows aligned to the bit clock and overlays them (\cref{fig:ch2-eye-fold}). Every bit pattern draws its own trace, so the overlay shows at a glance how much the patterns differ at each instant.

\begin{figure}[htbp]
  \centering
  \includegraphics[width=\linewidth]{ch02_eye_fold}
  \caption{\textbf{An eye diagram folds the waveform every 2~UI.} (a)~Twelve UI of PRBS7 after a one-pole channel ($\tau = 0.45\,\UI$), cut into six 2-UI windows. (b)~The same segments, same colors, overlaid on 600 others (grey).}
  \label{fig:ch2-eye-fold}
\end{figure}

\Cref{fig:eye-anatomy} names the parts you will measure throughout this book.
\begin{itemize}
  \item \textbf{Eye height} (EH): the vertical opening at the sampling instant, between the lowest 1 and the highest 0. It must beat noise, offset and sampler sensitivity.
  \item \textbf{Eye width} (EW): the horizontal opening at the threshold, between adjacent crossings. It must beat clock jitter and the sampler aperture.
  \item \textbf{Crossing points}: where edges cross the threshold. Their spread in time is \textbf{jitter}; their mean is what the CDR locks to.
  \item \textbf{Optimal sampling point}: the instant and threshold with the most margin both ways.
  \item \textbf{Slope} near the crossing: it converts timing error into voltage error, $\Delta v \approx (dV/dt)\,\Delta t$.
\end{itemize}

\begin{figure}[htbp]
  \centering
  \includegraphics[width=\linewidth]{ch02_eye_anatomy}
  \caption{\textbf{Anatomy of an eye diagram.} A PRBS7 stream through a smooth band-limited channel, with \qty{6}{\pico\second} rms edge jitter and \qty{15}{\milli\volt} rms noise. Eye height and eye width are measured between the 0.1\% tails of the traces. The right histogram shows the spread of the sampled levels (noise plus ISI) at $\ts$; the bottom histogram shows the spread of the crossing times (jitter). The orange line is the edge slope, here about \qty{2.8}{\milli\volt\per\pico\second}.}
  \label{fig:eye-anatomy}
\end{figure}

With the slope of \cref{fig:eye-anatomy}, sampling \qty{4}{\pico\second} off near an edge costs about $2.8 \times 4 \approx \qty{11}{\milli\volt}$. Noise eats the eye vertically and jitter horizontally; the next section asks how much is tolerable.

\begin{pitfall}[An eye is only as deep as the number of bits behind it]
A simulated eye of $10^4$ bits only shows events more frequent than about $10^{-4}$; a $7\sigma$ noise peak will almost never appear. An open eye in a short simulation says nothing about $10^{-12}$. \Cref{sec:stateye} computes the eye at that probability directly.
\end{pitfall}

\begin{tryit}[Lab 2: PRBS7 eye through the REF channel]
\code{lab2\_prbs\_eye.py} draws the PRBS7 eye before and after the REF channel. The output eye is closed ($\qty{-78}{\milli\volt}$): the motivation for Part~II (\cref{sec:lab2}).
\end{tryit}

\section{Noise and the bit-error rate}\label{sec:ber}
The receiver compares each sample $y_n$ with a threshold at 0, and an error happens when noise pushes a sample across it. With a worst-case eye of height EH, a 1 arrives at least $\mathrm{EH}/2$ above the threshold and a 0 at least $\mathrm{EH}/2$ below it. Add Gaussian noise $n \sim \mathcal{N}(0, \sign^2)$ from the termination, the front end and the sampler (\cref{fig:ch2-ber}a). A 1 is read as a 0 when $n < -\mathrm{EH}/2$:
\begin{align}
  P(\text{error}\mid 1) &= \int_{-\infty}^{-\mathrm{EH}/2} \frac{1}{\sqrt{2\pi}\,\sign}\,e^{-v^2/(2\sign^2)}\,dv = \Qf\!\left(\frac{\mathrm{EH}}{2\sign}\right), \notag\\
  \Qf(x) &= \int_x^{\infty}\frac{e^{-u^2/2}}{\sqrt{2\pi}}\,du = \tfrac12\erfc\!\left(\frac{x}{\sqrt2}\right). \label{eq:ch2-qdef}
\end{align}
By symmetry $P(\text{error}\mid 0)$ is the same. With 1s and 0s equally likely, the BER is the average of the two:

\begin{keyeq}[Bit-error rate with Gaussian noise]
\begin{equation}
  \BER = \Qf\!\left(\frac{\mathrm{EH}}{2\,\sign}\right), \qquad \SNR \equiv \left(\frac{\mathrm{EH}}{2\,\sign}\right)^{2} \;\Rightarrow\; \BER = \Qf\bigl(\sqrt{\SNR}\bigr)
  \label{eq:ber-q}
\end{equation}
\end{keyeq}

The argument of $\Qf$ is half the eye in units of rms noise. With the worst-case EH, \cref{eq:ber-q} is an upper bound, since most patterns leave more margin (\cref{sec:stateye} computes the average). \Cref{tab:ch2-qtable} lists common values. Remember $\Qf = 7.03$ for $10^{-12}$: the eye must be at least $2\Qf = 14.07$ times the rms noise.

\begin{figure}[htbp]
  \centering
  \includegraphics[width=\linewidth]{ch02_ber}
  \caption{\textbf{BER is the area of a Gaussian tail beyond half the eye.} (a)~Received 1s and 0s spread by noise of rms $\sign$ around $\pm\mathrm{EH}/2$; the shaded tails beyond the threshold are errors. (b)~BER against $\SNR = 20\log_{10}(\mathrm{EH}/2\sign)$. The curve is a cliff: near $10^{-12}$ (Q = 7.03, \qty{16.9}{\deci\bel}), one extra dB of SNR buys about three decades of BER.}
  \label{fig:ch2-ber}
\end{figure}

\begin{example}{How much noise can the eye tolerate?}\label{we:ch2-ber}
The target is $\BER \le 10^{-12}$, so $\mathrm{EH}/(2\sign) \ge 7.03$ and $\sign \le \mathrm{EH}/14.07$.
\begin{itemize}
  \item A sampler with $\sign = \qty{2}{\milli\volt}$ rms needs an eye of at least $14.07\times\qty{2}{\milli\volt} = \qty{28}{\milli\volt}$.
  \item Lab~5 opens the REF channel to $\mathrm{EH} = \qty{143}{\milli\volt}$ with an \qty{8}{\deci\bel} CTLE and a 3-tap DFE. The noise budget is then $\sign \le 143/14.07 = \qty{10.2}{\milli\volt}$ rms, for all noise sources together.
  \item At $\sign = \qty{15}{\milli\volt}$, $\Qf = 4.77$ and $\BER \approx 9\times10^{-7}$: 1.5$\times$ the noise costs six decades.
\end{itemize}
\end{example}

\paragraph{A first look at the bathtub} Jitter moves the crossings, so the BER also depends on \emph{when} we sample. Sweeping the sampling instant across the UI gives a \textbf{bathtub curve} (\cref{fig:ch2-bathtub}); its width at $10^{-12}$ is the usable eye width. With \qty{4}{\pico\second} rms random jitter (0.02~UI), 0.72~UI remains; doubling the jitter leaves 0.45~UI. \Cref{ch:jitter} covers this in full (\cref{sec:dualdirac}).

\begin{figure}[htbp]
\centering
\begin{minipage}[b]{0.45\linewidth}
  \centering\small
  \captionof{table}{\textbf{Q values for common BER targets.} $\BER = \Qf(x)$; $2\Qf$ is the eye-to-noise ratio needed.}
  \label{tab:ch2-qtable}
  \begin{tabular}{@{}rccc@{}}
  \toprule
  BER & $\Qf$ & $2\Qf$ & SNR (dB) \\
  \midrule
  $10^{-3}$  & 3.09 & 6.18  & 9.8 \\
  $10^{-6}$  & 4.75 & 9.51  & 13.5 \\
  $10^{-9}$  & 6.00 & 12.00 & 15.6 \\
  $10^{-10}$ & 6.36 & 12.72 & 16.1 \\
  \rowcolor{gtint} $10^{-12}$ & \textbf{7.03} & \textbf{14.07} & \textbf{16.9} \\
  $10^{-15}$ & 7.94 & 15.88 & 18.0 \\
  \bottomrule
  \end{tabular}
  \vspace{4mm}
\end{minipage}\hfill
\begin{minipage}[b]{0.52\linewidth}
  \centering
  \includegraphics{ch02_bathtub}
  \caption{\textbf{Bathtub curves.} BER against sampling instant for Gaussian edge jitter of 0.02 and 0.04~UI rms (half the edges are transitions). The opening at $10^{-12}$ shrinks from 0.72 to 0.45~UI.}
  \label{fig:ch2-bathtub}
\end{minipage}
\end{figure}

\begin{pitfall}[You cannot measure $10^{-12}$ in a coffee break]
Showing $\BER < 10^{-12}$ at 95\% confidence takes $3/\BER = 3\times10^{12}$ error-free bits: ten minutes on the bench, far beyond any simulation. Design with statistical models; verify with PRBS checkers.
\end{pitfall}

\section{Exercises}\label{sec:ch2-exercises}
\begin{enumerate}[label=\textbf{2.\arabic*},leftmargin=2.2em]
  \item \textbf{Rates.} A 2.5GBASE-KX lane runs at \qty{3.125}{\giga\bit\per\second} with 8b/10b. Find $\Tb$, $\fN$, the first PSD null and the payload rate.
  \item \textbf{PRBS9} ($x^9 + x^5 + 1$). Draw its LFSR; give its period, the counts of 1s and 0s, the longest runs and the line spacing at \qty{5}{\giga\bit\per\second}.
  \item \textbf{A CML stage as a pole.} For $\RL = \qty{200}{\ohm}$, $\CL = \qty{150}{\femto\farad}$, find $f_\mathrm{3dB}$, $t_r$ and $1 - 2e^{-\Tb/\tau}$. What $\CL$ costs 10\% of the eye?
  \item \textbf{Cursors by hand.} With the cursors of \cref{fig:ch2-pulse} ($\pk{-1}..\pk{3}$ = 0.0205, 0.5561, 0.1605, 0.0629, 0.0340) and $A = \qty{0.5}{\volt}$, find the sample of the bracketed 1 in \code{1 1 1 [1] 0} and in \code{0 0 0 [1] 1}. Compare the worst-case eye from these cursors with the \qty{230}{\milli\volt} of the full tail.
  \item \textbf{BER.} A receiver sees $\mathrm{EH} = \qty{60}{\milli\volt}$ with $\sign = \qty{4}{\milli\volt}$. Find Q and the BER. What eye height or noise gives $10^{-15}$, and how long must the link run error-free to show it?
\end{enumerate}

\begin{takeaways}
\begin{itemize}
  \item At \qty{5}{\giga\bit\per\second}, $1\,\UI = \qty{200}{\pico\second}$ and $\fN = \qty{2.5}{\giga\hertz}$; 77\% of NRZ power lies below $\fN$.
  \item PRBS-$n$ (an $n$-stage LFSR) repeats every $2^n - 1$ bits with runs up to $n$; PRBS7 for bring-up, PRBS31 for compliance.
  \item A one-pole block has $t_r \approx 0.35/f_\mathrm{3dB}$ and a worst-case eye of $2A(1 - 2e^{-\Tb/\tau})$; keeping ISI small needs $f_\mathrm{3dB} \gtrsim \fN$.
  \item The received signal is $y(t) = A\sum_n a_n\,p(t - n\Tb)$: main cursor plus ISI from the cursors $\pk{k}$, with worst-case eye $2A(\pk{0} - \sum_{k\ne0}\lvert\pk{k}\rvert)$.
  \item With Gaussian noise, $\BER = \Qf(\mathrm{EH}/2\sign)$, so $10^{-12}$ needs $\mathrm{EH} \ge 14.07\,\sign$.
\end{itemize}
\end{takeaways}

\tikzsetfigurename{03_channel-}
\chapter{The Channel}\label{ch:channel}
\chapterlead{Copper is not a wire at 5\,Gb/s: it is a lossy, reflective transmission line, and its pulse response sets almost every number in the rest of this book.}
% Omega for sans-serif contexts (captions, TikZ, sf tables): the sans text font has no Omega glyph.
\providecommand{\chThreeOhm}{\ensuremath{\Omega}}

\begin{objectives}
\begin{itemize}
  \item Decide when a trace is a transmission line, and compute $\Zo$ and the delay from its RLGC parameters.
  \item Explain why conductor loss grows as $\sqrt f$ and dielectric loss as $f$, and estimate FR-4 loss in dB per inch.
  \item Predict reflections from terminations, pads, bondwires and via stubs, and read them in a TDR plot.
  \item Use S-parameters, including the mixed-mode $S_{dd21}$ of a differential pair.
  \item Turn $S_{21}(f)$ into a causal impulse response and a pulse response.
  \item Know the \OS{} reference channel REF and its cursors.
\end{itemize}
\end{objectives}

%% =====================================================================
\section{What is the channel?}\label{sec:ch3-what}

The \textbf{channel} is everything between the transmitter's output transistors and the receiver's input transistors. \Cref{fig:ch3-path} follows a bit across a typical board-to-board link. Long traces and cables attenuate high frequencies (\textbf{loss}). Short lumps such as bondwires, pads and vias change the impedance and send part of the wave back (\textbf{reflections}). Dense pins and connectors couple neighboring lanes (\textbf{crosstalk}). We treat each mechanism in turn and then combine them into one transfer function $\Hch(f)$ and one pulse response $p(t)$, the two objects that the link-analysis tools of \cref{ch:linkanalysis} work with.

\begin{figure}[tbp]
\centering
\begin{tikzpicture}[
  seg/.style={blkb,minimum height=10mm,minimum width=14mm,inner sep=1pt,font=\sffamily\scriptsize,text width=13mm},
  die/.style={seg,blkn,text width=9mm,minimum width=10mm},
  imp/.style={font=\sffamily\scriptsize,text width=15mm,align=center,anchor=north,inner sep=0.5pt},
  node distance=2.4mm]
  \node[die] (tx) {TX\\die};
  \node[seg,right=of tx] (pk) {bondwire,\\package};
  \node[seg,right=of pk] (pcb) {PCB\\trace};
  \node[seg,right=of pcb] (via) {via};
  \node[seg,right=of via] (con) {connector};
  \node[seg,right=of con,fill=gtint,draw=golddk,minimum width=18mm,text width=17mm] (bp) {backplane\\or cable};
  \node[seg,right=of bp,text width=14mm,minimum width=15mm] (mir) {connector,\\via, PCB, pkg};
  \node[die,right=of mir] (rx) {RX\\die};
  \foreach \a/\b in {tx/pk,pk/pcb,pcb/via,via/con,con/bp,bp/mir,mir/rx}{\draw[sig] (\a) -- (\b);}
  \node[imp,below=1.5mm of tx]  {\textcolor{c4}{pad, ESD $C$;}\\\textcolor{c4}{$R_T$ error}};
  \node[imp,below=1.5mm of pk]  {\textcolor{c4}{series $L$}\\\textcolor{c4}{$\approx$1\,nH/mm}\\\textcolor{c5}{crosstalk}};
  \node[imp,below=1.5mm of pcb] {\textcolor{c2}{skin and}\\\textcolor{c2}{dielectric}\\\textcolor{c2}{loss}};
  \node[imp,below=1.5mm of via] {\textcolor{c4}{shunt $C$,}\\\textcolor{c4}{stub notch}};
  \node[imp,below=1.5mm of con] {\textcolor{c4}{$\Zo$ step}\\\textcolor{c5}{NEXT, FEXT}};
  \node[imp,below=1.5mm of bp,text width=18mm]  {\textcolor{c2}{most of the loss}\\\textcolor{c2}{(REF)}};
  \node[imp,below=1.5mm of mir] {the same\\again};
  \node[imp,below=1.5mm of rx]  {\textcolor{c4}{pad, ESD $C$;}\\\textcolor{c4}{$R_T$ error}};
  \draw[decorate,decoration={brace,amplitude=3pt},muted] ($(pk.north west)+(0,1.5mm)$) -- node[lbl,above=3pt]{package} ($(pk.north east)+(0,1.5mm)$);
  \draw[decorate,decoration={brace,amplitude=3pt},muted] ($(pcb.north west)+(0,1.5mm)$) -- node[lbl,above=3pt]{line card} ($(con.north east)+(0,1.5mm)$);
  \draw[decorate,decoration={brace,amplitude=3pt},muted] ($(bp.north west)+(0,1.5mm)$) -- node[lbl,above=3pt]{long interconnect} ($(bp.north east)+(0,1.5mm)$);
  \draw[decorate,decoration={brace,amplitude=3pt},muted] ($(mir.north west)+(0,1.5mm)$) -- node[lbl,above=3pt]{mirror image} ($(mir.north east)+(0,1.5mm)$);
  \node[lbl,above=3.2mm of tx] {on chip};
  \node[lbl,above=3.2mm of rx] {on chip};
  \node[font=\sffamily\scriptsize,anchor=north] at ($(via.south)!0.5!(con.south)+(0,-12mm)$)
    {\textcolor{c2}{$\blacksquare$}\,loss\qquad \textcolor{c4}{$\blacksquare$}\,reflection\qquad \textcolor{c5}{$\blacksquare$}\,crosstalk};
\end{tikzpicture}
\caption{\textbf{The physical path of one bit, with the main impairment of each segment.} Long, uniform segments (gold) cause most of the loss; short lumps cause reflections; dense pin fields cause crosstalk. The \OS{} reference channel REF lumps the long interconnect into 1\,m of FR-4 stripline with a 300\,fF pad at each end.}
\label{fig:ch3-path}
\end{figure}

%% =====================================================================
\section{Transmission lines}\label{sec:tline}

\subsection{When is a wire a transmission line?}
A wire is \textbf{lumped} when its voltage is the same everywhere at any instant. It becomes a \textbf{transmission line} when the signal changes noticeably while the wave travels along it. The usual rule compares the one-way delay $t_d$ with the rise time $t_r$~\citep{p2:johnson1993,hall2009}:
\begin{equation}
  t_d \gtrsim \frac{t_r}{6}
  \quad\Longleftrightarrow\quad
  \ell \gtrsim \frac{t_r\, v}{6},
  \label{eq:ch3-tline-rule}
\end{equation}
where $v$ is the propagation velocity. Equivalently, a line longer than about a tenth of a wavelength at the highest frequency of interest is a transmission line.

\subsection{The telegrapher model}
A transmission line is a ladder of tiny identical sections (\cref{fig:ch3-rlgc}). Each section of length $\Delta z$ has a series resistance $R\,\Delta z$ (the copper), a series inductance $L\,\Delta z$ (the magnetic field around the current), a shunt capacitance $C\,\Delta z$ (the electric field to the return plane) and a shunt conductance $G\,\Delta z$ (leakage through the dielectric). KVL and KCL around one section, in the limit $\Delta z\to0$, give the \textbf{telegrapher equations}
\begin{equation}
  \frac{\partial V}{\partial z} = -(R + j\omega L)\, I, \qquad
  \frac{\partial I}{\partial z} = -(G + j\omega C)\, V .
\end{equation}
Eliminating $I$ gives $\partial^2 V/\partial z^2 = \gamma^2 V$, whose solution is a forward and a backward wave, $V(z) = V^+e^{-\gamma z} + V^-e^{+\gamma z}$.

\begin{figure}[tbp]
\centering
\begin{tikzpicture}[scale=0.92, transform shape, font=\small]
  \draw (0,2) node[ocirc]{} to[short, i=$I(z)$] (0.9,2) to[R, l=$R\,\Delta z$] (2.9,2) to[L, l=$L\,\Delta z$] (4.9,2) -- (7.6,2) node[ocirc]{};
  \draw (5.4,2) to[R, l_=$G\,\Delta z$, *-*] (5.4,0);
  \draw (6.9,2) to[C, l=$C\,\Delta z$, *-*] (6.9,0);
  \draw (0,0) node[ocirc]{} -- (7.6,0) node[ocirc]{};
  \node[navy] at (-0.4,1) {$V(z)$};
  \node[navy] at (0,1.65) {\scriptsize $+$}; \node[navy] at (0,0.35) {\scriptsize $-$};
  \node[navy] at (8.3,1) {$V(z{+}\Delta z)$};
  \node[navy] at (7.6,1.65) {\scriptsize $+$}; \node[navy] at (7.6,0.35) {\scriptsize $-$};
  \draw[decorate,decoration={brace,mirror,amplitude=4pt},muted] (0,-0.25) -- node[lbl,below=4pt]{one section of length $\Delta z$; a line is many of these in cascade} (7.6,-0.25);
  \node[lbl,anchor=west] at (9.3,1.6) {$R$: copper, skin effect};
  \node[lbl,anchor=west] at (9.3,1.15) {$L$: magnetic field};
  \node[lbl,anchor=west] at (9.3,0.7) {$G$: dielectric loss};
  \node[lbl,anchor=west] at (9.3,0.25) {$C$: electric field};
\end{tikzpicture}
\caption{\textbf{One RLGC section of a transmission line.} The series elements carry the current and the shunt elements hold the voltage. On a board at 5\,Gb/s, $\omega L \gg R$ and $\omega C \gg G$: the line is an $LC$ ladder with small loss terms.}
\label{fig:ch3-rlgc}
\end{figure}

\begin{keyeq}[Propagation constant, impedance and velocity]
\begin{equation}
  \begin{gathered}
  \gamma = \alpha + j\beta = \sqrt{(R + j\omega L)(G + j\omega C)}, \qquad
  \Zo = \sqrt{\frac{R + j\omega L}{G + j\omega C}} \;\xrightarrow{\ \text{low loss}\ }\; \sqrt{\frac{L}{C}},\\[2pt]
  v = \frac{1}{\sqrt{LC}} = \frac{c}{\sqrt{\varepsilon_r}} \quad\text{(lossless)}.
  \end{gathered}
  \label{eq:ch3-z0}
\end{equation}
\end{keyeq}
$\alpha$ (Np/m) is the attenuation and $\beta$ (rad/m) the phase constant. The \textbf{characteristic impedance} $\Zo$ is the ratio $V^+/I^+$ of one traveling wave: the impedance the wave sees before it meets the far end, not something an ohmmeter can measure. Inductance takes over from resistance above $f = R/(2\pi L)$. For the REF line ($R_\mathrm{DC} = \qty{7}{\ohm\per\meter}$, $L = \qty{334}{\nano\henry\per\meter}$) that is only \qty{3.3}{\mega\hertz}, so every board trace is an LC line at gigahertz frequencies. Thin on-chip wires have far more resistance per meter and stay RC lines.

\begin{example}{FR-4 delay and the transmission-line length}\label{we:ch3-delay}
\textbf{Given:} FR-4 stripline, $\varepsilon_r = 4.0$, $\Zo = \qty{50}{\ohm}$, and a 5\,Gb/s CML edge with $t_r \approx \qty{50}{\pico\second}$ (0.25\,UI, a reasonable value for an \sgx{} HBT driver).

\textbf{Solution:} $v = c/\sqrt4 = \qty{1.50e8}{\meter\per\second}$, so the delay is \qty{6.67}{\pico\second\per\milli\meter} (169\,ps per inch); one meter of REF takes 6.67\,ns, or 33.4\,UI. The line has $L = \Zo/v = \qty{334}{\nano\henry\per\meter}$ and $C = 1/(\Zo v) = \qty{133}{\pico\farad\per\meter}$. By \cref{eq:ch3-tline-rule}, a trace is a transmission line beyond $t_r v/6 = \qty{1.25}{\milli\meter}$. The wavelength rule agrees: $\lambda/10$ at 7.5\,GHz, the third harmonic of $\fN$, is \qty{2.0}{\milli\meter}.

\textbf{Conclusion:} everything off-chip longer than about a millimeter, package traces included, is a transmission line at 5\,Gb/s. On-chip routes of a few hundred micrometers are lumped.
\end{example}

%% =====================================================================
\section{Loss mechanisms}\label{sec:loss}

A line loses energy in the copper and in the dielectric. In the low-loss limit ($R\ll\omega L$, $G\ll\omega C$), a first-order expansion of \cref{eq:ch3-z0} splits the attenuation into two independent terms, $\alpha \approx \alpha_R + \alpha_D$.

\paragraph{Skin effect: conductor loss grows as $\sqrt f$}
At DC the current fills the whole cross-section of the trace. At high frequency, the changing magnetic field inside the conductor induces eddy currents that cancel the current in the middle, and the current crowds into a surface layer of thickness
\begin{equation}
  \delta = \sqrt{\frac{\rho}{\pi f \mu_0}},
  \label{eq:ch3-skin}
\end{equation}
the \textbf{skin depth} ($\rho$ is the resistivity, $\mu_0 = 4\pi\times10^{-7}$\,H/m). For copper ($\rho = \qty{1.72e-8}{\ohm\meter}$) at 2.5\,GHz, $\delta = \qty{1.32}{\micro\meter}$. Once $\delta$ is below half the conductor thickness, the current flows in a sheet $\delta$ thick around the effective perimeter $P$, so
\begin{equation}
  R(f) \approx \frac{\rho}{P\,\delta} = \frac{\sqrt{\pi f \mu_0 \rho}}{P} \;\propto\; \sqrt f .
\end{equation}
The numerator $R_s = \rho/\delta$ is the \textbf{surface resistance} in ohms per square. \Cref{fig:ch3-skin} shows both effects for a 5-mil (\qty{127}{\micro\meter}) wide, half-ounce (\qty{17.5}{\micro\meter}) copper trace. On chip, the \sgx{} thick top metals are aluminum, with $\delta = \qty{1.6}{\micro\meter}$ at 2.5\,GHz, close to half the \qty{3}{\micro\meter} TopMetal2 thickness. That matters for inductors and long clock lines, not for the board channel.

\begin{figure}[htbp]
\centering
\includegraphics{ch03_skin}
\caption{\textbf{Skin effect makes resistance grow as $\sqrt f$.} (a) Skin depth falls as $1/\sqrt f$; below half the conductor thickness (dotted lines) the current no longer uses the whole cross-section. (b) Resistance of a 5-mil, half-ounce stripline: flat at $R_\mathrm{DC} = 7.7\,\chThreeOhm$/m up to about 57\,MHz, then $\propto\sqrt f$, reaching about 46\,\chThreeOhm/m at the 2.5\,GHz Nyquist frequency.}
\label{fig:ch3-skin}
\end{figure}

\paragraph{Dielectric loss grows as $f$}
The alternating field also rotates the dipoles in the board material, and each rotation dissipates a little heat. A lossy dielectric acts like a capacitor with a conductance in parallel, $G = \omega C \tan\delta$, where the \textbf{loss tangent} $\tan\delta$ is roughly constant over frequency (0.02 for FR-4, 0.002--0.005 for low-loss laminates). Because $G$ grows with $\omega$, dielectric loss grows linearly with frequency. (The loss angle $\delta$ is unrelated to the skin depth $\delta$.)

\begin{keyeq}[Conductor and dielectric attenuation]
\begin{equation}
  \alpha_R = \frac{R(f)}{2\Zo}, \qquad
  \alpha_D = \frac{G\Zo}{2} = \frac{\pi f\sqrt{\varepsilon_r}\,\tan\delta}{c}
  \ [\mathrm{Np/m}], \qquad
  \mathrm{IL}(f) = 8.686\,(\alpha_R+\alpha_D)\,\ell\ [\mathrm{dB}].
  \label{eq:ch3-alpha}
\end{equation}
\end{keyeq}
\Cref{fig:ch3-loss} plots both terms in dB per inch (1\,Np = 8.686\,dB). Conductor loss dominates only below about 0.5\,GHz. At the 2.5\,GHz Nyquist frequency of \OS{}, dielectric loss is already more than twice as large, and the gap widens because $f$ outruns $\sqrt f$.

\begin{figure}[htbp]
\centering
\includegraphics[width=\linewidth]{ch03_loss}
\caption{\textbf{Above about 0.5\,GHz, dielectric loss dominates FR-4.} Conductor loss (dashed) grows as $\sqrt f$, dielectric loss (dash-dot) as $f$. At the NRZ Nyquist frequency a smooth 5-mil stripline loses 0.33\,dB/in; a low-loss laminate cuts that to 0.15\,dB/in. The REF model (dotted) has a rougher, lossier conductor and loses 0.44\,dB/in.}
\label{fig:ch3-loss}
\end{figure}

\begin{example}{Loss per inch and reach for a 10\,dB budget}\label{we:ch3-loss}
\textbf{Given:} the 5-mil, half-ounce, \qty{50}{\ohm} FR-4 stripline of \cref{fig:ch3-skin} ($\varepsilon_r = 4.0$, $\tan\delta = 0.02$) at $\fN = 2.5$\,GHz.

\textbf{Dielectric:} $\alpha_D = \pi(2.5\times10^9)(2)(0.02)/(3\times10^8) = \qty{1.047}{Np/m}$: 9.10\,dB/m or \textbf{0.231\,dB/in}.

\textbf{Conductor:} $R_s = \rho/\delta = \qty{13.0}{\milli\ohm}$ per square. With $P \approx 2(w+t) = \qty{289}{\micro\meter}$, $R \approx \qty{45}{\ohm\per\meter}$, so $\alpha_R = 45/(2\cdot50) = \qty{0.45}{Np/m}$ or \textbf{0.10\,dB/in}. Copper roughness and return current in the planes add tens of percent in real boards.

\textbf{Total and reach:} $\alpha \approx 0.33$\,dB/in. A 10\,dB budget allows $10/0.33 = 30$\,in (76\,cm) of bare trace, or 21\,in (54\,cm) if 3\,dB is set aside for two packages and a connector pair. The REF model reaches 10\,dB at 0.572\,m (22.5\,in, the 10\,dB entry of \code{LOSS\_FAMILY}) because its conductor term is twice as large (0.21\,dB/in) and it includes the pads.
\end{example}

\begin{insight}[Why everyone quotes the loss at Nyquist]
NRZ puts most of its power below $\fN = \fb/2$ (\cref{sec:nrz}), and the fastest pattern, 1010\ldots, is a square wave whose fundamental is exactly $\fN$. The loss at $\fN$ therefore predicts how much the channel closes the eye, and how much equalization you need, far better than the length or the loss at $\fb$. For \OS{}, budget \emph{insertion loss at 2.5\,GHz}: REF has 17.4\,dB. Up to about 10\,dB, a small TX FFE or a CTLE is enough; from 10 to 25\,dB you need a CTLE and a DFE (\cref{ch:eqfund}).
\end{insight}

%% =====================================================================
\section{Reflections and termination}\label{sec:reflections}

A wave on a line of impedance $\Zo$ that meets a load $Z_L$ must satisfy $V = Z_L I$ at the load and $V/I = \Zo$ in the wave. One wave cannot do both, so a reflected wave is launched back toward the source, with relative amplitude given by the \textbf{reflection coefficient}.

\begin{keyeq}[Reflection coefficient and return loss]
\begin{equation}
  \Gamma_L = \frac{Z_L - \Zo}{Z_L + \Zo}, \qquad
  \Gamma_S = \frac{R_S - \Zo}{R_S + \Zo}, \qquad
  V_i = V_S\frac{\Zo}{R_S + \Zo}, \qquad
  \mathrm{RL} = -20\log_{10}\lvert\Gamma\rvert .
  \label{eq:ch3-gamma}
\end{equation}
\end{keyeq}
A matched load gives $\Gamma = 0$, an open $+1$ and a short $-1$. $V_i$ is the wave launched by a source of open-circuit voltage $V_S$ and resistance $R_S$. The \textbf{return loss} RL is $\lvert\Gamma\rvert$ in positive decibels: larger is better. With both ends mismatched, the wave bounces back and forth. The \textbf{lattice diagram} of \cref{fig:ch3-bounce} tracks each bounce, with position across and time down. Each arrival at the load adds $(1+\Gamma_L)$ times the incident wave, and the load voltage converges to the DC divider value only after several round trips of $2t_d$.

\begin{figure}[tbp]
\centering
\begin{tikzpicture}[font=\sffamily\scriptsize,x=1cm,y=0.6cm]
  \draw[navy,line width=1pt] (0,0.3) -- (0,-5.4);
  \draw[navy,line width=1pt] (4.6,0.3) -- (4.6,-5.4);
  \node[lblb,above] at (0,0.3) {source};
  \node[lblb,above] at (4.6,0.3) {load};
  \node[lbl,align=center,below] at (0,-5.4) {$R_S = 25\,\chThreeOhm$\\$\Gamma_S = -1/3$};
  \node[lbl,align=center,below] at (4.6,-5.4) {$R_L = 100\,\chThreeOhm$\\$\Gamma_L = +1/3$};
  \node[lbl,align=center] at (2.3,-6.0) {$\Zo = 50\,\chThreeOhm$\\delay $t_d$};
  \foreach \y/\l in {0/0,-1/t_d,-2/2t_d,-3/3t_d,-4/4t_d,-5/5t_d}{\draw[muted] (-0.1,\y) -- (0.1,\y); \node[lbl,left] at (-0.12,\y) {$\l$};}
  \draw[sigb] (0,0) -- node[sloped,above,text=bblue]{$V_i = +0.667$} (4.6,-1);
  \draw[sigb] (4.6,-1) -- node[sloped,above,text=bblue]{$\Gamma_L V_i = +0.222$} (0,-2);
  \draw[sigb] (0,-2) -- node[sloped,above,text=bblue]{$-0.074$} (4.6,-3);
  \draw[sigb] (4.6,-3) -- node[sloped,above,text=bblue]{$-0.025$} (0,-4);
  \draw[sigb] (0,-4) -- node[sloped,above,text=bblue]{$+0.008$} (4.6,-5);
  \foreach \y in {-1,-3,-5}{\node[dot,fill=golddk] at (4.6,\y) {};}
  \begin{scope}[shift={(6.8,-4.6)},x=0.62cm,y=4.2cm]
    \draw[->,muted] (0,0) -- (7.6,0) node[right,text=navy]{$t/t_d$};
    \draw[->,muted] (0,0) -- (0,1.05) node[above,text=navy]{$V_L$ (V)};
    \foreach \x in {1,3,5,7}{\draw[muted] (\x,0) -- (\x,-0.03) node[below,text=muted]{\x};}
    \foreach \v in {0.5,0.8}{\draw[muted] (0,\v) -- (-0.12,\v) node[left,text=muted]{\v};}
    \draw[gridc,dashed] (0,0.8) -- (7.5,0.8);
    \draw[golddk,line width=1.2pt] (0,0) -- (1,0) -- (1,0.889) -- (3,0.889) -- (3,0.790) -- (5,0.790) -- (5,0.801) -- (7.4,0.801);
    \node[lbl,anchor=south west,text=golddk] at (1.05,0.89) {0.889};
    \node[lbl,anchor=north west,text=golddk] at (3.05,0.785) {0.790};
    \node[lbl,anchor=south west] at (4.4,0.82) {DC: $\frac{R_L}{R_S+R_L} = 0.8$};
  \end{scope}
\end{tikzpicture}
\caption{\textbf{A lattice (bounce) diagram tracks every reflection.} Left: a 1\,V step from a 25\,\chThreeOhm{} source into a 50\,\chThreeOhm{} line ending in 100\,\chThreeOhm. Each arrow is one traveling wave with its amplitude. Right: the load voltage overshoots to 0.889\,V, undershoots to 0.790\,V and settles to 0.8\,V, one step every $2t_d$. In a link, the same echoes appear as bumps in the pulse response many UI after the main cursor.}
\label{fig:ch3-bounce}
\end{figure}

\paragraph{Termination choices}
\textbf{Source termination} ($R_S = \Zo$, open load) absorbs reflections at the driver. \textbf{Load termination} ($R_L = \Zo$) absorbs the wave at the receiver but only works if nothing in between reflects. \textbf{Double termination} ($R_S = R_L = \Zo$) absorbs every returning echo at either end. It halves the received swing, but echoes are worse than a smaller signal, so serial links use it. \OS{} has \qty{50}{\ohm} per side (\qty{100}{\ohm} differential) at TX and RX; \cref{sec:termination} builds these resistors in \sgx{}. A $\pm20\%$ resistor error (40--\qty{60}{\ohm}) gives $\lvert\Gamma\rvert \le 0.11$, i.e.\ RL $\ge 19$\,dB, which is acceptable. The parasitics around the resistor are the bigger problem.

\paragraph{Pad capacitance}
The bond pad, the ESD diodes and the input devices add a shunt capacitance $C_p$ across the termination. With $R_T = \Zo$, the load is $\Zo\parallel 1/(j\omega C_p)$ and \cref{eq:ch3-gamma} becomes
\begin{equation}
  \Gamma(f) = \frac{-j\omega \Zo C_p}{2 + j\omega \Zo C_p}, \qquad
  \lvert\Gamma\rvert = \frac{\omega \Zo C_p}{\sqrt{4 + (\omega \Zo C_p)^2}} ,
  \label{eq:ch3-padgamma}
\end{equation}
plotted in \cref{fig:ch3-pad}(a) for several values of $C_p$.

\begin{example}{Pad capacitance and return loss}\label{we:ch3-pad}
\textbf{Given:} $R_T = \qty{50}{\ohm}$ in parallel with $C_p = \qty{300}{\femto\farad}$ (the REF pad), at 2.5\,GHz.

\textbf{Solution:} $\omega\Zo C_p = 2\pi(2.5\times10^9)(50)(300\times10^{-15}) = 0.236$, so $\lvert\Gamma\rvert = 0.236/\sqrt{4.056} = 0.117$ and $\mathrm{RL} = \qty{18.6}{\deci\bel}$ (12.8\,dB at 5\,GHz).

\textbf{Budget:} inverting \cref{eq:ch3-padgamma}, $\omega\Zo C_p = 2\lvert\Gamma\rvert/\sqrt{1-\lvert\Gamma\rvert^2}$. RL $\ge 10$\,dB at 2.5\,GHz allows $C_p \le \qty{0.85}{\pico\farad}$; RL $\ge 15$\,dB allows $C_p \le \qty{0.46}{\pico\farad}$. Pad, ESD and input devices share this budget (\cref{sec:esd-pad}).
\end{example}

\begin{figure}[tbp]
\centering
\includegraphics[width=\linewidth]{ch03_pad}
\caption{\textbf{Pad capacitance spoils the termination at high frequency, and the reflections return as echoes.} (a) Return loss of a 50\,\chThreeOhm{} termination with shunt $C_p$; a 2.5\,nH bondwire (dashed) is as bad as 1\,pF. (b) Pulse response of a 10\,cm REF-type line. With 1\,pF pads, the wave reflected at the RX pad returns from the TX pad and arrives $2t_d = 6.7$\,UI after the main cursor. A short, low-loss line does not attenuate this echo.}
\label{fig:ch3-pad}
\end{figure}

\begin{example}{Bondwire inductance}\label{we:ch3-bond}
\textbf{Given:} a 2.5\,mm bondwire, $L \approx \qty{2.5}{\nano\henry}$ (about 1\,nH/mm), in series between two \qty{50}{\ohm} environments, at 2.5\,GHz.

\textbf{Solution:} a series impedance $Z = j\omega L$ between matched lines gives $S_{11} = (Z/2\Zo)/(1 + Z/2\Zo)$ and $S_{21} = 1/(1 + Z/2\Zo)$. Here $\omega L = \qty{39.3}{\ohm}$ and $\omega L/2\Zo = 0.393$. So $\lvert S_{21}\rvert = 1/\sqrt{1.154} = 0.931$ ($-0.62$\,dB) but $\lvert S_{11}\rvert = 0.393/1.075 = 0.366$: \textbf{RL = 8.7\,dB}, failing a 10\,dB target on its own.

\textbf{Fix:} keep bondwires below about 1\,mm, use flip-chip bumps, or absorb the inductance into a matched section: 2.5\,nH with 1.0\,pF of pad capacitance forms an artificial line with $\sqrt{L/C} = \qty{50}{\ohm}$ and a cutoff of $1/(\pi\sqrt{LC}) = 6.4$\,GHz.
\end{example}

\paragraph{Via stubs}
A through-hole via connects traces on different layers. If the signal enters on the top layer and leaves on an inner layer, the rest of the barrel is an unused \textbf{stub}: a short open-ended line hanging off the signal path (\cref{fig:ch3-via}a). Where the stub is a quarter wavelength long, its open end is transformed into a short at the signal path, and almost nothing gets through:
\begin{equation}
  f_\mathrm{null} \approx \frac{v}{4\,\ell_\mathrm{stub}} .
  \label{eq:ch3-stub}
\end{equation}

\begin{figure}[tbp]
\centering
\begin{minipage}[b]{0.43\linewidth}\centering
\begin{tikzpicture}[font=\sffamily\scriptsize,x=1cm,y=1cm]
  \fill[gtint] (-2.0,0) rectangle (2.0,-3.6);
  \foreach \y in {0,-0.9,-1.8,-2.7,-3.6}{\draw[muted,line width=0.4pt] (-2.0,\y) -- (2.0,\y);}
  \foreach \y/\n in {0/L1,-0.9/L2,-1.8/L3,-2.7/L4,-3.6/L5}{\node[lbl,left] at (-2.0,\y) {\n};}
  \fill[c1] (-0.18,0) rectangle (0.18,-1.8);
  \fill[c4] (-0.18,-1.8) rectangle (0.18,-3.6);
  \draw[navy,line width=2pt] (-1.9,0.04) -- (-0.18,0.04);
  \draw[navy,line width=2pt] (0.18,-1.8) -- (1.9,-1.8);
  \draw[sig] (-1.6,0.35) -- (-0.6,0.35) node[right,text=navy]{signal in (L1)};
  \draw[sig] (0.6,-1.45) -- (1.7,-1.45);
  \node[lbl,text=navy,anchor=south] at (1.25,-1.4) {out (L3)};
  \draw[decorate,decoration={brace,amplitude=3pt},c4] (0.3,-1.85) -- node[right=3pt,text=c4,align=left]{stub\\$\ell_\mathrm{stub}$} (0.3,-3.55);
  \node[lbl,align=center,fill=gtint,inner sep=1pt] at (-1.05,-3.15) {open end\\reflects};
  \node[lbl,align=center] at (0,-4.05) {backdrilling removes the red part};
\end{tikzpicture}
\subcaption{Via with an unused stub}
\end{minipage}\hfill
\begin{minipage}[b]{0.54\linewidth}\centering
\includegraphics{ch03_via}
\subcaption{Notch at $v/(4\ell_\mathrm{stub})$ in FR-4}
\end{minipage}
\caption{\textbf{A via stub is a quarter-wave resonator that notches $\lvert S_{21}\rvert$.} (a) The signal turns at L3, so the barrel below is an open stub. (b) Through-path $\lvert S_{21}\rvert$ with 2, 4 and 6\,mm stubs, labeled with their null frequencies. A 6\,mm stub puts the notch at 6.2\,GHz, close enough to the 0--5\,GHz NRZ main lobe (shaded) to distort the pulse.}
\label{fig:ch3-via}
\end{figure}

\begin{example}{Via-stub notch}\label{we:ch3-via}
\textbf{Given:} FR-4 ($v = \qty{1.5e8}{\meter\per\second}$), stubs of 1.5, 3 and 5\,mm.

\textbf{Solution:} \cref{eq:ch3-stub} gives $f_\mathrm{null} = 25$, 12.5 and 7.5\,GHz. To keep the first null above $3\fN = 7.5$\,GHz, stubs must be shorter than $v/(4\times7.5\,\mathrm{GHz}) = \qty{5}{\milli\meter}$. The barrel's extra capacitance slows the wave, so in practice keep stubs below about 3\,mm, route on layers near the bottom, or backdrill. The \OS{} test board should follow this rule for every high-speed via (\cref{sec:layout-hs}).
\end{example}

\begin{insight}
Every discontinuity sends back a small, filtered copy of the signal. Two discontinuities form a cavity, and the copy that bounces between them reaches the receiver $2t_d$ late, about 1.7\,UI per inch of FR-4 between them. In the pulse response this shows up as \emph{ripples and bumps in the tail}, often 5--20\,UI after the main cursor (\cref{fig:ch3-pad}b). A two-tap DFE cannot cancel ISI that far out, so control reflections at the source: accurate termination, small pad capacitance, short bondwires and no long via stubs. A short, low-loss channel does not attenuate its echoes, so it can be harder to equalize than a longer one.
\end{insight}

%% =====================================================================
\section{S-parameters}\label{sec:sparams}

At gigahertz frequencies we cannot measure with opens and shorts, because a short wire is already a transmission line. A vector network analyzer instead drives each port from \qty{50}{\ohm}, terminates the others in \qty{50}{\ohm}, and measures traveling waves: $a_j$ into port $j$ and $b_i$ out of port $i$. The \textbf{scattering parameters} are $S_{ij} = b_i/a_j$ with all other inputs matched. For a two-port, $S_{11}$ is the input reflection (return loss $-20\log_{10}\lvert S_{11}\rvert$) and $S_{21}$ the transmission; the \textbf{insertion loss} is $\mathrm{IL} = -20\log_{10}\lvert S_{21}\rvert$, a positive number. Because a doubly terminated link has the same \qty{50}{\ohm} reference, $S_{21}$ is the transfer function $\Hch(f) = 2V_\mathrm{out}/V_S$ that we need. To cascade elements, convert each to an ABCD matrix, multiply, and convert back; \code{link.py} builds REF this way.

\paragraph{Differential channels and mixed-mode S-parameters}
\OS{} is differential: the information is $V_P - V_N$. A differential channel has four single-ended ports, numbered as in \cref{fig:ch3-ports}(a): ports 1 and 3 are the P and N inputs, 2 and 4 the P and N outputs. Driving the input differentially ($a_1 = -a_3$) and reading the differential output $(b_2 - b_4)/\sqrt2$ gives the \textbf{mixed-mode} parameters, referenced to \qty{100}{\ohm} differential:
\begin{equation}
  S_{dd21} = \tfrac12\left(S_{21} - S_{23} - S_{41} + S_{43}\right), \qquad
  S_{dd11} = \tfrac12\left(S_{11} - S_{13} - S_{31} + S_{33}\right).
  \label{eq:ch3-sdd}
\end{equation}
$S_{cd21} = \tfrac12(S_{21} - S_{23} + S_{41} - S_{43})$ measures how much differential signal turns into common mode, for example through a P/N length mismatch. For a symmetric, weakly coupled pair, $S_{dd21}\approx S_{21}$ of one leg, which is why REF can be modeled as a single-ended \qty{50}{\ohm} line.

\begin{figure}[tbp]
\centering
\begin{minipage}[b]{0.49\linewidth}\centering
\begin{tikzpicture}[font=\sffamily\scriptsize]
  \node[blkb,minimum width=30mm,minimum height=20mm] (b) {4-port\\channel};
  \foreach \y/\p/\s in {0.55/1/P,-0.55/3/N}{
    \draw[navy,line width=0.8pt] ($(b.west)+(0,\y)$) -- ++(-0.9,0) node[ocirc]{};
    \node[lblb,above] at ($(b.west)+(-0.45,\y)$) {\p};
    \node[lbl,left] at ($(b.west)+(-1.0,\y)$) {in \s};}
  \foreach \y/\p/\s in {0.55/2/P,-0.55/4/N}{
    \draw[navy,line width=0.8pt] ($(b.east)+(0,\y)$) -- ++(0.9,0) node[ocirc]{};
    \node[lblb,above] at ($(b.east)+(0.45,\y)$) {\p};
    \node[lbl,right] at ($(b.east)+(1.0,\y)$) {out \s};}
  \draw[sigb,dashed] ($(b.west)+(0.15,0.55)$) -- ($(b.east)+(-0.15,0.55)$);
  \draw[sigb,dashed] ($(b.west)+(0.15,-0.55)$) -- ($(b.east)+(-0.15,-0.55)$);
  \node[lbl,align=center,below=1mm of b] {$S_{dd21} = \tfrac12(S_{21}-S_{23}-S_{41}+S_{43})$\\through paths $1\to2$ and $3\to4$};
\end{tikzpicture}
\subcaption{Port numbering used in this book}
\end{minipage}\hfill
\begin{minipage}[b]{0.49\linewidth}\centering
\begin{tikzpicture}[font=\sffamily\scriptsize]
  \draw[c4,line width=1.4pt] (0,1) -- (4.4,1);
  \draw[navy,line width=1.4pt] (0,0) -- (4.4,0);
  \node[lbl,left,text=c4,align=right] at (-0.1,1) {aggressor\\(neighbor TX)};
  \node[lbl,left,text=navy,align=right] at (-0.1,0) {victim\\(our RX lane)};
  \draw[sigg] (0.9,1.3) -- (2.1,1.3) node[right,text=golddk]{signal travels};
  \foreach \x in {1.0,1.8,2.6,3.4}{\draw[c5,densely dotted,line width=0.7pt] (\x,0.92) -- (\x,0.08);}
  \draw[c5,-{Stealth[length=2mm]},line width=0.8pt] (0.6,-0.12) -- (0.0,-0.55) node[below,text=c5,align=center]{NEXT\\near end};
  \draw[c5,-{Stealth[length=2mm]},line width=0.8pt] (3.8,-0.12) -- (4.4,-0.55) node[below,text=c5,align=center]{FEXT\\far end};
  \node[lbl,align=center] at (2.2,-0.75) {coupling $C_m$, $L_m$};
\end{tikzpicture}
\subcaption{Near- and far-end crosstalk}
\end{minipage}
\caption{\textbf{Differential channels are 4-ports, and neighboring lanes couple.} (a) Ports 1 and 3 are the P and N inputs, 2 and 4 the outputs. Check the port map of every Touchstone file: many tools number the pair (1,2) in and (3,4) out. (b) An aggressor couples into the victim at the near end (NEXT) and the far end (FEXT).}
\label{fig:ch3-ports}
\end{figure}

\paragraph{Crosstalk}
Mutual capacitance $C_m$ and inductance $L_m$ couple an \textbf{aggressor} line into a \textbf{victim} (\cref{fig:ch3-ports}b). \textbf{Near-end crosstalk} (NEXT) travels back toward the aggressor's source and is the dangerous kind in a transceiver: a full-swing local TX lane couples into a local RX lane whose own signal has already lost 17\,dB. A 0.8\,Vppd aggressor at $-30$\,dB NEXT injects about 25\,mVpp, comparable with the residual ISI after equalization. \textbf{Far-end crosstalk} (FEXT) travels with the victim's signal, is attenuated with it, and nearly cancels in a homogeneous stripline. Route TX and RX on different layers or connector rows, and prefer stripline.

\paragraph{Time-domain reflectometry}
A \textbf{TDR} launches a fast step and records the reflected voltage at the source. A reflection at round-trip time $t$ comes from distance $x = vt/2$, and its amplitude gives the local impedance through $Z = \Zo(1+\rho)/(1-\rho)$, with $\rho = V_r/V_i$. \Cref{fig:ch3-tdr} shows the two signature shapes: a series inductance bumps above \qty{50}{\ohm} and a shunt capacitance dips below. For a lumped element, the size of the feature depends on the rise time: with a 25\,ps step a 1\,nH pin reads \qty{97}{\ohm}, with a 100\,ps step only \qty{61}{\ohm}. A TDR result always comes with its rise time, which sets its resolution.

\begin{figure}[tbp]
\centering
\includegraphics[width=\linewidth]{ch03_tdr}
\caption{\textbf{TDR: inductance bumps up, capacitance dips down.} Simulated TDR of a 50\,\chThreeOhm{} line with a 1\,nH series inductance at 15\,mm and a 0.5\,pF shunt capacitance at 68\,mm, ending in a matched load. The top axis converts round-trip time to distance. A slower step (amber) blurs and shrinks both features. The wiggles near 1.4\,ns are the wave re-reflected between the two discontinuities.}
\label{fig:ch3-tdr}
\end{figure}

%% =====================================================================
\section{From frequency to time}\label{sec:ch3-f2t}

The tools of \cref{ch:linkanalysis} need the pulse response $p(t)$ of \cref{sec:pulse-response}. We get it from $S_{21}(f)$ in two steps.

\paragraph{Causality}
A physical channel cannot respond before it is driven. In the frequency domain this means the magnitude and phase of $\Hch(f)$ are linked: a given loss curve comes with a specific phase (dispersion). A model built from $\lvert S_{21}\rvert$ alone with a pure delay gives a symmetric impulse response that starts long before the wave could arrive (\cref{fig:ch3-causal}). A dielectric model with $\varepsilon_r$ and $\tan\delta$ held exactly constant has the same flaw in a milder form. REF uses a causal dielectric whose capacitance varies as $(j\omega/\omega_0)^{-2\tan\delta/\pi}$, the simplest of the causal FR-4 models~\citep{p2:djordjevic2001}. A non-causal model produces \emph{fake pre-cursors}, which then mislead the design of the TX FFE pre-tap.

\paragraph{Impulse and pulse response}
The impulse response is the inverse Fourier transform of the conjugate-symmetric spectrum, $h(t) = \mathcal F^{-1}\{\Hch(f)\}$ with $\Hch(-f) = \Hch^*(f)$. The pulse response is $h(t)$ convolved with a 1-UI rectangle.

\begin{keyeq}[From $S_{21}$ to the pulse response]
\begin{equation}
  p(t) = \int_{t-\Tb}^{t} h(\tau)\,d\tau
  \;\;\Longleftrightarrow\;\;
  P(f) = \Hch(f)\,\Tb\,\mathrm{sinc}(f\Tb)\,e^{-j\pi f\Tb}
  \label{eq:ch3-pulse}
\end{equation}
\end{keyeq}
On a grid of $N$ points spaced $\Delta f$ up to $f_\mathrm{max}$, the IFFT gives a time step $\Delta t = 1/(2f_\mathrm{max})$ and a window $T_w = 1/\Delta f$. \code{common.py} uses 0--160\,GHz in 10\,MHz steps: $\Delta t = \qty{3.125}{\pico\second}$ (64 samples per UI) and $T_w = \qty{100}{\nano\second}$, far longer than the REF delay plus tail. \Cref{fig:ch3-imp-pulse} compares the two responses: the 1-UI rectangle averages $h$ over 200\,ps, so $p(t)$ is smoother and lower, and its tail coincides with that of $h(t)\cdot\Tb$.

\begin{figure}[htbp]
\centering
\begin{minipage}[t]{0.49\linewidth}\centering
\includegraphics{ch03_causal}
\caption{\textbf{A magnitude-only model is not causal.} Keeping $\lvert S_{21}\rvert$ of REF but replacing its phase by a pure delay gives a symmetric response with a long precursor before the first arrival.}
\label{fig:ch3-causal}
\end{minipage}\hfill
\begin{minipage}[t]{0.49\linewidth}\centering
\includegraphics{ch03_imp_pulse}
\caption{\textbf{The pulse response is the impulse response averaged over one UI.} REF, with $h(t)$ scaled by $\Tb$ so both curves are in volts. The peak $p_0 = 0.340$\,V arrives 6.82\,ns after launch.}
\label{fig:ch3-imp-pulse}
\end{minipage}
\end{figure}

\Cref{fig:ch3-family} shows how the pulse response changes with loss, using the \code{LOSS\_FAMILY} lines of \code{link.py}. From 5 to 25\,dB of Nyquist loss, the main cursor falls from 0.77 to 0.22\,V while the first post-cursor stays near 0.16--0.18\,V, so $p_1/p_0$ grows from 0.12 to 0.71. That ratio, more than the loss itself, is what the equalizer must fight.

\begin{figure}[htbp]
\centering
\includegraphics[width=\linewidth]{ch03_family}
\caption{\textbf{More loss means a smaller main cursor and a relatively larger tail.} Pulse responses of the \code{LOSS\_FAMILY} channels, aligned at their peaks; open dots mark the post-cursors $p_1\ldots p_{10}$. At 15\,dB (close to REF) the first post-cursor is already 46\% of the main cursor.}
\label{fig:ch3-family}
\end{figure}

\begin{pitfall}[IFFT mistakes that change your cursors]
Forgetting conjugate symmetry gives a complex $h(t)$. A coarse $\Delta f$ makes the window shorter than the delay plus tail, so the tail wraps around and looks like pre-cursor ISI. Data cut abruptly at its last frequency rings (Gibbs), and data without a DC point must be extrapolated. And the cursors are samples of $p(t)$, not of $h(t)$.
\end{pitfall}

%% =====================================================================
\section{The Open-SERDES reference channel}\label{sec:refchannel}

Every model, figure and lab in this book uses one synthetic channel, \textbf{REF}: a \qty{50}{\ohm} source, a 300\,fF pad, 1\,m of lossy FR-4 stripline, another 300\,fF pad and a \qty{50}{\ohm} load. It is built from serdespy's ABCD functions~\citep{serdespy} by \code{channel(1.0)} in \code{tutorial/figsrc/link.py} and \code{tutorial/code/common.py}. \Cref{tab:ch3-refparams} lists the parameters, and \cref{fig:ch3-ref-s21} shows the insertion loss: \qty{-17.4}{\deci\bel} at 2.5\,GHz. The pads cost only 0.1\,dB of that; their real damage is the reflections they add.

\begin{figure}[tbp]
\centering
\begin{minipage}[b]{0.47\linewidth}\centering
\captionof{table}{\textbf{REF model parameters} and the matching \code{channel()} arguments.}
\label{tab:ch3-refparams}
\footnotesize\sffamily
\begin{tabular}{@{}lll@{}}
\toprule
Quantity & Value & Arg.\\
\midrule
Length & 1.0\,m & \code{length}\\
$\Zo$ per side & 50\,\chThreeOhm & \code{z0}\\
$\varepsilon_r$ & 4.0 (6.67\,ns/m) & \code{er}\\
$\tan\delta$ & 0.02 at 1\,GHz, causal & \code{tand}\\
$R_\mathrm{DC}$ & 7\,\chThreeOhm/m & \code{rdc}\\
Skin $R$ & 60\,\chThreeOhm/m at 1\,GHz, $\propto\!\sqrt f$ & \code{kr}\\
Pads & 300\,fF each end & \code{c\_pad}\\
Source, load & 50\,\chThreeOhm{} each & \\
Grid & 0--160\,GHz, 10\,MHz & \code{F}\\
\midrule
Loss & 0.59 / 10.4 / 17.4 / 30.1\,dB & \\
 & at 0 / 1.25 / 2.5 / 5\,GHz & \\
\bottomrule
\end{tabular}
\end{minipage}\hfill
\begin{minipage}[b]{0.50\linewidth}\centering
\includegraphics{ch03_ref_s21}
\caption{\textbf{REF loses 17.4\,dB at Nyquist.} $\lvert S_{21}\rvert$ with and without the 300\,fF pads, which barely change the loss.}
\label{fig:ch3-ref-s21}
\end{minipage}
\end{figure}

\Cref{fig:ref-pulse} shows the REF pulse response sampled every UI from its peak, and \cref{tab:ref-cursors} lists the cursors. These thirteen numbers are the most-used data in this book.

\begin{figure}[tbp]
\centering
\includegraphics[width=\linewidth]{ch03_ref_pulse}
\caption{\textbf{The REF pulse response and its cursors.} A 1\,V, 1-UI pulse arrives as a 0.340\,V peak with a long tail. Sampling every UI from the peak gives the pre-cursor $p_{-1}$ (red), the main cursor $p_0$ (gold) and the post-cursors $p_1, p_2, \ldots$ (blue). The first post-cursor alone is 53.5\% of the main cursor.}
\label{fig:ref-pulse}
\end{figure}

\begin{table}[tbp]
\centering
\caption{\textbf{REF cursors} in volts for a 1\,V, 1-UI pulse, sampled at the pulse peak ($\ts = 6.82$\,ns), from \code{link.cursors}.}
\label{tab:ref-cursors}
\small\sffamily
\setlength{\tabcolsep}{3.6pt}
\begin{tabular}{@{}l*{13}{c}@{}}
\toprule
$k$ & $-2$ & $-1$ & $\mathbf{0}$ & 1 & 2 & 3 & 4 & 5 & 6 & 7 & 8 & 9 & 10\\
\midrule
$p_k$ & 0.000 & 0.0486 & \textbf{0.3396} & 0.1817 & 0.0909 & 0.0537 & 0.0356 & 0.0255 & 0.0192 & 0.0150 & 0.0121 & 0.0100 & 0.0084\\
\bottomrule
\end{tabular}
\end{table}

\paragraph{Why this channel is hard}
Three numbers tell the story. First, $p_0 = 0.34$: only a third of the transmitted pulse arrives in its own UI. Second, $p_1/p_0 = 0.535$: the previous bit contributes more than half as much as the current one. Third, the ISI magnitudes $\sum_{k\ne0}\lvert p_k\rvert$ of these cursors add up to 0.501\,V, more than $p_0$. A worst-case pattern can push a ``1'' below zero, and the peak-distortion eye of \cref{sec:pda} is $1\,\mathrm{V}\times(0.340 - 0.501) = -0.161$\,V: \emph{closed}. A PRBS7 simulation agrees (\cref{sec:lab2}). REF cannot be received without equalization, and the \OS{} link must still reach BER $\le 10^{-12}$ over it, with margin up to the 20--25\,dB end of \code{LOSS\_FAMILY} (\cref{sec:eq-partition}, \cref{tab:specs}).

\paragraph{Replacing REF with real S-parameters}
REF is a teaching model. Once the \OS{} test board is measured, or its traces are simulated with \tool{openEMS} or \tool{Palace}, you will swap in a 4-port Touchstone file. \tool{scikit-rf}~\citep{p2:skrf2022} reads it, converts it to mixed mode and extrapolates it to DC; the rest of the flow is unchanged:
\begin{lstlisting}
import skrf as rf
from common import *                       # grid F, impulse(), pulse(), SPS
ch = rf.Network("board.s4p")               # measured or EM-simulated 4-port
ch.renumber([1, 2], [2, 1])                # book map (1,3 in; 2,4 out) -> skrf (0,1 in)
ch.se2gmm(p=2)                             # mixed mode: ports d-in, d-out, c-in, c-out
ch = ch.extrapolate_to_dc(kind="linear")   # add 0 Hz, uniform frequency step
sdd21 = ch.s[:, 1, 0]
H = np.interp(F, ch.f, sdd21.real, right=0) + 1j * np.interp(F, ch.f, sdd21.imag, right=0)
h, t = impulse(H); p = pulse(h)            # then cursors(), pda_eye(), ... as for REF
\end{lstlisting}
Check the port map first (\cref{fig:ch3-ports}), and taper the data toward its highest frequency instead of cutting it to zero. serdespy's \code{four\_port\_to\_diff} does the same job but assumes through paths $1\to2$ and $3\to4$ and ignores its \code{port\_def} argument.

\begin{tryit}[Lab 1: the channel]
Run \code{lab1\_channel.py} (\cref{sec:lab1}). It plots $\lvert S_{21}\rvert$ of REF, converts it to $h(t)$ and $p(t)$, and prints the cursors of \cref{tab:ref-cursors}. Then set \code{c\_pad} to 1\,pF and \code{length} to 0.3\,m and watch the reflections of \cref{fig:ch3-pad}(b) appear in the tail.
\end{tryit}

%% =====================================================================
\section{Exercises}\label{sec:ch3-exercises}
\begin{enumerate}[label=\textbf{3.\arabic*},leftmargin=2.2em]
  \item \textbf{Is it a line?} A package trace is 8\,mm long on a substrate with $\varepsilon_r = 3.5$, driven with $t_r = \qty{40}{\pico\second}$. Compute its delay. Must it be modeled as a transmission line?
  \item \textbf{Better laminate.} Repeat \cref{we:ch3-loss} for a low-loss laminate ($\varepsilon_r = 3.7$, $\tan\delta = 0.004$) with the same copper. How long can the trace be for 17.4\,dB at 2.5\,GHz? Which loss term now dominates?
  \item \textbf{ESD budget.} The RX must have RL $\ge 15$\,dB at 2.5\,GHz. The pad is 120\,fF and the input devices add 60\,fF. How much capacitance is left for ESD protection?
  \item \textbf{Via stub.} A 3.2\,mm thick board routes a signal from L1 to an inner layer 0.8\,mm deep, and the via's capacitance gives $\varepsilon_{r,\mathrm{eff}} = 6$. Where is the stub null? Is it acceptable for \OS{}? What if the signal exits on the bottom layer instead?
  \item \textbf{Cursor ratio.} With \code{channel()} from \code{link.py}, find the length at which $p_1/p_0 = 0.5$ and its Nyquist loss. Repeat with \code{tand=0.01} and explain why the answer moves.
\end{enumerate}

\begin{takeaways}
\begin{itemize}
  \item Off-chip, every trace longer than about 1\,mm is a transmission line at 5\,Gb/s. FR-4 delay is 6.67\,ps/mm (169\,ps/in); REF takes 6.67\,ns, or 33\,UI.
  \item Conductor loss grows as $\sqrt f$ (skin depth 1.32\,\textmu m at 2.5\,GHz); dielectric loss grows as $f\tan\delta$ and dominates FR-4 above 0.5\,GHz. A smooth FR-4 stripline loses about 0.33\,dB/in at 2.5\,GHz.
  \item Quote and budget insertion loss at Nyquist, 2.5\,GHz. REF loses 17.4\,dB there.
  \item Discontinuities reflect: 300\,fF of pad gives RL = 18.6\,dB, a 2.5\,nH bondwire only 8.7\,dB, and a via stub notches $S_{21}$ at $v/(4\ell)$. Echoes return as late bumps in the pulse response.
  \item Use mixed-mode $S_{dd21}$ for differential channels and check the port map. On a TDR, inductance bumps up and capacitance dips down.
  \item A causal model and a correct IFFT turn $S_{21}$ into $h(t)$ and $p(t)$. REF has $p_{-1} = 0.049$, $p_0 = 0.340$ and $p_1 = 0.182$ ($p_1/p_0 = 0.535$), and its unequalized peak-distortion eye is closed ($-0.161$\,V).
\end{itemize}
\end{takeaways}

\tikzsetfigurename{04_linkanalysis-}
\chapter{Link Budgets and Statistical Analysis}\label{ch:linkanalysis}
\chapterlead{The pulse response is enough to predict the eye: its worst case, its probability at $10^{-12}$, and how much margin is left after noise, jitter and the sampler take their share.}
% Omega for sans-serif contexts (captions, TikZ): the sans text font has no Omega glyph.
\providecommand{\chFourOhm}{\ensuremath{\Omega}}

\begin{objectives}
\begin{itemize}
  \item Explain why superposition lets us predict every received sample from the cursors $p_k$.
  \item Derive the peak-distortion (worst-case) eye and the worst-case bit pattern, and apply them to REF.
  \item Build a statistical eye by convolving ISI, noise and jitter distributions, and read eye height at $10^{-12}$.
  \item Write a top-down margin budget for the \OS{} receiver and decide what to spend margin on.
  \item Decide between NRZ and PAM4 with the 9.54\,dB rule.
\end{itemize}
\end{objectives}

%% =====================================================================
\section{Pulse-response link analysis}\label{sec:ch4-pulse}

The channel, the TX FFE and the CTLE are linear and time-invariant, so the received waveform is a sum of shifted, scaled copies of one pulse response (\cref{sec:lti,sec:pulse-response}): $y(t) = A\sum_n a_n\,p(t-n\Tb)$. Sampling at $\ts + n\Tb$ and writing $p_k = p(\ts + k\Tb)$ gives
\begin{equation}
  y_n = \underbrace{A\,a_n\,p_0}_{\text{signal}} \;+\; \underbrace{A\sum_{k\neq0} a_{n-k}\,p_k}_{\text{ISI}} \;+\; \text{noise}.
  \label{eq:ch4-yn}
\end{equation}
At the sampler, the channel is a finite-impulse-response (FIR) filter whose taps are the cursors (\cref{fig:ch4-fir}). Positive $k$ are \textbf{post-cursors}, which come from \emph{earlier} bits; negative $k$ are \textbf{pre-cursors}, from later bits. Every question in this chapter (how bad the eye can get, how often, and how much margin is left) is answered from the list of $p_k$, without simulating long bit streams.

\begin{figure}[tbp]
\centering
\begin{tikzpicture}[font=\sffamily\small,scale=0.95,transform shape,
  dl/.style={blk,minimum width=10mm,minimum height=7mm,inner sep=1pt},
  mul/.style={circle,draw=navy,line width=0.7pt,minimum size=8.5mm,inner sep=0pt,font=\sffamily\scriptsize}]
  \node[navy] (in) at (-0.3,0) {$A\,a_{n+1}$};
  \foreach \i/\x in {1/2.35,2/4.75,3/7.15}{\node[dl] (d\i) at (\x,0) {$z^{-1}$};}
  \node (dots) at (9.15,0) {$\cdots$};
  \foreach \j/\x/\l in {0/1.15/a_{n+1},1/3.55/a_n,2/5.95/a_{n-1},3/8.35/a_{n-2},4/10.2/a_{n-K}}{
    \coordinate (t\j) at (\x,0); \node[dot] at (t\j) {}; \node[lbl,above=1pt] at (\x,0.12) {$\l$};}
  \draw[sig] (in) -- (d1); \draw[sig] (d1) -- (d2); \draw[sig] (d2) -- (d3); \draw[navy,line width=0.75pt] (d3) -- (dots); \draw[navy,line width=0.75pt] (dots) -- (t4);
  \node[mul,fill=c4!15,draw=c4] (m0) at (1.15,-1.15) {$p_{-1}$};
  \node[mul,fill=gtint,draw=golddk,line width=1pt] (m1) at (3.55,-1.15) {$p_0$};
  \node[mul,fill=btint] (m2) at (5.95,-1.15) {$p_1$};
  \node[mul,fill=btint] (m3) at (8.35,-1.15) {$p_2$};
  \node[mul,fill=btint] (m4) at (10.2,-1.15) {$p_K$};
  \foreach \j in {0,...,4}{\draw[sig] (t\j) -- (m\j); \draw[navy,line width=0.75pt] (m\j) -- ++(0,-0.85);}
  \draw[navy,line width=0.75pt] (1.15,-2.0) -- (11.3,-2.0);
  \node[sumnode] (s) at (11.6,-2.0) {$+$};
  \draw[sig] (11.3,-2.0) -- (s);
  \node[lbl,align=center] (nz) at (11.6,-0.85) {noise};
  \draw[sig] (nz) -- (s);
  \node[blkg,minimum width=14mm,minimum height=8mm] (sl) at (13.6,-2.0) {sampler};
  \draw[sig] (s) -- node[above,font=\sffamily\scriptsize,text=navy]{$y_n$} (sl);
  \draw[sig] (sl) -- ++(1.25,0) node[right,text=navy]{$\hat a_n$};
  \node[lbl,text=c4] at (1.15,-2.35) {pre-cursor ISI};
  \node[lbl,text=golddk] at (3.55,-2.35) {signal};
  \node[lbl,text=bblue] at (7.3,-2.35) {post-cursor ISI from earlier bits};
\end{tikzpicture}
\caption{\textbf{Sampled once per UI, the channel is an FIR filter whose taps are the pulse-response cursors.} Each received sample is the current symbol times $p_0$ plus every other symbol times its cursor. REF has one significant pre-cursor ($p_{-1} = 0.049$) and a long post-cursor tail.}
\label{fig:ch4-fir}
\end{figure}

%% =====================================================================
\section{Peak distortion analysis}\label{sec:pda}

\paragraph{The worst case}
Take a transmitted ``1'' ($a_n = +1$). By \cref{eq:ch4-yn}, its sample is smallest when every ISI term pulls down, $a_{n-k}\,p_k = -\lvert p_k\rvert$ for all $k \ne 0$. The lowest ``1'' is then $A(p_0 - \sum_{k\neq0}\lvert p_k\rvert)$, and by symmetry the highest ``0'' is its negative. The gap between them is the worst-case inner eye height.

\begin{keyeq}[Peak-distortion eye and worst-case pattern]
\begin{equation}
  \mathrm{EH}(\ts) = 2A\Big(p_0 - \sum_{k\neq0}\lvert p_k\rvert\Big),
  \qquad
  a_{n-k} = -\sgn(p_k)\,a_n .
  \label{eq:pda}
\end{equation}
\end{keyeq}
This is \textbf{peak distortion analysis} (PDA)~\citep{p2:casper2002,palermo720ch}. It needs only the cursors, it bounds the eye for \emph{any} bit stream, and the second half of \cref{eq:pda} gives the \textbf{worst-case bit pattern}. Because all cursors depend on the sampling instant, $\mathrm{EH}$ is a function of $\ts$. Sweeping $\ts$ across the UI and keeping the range where $\mathrm{EH}(\ts) > 0$ gives the worst-case \textbf{eye width}.

\begin{example}{PDA of the REF channel}\label{we:ch4-pda}
\textbf{Given:} the REF cursors of \cref{tab:ref-cursors}, $2A = 1$\,V, $\ts$ at the pulse peak.

\textbf{Thirteen cursors:} $\sum_{k\neq0}\lvert p_k\rvert = 0.0486 + 0.1817 + 0.0909 + \cdots + 0.0084 = 0.501$\,V, so $\mathrm{EH} = 1\,\mathrm{V}\times(0.340 - 0.501) = \mathbf{-0.161\,V}$: closed. \Cref{fig:ch4-pda}(a) shows that the ISI already exceeds $p_0$ after three post-cursors.

\textbf{Full tail:} the cursors beyond $k = 10$ are small but many, and they add another 0.093\,V. Summing out to 300\,UI gives $\mathrm{EH} = -0.255$\,V. A check: all REF cursors are positive, so $\sum_k p_k = \Hch(0) = 0.935$ (the DC gain) and $\mathrm{EH} = 2A(2p_0 - \Hch(0)) = -0.255$\,V.

\textbf{Worst-case pattern:} every $p_k$ with $k \neq 0$ is positive, so for $a_n = +1$ all other bits are $-1$: a lone ``1'' after a long run of ``0''s, followed by a ``0''. \Cref{fig:ch4-pda}(b) simulates it: the ``1'' arrives at $-126$\,mV, below the threshold.
\end{example}

\begin{figure}[tbp]
\centering
\includegraphics[width=\linewidth]{ch04_pda}
\caption{\textbf{Peak distortion on REF: the ISI outweighs the main cursor.} (a) Cursor magnitudes and the cumulative ISI (navy). It passes $p_0$ at $k = 3$ and reaches 0.501\,V by $k = 10$; the full tail adds 0.093\,V more. (b) The worst-case pattern, a lone ``1'' after a long run of ``0''s. The sample at the main cursor (red dot) is $-126$\,mV, so this ``1'' is read as a ``0''.}
\label{fig:ch4-pda}
\end{figure}

\paragraph{PDA versus a PRBS simulation}
PDA is a bound, not a probability. The worst-case pattern for REF needs dozens of identical bits before the ``1''; PRBS7 never has more than six ``0''s in a row, so a PRBS7 simulation cannot find it. \Cref{fig:ch4-pda-prbs} compares eye height versus sampling phase from PDA, from PRBS simulations and from the statistical eye of \cref{sec:stateye}. On REF alone, PRBS7 reports $-78$\,mV (\cref{sec:lab2}), PRBS13 $-181$\,mV and full-tail PDA $-252$\,mV: closed in every view. After a 2-tap TX de-emphasis, the ranking is the same: PRBS7 163\,mV, PRBS13 148\,mV, PDA 126\,mV. A short PRBS is optimistic because it never visits the rare patterns that set the eye at $10^{-12}$.

\begin{figure}[tbp]
\centering
\includegraphics[width=\linewidth]{ch04_pda_prbs}
\caption{\textbf{PDA is the most pessimistic estimate; short PRBS runs are the most optimistic.} Inner eye height versus sampling phase. (a) REF is closed by every measure; truncating PDA at 13 cursors (dotted) hides 90\,mV of tail ISI. (b) REF with the 2-tap TX de-emphasis $c_0 = 0.652$, $c_1 = -0.348$ (Lab~3). The statistical eye at $10^{-12}$ sits only 5\,mV above PDA, while PRBS7 overestimates it by 32\,mV.}
\label{fig:ch4-pda-prbs}
\end{figure}

\begin{pitfall}[Truncating the tail]
PDA sums $\lvert p_k\rvert$ over \emph{all} $k$. Many small cursors add up: for REF, the cursors beyond $k = 10$ are each below 0.008\,V yet add 0.093\,V. Always extend the cursor list until the tail is negligible, and simulate with a PRBS whose period is longer than the channel memory (PRBS13 or longer for REF), or the eye will look better than it is.
\end{pitfall}

%% =====================================================================
\section{The statistical eye}\label{sec:stateye}

PDA tells us the worst case but not how often it happens. A BER target is a statement about probability, so we want the \emph{distribution} of $y_n$. In \cref{eq:ch4-yn} each ISI term $A\,a_{n-k}p_k$ is an independent random variable that takes $+A p_k$ or $-A p_k$ with probability $\tfrac12$. The PDF of a sum of independent variables is the convolution of their PDFs, so the ISI PDF is built one cursor at a time (\cref{fig:ch4-isipdf}a): one cursor gives two spikes, two give four, $m$ give up to $2^m$, and the shape fills in~\citep{p2:casper2002,p2:sanders2004}.

\begin{keyeq}[Statistical eye]
\begin{gather}
  f_\mathrm{ISI}(v) = \mathop{\scalebox{1.4}{$\ast$}}_{k\neq0}\Big[\tfrac12\delta(v - A p_k) + \tfrac12\delta(v + A p_k)\Big],
  \qquad
  f_X = f_\mathrm{ISI} \ast \mathcal N(0,\sign^2),
  \label{eq:ch4-isipdf}\\
  \BER(\ts, v) = \tfrac12 P(A p_0 + X < v) + \tfrac12 P(-A p_0 + X > v),
  \label{eq:ch4-ber}\\
  \BER_J(\ts, v) = \int \BER(\ts - \tau, v)\, f_J(\tau)\,d\tau .
  \label{eq:ch4-berj}
\end{gather}
\end{keyeq}
Gaussian noise of rms $\sign$ is one more convolution in voltage. Sampling-clock jitter moves $\ts$, so it is a convolution in \emph{time} with the jitter PDF $f_J$. Here we use the dual-Dirac model~\cite{p5:agilentdualdirac} (a Gaussian of rms $\sigma_\mathrm{RJ}$ at $\pm\mathrm{DJ}/2$); \cref{sec:dualdirac} explains where it comes from. Evaluating \cref{eq:ch4-ber} over a grid of $(\ts, v)$ gives the \textbf{statistical eye}: contours of constant BER. The \textbf{vertical bathtub} is one column of it, BER versus threshold at a fixed phase. The eye height at a BER target is the width of the bathtub at that BER.

\begin{figure}[tbp]
\centering
\includegraphics[width=\linewidth]{ch04_isipdf}
\caption{\textbf{The ISI distribution is a convolution of two-point distributions.} (a) REF cursors convolved one at a time: each new cursor splits every spike in two. (b) Distribution of the received sample for a ``1'', with all 153 cursors and 1.5\,mV rms noise, on a log scale. Without equalization (red), 6\% of all ``1''s fall below the threshold. With the 2-tap TX FFE (teal), the distribution ends at 61\,mV at $10^{-12}$, just inside the PDA worst case (dashed).}
\label{fig:ch4-isipdf}
\end{figure}

The statistical eye needs a channel that is open, so from here on we use \textbf{REF with the 2-tap TX de-emphasis} of Lab~3, $c_0 = 0.652$ and $c_1 = -0.348$ (\cref{sec:ffe-taps,sec:lab3}). It is the simplest equalizer that opens REF, not the final \OS{} design. \Cref{fig:ch4-stateye} shows its statistical eye with noise and jitter. The code is \code{stat\_eye()} in \code{figsrc/ch04\_stateye.py}: a few dozen lines of NumPy that shift and add the two-point PDFs on a 0.05\,mV grid.

\begin{figure}[tbp]
\centering
\includegraphics[width=\linewidth]{ch04_stateye}
\caption{\textbf{The statistical eye gives BER contours instead of one worst case.} REF with 2-tap TX de-emphasis, $\sign = 1.5$\,mV rms, sampling jitter of 2\,ps RJ and 20\,ps DJ (dual-Dirac). (a) Signal density (shading) and contours at BER $10^{-3}$ to $10^{-12}$; the $10^{-12}$ eye is 119\,mV high and 0.59\,UI wide. (b) Vertical bathtubs at the best phase. Noise and jitter each shave a few millivolts off the 131\,mV ISI-only eye.}
\label{fig:ch4-stateye}
\end{figure}

\begin{example}{PDA versus statistical eye height}\label{we:ch4-stat}
\textbf{Given:} REF with 2-tap de-emphasis. The equalized pulse peaks at $q_0 = 0.209$\,V, 22\,ps earlier than REF. Its residual cursors are $q_{-1} = 0.0155$, $q_1 = 0.0125$, $q_2 = -0.0056$, then a slowly decaying tail of 0.002--0.0045\,V per UI. $\sum_{k\neq0}\lvert q_k\rvert = 0.084$\,V.

\textbf{PDA:} $\mathrm{EH} = 1\,\mathrm{V}\times(0.209 - 0.084) = \mathbf{126}$\,mV, open over 0.75\,UI.

\textbf{Statistical eye, ISI only:} 139\,mV at $10^{-6}$ and \textbf{131}\,mV at $10^{-12}$, open over 0.77\,UI.

\textbf{With noise and jitter:} $\sign = 1.5$\,mV gives 123\,mV; adding 2\,ps RJ and 20\,ps DJ gives \textbf{119}\,mV and cuts the width to 0.59\,UI.

\textbf{Interpretation:} at $10^{-12}$ the statistical eye is only 4\% larger than PDA. A probability of $10^{-12} \approx 2^{-40}$ means roughly 40 cursors line up in their worst direction, and here the cursors beyond the first 40 add only 0.006\,V. PDA is far too pessimistic only when the ISI is spread over many comparable cursors. A large gap between PDA and the statistical eye should make you check your tail, not celebrate.
\end{example}

\begin{insight}
Use the three tools for different jobs. PDA is instant and gives the worst-case pattern, which is the best stimulus for a transistor-level simulation. The statistical eye answers the actual specification, eye height and width at $10^{-12}$, and is where noise and jitter enter. A PRBS7 eye from a short transient is a sanity check, never a sign-off. For REF-like channels, PDA is within a few percent of the $10^{-12}$ eye, so design to PDA early and confirm with the statistical eye.
\end{insight}

\begin{tryit}[Statistical eye in serdespy]
Import \code{stat\_eye} and \code{eye\_metrics} from \code{ch04\_stateye.py}, and feed them the pulse response from any lab (\cref{sec:lab2,sec:lab3}). Change \code{sigma\_n}, \code{rj} and \code{dj} and watch which one moves the height and which the width.
\end{tryit}

%% =====================================================================
\section{The system noise and margin budget}\label{sec:noisebudget}

A \textbf{margin budget} starts from the transmitted swing and subtracts, step by step, everything that eats the eye before the sampler decides. Whatever is left is the \textbf{margin} that covers what the model does not know: process corners, crosstalk, adaptation error. The terms are:
\begin{enumerate}[label=(\roman*)]
  \item \textbf{TX swing} $2A$ = 1\,Vppd.
  \item \textbf{Channel loss}, which shrinks the main cursor to $2A p_0$.
  \item \textbf{Equalization}: a TX FFE costs main cursor (\cref{ch:txffe}); a CTLE and DFE do not, to first order.
  \item \textbf{Residual ISI} at the BER target, from the statistical eye.
  \item \textbf{RX noise}: for Gaussian noise of rms $\sign$, a BER of $10^{-12}$ needs $\mathrm{EH} \ge 2Q\sign = 14.07\,\sign$ (\cref{eq:ber-q}).
  \item \textbf{Jitter}: sampling away from the best phase lowers the eye.
  \item \textbf{Sampler offset} $V_\mathrm{os}$ and \textbf{sensitivity} $V_\mathrm{sens}$ (\cref{sec:offset,sec:sampler}).
\end{enumerate}

\begin{keyeq}[Linear margin budget at the sampler]
\begin{equation}
  M = \mathrm{EH}_\mathrm{ISI}(10^{-12}) - 2Q\,\sign - \Delta V_\mathrm{jitter} - 2V_\mathrm{os} - V_\mathrm{sens} \;>\; 0 .
  \label{eq:ch4-margin}
\end{equation}
\end{keyeq}
An offset moves the threshold toward one rail, so the worse half-eye shrinks by $V_\mathrm{os}$; that costs $2V_\mathrm{os}$ of symmetric eye height. The sensitivity is the minimum peak-to-peak input the sampler resolves in time. All terms are referred to the RX input.

\begin{example}{Margin budget for REF with 2-tap de-emphasis}\label{we:ch4-budget}
\textbf{Assumptions:} $\sign = 1.5$\,mV rms input-referred (CTLE plus samplers); sampling jitter of 2\,ps RJ and 20\,ps DJ, the latter from PI quantization dither (6.25\,ps steps) and CDR wander (\cref{sec:jitter-budget}); offset after calibration $V_\mathrm{os} = 5$\,mV; sampler sensitivity $V_\mathrm{sens} = 20$\,mV.

\textbf{Stack-up} (\cref{fig:ch4-waterfall}): 1000\,mV $\to$ channel $\to$ 340\,mV main cursor $\to$ de-emphasis $\to$ 209\,mV $\to$ residual ISI at $10^{-12}$ $\to$ 131\,mV. Then noise $2Q\sign = 14.07\times1.5 = 21$\,mV, jitter 3.5\,mV, offset $2\times5 = 10$\,mV and sensitivity 20\,mV leave \textbf{$M \approx 76$\,mV}.

\textbf{Statistical check:} the statistical eye with noise and jitter included is 119\,mV, which leaves $119 - 10 - 20 = 89$\,mV. Noise really costs only 8\,mV here, not 21\,mV, because the $7\sigma$ noise peak rarely coincides with a worst-case ISI pattern. The linear budget is conservative by 13\,mV.
\end{example}

\begin{figure}[tbp]
\centering
\includegraphics[width=\linewidth]{ch04_waterfall}
\caption{\textbf{A top-down margin budget for REF with 2-tap de-emphasis.} (a) The channel and the de-emphasis take 79\% of the transmitted swing before any circuit impairment. (b) Zoom on the 209\,mV main cursor: residual ISI is the largest single term. Adding the rest linearly leaves 76\,mV; combining noise and jitter statistically (dashed) leaves 89\,mV.}
\label{fig:ch4-waterfall}
\end{figure}

\begin{insight}[What to spend margin on]
The largest term is residual ISI (79\,mV), not noise (8--21\,mV). The best use of design effort is therefore equalization: the CTLE and 2-tap DFE of \cref{ch:ctle,ch:dfe} remove most of the residual ISI without the main-cursor cost of a TX FFE (\cref{sec:eq-partition}). Keep the margin for things the model leaves out: the 20--25\,dB end of \code{LOSS\_FAMILY} (each extra dB of Nyquist loss shrinks the main cursor by about 6\%), crosstalk (25\,mVpp of NEXT would take a third of it), termination and CTLE drift over corners, and DFE tap error. Do not spend power to push $\sign$ below 1\,mV: halving it buys only a few millivolts.
\end{insight}

%% =====================================================================
\section{Modulation choice: NRZ or PAM4?}\label{sec:pam4}

\textbf{PAM4} sends two bits per symbol on four levels, $\{-A, -A/3, +A/3, +A\}$. At the same bit rate the symbol rate halves (2.5\,GBd for 5\,Gb/s), so the Nyquist frequency halves to 1.25\,GHz and the channel loss there is smaller. The price is that the three stacked eyes are each $2A/3$ high instead of $2A$ (\cref{fig:ch4-pam4}a,b).

\begin{keyeq}[The PAM4 rule]
\begin{equation}
  20\log_{10}3 = 9.54\,\mathrm{dB}
  \qquad\Longrightarrow\qquad
  \text{PAM4 wins only if}\quad \Delta = \mathrm{IL}(\fN) - \mathrm{IL}(\fN/2) > 9.54\,\mathrm{dB}.
  \label{eq:ch4-pam4}
\end{equation}
\end{keyeq}
PAM4 recovers its 9.54\,dB eye penalty only if the channel loses more than that between the two Nyquist frequencies, which are an octave apart. If the loss were proportional to $f$ (pure dielectric), $\Delta$ would be half the NRZ Nyquist loss; if it were proportional to $\sqrt f$ (pure skin effect), $\Delta = 0.29\,\mathrm{IL}(\fN)$. Real FR-4 sits in between.

\begin{example}{NRZ or PAM4 for Open-SERDES?}\label{we:ch4-pam4}
\textbf{Given:} \code{channel()} from \code{link.py} for REF and the \code{LOSS\_FAMILY} lengths.

\textbf{Solution:} \cref{tab:ch4-pam4} lists the losses at both Nyquist frequencies. For our model $\Delta \approx 0.40\,\mathrm{IL}(\fN)$. REF has $\Delta = 17.4 - 10.4 = 7.0$\,dB, so NRZ wins by $9.54 - 7.0 = 2.6$\,dB. $\Delta$ reaches 9.54\,dB at 1.37\,m, where the NRZ Nyquist loss is 23.8\,dB (\cref{fig:ch4-pam4}c). Only the 25 and 30\,dB channels pass the rule, and only by 0.5 and 2.5\,dB.
\end{example}

\begin{table}[tbp]
\centering
\caption{\textbf{Loss at the NRZ and PAM4 Nyquist frequencies} for REF and the \code{LOSS\_FAMILY} channels, from \code{link.py}.}
\label{tab:ch4-pam4}
\small\sffamily
\begin{tabular}{@{}lccccc@{}}
\toprule
Channel & Length (m) & IL at 2.5\,GHz (dB) & IL at 1.25\,GHz (dB) & $\Delta$ (dB) & Rule favors\\
\midrule
5\,dB  & 0.285 & 5.0  & 3.0  & 2.0  & NRZ\\
10\,dB & 0.572 & 10.0 & 6.0  & 4.0  & NRZ\\
15\,dB & 0.862 & 15.0 & 9.0  & 6.0  & NRZ\\
\textbf{REF} & \textbf{1.000} & \textbf{17.4} & \textbf{10.4} & \textbf{7.0} & \textbf{NRZ}\\
20\,dB & 1.152 & 20.0 & 12.0 & 8.0  & NRZ\\
25\,dB & 1.441 & 25.0 & 15.0 & 10.0 & PAM4 (barely)\\
30\,dB & 1.731 & 30.0 & 18.0 & 12.0 & PAM4\\
\bottomrule
\end{tabular}
\end{table}

\begin{figure}[tbp]
\centering
\includegraphics[width=\linewidth]{ch04_pam4}
\caption{\textbf{PAM4 pays 9.54\,dB of eye height to halve the Nyquist frequency.} (a, b) Ideal NRZ and PAM4 eyes at the same peak swing: each PAM4 eye is one third as tall. (c) REF-model loss at 2.5 and 1.25\,GHz versus length, and their difference $\Delta$. PAM4 is favored only where $\Delta$ exceeds 9.54\,dB, beyond 1.37\,m (23.8\,dB at 2.5\,GHz). REF (gold) is 2.6\,dB short of that.}
\label{fig:ch4-pam4}
\end{figure}

The rule is a first-order comparison, and it flatters PAM4. A DFE cancels post-cursors for either format, but whatever ISI is left (pre-cursors, tap errors) weighs three times more against a PAM4 eye. PAM4 also needs three comparators per sampling phase (up to twelve samplers instead of the two data and two edge samplers of our half-rate receiver), reference levels that track the signal amplitude, a linear TX driver and CTLE, and a CDR that copes with several zero-crossing times per transition. It is also three times more sensitive to residual ISI, crosstalk and sampler offset.

\begin{specnote}[Open-SERDES decision]
\OS{} uses \textbf{NRZ}. Over REF and every channel up to 20\,dB, NRZ wins by the 9.54\,dB rule, and it costs a third of the slicers. PAM4 is kept as a Phase-4 option, worth revisiting only for channels beyond about 24\,dB at 2.5\,GHz or for a 10\,Gb/s follow-on (\cref{tab:decisions}).
\end{specnote}

%% =====================================================================
\section{Exercises}\label{sec:ch4-exercises}
\begin{enumerate}[label=\textbf{4.\arabic*},leftmargin=2.2em]
  \item \textbf{PDA by hand.} A channel has $p_{-1} = 0.05$, $p_0 = 0.50$, $p_1 = 0.20$, $p_2 = -0.04$, $p_3 = 0.02$. Find the PDA eye for $2A = 0.8$\,V and the worst-case bit pattern in time order.
  \item \textbf{The tail matters.} Using \code{link.cursors} with \code{n\_post} = 10, 30, 100 and 300, plot the REF PDA eye against \code{n\_post}. How many cursors do you need before the answer changes by less than 5\,mV?
  \item \textbf{Two cursors, exactly.} For ISI from only $p_1$ and $p_2$, write the ISI PDF and the BER versus threshold in closed form, with Gaussian noise $\sign$. Compare with \code{stat\_eye()}.
  \item \textbf{Budget.} Redo \cref{we:ch4-budget} with $\sign = 3$\,mV and an uncalibrated offset of 15\,mV. Is the margin still positive? Which single improvement restores the most margin?
  \item \textbf{PAM4 at 10\,Gb/s.} For a 10\,Gb/s follow-on over REF, compare NRZ ($\fN = 5$\,GHz) with PAM4 ($\fN = 2.5$\,GHz) using \code{loss\_db}. Which wins by the rule?
\end{enumerate}

\begin{takeaways}
\begin{itemize}
  \item At the sampler, the channel is an FIR filter with taps $p_k$, so the cursors alone predict every received sample.
  \item PDA gives the worst-case eye $2A(p_0 - \sum_{k\neq0}\lvert p_k\rvert)$ and the worst-case pattern $a_{n-k} = -\sgn(p_k)\,a_n$. REF is closed: $-0.161$\,V with 13 cursors, $-0.255$\,V with the full tail.
  \item The statistical eye convolves two-point ISI distributions with noise (in voltage) and jitter (in time) to give BER contours. For REF with 2-tap de-emphasis: PDA 126\,mV, statistical 131\,mV at $10^{-12}$, 119\,mV with noise and jitter.
  \item Short PRBS simulations are optimistic (PRBS7: 163\,mV) because they never meet the rare worst-case patterns.
  \item A linear margin budget for that link leaves about 76\,mV (89\,mV statistically). Residual ISI, not noise, is the largest term, so spend design effort on equalization.
  \item PAM4 pays 9.54\,dB of eye height and wins only if the loss between 1.25 and 2.5\,GHz exceeds that. REF has 7.0\,dB, so \OS{} uses NRZ.
\end{itemize}
\end{takeaways}


\tutpart{Equalization}{Undoing the channel. We look at where equalization can live in the link (TX FFE, CTLE, DFE), how much each can do, what it costs, and how to split the work between them for the Open-SERDES channel.}
\tikzsetfigurename{05_eqfund-}
\chapter{Equalization Fundamentals}\label{ch:eqfund}
\chapterlead{The channel is a low-pass filter. An equalizer is the filter we add so that channel and equalizer together look flat out to the Nyquist frequency.}
\suppressfloats[t]

\begin{objectives}
\begin{itemize}
  \item Explain how loss at $\fN$ turns into inter-symbol interference and closes the REF eye.
  \item State the zero-forcing idea and quantify why an ideal channel inverse amplifies noise.
  \item Use Liu's cable model to connect a loss in dB to the size of the main cursor.
  \item Name the four places an equalizer can sit in a link, and what each one can and cannot fix.
  \item Frame the partition question that \cref{sec:eq-partition} answers for \OS{}.
\end{itemize}
\end{objectives}

\section{Why equalize?}\label{sec:ch5-why}

\Cref{sec:refchannel} showed that the REF channel loses \qty{17.4}{\deci\bel} at $\fN=\qty{2.5}{\giga\hertz}$ but only \qty{0.6}{\deci\bel} at DC. A single 1-UI pulse therefore comes out smeared. Its main cursor is only $\pk{0}=\qty{0.340}{\volt}$ instead of \qty{1}{\volt}, and its energy leaks into the neighbouring bits: $\pk{1}=0.182$, $\pk{2}=0.091$ and a long tail after it, plus a small pre-cursor $\pk{-1}=0.049$ (\cref{tab:ref-cursors}). Summed over all bits, this \textbf{inter-symbol interference} (ISI) exceeds the main cursor, so the peak-distortion eye of \cref{eq:pda} is $-0.161$\,V. The eye is closed.

An \textbf{equalizer} is a filter $H_\mathrm{eq}(f)$ placed somewhere in the link so that the cascade $\Hch(f)\,H_\mathrm{eq}(f)$ is flat, at least up to $\fN$~\cite{liuEq,palermo720eq}. Flat to $\fN$ is enough. A pulse sampled once per UI is free of ISI when its spectrum, folded around multiples of $1/\Tb$, adds up to a constant. A response that is flat to $\fN$ and rolls off smoothly above it comes close to meeting that condition. In the time domain, the same goal reads: \emph{make every cursor except $\pk{0}$ zero}. \Cref{fig:ch5-inverse}(a) shows the most direct way to do it, which is to build the exact inverse of the channel.

\begin{figure}[tbp]
  \centering
  \includegraphics[width=\linewidth]{ch05_inverse}
  \caption{\textbf{Inverting the channel removes ISI but amplifies noise.} (a)~REF $|\Hch|$, its ideal inverse and their flat product. The inverse needs $+17.4$\,dB at $\fN$ and keeps rising beyond it. (b)~For white noise at the equalizer input, the noise power spectrum is multiplied by $|H_\mathrm{eq}|^2$. The rms noise gain is \qty{12.1}{\deci\bel} if equalization stops at $\fN$ and \qty{22.5}{\deci\bel} if it extends to $2\fN$.}
  \label{fig:ch5-inverse}
\end{figure}

\section{The zero-forcing idea}\label{sec:zf}

A \textbf{zero-forcing} (ZF) equalizer chooses $H_\mathrm{eq}=1/\Hch$ over the band of interest, which forces all ISI to zero. It ignores noise, and that is its weakness. Thermal noise from the termination and from the first amplifier enters \emph{before} the equalizer and is filtered by it. By the Wiener--Khinchin theorem, a noise process with power spectral density $S_n(f)$ passing through $H_\mathrm{eq}$ comes out with variance:

\begin{keyeq}[Noise at the equalizer output]
\begin{equation}
  \sigma_\mathrm{out}^2=\int_0^{B}\lvert H_\mathrm{eq}(f)\rvert^2\,S_n(f)\,\mathrm{d}f
  \quad\xrightarrow{\;H_\mathrm{eq}=1/\Hch,\ S_n=N_0\;}\quad
  \sigma_\mathrm{out}^2=N_0B\cdot\underbrace{\frac{1}{B}\int_0^{B}\frac{\mathrm{d}f}{\lvert\Hch(f)\rvert^2}}_{\text{noise enhancement}}
  \label{eq:ch5-noise}
\end{equation}
\end{keyeq}

For REF the enhancement factor is \qty{12.1}{\deci\bel} for $B=\fN$ (\cref{fig:ch5-inverse}b). The signal is not boosted by that much on average. The main cursor recovers, but the noise integrated over the whole band grows faster than the signal does. Extending the inverse to $2\fN$ costs another \qty{10}{\deci\bel} of noise for almost no gain in eye opening. Practical equalizers are therefore \emph{partial} inverses. They stop near $\fN$ and limit their peak boost, or, like the DFE of \cref{ch:dfe}, they avoid linear filtering of the noise altogether.

\paragraph{Liu's cable model} A model with a closed-form pulse helps build intuition~\cite{liuEq}. In a cable, dielectric loss dominates, and the transfer function is close to $C(f)=\exp(-k_d\,l\,\lvert f\rvert)$, with $l$ the length and $k_d$ a loss coefficient. Write $a=k_dl$. The inverse Fourier transform is a Lorentzian,
\begin{equation}
  c(t)=\int_{-\infty}^{\infty}e^{-a\lvert f\rvert}e^{j2\pi ft}\,\mathrm{d}f=2\int_0^\infty e^{-af}\cos(2\pi ft)\,\mathrm{d}f=\frac{2a}{a^2+4\pi^2t^2}.
  \label{eq:ch5-lorentz}
\end{equation}
Convolving with a pulse of height $V_0$ on $0\le t\le\Tb$, and using the antiderivative $\int c(u)\,\mathrm{d}u=\tfrac1\pi\arctan(2\pi u/a)$, gives
\begin{equation}
  y(t)=\frac{V_0}{\pi}\left[\arctan\frac{2\pi t}{a}-\arctan\frac{2\pi(t-\Tb)}{a}\right].
  \label{eq:ch5-cablepulse}
\end{equation}
The pulse is symmetric about $t=\Tb/2$, so setting $\mathrm{d}y/\mathrm{d}t=0$ puts the peak there:

\begin{keyeq}[Main cursor of a dielectric-loss cable]
\begin{equation}
  y_{\max}=\frac{2V_0}{\pi}\arctan\frac{\pi\Tb}{k_d\,l},\qquad
  \text{loss at }\fN\ (\mathrm{dB})=20\log_{10}(e)\,k_dl\,\fN=8.686\,k_dl\,\fN .
  \label{eq:ch5-ymax}
\end{equation}
\end{keyeq}

\begin{example}{How much does \qty{15}{\deci\bel} of cable loss cost?}\label{we:ch5-cable}
Take a cable with \qty{15}{\deci\bel} of loss at $\fN=\qty{2.5}{\giga\hertz}$. Then $k_dl=15/(8.686\times\qty{2.5}{\giga\hertz})=\qty{691}{\pico\second}$, and $\pi\Tb/(k_dl)=\pi\cdot200/691=0.909$. \Cref{eq:ch5-ymax} gives $y_{\max}=(2/\pi)\arctan(0.909)=(2/\pi)(0.738)=0.470\,V_0$, so the main cursor has fallen to less than half.

The neighbours follow from \cref{eq:ch5-cablepulse} at $t=\Tb/2+k\Tb$: $y_{\pm1}=0.153$, $y_{\pm2}=0.043$ and $y_{\pm3}=0.019$, in units of $V_0$. Summing all cursors out to $\pm100$\,UI gives $\sum\lvert y_k\rvert=0.527>0.470$, so the PDA eye is closed. \Cref{fig:ch5-cable}(b) shows that the eye of this model closes at about \qty{14}{\deci\bel}.

A zero-forcing inverse needs $+15$\,dB at $\fN$. By \cref{eq:ch5-noise} with $\lvert C\rvert^{-2}=e^{2af}$, its noise enhancement is $(e^{2a\fN}-1)/(2a\fN)=$ \qty{9.5}{\deci\bel} up to $\fN$. Up to $2\fN$ it would be \qty{21.6}{\deci\bel}.
\end{example}

\begin{figure}[!tbp]
  \centering
  \includegraphics[width=\linewidth]{ch05_cable}
  \caption{\textbf{Liu's cable model ties loss in dB to the pulse.} (a)~Pulse responses from \cref{eq:ch5-cablepulse} for 5--20\,dB of loss at $\fN$. The dots mark the cursors of the 15\,dB case. The model has zero phase, so the pre-cursors equal the post-cursors. (b)~Main cursor from \cref{eq:ch5-ymax} and PDA half-eye against loss. The tails close the eye at about \qty{14}{\deci\bel}, while the main cursor is still about half its lossless value.}
  \label{fig:ch5-cable}
\end{figure}

\begin{pitfall}[A zero-phase model is not a real channel]
$C(f)=e^{-k_dl\lvert f\rvert}$ is real and even, so its impulse response starts before $t=0$ and is not causal. That is why its pre- and post-cursors are equal in \cref{fig:ch5-cable}. A real channel such as REF is causal, and its post-cursors dominate ($\pk{1}=0.182$ against $\pk{-1}=0.049$). Use the cable model for intuition about loss and the main cursor, and use REF (\cref{sec:refchannel}) for design numbers.
\end{pitfall}

\begin{insight}
\begin{itemize}
  \item Every 5\,dB of loss at $\fN$ costs a large share of the main cursor and adds tail ISI. Near 15\,dB the tails, not the main cursor, are what close the eye.
  \item Equalize to about $\fN$ and no further. A full inverse to $2\fN$ roughly doubles the noise penalty in dB (\qty{12.1}{\deci\bel} becomes \qty{22.5}{\deci\bel} on REF) and buys almost no eye.
  \item Zero forcing ignores noise and swing. Every practical equalizer in this part of the book is a constrained version of it.
\end{itemize}
\end{insight}

\section{Where equalizers live}\label{sec:eq-taxonomy}

An equalizer can be placed at the transmitter, at the receiver input, or after the receiver's decision. \Cref{fig:ch5-eqmap} places the options in the \OS{} link, and \cref{tab:ch5-eqcompare} compares them.

\begin{figure}[!tbp]
\centering
\begin{tikzpicture}[font=\sffamily\small, node distance=4mm and 5mm]
  \node[blkg, minimum width=14mm] (ser) {SER\\[-1pt]\scriptsize 16:1};
  \node[blkb, right=5mm of ser, minimum width=19mm] (ffe) {TX FFE\\[-1pt]\scriptsize + CML driver};
  \node[blk, right=7mm of ffe, minimum width=17mm, fill=bgpale] (ch) {channel\\[-1pt]\scriptsize pkg, PCB, pkg};
  \node[blkb, right=7mm of ch, minimum width=14mm] (ctle) {CTLE};
  \node[sumnode, right=5mm of ctle] (sum) {$\Sigma$};
  \node[blkg, right=5mm of sum, minimum width=15mm] (slc) {slicers\\[-1pt]\scriptsize half-rate};
  \node[blkg, right=5mm of slc, minimum width=13mm] (des) {DES\\[-1pt]\scriptsize 1:16};
  \node[blkb, below=7mm of slc, minimum width=15mm, minimum height=7mm] (dfe) {DFE taps};
  \draw[sig] (ser) -- (ffe);
  \draw[sig] (ffe) -- (ch);
  \draw[sig] (ch) -- (ctle);
  \draw[sig] (ctle) -- (sum);
  \draw[sig] (sum) -- (slc);
  \draw[sig] (slc) -- (des);
  \draw[sig] ($(slc.east)+(2.5mm,0)$) |- (dfe.east);
  \draw[sig] (dfe.west) -| node[pos=0.92, right, font=\sffamily\footnotesize, text=navy]{$-$} (sum.south);
  \node[blk, dashed, draw=muted, text=muted, above=4mm of ctle, minimum width=14mm, minimum height=6mm, font=\sffamily\scriptsize] (rxffe) {RX FFE (not used)};
  % clocks
  \node[blkg, below=7mm of ser, minimum width=12mm, minimum height=7mm] (pll) {PLL};
  \node[blkg, above=4mm of slc, minimum width=12mm, minimum height=6mm] (cdr) {CDR};
  \draw[clk] (pll) -- (ser);
  \draw[clk] (cdr) -- (slc);
  % tags
  \node[lbl, align=center, below=1mm of ch] {loss turns\\into ISI};
  \node[font=\sffamily\scriptsize, text=c3, align=left, anchor=north] (txtag) at ($(ffe.south)+(0,-1.5mm)$) {linear, pre + post\\no RX noise gain\\costs launch swing};
  \node[font=\sffamily\scriptsize, text=c4, align=left, anchor=north] at ($(ctle.south)+(0,-1mm)$) {linear, analog\\boosts HF noise\\and crosstalk};
  \node[font=\sffamily\scriptsize, text=c3, align=left, anchor=west] (dfetag) at ($(dfe.east)+(3mm,0)$) {nonlinear\\post-cursors only\\no noise gain};
  % regions
  \begin{scope}[on background layer]
    \node[region, fit=(ser)(ffe)(pll)(txtag), inner sep=2.5mm, label={[lbl]above:transmitter}] {};
    \node[region, fit=(ctle)(sum)(slc)(des)(dfe)(cdr)(rxffe)(dfetag), inner sep=2.5mm, label={[lbl]above:receiver}] {};
  \end{scope}
\end{tikzpicture}
\caption{\textbf{Equalization can sit at the TX, at the RX input, or after the decision.} The \OS{} link uses a TX FFE in the driver, a CTLE as the first receiver stage and a DFE that feeds back from the half-rate slicers. A receiver FIR (RX FFE) is the fourth option. We do not use it.}
\label{fig:ch5-eqmap}
\end{figure}

\paragraph{TX FFE} A FIR filter in the driver pre-distorts the waveform (\cref{ch:txffe}). It handles pre- and post-cursors and cannot amplify receiver noise, but with a fixed peak swing it can only \emph{lower} the low-frequency level, so the eye arrives smaller.

\paragraph{RX CTLE} An analog high-pass stage at the receiver input (\cref{ch:ctle}). It is cheap, fast and adaptable, and reshapes the long tail and part of the pre-cursor, but it boosts high-frequency noise and crosstalk with the signal.

\paragraph{RX FFE} A receiver FIR, from analog delay cells or after an ADC. It amplifies noise like a CTLE, and needs accurate \qty{200}{\pico\second} analog delays or a power-hungry ADC. \OS{} does not use it.

\paragraph{DFE} Subtracts the ISI that already-decided bits will cause (\cref{ch:dfe}). Its feedback is a clean $\pm1$ decision, so it adds no noise, but it cancels only \emph{post}-cursors and its first tap must close a loop within one UI.

\Cref{fig:ch5-anatomy} reads the REF pulse in these terms. Each zone of cursors has a natural owner.

\begin{figure}[!tbp]
  \centering
  \includegraphics[width=\linewidth]{ch05_anatomy}
  \caption{\textbf{Each part of the pulse response has a natural owner.} The pre-cursor (red) can only be removed by linear equalization, either the TX pre-tap or the CTLE. The first post-cursors (blue) are the DFE's job, which it does without noise. The long tail (teal) is best flattened by the CTLE or the TX post-tap, because removing it with DFE taps would take many taps. Shown is the unequalized REF pulse with $\pk{-1}..\pk{10}$.}
  \label{fig:ch5-anatomy}
\end{figure}

\begin{table}[!tbp]
\centering\small
\caption{\textbf{Four equalizers, four sets of costs.} Power and complexity are relative, for a 5\,Gb/s link.}
\label{tab:ch5-eqcompare}
\begin{tabularx}{\linewidth}{@{}l L L L l L@{}}
\toprule
 & \textbf{Cursors it fixes} & \textbf{Frequency view} & \textbf{Noise, crosstalk} & \textbf{Power} & \textbf{Main limit} \\
\midrule
TX FFE  & pre + post (as many as taps) & lowers LF, Nyquist kept & no RX noise gain; launch swing lost & low & peak swing; not adaptive without a back-channel \\
RX CTLE & long tail, some pre-cursor & boosts HF over LF & amplifies HF noise and crosstalk & low & peaking range and linearity \\
RX FFE  & pre + post & like a CTLE, discrete taps & amplifies noise & high & analog delays or ADC power \\
DFE     & post only (one per tap) & nonlinear (no $H(f)$) & no noise gain & medium & tap-1 timing, error propagation \\
\bottomrule
\end{tabularx}
\end{table}

\section{Preview: who does how much?}\label{sec:ch5-preview}

None of the four is free: the TX FFE costs swing, the CTLE noise and linearity, the DFE timing margin and loop power. The design question is how to \emph{split} the work of undoing \qty{17.4}{\deci\bel} between them. Chapters~\ref{ch:txffe}--\ref{ch:dfe} build each equalizer, and \cref{sec:eq-partition} answers the split with numbers for REF and for 5--30\,dB channels.

\begin{insight}[Design rule of thumb]
Give each equalizer the job it does most cheaply. The TX pre-tap takes the pre-cursor, the DFE takes the first one or two post-cursors without adding noise, and a linear equalizer (CTLE or TX post-tap) flattens the long tail. Any equalizer pushed far beyond its natural job costs more than it returns.
\end{insight}

\begin{takeaways}
\begin{itemize}
  \item Loss at $\fN$ spreads each bit into its neighbours. On REF the ISI sum exceeds the main cursor, so the PDA eye is $-0.161$\,V.
  \item Equalization aims for a cascade that is flat up to $\fN$. In the time domain that means zero ISI cursors.
  \item A zero-forcing inverse multiplies the noise by $\int\lvert H_\mathrm{eq}\rvert^2$. That is \qty{12.1}{\deci\bel} on REF up to $\fN$ and \qty{22.5}{\deci\bel} up to $2\fN$.
  \item In Liu's cable model, $y_{\max}=(2V_0/\pi)\arctan(\pi\Tb/k_dl)$. At 15\,dB of loss the main cursor is 0.47 and the eye is already closed.
  \item The TX FFE, CTLE, RX FFE and DFE differ in which cursors they fix, whether they amplify noise, and what limits them.
  \item The partition between them is a design decision. \Cref{sec:eq-partition} makes it for \OS{}.
\end{itemize}
\end{takeaways}

\tikzsetfigurename{06_txffe-}
\chapter{Transmit FFE}\label{ch:txffe}
\chapterlead{The transmitter knows every bit before it leaves the chip, so it can pre-distort the waveform to cancel the ISI the channel is about to add.}
\suppressfloats[t]

\begin{objectives}
\begin{itemize}
  \item Describe the TX feed-forward equalizer (FFE) as a 3-tap FIR filter, and write its output and its effect on the pulse response.
  \item Derive the FFE frequency response and the de-emphasis in dB, and explain the peak-swing constraint $\sum\lvert c_k\rvert=1$.
  \item Solve for zero-forcing taps as a small matrix equation, and say when least squares does better.
  \item Show on the REF channel why zero-forcing taps are not eye-optimal (Lab~3).
  \item Pick practical settings: standard presets, tap resolution, and the swing the FFE costs.
\end{itemize}
\end{objectives}

\section{A FIR filter in the driver}\label{sec:ch6-fir}

The TX FFE is a finite-impulse-response filter that acts on the symbol stream $a_n\in\{-1,+1\}$ before the driver (\cref{fig:ch6-fir}). With three taps, the driver output during UI $n$ is
\begin{equation}
  x_n=A\left(c_{-1}\,a_{n+1}+c_0\,a_n+c_1\,a_{n-1}\right).
  \label{eq:ch6-fir}
\end{equation}
The \textbf{pre-tap} $c_{-1}$ weights the \emph{next} bit. The transmitter already knows that bit, because in hardware the whole chain is delayed by one UI. The \textbf{main tap} is $c_0$ and the \textbf{post-tap} is $c_1$. Liu writes the general $(2M+1)$-tap version as $H(z)=\sum_{k=0}^{2M}\alpha_{-M+k}z^{-k}$ with $z=e^{j\omega\Tb}$~\cite{liuEq}. It has no feedback, so it is unconditionally stable.

Since the channel is linear, sending $x_n$ is the same as sending three scaled, shifted copies of each bit. The pulse response seen at the receiver is therefore the convolution of the tap list with the cursor list:
\begin{equation}
  q_k=\sum_{j=-1}^{1}c_j\,\pk{k-j}=c_{-1}\pk{k+1}+c_0\pk{k}+c_1\pk{k-1}.
  \label{eq:ch6-conv}
\end{equation}
\Cref{eq:ch6-conv} is all we need. Every FFE question in this chapter is a question about the new cursors $q_k$.

\begin{figure}[tbp]
\centering
\begin{tikzpicture}[font=\sffamily\small,
    tap/.style={draw=navy, line width=0.7pt, fill=btint, isosceles triangle, isosceles triangle apex angle=60,
                shape border rotate=-90, inner sep=0pt, minimum size=7mm}]
  \node[text=navy] (in) at (0,0) {$a_{n+1}$};
  \node[blk, minimum width=11mm, minimum height=7mm] (d1) at (3.0,0) {$z^{-1}$};
  \node[blk, minimum width=11mm, minimum height=7mm] (d2) at (6.0,0) {$z^{-1}$};
  \coordinate (n0) at (1.4,0); \coordinate (n1) at (4.5,0); \coordinate (n2) at (7.5,0);
  \draw[sig] (in) -- (d1);
  \draw[sig] (d1) -- node[above, lbl, pos=0.45]{$a_n$} (d2);
  \draw[line width=0.75pt, navy] (d2.east) -- (n2) node[above, lbl, pos=0.4]{$a_{n-1}$};
  \foreach \n/\c/\i in {n0/{c_{-1}}/0, n1/{c_0}/1, n2/{c_1}/2}{
    \node[dot] at (\n) {};
    \node[tap] (t\i) at ($(\n)+(0,-1.25)$) {};
    \draw[line width=0.75pt, navy] (\n) -- (t\i.north);
    \node[right=1mm of t\i, text=navy] {$\c$};
  }
  \node[sumnode] (s) at (4.5,-2.75) {$\Sigma$};
  \draw[sig] (t0.south) -- (s.north west);
  \draw[sig] (t1.south) -- (s.north);
  \draw[sig] (t2.south) -- (s.north east);
  \node[blkb, minimum width=16mm, minimum height=8mm] (drv) at (8.0,-2.75) {CML\\[-2pt]\scriptsize driver};
  \draw[sig] (s.east) -- (drv.west);
  \draw[sig] (drv.east) -- ++(0.9,0) node[right, text=navy]{$x(t)$};
  \node[align=left, anchor=west, font=\footnotesize, text=muted] at (9.1,-1.0) {$x_n=A\sum_j c_j\,a_{n-j}$\\[1pt] $\sum_j\lvert c_j\rvert=1$};
  \node[lbl, anchor=south] at (1.4,0.25) {from serializer};
\end{tikzpicture}
\caption{\textbf{A 3-tap TX FFE is a short delay line with weighted taps.} Each bit is sent three times: once early with weight $c_{-1}$, once on time with $c_0$, and once late with $c_1$. In a CML driver the taps are parallel current-steering slices that share the output load (\cref{sec:ffe-impl}).}
\label{fig:ch6-fir}
\end{figure}

\section{De-emphasis, pre-emphasis and the swing limit}\label{sec:ffe-taps}

The driver's output swing is fixed by its tail current and load at $\pm A$. By \cref{eq:ch6-fir}, the largest output occurs when all terms add, so it equals $A\sum\lvert c_k\rvert$. To keep the driver within its swing we require the \textbf{peak-power constraint}:
\begin{equation}
  \lvert c_{-1}\rvert+\lvert c_0\rvert+\lvert c_1\rvert=1 .
  \label{eq:ch6-peak}
\end{equation}
For a low-pass channel the side taps are negative. Consider a 2-tap FFE with $c_1<0$. A \emph{transition} bit ($a_n=-a_{n-1}$) gets $\lvert x_n\rvert=A(c_0+\lvert c_1\rvert)=A$, the full swing. A \emph{repeated} bit gets $A(c_0-\lvert c_1\rvert)$, which is smaller. This is \textbf{de-emphasis}: transitions keep full swing and runs of identical bits are attenuated (\cref{fig:ch6-deemph}).

\textbf{Pre-emphasis} has exactly the same filter shape, but it is scaled the other way. It keeps $c_0=1$ and lifts the transitions to $A(1+\lvert c_1\rvert)$. A driver that already runs at its maximum swing has no headroom for that, so a practical TX de-emphasizes.

\begin{figure}[!tbp]
  \centering
  \includegraphics[width=\linewidth]{ch06_deemph}
  \caption{\textbf{De-emphasis keeps transitions at full swing and halves the repeated bits.} (a)~TX output for a 24-bit pattern without FFE, and with 6\,dB of 2-tap de-emphasis ($c_0=0.751$, $c_1=-0.249$). (b)~The same signals after the REF channel. Without FFE a long run charges the line far beyond the transition levels, so the next transition cannot cross in time. With de-emphasis the run and transition levels are much closer, and the dots (one sample per UI) form a cleaner pattern.}
  \label{fig:ch6-deemph}
\end{figure}

\paragraph{Frequency response} The FFE transfer function is $H_\mathrm{FFE}(f)=\sum_k c_k\,e^{-j2\pi fk\Tb}$. For two taps, $\lvert H\rvert^2=c_0^2+c_1^2+2c_0c_1\cos(2\pi f\Tb)$. In Liu's normalization $c_0=1$, $c_1=\alpha$, with $-1<\alpha<0$:

\begin{keyeq}[2-tap FFE response and de-emphasis]
\begin{equation}
  \lvert H_\mathrm{FFE}(f)\rvert=\sqrt{1+\alpha^2+2\alpha\cos(2\pi f\Tb)},\qquad
  \mathrm{DE}=20\log_{10}\frac{\lvert H(\fN)\rvert}{\lvert H(0)\rvert}=20\log_{10}\frac{1-\alpha}{1+\alpha}
  \label{eq:ffe-resp}
\end{equation}
With $\sum\lvert c\rvert=1$ the response is $c_0+c_1=1-2\lvert c_1\rvert$ at DC and $c_0-c_1=1$ at $\fN$, so $\mathrm{DE}=-20\log_{10}(1-2\lvert c_1\rvert)$.
\end{keyeq}

As a check, \cref{eq:ffe-resp} gives $1+\alpha$ at DC and $1-\alpha$ at $\fN$, matching Liu's sketch. (The handwritten notes contain a typo, $\sqrt{(1+\alpha)^2+2\alpha\cos\omega\Tb}$, which fails this DC check.) The response rises monotonically from DC to $\fN$ and repeats with period $1/\Tb=\qty{5}{\giga\hertz}$ (\cref{fig:ch6-ffefreq}a). It never rises above 0\,dB. It only pulls the low frequencies down.

\begin{example}{The PCIe de-emphasis presets in volts}\label{we:ch6-pcie}
PCIe Gen2 at 5\,GT/s specifies TX de-emphasis of \qty{-3.5}{\deci\bel} and \qty{-6}{\deci\bel}. For $-3.5$\,dB, $1-2\lvert c_1\rvert=10^{-3.5/20}=0.668$, so $c_1=-0.166$ and $c_0=0.834$. For $-6$\,dB, $1-2\lvert c_1\rvert=0.501$, so $c_1=-0.249$ and $c_0=0.751$.

With $A=\qty{0.5}{\volt}$ (1\,Vppd), transition bits swing $\pm\qty{500}{\milli\volt}$ in both cases. Repeated bits swing $\pm0.668\times500=\pm\qty{334}{\milli\volt}$ at $-3.5$\,dB and $\pm\qty{251}{\milli\volt}$ at $-6$\,dB. A 9.5\,dB setting would need $c_1=-1/3$ and leave repeated bits at $\pm\qty{167}{\milli\volt}$.

In a CML driver the taps steer a fixed total tail current, so the supply current is the same at every setting. De-emphasis does not save power. It only redistributes the swing.
\end{example}

\begin{figure}[!tbp]
  \centering
  \includegraphics[width=\linewidth]{ch06_ffefreq}
  \caption{\textbf{A TX FFE pulls the low frequencies down; it never lifts $\fN$.} (a)~\Cref{eq:ffe-resp} for 0, 3.5, 6 and 9.5\,dB of de-emphasis. The response peaks at $\fN$ and is periodic in $\fb$. (b)~REF channel times FFE. Each setting removes as many dB of the low-frequency excess as its de-emphasis, so the cascade gets flatter but stays at $-17.4$\,dB at $\fN$.}
  \label{fig:ch6-ffefreq}
\end{figure}

\section{Solving for the taps}\label{sec:ch6-solve}

\paragraph{Zero forcing} Choose the three taps so that the equalized pulse has $q_{-1}=0$, $q_0=1$ and $q_1=0$. Writing \cref{eq:ch6-conv} for $k=-1,0,1$ gives three linear equations in three unknowns. Each row is one cursor of the equalized pulse, and each column is the pulse shifted by one tap:

\begin{keyeq}[Zero-forcing FFE taps]
\begin{equation}
  \begin{bmatrix}\pk{0}&\pk{-1}&\pk{-2}\\ \pk{1}&\pk{0}&\pk{-1}\\ \pk{2}&\pk{1}&\pk{0}\end{bmatrix}
  \begin{bmatrix}c_{-1}\\ c_0\\ c_1\end{bmatrix}=
  \begin{bmatrix}0\\1\\0\end{bmatrix},
  \qquad c_k\leftarrow\frac{c_k}{\sum_j\lvert c_j\rvert}.
  \label{eq:ch6-zf}
\end{equation}
\end{keyeq}

The matrix has constant diagonals: it is a \textbf{Toeplitz} matrix. A $(2M+1)$-tap FFE gives a $(2M+1)\times(2M+1)$ system of the same form. The normalization step comes last. It scales the whole equalized pulse down, but it does not change where the zeros are.

\begin{example}{Liu's 3-tap zero-forcing FFE}\label{we:ch6-liu}
Liu's example pulse has $x_{-1}=0.2$, $x_0=0.85$, $x_1=-0.2$, $x_2=0.1$ and $x_3=0.05$~\cite{liuEq}. \Cref{eq:ch6-zf} reads
\[
\begin{bmatrix}0.85&0.2&0\\-0.2&0.85&0.2\\0.1&-0.2&0.85\end{bmatrix}
\begin{bmatrix}\alpha_{-1}\\\alpha_0\\\alpha_1\end{bmatrix}=\begin{bmatrix}0\\1\\0\end{bmatrix}
\;\Rightarrow\;\boldsymbol\alpha=[-0.248,\ 1.053,\ 0.277].
\]
Since $\sum\lvert\alpha\rvert=1.578$, the normalized taps are $[-0.157,\ 0.667,\ 0.176]$. The equalized cursors for $k=-2..4$ are $[-0.031,\ 0,\ 0.634,\ 0,\ 0.024,\ 0.051,\ 0.009]$. With $2A=1$, the PDA eye opens from $0.85-0.55=0.30$ to $0.634-0.115=0.519$ (\cref{fig:ch6-zf}). Three lessons:
\begin{itemize}
  \item The post-tap is \emph{positive}, because $x_1$ is negative. ZF inverts whatever ISI is there. This FFE is not even high-pass: it is $-3.3$\,dB at DC and $-3.8$\,dB at $\fN$.
  \item ISI is pushed outside the forced window. A new cursor appears at $k=-2$, and $k=3$ grows from 6\,\% to 8\,\% of the main cursor.
  \item Normalization costs main cursor. It drops from 0.85 to 0.634.
\end{itemize}
\end{example}

\begin{figure}[!tbp]
  \centering
  \includegraphics[width=\linewidth]{ch06_zf}
  \caption{\textbf{Zero forcing zeros only the cursors it is told to.} (a)~Liu's example pulse. (b)~After the normalized 3-tap ZF FFE, the cursors at $k=\pm1$ are exactly zero, but ISI appears at $k=-2$ and grows at $k=+3$ (circled). Normalization to $\sum\lvert c\rvert=1$ lowers the main cursor to 0.634.}
  \label{fig:ch6-zf}
\end{figure}

\paragraph{Least squares and MMSE} Zero forcing looks at only $2M+1$ cursors. \textbf{Least squares} (LS) uses the whole pulse. Stack shifted copies of the full pulse as the columns of a matrix $\mathbf H$ (one row per output cursor), then minimize $\lVert\mathbf H\mathbf c-\mathbf y_\mathrm{des}\rVert^2$, where $\mathbf y_\mathrm{des}$ is 1 at the main cursor and 0 elsewhere. The solution is $\mathbf c=(\mathbf H^\mathsf{T}\mathbf H)^{-1}\mathbf H^\mathsf{T}\mathbf y_\mathrm{des}$, which is then normalized. For a receiver FIR, where noise passes through the filter, the noise-aware \textbf{MMSE} form adds $(\sign^2/\sigma_a^2)\mathbf I$ to $\mathbf H^\mathsf{T}\mathbf H$. Palermo's 10\,Gb/s example gives $\mathbf c=[-0.131,\ 0.595,\ -0.274]$~\cite{palermo720eq}. Its DC gain is $\sum c=0.190$ ($-14.4$\,dB) and its Nyquist gain is $\lvert\sum(-1)^kc_k\rvert=1$, so it applies \qty{14.4}{\deci\bel} of de-emphasis, all of it paid in swing.

\paragraph{The REF channel} Lab~3 applies \cref{eq:ch6-zf} to the REF cursors $\pk{-1}=0.0486$, $\pk{0}=0.3396$ and $\pk{1}=0.1817$. The 3-tap ZF taps are $[-0.0853,\ 0.5959,\ -0.3188]$, which is \qty{14.3}{\deci\bel} of boost. The 2-tap ZF taps, forcing only $q_1=0$, are $[0.6515,\ -0.3485]$, or \qty{10.4}{\deci\bel}. \Cref{fig:ch6-ref} shows the result for PRBS7. The eye goes from \qty{-77.8}{\milli\volt} without FFE to \qty{162.7}{\milli\volt} (0.89\,UI) with the 2-tap FFE, but only to \qty{144.8}{\milli\volt} (0.83\,UI) with the 3-tap ZF FFE. (With PRBS13 the 2-tap eye is \qty{148}{\milli\volt}. Its PDA and statistical eyes, \qty{126}{\milli\volt} and \qty{131}{\milli\volt} at $10^{-12}$, are worked out in \cref{sec:pda,sec:stateye}.)

So the 3-tap zero-forcing FFE is \emph{worse} than the 2-tap one. Under $\sum\lvert c\rvert=1$, every unit of tap weight comes out of the main cursor, which drops from 0.340 to 0.175\,V with three taps. The pre-tap spends 0.085 of the budget to null $\pk{-1}=0.049$ exactly, and it leaks $\pk{2}$ into $q_1$. A grid search for the 3-tap setting with the largest PDA eye finds $[-0.04,\ 0.61,\ -0.35]$. Its PRBS7 eye is \qty{163.6}{\milli\volt}, only \qty{1}{\milli\volt} better than the 2-tap FFE. \emph{Zero forcing is a recipe, not an optimum.} Judge taps by the eye they produce.

\begin{figure}[!tbp]
  \centering
  \includegraphics[width=\linewidth]{ch06_ref}
  \caption{\textbf{On REF, a 2-tap FFE opens the eye; the 3-tap ZF FFE does slightly worse.} (a)~Pulse responses without FFE, with the 2-tap ZF taps $[0.652,-0.349]$ and with the 3-tap ZF taps $[-0.085,0.596,-0.319]$. Dots mark the cursors. (b,\,c)~PRBS7 eyes at the RX input, plotted on the same scale. The open eye is \qty{163}{\milli\volt} but only about a third of the unequalized swing. That is the price of equalizing at the TX.}
  \label{fig:ch6-ref}
\end{figure}

\begin{tryit}[Lab 3: TX FFE]
\code{lab3\_tx\_ffe.py} computes the ZF taps with \code{sdp.forcing\_ffe}, normalizes them, and measures the eyes quoted above. Change the pre-tap by hand and watch the eye move in both directions (\cref{sec:lab3}).
\end{tryit}

\section{Practical notes}\label{sec:ch6-practical}

\paragraph{Presets} Standards define a few TX presets rather than free taps, because a transmitter cannot measure the channel. PCIe Gen2 uses \qty{-3.5}{\deci\bel} and \qty{-6}{\deci\bel} (\cref{we:ch6-pcie}). \OS{} offers the same two presets plus a small pre-tap. \Cref{sec:eq-partition} shows that $-6$\,dB is the right preset for REF once a DFE follows.

\paragraph{Tap resolution} A CML driver builds each tap from $N$ equal current slices, so the taps move in steps of $1/N$. For the REF 2-tap FFE, 8 slices give $[0.625,-0.375]$ and a \qty{162.6}{\milli\volt} eye. 32 slices give $[0.656,-0.344]$ and \qty{162.3}{\milli\volt}. The ideal setting gives \qty{162.7}{\milli\volt}. The eye is flat near its optimum, so 4--5 bits (16--32 slices) are plenty at this data rate. Palermo quotes 4--6 bits as typical~\cite{palermo720eq}.

\begin{insight}
A TX FFE does not boost the high frequencies. It \emph{lowers the low frequencies}, because the driver is already at its peak swing. The 10.4\,dB ZF setting on REF leaves repeated bits at 30\,\% of full swing, and the eye at the receiver is about a third of the unequalized swing. Receiver noise, crosstalk and sampler offset do not shrink with it, so the SNR falls by the de-emphasis. Keep the TX FFE at moderate de-emphasis (3.5--6\,dB) with a small pre-tap, and let the receiver equalizers do the rest. How the driver builds the taps is covered in \cref{sec:ffe-impl}.
\end{insight}

\begin{takeaways}
\begin{itemize}
  \item A TX FFE is a short FIR filter. Its taps convolve with the cursor list: $q_k=\sum_jc_j\pk{k-j}$.
  \item The peak-swing constraint $\sum\lvert c_k\rvert=1$ forces de-emphasis. Transitions keep full swing and repeated bits are attenuated.
  \item A 2-tap FFE has $\lvert H\rvert=\sqrt{1+\alpha^2+2\alpha\cos2\pi f\Tb}$, with de-emphasis $-20\log_{10}(1-2\lvert c_1\rvert)$. PCIe's $-3.5$ and $-6$\,dB correspond to $c_1=-0.166$ and $-0.249$.
  \item Zero-forcing taps solve a small Toeplitz system. Least squares uses the whole pulse, and MMSE adds noise.
  \item On REF, 2-tap ZF gives \qty{162.7}{\milli\volt} and 3-tap ZF gives \qty{144.8}{\milli\volt}. Zero forcing is not eye-optimal.
  \item The FFE costs launch swing and SNR, not power. Use moderate presets and 4--5-bit taps.
\end{itemize}
\end{takeaways}

\tikzsetfigurename{07_ctle-}
\chapter{Continuous-Time Linear Equalizer}\label{ch:ctle}
\chapterlead{A CTLE is a high-pass filter built into the first receiver stage. It lifts the frequencies the channel attenuated by giving up gain at low frequencies.}
\suppressfloats[t]

\begin{objectives}
\begin{itemize}
  \item Compare passive RC, RLC and degenerated-pair CTLEs and know why the degenerated pair is the workhorse.
  \item Derive the degenerated-pair transfer function, with its zero, two poles, DC gain and peaking.
  \item Turn a peaking target into $R_S$, $C_S$, $R_D$ and $\CL$, and respect the gain--bandwidth limit.
  \item Choose the peaking for the REF channel, and recognize over-peaking.
  \item Design an \sgx{} HBT CTLE, explain why its realized boost falls short of the ideal, and make it programmable from 0 to 12\,dB.
\end{itemize}
\end{objectives}

\section{From RC to RLC: three passive ideas}\label{sec:ch7-passive}

The channel is a low-pass filter, so its inverse must be a high-pass filter. Liu's notes walk through three ways to build one, shown in \cref{fig:ch7-passive}~\cite{liuEq}.

\begin{figure}[tbp]
\centering
\begin{tikzpicture}[american, scale=0.82, transform shape, font=\small]
  % (a) RC high-pass
  \begin{scope}
    \draw (0,0) node[left]{$v_\mathrm{in}$} to[C, l=$C$, o-*] (2,0) to[short,-o] (2.8,0) node[right]{$v_\mathrm{out}$};
    \draw (2,0) to[R, l_=$R$] (2,-1.7) node[ground]{};
    \node[font=\sffamily\small, text=navy] at (1.3,-2.6) {(a) RC high-pass};
  \end{scope}
  % (b) compensated divider
  \begin{scope}[xshift=5.0cm]
    \draw (0,0) node[left]{$v_\mathrm{in}$} to[short,o-*] (0.4,0);
    \draw (0.4,0) -- (0.4,0.65) to[C, l=$C_1$] (2.2,0.65) -- (2.2,0);
    \draw (0.4,0) -- (0.4,-0.65) to[R, l_=$R_1$] (2.2,-0.65) -- (2.2,0);
    \draw (2.2,0) to[short,*-*] (3.1,0) to[short] (4.1,0) to[short,*-o] (4.8,0) node[right]{$v_\mathrm{out}$};
    \draw (3.1,0) to[R, l_=$R_2$] (3.1,-1.7) node[ground]{};
    \draw (4.1,0) to[C, l=$C_2$] (4.1,-1.7) node[ground]{};
    \node[font=\sffamily\small, text=navy] at (2.2,-2.6) {(b) $R_1\Vert C_1$ / $R_2\Vert C_2$ divider};
  \end{scope}
  % (c) RLC load (half circuit)
  \begin{scope}[xshift=12.2cm, yshift=-0.3cm]
    \draw (0,2.5) node[vcc]{$V_\mathrm{DD}$} to[L, l_=$L$] (0,1.3) to[R, l_=$R_S$] (0,0.1) node[circ]{};
    \draw (0,0.1) -- (1.0,0.1) node[circ]{} to[short,-o] (1.7,0.1) node[right]{$v_\mathrm{out}$};
    \draw (0,0.1) to[I, l_=$g_mv_\mathrm{in}$] (0,-1.4) node[ground]{};
    \draw (1.0,0.1) to[C, l=$C$] (1.0,-1.4) node[ground]{};
    \node[font=\sffamily\small, text=navy] at (0.7,-2.3) {(c) RLC load};
  \end{scope}
\end{tikzpicture}
\caption{\textbf{Three passive ways to make a high-pass equalizer.} (a)~A plain RC high-pass has no DC gain. (b)~A deliberately over-compensated divider has a finite DC gain but attenuates at every frequency. (c)~An inductively peaked load gives real boost, but with complex poles that ring.}
\label{fig:ch7-passive}
\end{figure}

\paragraph{RC high-pass} $H(s)=sRC/(1+sRC)$ has zero gain at DC. A long run of identical bits decays toward 0\,V, so the baseline wanders with the data. On its own it is unusable.

\paragraph{Over-compensated divider} With a series $R_1\Vert C_1$ and a shunt $R_2\Vert C_2$,
\begin{equation}
  H(s)=\frac{R_2}{R_1+R_2}\cdot\frac{1+sR_1C_1}{1+s\,(R_1\Vert R_2)(C_1+C_2)},\qquad
  \frac{\omega_p}{\omega_z}=\frac{R_1+R_2}{R_2}\cdot\frac{C_1}{C_1+C_2}.
  \label{eq:ch7-passive}
\end{equation}
The DC gain is $R_2/(R_1+R_2)$ and the high-frequency gain is $C_1/(C_1+C_2)$. With $R_1C_1=R_2C_2$ it is a flat compensated probe; a larger $R_1C_1$ makes it an equalizer. It is very linear, but its gain is below 1 everywhere (``no gain'', in Liu's words), which degrades the SNR at the next stage.

\paragraph{RLC load} Put an inductor in series with the load resistor $R_S$ of a differential pair and let $C$ be the node capacitance. With $\omega_n=1/\sqrt{LC}$ and $\zeta=\tfrac{R_S}{2}\sqrt{C/L}$, the load impedance is
\begin{equation}
  Z(s)=R_S\,\frac{\omega_n}{2\zeta}\cdot\frac{s+2\zeta\omega_n}{s^2+2\zeta\omega_n s+\omega_n^2},\qquad
  \frac{\lvert Z(j\omega_n)\rvert}{R_S}=\frac{\sqrt{1+4\zeta^2}}{4\zeta^2}.
  \label{eq:ch7-rlc}
\end{equation}
With $\zeta=0.2$ the boost is 6.7, or \qty{16.6}{\deci\bel} (\cref{fig:ch7-rlc}). Liu calls the result a ``fatal drawback'': the complex poles sit close to the $j\omega$ axis, so the step response rings as $e^{-\zeta\omega_nt}\cos(\sqrt{1-\zeta^2}\,\omega_nt)$, and the ringing is ISI of its own. Mild inductive peaking with $\zeta\approx0.5$--0.7 is still useful. It extends bandwidth without much ringing, and we return to it in \cref{sec:ctle-hbt}.

\begin{figure}[!tbp]
  \centering
  \includegraphics[width=\linewidth]{ch07_rlc}
  \caption{\textbf{An RLC load buys boost with ringing.} (a)~Load impedance from \cref{eq:ch7-rlc} for several damping factors. $\zeta=0.2$ gives \qty{16.6}{\deci\bel} of peaking at $\omega_n$. (b)~Step responses. At $\zeta=0.2$ the output overshoots to 2.7 times its final value and rings for several periods inside the $e^{-\zeta\omega_nt}$ envelope (dotted).}
  \label{fig:ch7-rlc}
\end{figure}

\section{The source-degenerated differential pair}\label{sec:ctle-degen}

The workhorse CTLE is a differential pair with a parallel $R_S\Vert C_S$ between its sources or emitters (\cref{fig:ch7-degen}). At low frequency $R_S$ degenerates the pair and lowers its gain. At high frequency $C_S$ shorts $R_S$, and the full gain $\gm R_D$ returns. The zero and first pole are set by $R_S$ and $C_S$, so $R_S$ can be a triode MOSFET and $C_S$ a varactor, giving a tunable CTLE~\cite{liuEq,palermo720eq}.

\begin{figure}[!tbp]
\centering
\begin{tikzpicture}[american, scale=0.8, transform shape, font=\small]
  % differential pair
  \draw (0,0) node[nmos](M1){};
  \draw (4,0) node[nmos, xscale=-1](M2){};
  \node[left=1pt, font=\small] at ($(M1.center)+(0.55,0.35)$) {};
  \draw (M1.D) -- (0,1.3) coordinate(o1) to[R, l=$R_D$] (0,3.1);
  \draw (M2.D) -- (4,1.3) coordinate(o2) to[R, l_=$R_D$] (4,3.1);
  \draw (0,3.1) -- (4,3.1); \draw (2,3.1) node[vcc]{$V_\mathrm{DD}$};
  \draw (o1) to[short,*-] (1.1,1.3) to[C, l_=$\CL$] (1.1,0.2) node[ground]{};
  \draw (o2) to[short,*-] (2.9,1.3) to[C, l=$\CL$] (2.9,0.2) node[ground]{};
  \draw (o1) to[short,-o] (-1.2,1.3) node[left]{$v_\mathrm{o}^-$};
  \draw (o2) to[short,-o] (5.2,1.3) node[right]{$v_\mathrm{o}^+$};
  \draw (M1.G) to[short,-o] (-1.2,0) node[left]{$v_\mathrm{in}^+$};
  \draw (M2.G) to[short,-o] (5.2,0) node[right]{$v_\mathrm{in}^-$};
  \draw (M1.S) -- (0,-1.3) coordinate(s1);
  \draw (M2.S) -- (4,-1.3) coordinate(s2);
  \draw (s1) to[short,*-] (0.9,-1.3) to[R, l=$R_S$, *-*] (3.1,-1.3) to[short,-*] (s2);
  \draw (0.9,-1.3) -- (0.9,-2.2) to[C, l_=$C_S$] (3.1,-2.2) -- (3.1,-1.3);
  \draw (s1) to[I, l_=$I/2$] (0,-3.4) node[ground]{};
  \draw (s2) to[I, l=$I/2$] (4,-3.4) node[ground]{};
  % half circuit
  \begin{scope}[xshift=8.6cm]
    \draw (0,0) node[nmos](M){};
    \draw (M.D) -- (0,1.3) coordinate(oh) to[R, l=$R_D$] (0,3.1) node[vcc]{$V_\mathrm{DD}$};
    \draw (oh) to[short,*-] (0.9,1.3) node[circ]{} to[C, l=$\CL$] (0.9,0.1) node[ground]{};
    \draw (0.9,1.3) to[short,-o] (1.7,1.3) node[right]{$v_\mathrm{o}$};
    \draw (M.G) to[short,-o] (-0.9,0) node[left]{$v_\mathrm{in}$};
    \draw (M.S) -- (0,-1.3) coordinate(sh);
    \draw (sh) to[short,*-] (-0.6,-1.3) to[R, l_=$R_S/2$] (-0.6,-3.0) node[ground]{};
    \draw (sh) -- (0.6,-1.3) to[C, l=$2C_S$] (0.6,-3.0) node[ground]{};
    \node[font=\sffamily\small, text=navy] at (0.3,-3.9) {half circuit};
  \end{scope}
\end{tikzpicture}
\caption{\textbf{The degenerated differential pair is the workhorse CTLE.} $R_S\Vert C_S$ between the sources sets the zero and the first pole. $R_D\Vert\CL$ at the drains sets the second pole. The split tail ($2\times I/2$) keeps DC current out of $R_S$. In the half circuit, the degeneration appears as $R_S/2$ in parallel with $2C_S$.}
\label{fig:ch7-degen}
\end{figure}

\begin{derivation}[Transfer function of the degenerated pair]
In the half circuit the source sees $Z_S=\tfrac{R_S}{2}\big\Vert\tfrac{1}{2sC_S}=\dfrac{R_S/2}{1+sR_SC_S}$. Series feedback divides the transconductance by $1+\gm Z_S$:
\[
G_m(s)=\frac{\gm}{1+\gm Z_S}=\frac{\gm\,(1+sR_SC_S)}{1+\gm R_S/2+sR_SC_S}
=\frac{\gm}{1+\gm R_S/2}\cdot\frac{1+sR_SC_S}{1+\dfrac{sR_SC_S}{1+\gm R_S/2}} .
\]
The drain load is $Z_D=R_D/(1+sR_D\CL)$, and $H=G_mZ_D$ (the sign inversion is absorbed by the output polarity).
\end{derivation}

\begin{keyeq}[Degenerated-pair CTLE]
\begin{equation}
  \begin{gathered}
  \Hctle(s)=\frac{\gm R_D}{1+\gm R_S/2}\cdot\frac{1+s/\omega_z}{(1+s/\omega_{p1})(1+s/\omega_{p2})},\\[2pt]
  \omega_z=\frac{1}{R_SC_S},\qquad
  \omega_{p1}=\frac{1+\gm R_S/2}{R_SC_S},\qquad
  \omega_{p2}=\frac{1}{R_D\CL}
  \end{gathered}
  \label{eq:ctle-tf}
\end{equation}
DC gain $=\gm R_D/(1+\gm R_S/2)$. Gain between $\omega_{p1}$ and $\omega_{p2}$ $\approx\gm R_D$. Ideal peaking $=\omega_{p1}/\omega_z=1+\gm R_S/2$.
\end{keyeq}

$R_S$ sets the peaking and the DC gain, and barely moves $\omega_{p1}\approx\gm/(2C_S)$ when $\gm R_S/2\gg1$. $C_S$ moves $\omega_z$ and $\omega_{p1}$ together at constant peaking (\cref{fig:ch7-bode}a). In this book the ``peaking'' of a setting means the \emph{ideal} ratio $\omega_{p1}/\omega_z$, as in \code{ctle()} of \code{link.py}, which fixes $f_{p1}=\fN$, $f_{p2}=\qty{10}{\giga\hertz}$ and 0\,dB high-frequency gain. The boost realized at $\fN$ is smaller: the 12\,dB setting gives \qty{9.0}{\deci\bel}, because the first pole already costs 3\,dB there (\cref{fig:ch7-bode}b).

\begin{figure}[!tbp]
  \centering
  \includegraphics[width=\linewidth]{ch07_bode}
  \caption{\textbf{Zero, poles and peaking of the degenerated pair.} (a)~\Cref{eq:ctle-tf} for the Lab~4 example circuit ($\gm=\qty{20}{\milli\siemens}$, $R_S=\qty{401}{\ohm}$, $C_S=\qty{0.795}{\pico\farad}$, $R_D=\qty{50}{\ohm}$, $\CL=\qty{318}{\femto\farad}$), with its asymptotes. The ideal peaking is 14\,dB, but the boost at $\fN$ is only 10.9\,dB. (b)~The \code{ctle()} family from 0 to 12\,dB, as used in the REF studies. The legend gives the boost realized at $\fN$.}
  \label{fig:ch7-bode}
\end{figure}

\begin{example}{From a peaking target to component values}\label{we:ch7-design}
Lab~4 finds 14\,dB best for REF without a DFE (\cref{sec:ch7-ref}). Target $a=10^{14/20}=5.01$, $f_{p1}=\fN$, $f_{p2}=\qty{10}{\giga\hertz}$ and 0\,dB high-frequency gain, with $\gm=\qty{20}{\milli\siemens}$ per side. Peaking: $R_S=2(a-1)/\gm=\qty{401}{\ohm}$. Zero: $f_z=f_{p1}/a=\qty{0.499}{\giga\hertz}$, so $C_S=1/(2\pi f_zR_S)=\qty{0.795}{\pico\farad}$. Gain: $\gm R_D=1$ gives $R_D=\qty{50}{\ohm}$, and $\CL=1/(2\pi f_{p2}R_D)=\qty{318}{\femto\farad}$. The DC gain is $\gm R_D/a=0.20$ (\qty{-14}{\deci\bel}). This circuit only \emph{loses} gain at low frequency; a real front end adds gain with more $\gm R_D$ or a following stage.
\end{example}

\paragraph{Gain--bandwidth} Multiplying the mid-band gain by the second pole gives
\begin{equation}
  (\gm R_D)\cdot\omega_{p2}=\frac{\gm}{\CL}.
  \label{eq:ch7-gbw}
\end{equation}
This is the gain--bandwidth product of a plain differential pair; degeneration does not change it. A high $\omega_{p2}$ together with a large peak gain needs more $\gm$ (current) or a smaller $\CL$ (smaller next stage, inductive peaking).

\section{The CTLE on the REF channel}\label{sec:ch7-ref}

Lab~4 cascades REF with \code{ctle(pk)} and sweeps the peaking (\cref{fig:ch7-ref}). Without a CTLE the PRBS7 eye is \qty{-84}{\milli\volt}. It opens at about 4\,dB, reaches \qty{104}{\milli\volt} at 12\,dB and peaks at \qty{113.3}{\milli\volt} (0.94\,UI) at 14\,dB. Beyond 16\,dB it closes again. Meanwhile the main cursor $2A\pk{0}$ falls steadily, from \qty{337}{\milli\volt} to \qty{140}{\milli\volt} at 14\,dB. This is the CTLE ``boosting'' by suppressing the low frequencies.

\begin{figure}[!tbp]
  \centering
  \includegraphics[width=\linewidth]{ch07_ref}
  \caption{\textbf{More peaking is better only up to a point.} (a)~REF pulse response alone and after 8, 14 and 20\,dB CTLEs. At 20\,dB the first post-cursors go negative, which is over-equalization. (b)~PRBS7 eye height and main cursor against peaking (Lab~4). The eye is best at 14\,dB (\qty{113}{\milli\volt}), just outside the \OS{} range (shaded). The main cursor shrinks at every step.}
  \label{fig:ch7-ref}
\end{figure}

\begin{insight}
\begin{itemize}
  \item \textbf{Over-peaking returns the ISI with the opposite sign.} Above 16\,dB, $\pk{1}$ and $\pk{2}$ turn negative and the eye closes again.
  \item \textbf{Every dB of peaking costs main cursor.} From 0 to 14\,dB the main cursor drops by 7.6\,dB. Noise that enters after the CTLE, such as sampler and DFE-summer noise, does not drop with it, and the CTLE's high-frequency gain lifts noise and crosstalk relative to the signal.
  \item \textbf{Pre-cursors shrink but do not vanish.} In our model $\pk{-1}$ falls from 0.049 to 0.006 at 14\,dB. A larger pre-cursor, for example from a reflection, is a job for the TX pre-tap.
  \item With a DFE behind it, the CTLE should be tuned for the best eye \emph{after} the DFE. Let the DFE take $\pk{1}$ instead of peaking it away (\cref{sec:eq-partition}).
\end{itemize}
\end{insight}

\begin{tryit}[Lab 4: CTLE]
\code{lab4\_ctle.py} sweeps the peaking from 0 to 20\,dB, prints the eye and the cursors, and converts the best setting into $R_S$, $C_S$, $R_D$ and $\CL$ (\cref{sec:lab4}).
\end{tryit}

\section{An HBT CTLE in SG13G2}\label{sec:ctle-hbt}

\OS{} builds its CTLE from \sgx{} SiGe HBTs. A bipolar transistor gives $\gm=I_C/V_T=\qty{38.7}{\milli\siemens}$ per mA, several times more than a MOSFET at the same current, and its $f_T\approx\qty{300}{\giga\hertz}$ keeps device capacitance small. The circuit is \cref{fig:ch7-degen} with npn devices, emitter degeneration $R_E\Vert C_E$ and collector loads $R_C$. Each $V_{CE}$ must stay below the $\mathrm{BV}_{CEO}$ of about 1.6--1.7\,V on the 2.5\,V analog supply.

\begin{example}{An SG13G2 HBT CTLE with 10\,dB of ideal peaking}\label{we:ch7-hbt}
Bias with $I_\mathrm{tail}=\qty{2}{\milli\ampere}$ (1\,mA per side): $\gm=\qty{1}{\milli\ampere}/\qty{25.85}{\milli\volt}=\qty{38.7}{\milli\siemens}$.
\begin{itemize}
  \item 10\,dB ideal peaking: $1+\gm R_E/2=3.16$, so $R_E=2(2.16)/\gm=\qty{112}{\ohm}$.
  \item Put $f_{p1}$ at $\fN$: $f_z=2.5/3.16=\qty{0.79}{\giga\hertz}$, so $C_E=1/(2\pi f_zR_E)=\qty{1.8}{\pico\farad}$.
  \item $R_C=\qty{200}{\ohm}$ and $\CL=\qty{100}{\femto\farad}$ give $f_{p2}=\qty{7.96}{\giga\hertz}$.
  \item DC gain $=\gm R_C/3.16=2.45$ (\qty{7.8}{\deci\bel}). $\lvert H(j2\pi\fN)\rvert=\qty{14.8}{\deci\bel}$.
\end{itemize}
The \textbf{realized boost} at $\fN$ is $14.8-7.8=\qty{7.0}{\deci\bel}$, not 10\,dB. Each factor of \cref{eq:ctle-tf} evaluated at $\fN$ shows why: the zero contributes $20\log_{10}\lvert1+j\,2.5/0.79\rvert=+10.4$\,dB, the first pole $-3.0$\,dB because it sits at $\fN$, and the second pole $-0.4$\,dB. The response peaks at \qty{4.3}{\giga\hertz}, well above $\fN$.

\emph{Fix.} The boost at $\fN$ is set by $\fN/f_z$, minus whatever the poles take away. Move $f_{p1}$ up to $2\fN$ and lower $f_z$ to \qty{0.70}{\giga\hertz}. That needs $1+\gm R_E/2=7.1$ (17\,dB ideal), so $R_E=\qty{317}{\ohm}$ and $C_E=\qty{0.72}{\pico\farad}$. The realized boost is then \qty{10.0}{\deci\bel}. The price is DC gain, which falls to \qty{0.7}{\deci\bel} (\cref{fig:ch7-we6}). A second option is inductive (shunt) peaking in the collector loads. It pushes $\omega_{p2}$ out and adds a little gain near $\fN$ (\cref{sec:ch7-passive}).
\end{example}

\begin{figure}[!tbp]
  \centering
  \includegraphics[width=\linewidth]{ch07_we6}
  \caption{\textbf{The pole-zero ``ideal'' overstates the boost at $\fN$.} Blue: \cref{we:ch7-hbt} as designed. The asymptotes promise 10\,dB, but only 7.0\,dB is realized at $\fN$ (arrow). Teal: the fixed design, with $f_{p1}=2\fN$ and $f_z=\qty{0.70}{\giga\hertz}$, realizes 10.0\,dB at the cost of DC gain.}
  \label{fig:ch7-we6}
\end{figure}

\paragraph{Programmability} The \OS{} CTLE must cover 0--12\,dB of peaking for channels shorter and longer than REF. \Cref{fig:ch7-hbt} shows the usual way to do it: binary-weighted banks of $R_E$ and $C_E$ segments, switched in parallel by MOS switches under a digital code. Closing more $R_E$ segments lowers $R_{E,\mathrm{eff}}$, which reduces the peaking and raises the DC gain. A fully closed bank (0\,dB) makes it a flat amplifier. $C_E$ segments are switched together with $R_E$, so $\omega_{p1}$ stays near $2\fN$. Keep the switches small. Their on-resistance adds to $R_E$ and their junction capacitance adds to $C_E$, and both shift the corners. Include them in the \tool{ngspice} netlist.

\begin{figure}[!tbp]
\centering
\begin{tikzpicture}[american, scale=0.72, transform shape, font=\small]
  \draw (0,0) node[npn](Q1){};
  \draw (5,0) node[npn, xscale=-1](Q2){};
  \node[font=\small, text=navy] at ($(Q1.center)+(0.75,0)$) {$Q_1$};
  \node[font=\small, text=navy] at ($(Q2.center)+(-0.75,0)$) {$Q_2$};
  \draw (Q1.C) -- (0,1.2) coordinate(o1) to[R, l=$R_C$] (0,3.0);
  \draw (Q2.C) -- (5,1.2) coordinate(o2) to[R, l_=$R_C$] (5,3.0);
  \draw (0,3.0) -- (5,3.0); \draw (2.5,3.0) node[vcc]{$V_\mathrm{CC}=\qty{2.5}{\volt}$};
  \draw (o1) to[short,*-o] (-1.2,1.2) node[left]{$v_\mathrm{o}^-$};
  \draw (o2) to[short,*-o] (6.2,1.2) node[right]{$v_\mathrm{o}^+$};
  \draw (Q1.B) to[short,-o] (-1.2,0) node[left]{$v_\mathrm{in}^+$};
  \draw (Q2.B) to[short,-o] (6.2,0) node[right]{$v_\mathrm{in}^-$};
  \draw (Q1.E) -- (0,-1.2) coordinate(e1);
  \draw (Q2.E) -- (5,-1.2) coordinate(e2);
  % banks
  \draw (e1) to[short,*-*] (1.0,-1.2) to[R, l=$R_{E0}$] (4.0,-1.2) to[short,*-*] (e2);
  \draw (1.0,-1.2) to[short,-*] (1.0,-2.3) to[R, l_=$R_{E1}$] (2.6,-2.3) to[nos, l_=$b_1$] (4.0,-2.3) to[short,-*] (4.0,-1.2);
  \draw (1.0,-2.3) to[short,-*] (1.0,-3.4) to[C, l_=$C_{E0}$] (4.0,-3.4) to[short,-*] (4.0,-2.3);
  \draw (1.0,-3.4) -- (1.0,-4.5) to[C, l_=$C_{E1}$] (2.6,-4.5) to[nos, l_=$b_2$] (4.0,-4.5) -- (4.0,-3.4);
  \draw (e1) to[I, l_=$I/2$] (0,-5.2) node[ground]{};
  \draw (e2) to[I, l=$I/2$] (5,-5.2) node[ground]{};
\end{tikzpicture}
\caption{\textbf{A programmable HBT CTLE.} An \sgx{} npn pair with emitter degeneration banks. Closing switch $b_1$ puts $R_{E1}$ in parallel with $R_{E0}$, lowering the peaking. Switch $b_2$ adds $C_{E1}$ and lowers the corners. More binary-weighted segments (code $b_{3:0}$) cover 0--12\,dB. The split tail keeps DC current out of the banks, so the switches carry no bias current.}
\label{fig:ch7-hbt}
\end{figure}

\paragraph{Noise} The degeneration resistor adds $\sqrt{4kTR_E}=\qty{1.4}{\nano\volt\per\sqrt{\hertz}}$ for $R_E=\qty{112}{\ohm}$, against $\sqrt{2\cdot2kT/\gm}=\qty{0.65}{\nano\volt\per\sqrt{\hertz}}$ of collector shot noise from the pair. Over \qty{5}{\giga\hertz} that is about \qty{0.1}{\milli\volt} rms, far below the \qty{10}{\milli\volt} slicer noise assumed in Lab~5. Offset, crosstalk and linearity limit a 5\,Gb/s CTLE before thermal noise does.

\paragraph{Linearity and overload} With a split tail, the pair fully switches when the differential input exceeds about $I_\mathrm{tail}R_E/2$ plus a few $V_T$. For \cref{we:ch7-hbt} that is $112+\sim100\approx\qty{0.2}{\volt}$. Long runs on REF arrive at nearly the full DC level, $\pm0.93\times\qty{0.5}{\volt}=\pm\qty{0.47}{\volt}$ differential. The CTLE would compress them, and a compressed CTLE breaks the linear superposition that the DFE taps rely on. Three remedies stack: TX de-emphasis lowers the long-run level ($\pm\qty{0.23}{\volt}$ at $-6$\,dB); the fixed design's $I_\mathrm{tail}R_E/2=\qty{317}{\milli\volt}$ widens the linear range; and an input attenuator or gain control handles short channels (\cref{sec:rx-frontend}). At high frequency $C_E$ shorts $R_E$ and the range shrinks toward that of a bare pair, about $\pm2V_T$. That is tolerable, because the channel has already attenuated the high-frequency content by \qty{17}{\deci\bel}.

\begin{specnote}[Open-SERDES CTLE targets (derived here)]
Ideal peaking 0--12\,dB in 1--2\,dB steps ($\approx9$\,dB realized at $\fN$); $f_{p1}\approx1$--$2\,\fN$; $f_{p2}\ge\qty{7.5}{\giga\hertz}$; DC gain $\ge0$\,dB at 12\,dB; LF linear range $\ge\pm\qty{0.25}{\volt}$ differential. Verify the realized boost in \tool{ngspice}.
\end{specnote}

\begin{takeaways}
\begin{itemize}
  \item An RC high-pass has no DC gain and a passive divider has no gain at all. RLC peaking gives boost but rings when $\zeta$ is small (16.6\,dB at $\zeta=0.2$).
  \item The degenerated pair gives $\omega_z=1/R_SC_S$, $\omega_{p1}=(1+\gm R_S/2)/R_SC_S$ and $\omega_{p2}=1/R_D\CL$, with ideal peaking $1+\gm R_S/2$.
  \item The CTLE boosts by suppressing the low frequencies. Its gain--bandwidth is still $\gm/\CL$.
  \item On REF, the eye is best at 14\,dB without a DFE (\qty{113.3}{\milli\volt}). Over-peaking turns the post-cursors negative and costs main cursor.
  \item Realized boost at $\fN$ is well below the ideal when $f_{p1}=\fN$: 7.0\,dB instead of 10\,dB in \cref{we:ch7-hbt}. Move $f_{p1}$ up and $f_z$ down to recover it.
  \item Make the HBT CTLE programmable with switched $R_E$/$C_E$ banks, and check its linear range against the long-run input level.
\end{itemize}
\end{takeaways}

\tikzsetfigurename{08_dfe-}
\chapter[Decision-Feedback Equalization and Adaptation]{Decision-Feedback Equalization\newline and Adaptation}\label{ch:dfe}
\chapterlead{A DFE removes post-cursor ISI by subtracting what the past decisions predict. It equalizes without amplifying noise, and it can learn its own taps.}
\suppressfloats[t]

\begin{objectives}
\begin{itemize}
  \item Derive the DFE from a linear IIR equalizer, and explain why the slicer in the loop removes noise enhancement.
  \item Compute DFE taps from a pulse response and combine a DFE with a TX pre-tap.
  \item Budget the tap-1 timing loop, and choose between direct feedback, half-rate and loop unrolling.
  \item Write the LMS and sign-sign LMS updates, and recognize when adaptation wanders.
  \item Split the equalization work between TX FFE, CTLE and DFE for \OS{}, with numbers.
\end{itemize}
\end{objectives}

\section{The DFE principle}\label{sec:ch8-principle}

\paragraph{From IIR filter to DFE} Start from a one-tap \emph{linear} feedback equalizer, $y(t)=x(t)-\alpha_1y(t-\Tb)$. Its transfer function is $H(z)=1/(1+\alpha_1z^{-1})$, with
\begin{equation}
  \lvert H\rvert=\frac{1}{\sqrt{1+\alpha_1^2+2\alpha_1\cos\omega\Tb}} .
  \label{eq:ch8-iir}
\end{equation}
For $0<\alpha_1<1$ it rises from $1/(1+\alpha_1)$ at DC to $1/(1-\alpha_1)$ at $\fN$, so it could equalize, but it boosts noise like any linear equalizer and is unstable for $\lvert\alpha_1\rvert\ge1$.

The key step is to put a \textbf{slicer} \emph{inside} the loop. The fed-back value is then a clean decision $\hat a_{n-1}\in\{-1,+1\}$, the noise on $y$ is never recirculated, and the DFE adds no noise. With $N$ taps (\cref{fig:ch8-dfe}):

\begin{keyeq}[$N$-tap DFE]
\begin{equation}
  z_n=y_n-\sum_{k=1}^{N}w_k\,\hat a_{n-k},\qquad \hat a_n=\sgn(z_n),\qquad
  w_k=A\,\pk{k}
  \label{eq:ch8-dfe}
\end{equation}
\end{keyeq}

Here $y_n$ is the summer input in volts and the slicer threshold is 0. The DFE coefficients are the post-cursors: the ideal $w_k$ is the post-cursor \emph{at the DFE input} (after TX FFE and CTLE) times $A$.

\begin{figure}[tbp]
\centering
\begin{tikzpicture}[font=\sffamily\small]
  \node[text=navy] (in) at (-0.2,0) {$y_n$};
  \node[sumnode] (s) at (1.4,0) {$\Sigma$};
  \node[blkg, minimum width=15mm, minimum height=8mm] (slc) at (3.6,0) {slicer\\[-2pt]\scriptsize latch};
  \node[dot] (o) at (5.4,0) {};
  \node[blk, minimum width=10mm, minimum height=7mm] (d1) at (6.6,-1.3) {FF};
  \node[blk, minimum width=10mm, minimum height=7mm] (d2) at (8.8,-1.3) {FF};
  \node[dot] (n1) at (7.7,-1.3) {};
  \node[dot] (n2) at (9.9,-1.3) {};
  \node[blkb, minimum width=10mm, minimum height=7mm] (w1) at (7.7,-2.6) {$w_1$};
  \node[blkb, minimum width=10mm, minimum height=7mm] (w2) at (9.9,-2.6) {$w_2$};
  \draw[sig] (in) -- (s);
  \draw[sig] (s) -- node[above, lbl]{$z_n$} (slc);
  \draw[line width=0.75pt, navy] (slc) -- (o);
  \draw[sig] (o) -- (11.0,0) node[right, text=navy]{$\hat a_n$ to DES};
  \draw[sig] (o) |- (d1);
  \draw[line width=0.75pt, navy] (d1) -- (n1);
  \draw[sig] (n1) -- (d2);
  \draw[line width=0.75pt, navy] (d2) -- (n2);
  \node[lbl, above] at (n1) {$\hat a_{n-1}$};
  \node[lbl, above] at (n2) {$\hat a_{n-2}$};
  \draw[sig] (n1) -- (w1);
  \draw[sig] (n2) -- (w2);
  \draw[line width=0.75pt, navy] (w1.south) -- (7.7,-3.4);
  \draw[sig] (w2.south) -- (9.9,-3.4) -- (1.4,-3.4) -- node[pos=0.85, right, text=navy]{$-$} (s.south);
  \node[dot] at (7.7,-3.4) {};
  \draw[clk] (3.6,-1.5) node[below, lbl]{CK} -- (slc.south);
  \draw[clk] (6.6,-2.3) -- (d1.south);
  \draw[clk] (8.8,-2.3) -- (d2.south);
  \draw[sigg, line width=1.2pt, -] ($(slc.east)+(0,0.08)$) -- ($(o)+(0,0.08)$);
  \node[font=\sffamily\scriptsize, text=golddk, align=left, anchor=west] at (2.2,-2.35) {tap-1 loop: slicer $\to$ $w_1$ $\to$ summer\\must settle within 1\,UI};
\end{tikzpicture}
\caption{\textbf{A 2-tap DFE.} The slicer decision $\hat a_n$ is delayed by flip-flops, weighted by $w_1$ and $w_2$, and subtracted from the next samples at the summer. The fed-back values are clean $\pm1$ decisions, so noise is not recirculated. Tap~1 closes a loop that must settle within one UI (\cref{sec:dfe-timing}).}
\label{fig:ch8-dfe}
\end{figure}

\paragraph{Post-cursors only} The DFE cancels only the ISI of bits already decided: the pre-cursor of bit $n+1$ hits $y_n$ before $\hat a_{n+1}$ exists. \Cref{fig:ch8-cancel} shows the feedback staircase on the REF pulse after an 8\,dB CTLE. It matches the pulse at $k=1,2$ and leaves the pre-cursor and the tail untouched.

\begin{figure}[!tbp]
  \centering
  \includegraphics[width=\linewidth]{ch08_cancel}
  \caption{\textbf{A DFE subtracts the post-cursors at the sampling instants.} REF + 8\,dB CTLE (the Lab~5 setup). (a)~Pulse at the DFE input, the 2-tap feedback held for one UI per tap, and the result. Only the sampling instants (dots) matter: $h_1$ and $h_2$ become zero, while the pre-cursor and the tail beyond $k=2$ remain. (b)~Cursors before and after the 2-tap DFE. The main cursor is unchanged, which is the advantage over linear equalization.}
  \label{fig:ch8-cancel}
\end{figure}

\paragraph{Error propagation} Instability is a risk for the linear prototype only. With a slicer in the loop the output is bounded, and the real risk is \textbf{error propagation}: a wrong decision adds $2w_1$ of error to the next sample instead of removing $w_1$. In the Lab~5 setup ($w_1=\qty{27.6}{\milli\volt}$, half-eye \qty{64.7}{\milli\volt} after two taps) that leaves only \qty{9.5}{\milli\volt} of margin on the next bit. At a BER near $10^{-12}$ first errors are so rare that this at most doubles the BER~\cite{palermo720eq}. Keep $\pk{1}/\pk{0}\lesssim0.5$ at the DFE input.

\begin{example}{A DFE alone, and a DFE with a TX pre-tap}\label{we:ch8-pretap}
Take the pulse of \cref{we:ch6-zf}, $x_{-1}=0.2$, $x_0=0.85$, $x_1=-0.2$, $x_2=0.1$, $x_3=0.05$. Forcing the decision sequence $\hat y=\{\dots,0,1,0,0,\dots\}$ at $t=\Tb,2\Tb,3\Tb$ gives $-\alpha_k+x_k=0$, so the taps are the post-cursors: $\alpha=[-0.2,\ 0.1,\ 0.05]$. With $2A=1$ (so $w_k=x_k/2$ in volts), the PDA eye becomes $0.85-0.2=0.65$. Only the pre-cursor remains.

Now add a TX pre-tap with $c_{-1}/c_0=-0.2/0.85$, normalized to $c=[-0.190,\,0.810]$. By \cref{eq:ch6-conv} the cursors for $k=-2..3$ become $-0.038$, 0, 0.726, $-0.181$, 0.071 and 0.041. The DFE re-adapts to the last three, and the eye is $0.726-0.038=\mathbf{0.688}$. That beats both the 3-tap ZF FFE alone (0.519, \cref{we:ch6-zf}) and the DFE alone (0.65). This is the canonical split: the TX takes the pre-cursor and the DFE takes the post-cursors.
\end{example}

\section{The tap-1 timing loop}\label{sec:dfe-timing}

Decision $\hat a_{n-1}$ must be resolved, weighted and subtracted, and the summer settled, before $z_n$ is sampled one UI later:

\begin{keyeq}[DFE tap-1 critical path]
\begin{equation}
  t_\mathrm{CK\to Q}+t_\mathrm{tap}+t_\mathrm{settle}+t_\mathrm{setup}\le\Tb=\qty{200}{\pico\second}
  \label{eq:ch8-timing}
\end{equation}
\end{keyeq}

Taps 2 and higher have one or more extra UIs, so tap~1 sets the speed limit.

\paragraph{Half rate} \OS{} slices even bits on the rising and odd bits on the falling edge of the \qty{2.5}{\giga\hertz} clock, so each latch has 2\,UI to regenerate. Half rate does \emph{not} relax \cref{eq:ch8-timing}: the even decision must correct the odd sample one UI later, so the tap-1 path crosses between the halves and still has 1\,UI.

\paragraph{Loop unrolling} \textbf{Loop unrolling} (\textbf{speculation}) removes the summer from the tap-1 path. Two comparators per half precompute both answers, against $+w_1$ and $-w_1$, and a multiplexer selected by the previous decision picks one~\cite{palermo720eq}:
\begin{equation}
  \hat a_n=\begin{cases}\sgn(y_n-w_1), & \hat a_{n-1}=+1\\ \sgn(y_n+w_1), & \hat a_{n-1}=-1\end{cases}
  \qquad\Rightarrow\qquad t_\mathrm{CK\to Q}+t_\mathrm{mux}+t_\mathrm{setup}\le\Tb .
  \label{eq:ch8-unroll}
\end{equation}
\Cref{fig:ch8-unroll} shows the half-rate version. The cost is $2^N$ comparators for $N$ unrolled taps, twice the load on the CTLE, and taps that live in comparator thresholds.

\begin{figure}[!tbp]
\centering
\begin{tikzpicture}[font=\sffamily\small]
  \node[text=navy] (in) at (-0.4,0) {$y$};
  \draw[line width=0.75pt, navy] (in) -- (0.6,0);
  \draw[line width=0.75pt, navy] (0.6,-2.2) -- (0.6,2.2);
  \node[dot] at (0.6,0) {};
  \foreach \y/\n/\t in {2.2/ep/{+w_1}, 1.2/em/{-w_1}, -1.2/op/{+w_1}, -2.2/om/{-w_1}}{
    \node[blkg, minimum width=17mm, minimum height=7mm] (\n) at (2.4,\y) {$y\gtrless\t$};
    \draw[sig] (0.6,\y) -- (\n.west);
  }
  \node[lbl, anchor=east] at (0.5,1.7) {CK};
  \node[lbl, anchor=east] at (0.5,-1.7) {$\overline{\mathrm{CK}}$};
  \node[lblb, anchor=south west] at (1.5,2.55) {even half};
  \node[lblb, anchor=north west] at (1.5,-2.55) {odd half};
  \node[blkb, minimum width=11mm, minimum height=14mm] (me) at (5.0,1.7) {MUX};
  \node[blkb, minimum width=11mm, minimum height=14mm] (mo) at (5.0,-1.7) {MUX};
  \draw[sig] (ep.east) -- ($(me.west)+(0,0.5)$);
  \draw[sig] (em.east) -- ($(me.west)+(0,-0.5)$);
  \draw[sig] (op.east) -- ($(mo.west)+(0,0.5)$);
  \draw[sig] (om.east) -- ($(mo.west)+(0,-0.5)$);
  \node[blkg, minimum width=12mm, minimum height=7mm] (le) at (7.4,1.7) {latch};
  \node[blkg, minimum width=12mm, minimum height=7mm] (lo) at (7.4,-1.7) {latch};
  \draw[sig] (me) -- (le);
  \draw[sig] (mo) -- (lo);
  \draw[sig] (le.east) -- (10.3,1.7) node[right, text=navy]{$\hat a_\mathrm{even}$};
  \draw[sig] (lo.east) -- (10.3,-1.7) node[right, text=navy]{$\hat a_\mathrm{odd}$};
  % cross-coupled selects (critical path)
  \draw[sigg, line width=1.1pt] (8.6,1.7) node[dot, fill=golddk]{} -- (8.6,-0.4) -- (5.0,-0.4) -- (mo.north);
  \draw[sigg, line width=1.1pt] (9.4,-1.7) node[dot, fill=golddk]{} -- (9.4,0.4) -- (5.0,0.4) -- (me.south);
  \node[font=\sffamily\scriptsize, text=golddk, align=left, anchor=west] at (9.6,0.0) {select from the other half:\\$t_\mathrm{CK\to Q}+t_\mathrm{mux}+t_\mathrm{setup}<1\,\UI$};
\end{tikzpicture}
\caption{\textbf{Half-rate DFE with tap~1 unrolled.} Each half has two comparators with thresholds $\pm w_1$ and a multiplexer. The multiplexer is selected by the latest decision of the \emph{other} half (gold paths). These cross-coupled selects are the critical path and still have only 1\,UI, but there is no summer to settle. Tap~2 and higher would still be subtracted at a summer before the comparators.}
\label{fig:ch8-unroll}
\end{figure}

\begin{example}{Tap-1 timing budget at 5\,Gb/s}\label{we:ch8-timing}
Use illustrative \sgx{} CML numbers, to be confirmed post-layout: latch $t_\mathrm{CK\to Q}=\qtyrange{25}{35}{\pico\second}$; tap switch about \qty{10}{\pico\second}; summer with $R_L=\qty{250}{\ohm}$ and $C=\qty{80}{\femto\farad}$, so $\tau=\qty{20}{\pico\second}$ and 95\,\% settling takes $3\tau=\qty{60}{\pico\second}$; setup and aperture about \qty{20}{\pico\second}; skew and PVT margin about \qty{20}{\pico\second}. The total is about \qtyrange{135}{145}{\pico\second}, which is below \qty{200}{\pico\second}. Direct feedback works at 5\,Gb/s, including in the half-rate cross-coupled form, with 55--65\,ps to spare (\cref{fig:ch8-timing}).

It does \emph{not} meet the \qty{100}{\pico\second} half-UI budget that would apply if the $90^\circ$/$270^\circ$ edge samplers also had to see a fully corrected signal. In \OS{} the edge samplers therefore see a partly settled correction, which works because TX and CTLE keep the pre-DFE eye open (\cref{sec:eq-partition}). If post-layout corners push the direct path past about 0.7\,UI, unroll tap~1: with an assumed \qty{15}{\pico\second} mux the unrolled path is \qty{85}{\pico\second}. Halving $C$ (fewer taps) or an integrating summer saves about \qty{30}{\pico\second}.
\end{example}

\begin{figure}[!tbp]
\centering
\begin{tikzpicture}[x=0.052cm, y=1cm, font=\sffamily\scriptsize]
  \foreach \x in {0,50,100,150,200}{
    \draw[gridc] (\x,-0.25) -- (\x,1.75);
    \node[text=muted, below] at (\x,-0.25) {\x};
  }
  \node[text=navy, font=\sffamily\small] at (110,-0.95) {Time after the clock edge (ps)};
  % direct feedback bar (y = 1.1)
  \def\yy{1.15}
  \foreach \a/\b/\c/\t in {0/30/c1/{$t_\mathrm{CK\to Q}$}, 30/40/c5/{}, 40/100/c2/{settle $3\tau$}, 100/120/c3/{setup}, 120/140/pale/{margin}}{
    \fill[\c] (\a,\yy-0.22) rectangle (\b,\yy+0.22);
    \node[text=white, font=\sffamily\scriptsize] at ({(\a+\b)/2},\yy) {\t};
  }
  \node[text=c5, font=\sffamily\scriptsize, above] at (35,\yy+0.22) {tap};
  \node[text=navy, anchor=east, font=\sffamily\small] at (-4,\yy) {direct};
  \node[text=navy, anchor=west] at (143,\yy) {140 ps};
  % unrolled bar (y = 0.35)
  \def\yy{0.35}
  \foreach \a/\b/\c/\t in {0/30/c1/{$t_\mathrm{CK\to Q}$}, 30/45/c5/{mux}, 45/65/c3/{setup}, 65/85/pale/{margin}}{
    \fill[\c] (\a,\yy-0.22) rectangle (\b,\yy+0.22);
    \node[text=white, font=\sffamily\scriptsize] at ({(\a+\b)/2},\yy) {\t};
  }
  \node[text=navy, anchor=east, font=\sffamily\small] at (-4,\yy) {unrolled};
  \node[text=navy, anchor=west] at (102,\yy) {85 ps};
  \draw[c4, dashed, line width=0.8pt] (100,-0.25) -- (100,1.85) node[above, text=c4, align=center] {$\tfrac12$\,UI: edge samplers};
  \draw[c4, line width=0.9pt] (200,-0.25) -- (200,1.85) node[above, text=c4] {1\,UI deadline};
\end{tikzpicture}
\caption{\textbf{The tap-1 budget fits in one UI with direct feedback.} Numbers from \cref{we:ch8-timing}. Direct feedback (\qty{140}{\pico\second}) meets the 1\,UI deadline but not the half-UI one. Unrolling removes the summer settling and fits both.}
\label{fig:ch8-timing}
\end{figure}

\section{Adapting the taps}\label{sec:adaptation}

The ideal taps depend on channel, temperature and supply, so the receiver should find them itself. Define the slicer error $e_n=z_n-d_\mathrm{lev}\,\hat a_n$, where $d_\mathrm{lev}$ is the expected signal level (ideally $A\pk{0}$). Since $\partial E[e_n^2]/\partial w_k=-2E[e_n\hat a_{n-k}]$, stepping against the gradient gives the \textbf{LMS} update. Hardware prefers the \textbf{sign-sign} form (SS-LMS), which needs only signs~\cite{palermo720eq,stojanovic2005}:

\begin{keyeq}[LMS and sign-sign LMS]
\begin{equation}
  \begin{gathered}
  \text{LMS:}\quad w_k\leftarrow w_k+\mu\,e_n\,\hat a_{n-k}\\[2pt]
  \text{SS-LMS:}\quad w_k\leftarrow w_k+\mu\,\sgn(e_n)\,\hat a_{n-k},\qquad
  d_\mathrm{lev}\leftarrow d_\mathrm{lev}+\mu\,\sgn(e_n)\,\hat a_n
  \end{gathered}
  \label{eq:ch8-sslms}
\end{equation}
\end{keyeq}

$\sgn(e_n)$ comes from one extra comparator, the \textbf{error sampler}, with threshold $\pm d_\mathrm{lev}$ chosen by $\hat a_n$. The updates are digital counters decimated to the deserialized rate: slow compared with the data, fast compared with temperature drift.

\begin{example}{Sign-sign LMS on the REF channel}\label{we:ch8-sslms}
Simulate the Lab~5 setup (REF + 8\,dB CTLE) with PRBS13, \qty{4}{\milli\volt} rms noise and a 2-tap DFE adapted by \cref{eq:ch8-sslms} with $\mu=\qty{0.25}{\milli\volt}$, starting from $w=0$ and $d_\mathrm{lev}=\qty{50}{\milli\volt}$. The pre-DFE eye is open (\qty{60}{\milli\volt}), so decisions are correct from the first bit. $w_1$ comes within \qty{2}{\milli\volt} of $A\pk{1}=\qty{27.6}{\milli\volt}$ after about 310\,UI and $w_2$ reaches $A\pk{2}=\qty{11.8}{\milli\volt}$ after about 170\,UI (\cref{fig:ch8-sslms}a). Both then dither by a few $\mu$; after 6000\,UI they read 29.8 and \qty{11.5}{\milli\volt}, with $d_\mathrm{lev}=\qty{94.0}{\milli\volt}$ (target 91.6). A larger $\mu$ converges faster but dithers more, and dither is a tap error that costs eye.
\end{example}

\begin{figure}[!tbp]
  \centering
  \includegraphics[width=\linewidth]{ch08_sslms}
  \caption{\textbf{Sign-sign LMS finds the taps, unless a pre-cursor is left in.} (a)~REF + 8\,dB CTLE, 2-tap DFE ($\mu=\qty{0.25}{\milli\volt}$, \qty{4}{\milli\volt} noise). The taps and the data level settle within a few hundred UI onto their zero-forcing values (dashed: 91.6, 27.6 and \qty{11.8}{\milli\volt}). (b)~The pulse of \cref{we:ch8-pretap} with a 3-tap DFE ($\mu=\qty{2}{\milli\volt}$, \qty{10}{\milli\volt} noise). With its 0.2 pre-cursor removed, $w_2$ stays within 36--\qty{62}{\milli\volt} of its \qty{50}{\milli\volt} target. With the pre-cursor left in, it wanders between $-6$ and \qty{92}{\milli\volt}.}
  \label{fig:ch8-sslms}
\end{figure}

\begin{pitfall}[An uncancelled pre-cursor makes SS-LMS wander]
The 0.2 pre-cursor puts $\pm Ax_{-1}=\pm\qty{100}{\milli\volt}$ into every error sample. Once the tap errors are smaller than that, $\sgn(e_n)$ is set by $Ax_{-1}\hat a_{n+1}$ alone, which is uncorrelated with $\hat a_{n-k}$: the average update vanishes and the taps random-walk (\cref{fig:ch8-sslms}b). Cancel the pre-cursor first (TX pre-tap or CTLE), or use a smaller $\mu$ and average.
\end{pitfall}

\section{The DFE on the REF channel}\label{sec:ch8-ref}

Lab~5 adds DFE taps one at a time after an 8\,dB CTLE (\cref{fig:ch8-taps}). The taps are the post-cursors at the DFE input, 0.0552, 0.0235 and 0.0160, so $w_k=27.6$, 11.8 and \qty{8.0}{\milli\volt}. The first tap lifts the PRBS7 eye from 60.0 to \qty{110.7}{\milli\volt}, the second adds 19 and the third \qty{13.5}{\milli\volt}. Beyond four taps nothing more is gained (153.6, 152.1 and \qty{156.4}{\milli\volt} for 4, 5 and 8 taps). At \qty{10}{\milli\volt} rms slicer noise the estimated BER is $3.2\times10^{-14}$ with three taps and $4.5\times10^{-12}$ with two.

\begin{figure}[!tbp]
  \centering
  \includegraphics[width=\linewidth]{ch08_taps}
  \caption{\textbf{The first DFE taps do most of the work.} REF + 8\,dB CTLE, PRBS7 (Lab~5). (a)~Eye height against the number of ideal DFE taps. The first two taps add 70\,mV, and beyond four taps the gain is lost in the simulation's resolution. (b,\,c)~Summer output before and after a 2-tap DFE. The feedback switches half a UI after each decision, so the eye shows steps at $\pm0.5$\,UI. Only the sampling instant at 0\,UI matters.}
  \label{fig:ch8-taps}
\end{figure}

\begin{tryit}[Lab 5: DFE]
\code{lab5\_dfe.py} reads the taps off the pulse response, runs \code{sdp.Receiver.nrz\_DFE} with 0--8 taps and estimates the BER. Then try a $-6$\,dB TX preset with 2 taps (\cref{sec:lab5}).
\end{tryit}

\section{Splitting the work}\label{sec:eq-partition}

How should \OS{} divide the work between TX FFE, CTLE and DFE (\cref{sec:ch5-preview})? \Cref{tab:eq-partition} and \cref{fig:ch8-partition} compare seven strategies on the \code{link.py} channel family.

\paragraph{Method} Channels are \code{channel(LOSS\_FAMILY[d])} for 5--30\,dB at $\fN$, plus REF. The CTLE is \code{ctle(pk)}, 0--20\,dB in 1\,dB steps. The TX FFE is a grid search over $c_{-1}\in[-0.2,0]$ and $c_1\in[-0.5,0]$ (step 0.02) with $c_0=1-\lvert c_{-1}\rvert-\lvert c_1\rvert$. The DFE is ideal: it removes the first $N$ post-cursors exactly. Each strategy is scored by its PDA eye over $k=-4..40$, maximized over sampling phase and free settings, and also normalized to $2A\pk{0}$ of the \emph{unequalized} pulse, which exposes the swing cost of TX equalization. The ``\OS{} limits'' rows add $\lvert c_1\rvert\le0.25$ (6\,dB post-tap de-emphasis), $\lvert c_{-1}\rvert\le0.1$, CTLE $\le12$\,dB, and an eye open \emph{before} the DFE (for start-up adaptation and the edge samplers). Because the tail is summed to 40\,UI, the unequalized REF eye is $-0.23$\,V rather than the 13-cursor $-0.161$\,V of \cref{tab:ref-cursors}.

\begin{table}[!tbp]
\centering\footnotesize
\caption{\textbf{PDA eye height (mV) and eye normalized to the unequalized main cursor, by strategy and channel loss.} Regenerated with \code{link.py}; method in the text. $^\ast$\OS{} limits: $\lvert c_1\rvert\le0.25$, $\lvert c_{-1}\rvert\le0.1$, CTLE $\le12$\,dB, open pre-DFE eye.}
\label{tab:eq-partition}
\setlength{\tabcolsep}{4.2pt}
\begin{tabular}{@{}l rr rr rr rr rr@{}}
\toprule
 & \multicolumn{2}{c}{\textbf{10\,dB}} & \multicolumn{2}{c}{\textbf{15\,dB}} & \multicolumn{2}{c}{\textbf{REF 17.4\,dB}} & \multicolumn{2}{c}{\textbf{20\,dB}} & \multicolumn{2}{c}{\textbf{25\,dB}}\\
\cmidrule(lr){2-3}\cmidrule(lr){4-5}\cmidrule(lr){6-7}\cmidrule(lr){8-9}\cmidrule(l){10-11}
\textbf{Strategy} & mV & norm & mV & norm & mV & norm & mV & norm & mV & norm\\
\midrule
No equalization            & 165 & 0.30 & \multicolumn{2}{c}{\textcolor{c4}{closed}} & \multicolumn{2}{c}{\textcolor{c4}{closed}} & \multicolumn{2}{c}{\textcolor{c4}{closed}} & \multicolumn{2}{c}{\textcolor{c4}{closed}}\\
TX 2-tap FFE               & 341 & 0.61 & 189 & 0.48 & 141 & 0.42 & 100 & 0.34 & 44 & 0.20\\
CTLE (0--20\,dB)           & 236 & 0.42 & 132 & 0.33 & 99 & 0.29 & 70 & 0.24 & 32 & 0.15\\
CTLE + 2-tap DFE           & 395 & 0.71 & 146 & 0.37 & 110 & 0.32 & 84 & 0.29 & 51 & 0.23\\
TX 3-tap + CTLE + 2-tap DFE & 396 & 0.71 & 204 & 0.52 & 160 & 0.47 & 123 & 0.42 & 75 & 0.34\\
\midrule
\OS{} limits, 2-tap DFE$^\ast$ & 396 & 0.71 & 190 & 0.48 & 134 & 0.40 & 87 & 0.30 & 35 & 0.16\\
\OS{} limits, 5-tap DFE$^\ast$ & 467 & 0.84 & 239 & 0.60 & 170 & 0.50 & 120 & 0.41 & 43 & 0.19\\
\bottomrule
\end{tabular}
\end{table}


\paragraph{What the numbers say}
\begin{itemize}
  \item \textbf{No single equalizer is enough at REF.} The TX FFE alone reaches \qty{141}{\milli\volt} but needs 12\,dB of de-emphasis; the CTLE alone needs 16\,dB of peaking, outside our range, for \qty{99}{\milli\volt}.
  \item \textbf{With a DFE present, the noise-free optimum moves the linear work into the TX} (CTLE at 0\,dB, 11--16\,dB of de-emphasis). In a PDA our normalized CTLE and the TX FFE pay the same price, a lower low-frequency level, and the 1-UI FIR fits a sampled channel better. The PDA cannot see noise, crosstalk, sampler sensitivity or adaptivity, which are what limit large de-emphasis in practice.
  \item \textbf{Within practical limits} the REF eye is \qty{134}{\milli\volt} (0.40 of the unequalized main cursor). Three more DFE taps add \qty{36}{\milli\volt}.
  \item \textbf{Above about 22\,dB the 12\,dB CTLE limit binds.} At 25\,dB, CTLE + 2-tap DFE with 19\,dB of peaking (\qty{51}{\milli\volt}) beats both \OS{} rows; more DFE taps help less than more CTLE.
\end{itemize}

\begin{figure}[!tbp]
  \centering
  \includegraphics[width=\linewidth]{ch08_partition}
  \caption{\textbf{Every strategy loses about a factor of two in eye per 5\,dB of loss; combinations lose less.} PDA eye against loss at $\fN$ (log scale, data of \cref{tab:eq-partition}). At REF (dotted) the \OS{} partition gives \qty{134}{\milli\volt}. Above about 22\,dB, a CTLE with more than 12\,dB of peaking (purple) overtakes it, which signals when \OS{} would need more equalization.}
  \label{fig:ch8-partition}
\end{figure}

\begin{example}{The REF operating point in the time domain}\label{we:ch8-refop}
Check the partition with the Lab~5 machinery: PRBS7 through REF, a TX preset, \code{ctle(pk)} and an ideal 2-tap DFE. The BER is estimated at \qty{10}{\milli\volt} rms slicer noise.
\begin{center}\small
\begin{tabular}{@{}lccccc@{}}
\toprule
TX preset, CTLE & pre-DFE eye & + 2-tap DFE & $\pk{1}/\pk{0}$ & $w_1$, $w_2$ (mV) & BER \\
\midrule
none, 8\,dB (Lab 5) & 60.0 & 129.5 & 0.30 & 27.6, 11.8 & $4.5\times10^{-12}$\\
none, 0\,dB         & $-84.0$ & 167.1 & 0.54 & 91.2, 45.6 & $1.3\times10^{-15}$\\
$-3.5$\,dB, 0\,dB   & 35.4 & 169.5 & 0.37 & 51.1, 23.5 & $3.5\times10^{-17}$\\
$-6$\,dB, 0\,dB     & 96.9 & \textbf{171.8} & 0.26 & 31.1, 11.7 & $3.0\times10^{-18}$\\
$-6$\,dB, 4\,dB     & 103.0 & 132.6 & 0.15 & 12.7, 5.9 & $1.4\times10^{-12}$\\
\bottomrule
\end{tabular}
\end{center}
A DFE alone (row 2) opens the eye, but with a closed pre-DFE eye it cannot adapt from start-up, and $\pk{1}/\pk{0}=0.54$ invites error propagation. The $-6$\,dB preset with the CTLE near 0\,dB gives the best eye, an open pre-DFE eye (\qty{97}{\milli\volt}) for the edge samplers, and $\pk{1}/\pk{0}=0.26$: a wrong decision costs $2w_1=\qty{62}{\milli\volt}$ of the \qty{86}{\milli\volt} half-eye.
\end{example}

\paragraph{The \OS{} partition, and why}
\begin{itemize}
  \item \textbf{TX: 3-tap FFE, mostly de-emphasis} ($-3.5$/$-6$\,dB presets, pre-tap up to $\lvert c_{-1}\rvert=0.1$; REF uses $-6$\,dB). It is the only block that removes the pre-cursor without noise, and it lowers the long-run level that would overload the CTLE (\cref{sec:ctle-hbt}). It stays moderate because it cannot adapt without a back-channel and because lost launch swing is lost SNR.
  \item \textbf{RX: CTLE, 0--12\,dB.} On REF it sits near 0--4\,dB and acts mainly as the front-end gain stage. Its range covers longer channels and board-to-board and temperature variation, keeps the pre-DFE eye open for the CDR, and, unlike the TX presets, adapts at the receiver.
  \item \textbf{RX: 2-tap half-rate DFE, direct feedback.} It removes $\pk{1}$ and $\pk{2}$ without noise, with 55--65\,ps of tap-1 slack (\cref{we:ch8-timing}). Each extra tap adds summer capacitance and slows tap~1 for about \qty{10}{\milli\volt} more eye.
\end{itemize}
\textbf{When we would need more.} Beyond about 22\,dB of loss, add a second CTLE stage ($>12$\,dB), 3--5 DFE taps or an IIR tail tap, and a larger pre-tap. Unroll tap~1 if post-layout timing exceeds about 0.7\,UI.

\begin{insight}
The partition follows the pulse anatomy of \cref{fig:ch5-anatomy}: the TX pre-tap takes the pre-cursor, de-emphasis and the CTLE take the long tail and keep the pre-DFE eye open, and the DFE takes $\pk{1}$ and $\pk{2}$ without noise. A noise-free PDA would push all linear work into the TX; noise, crosstalk and adaptation pull some of it back into the receiver. Always report the result in millivolts against the noise budget of \cref{sec:noisebudget}.
\end{insight}

\section*{Exercises}\addcontentsline{toc}{section}{Exercises}
\begin{enumerate}[label=\textbf{8.\arabic*}, leftmargin=2.2em]
  \item The DFE of \cref{we:ch8-pretap} with one tap only: what is the PDA eye? If $\hat a_{n-1}$ is wrong, how much margin is left on bit $n$?
  \item Redo \cref{we:ch8-timing} with \qty{120}{\femto\farad} of summer capacitance (more taps). Does direct feedback still fit? What $R_L$ restores it, and at what cost in summer gain?
  \item Show that the SS-LMS update $E[\sgn(e_n)\hat a_{n-1}]$ is zero at $w_1=A\pk{1}$ if the remaining ISI plus noise is symmetric. Why does this not stop the wandering in \cref{fig:ch8-sslms}(b)?
  \item From the cursors in \cref{sec:ch8-ref} ($\pk{-1}=0.0138$, $\pk{0}=0.1832$, tail to $\pk{8}$), compute the PDA eye after 0, 1 and 2 DFE taps. Why is it lower than the PRBS7 numbers of \cref{fig:ch8-taps}? Then use \cref{tab:eq-partition}: at 25\,dB, which buys more, three extra taps or a 19\,dB CTLE?
\end{enumerate}

\begin{takeaways}
\begin{itemize}
  \item A DFE is an IIR equalizer with a slicer in the loop: it cancels post-cursors without amplifying noise, with taps equal to the post-cursors at its input. It cannot touch pre-cursors, so a TX pre-tap plus a DFE beats either alone (0.688 vs 0.519 and 0.65 in \cref{we:ch8-pretap}).
  \item Tap~1 must settle within 1\,UI even at half rate. At 5\,Gb/s direct feedback fits in about 140\,ps. Unrolling is the fallback if it does not.
  \item SS-LMS adapts taps with one error sampler and counters. It wanders if a large pre-cursor is left uncancelled.
  \item On REF + 8\,dB CTLE, DFE taps give $60\to111\to130\to143$\,mV for 0--3 taps. Beyond four taps the gain is negligible.
  \item The \OS{} partition is a TX 3-tap FFE ($-6$\,dB on REF), a CTLE with 0--12\,dB and a 2-tap half-rate DFE. On REF it gives \qty{172}{\milli\volt} (PRBS7) and \qty{134}{\milli\volt} (PDA).
  \item Beyond about 22\,dB of loss, more CTLE peaking and more DFE taps become necessary.
\end{itemize}
\end{takeaways}


\tutpart{Transceiver Circuits}{Turning the system model into transistors: the serializer and output driver on the transmit side, and the termination, front-end and samplers on the receive side, all with SG13G2 HBTs and CMOS.}
\tikzsetfigurename{09_tx-}
\chapter{The Transmitter}\label{ch:tx}
\chapterlead{Sixteen slow bits go in, one fast differential wave comes out: the serializer, the termination, the CML driver with its FFE, and the pad that connects them to the channel.}

\begin{objectives}
\begin{itemize}
  \item Explain with the FO4 rule why the last serializer stages and their clocks must be half-rate CML in 130\,nm.
  \item Design an HBT CML 2:1 multiplexer and read its timing, including duty-cycle errors.
  \item Compare series, parallel and double termination, and compute what double termination costs.
  \item Size the \OS{} CML driver (current, swing, power, headroom, $\mathrm{BV}_\mathrm{CEO}$) and compare it with a voltage-mode driver.
  \item Build the 3-tap FFE as a current-summing driver and choose its tap resolution.
  \item Budget the pad, ESD and driver capacitance, and decide when a T-coil is worth its area.
\end{itemize}
\end{objectives}

%% =====================================================================
\section{The transmitter at a glance}\label{sec:ch9-overview}

The transmitter takes a 16-bit word from the digital core every \qty{3.2}{\nano\second} (\qty{312.5}{\mega\hertz}) and launches it as a \qty{5}{\giga\bit\per\second} NRZ wave into the \qty{50}{\ohm}-per-side channel (\cref{fig:ch9-arch}). Each 2:1 stage doubles the lane rate and is clocked at \emph{half} its output bit rate. The clocks come from the \qty{2.5}{\giga\hertz} PLL (\cref{ch:pll}) through a chain of $\div2$ dividers. Static CMOS on \qty{1.2}{\volt} handles the slow, wide part. HBT current-mode logic (CML) on the \qty{2.5}{\volt} analog supply handles everything at \qty{1.25}{\giga\hertz} and above. The final 2:1 multiplexer is built three times, so that the driver receives $D_{n+1}$, $D_n$ and $D_{n-1}$ for its three FFE taps.

\begin{figure}[tbp]
\centering
\begin{tikzpicture}[x=0.98cm,y=1cm,font=\sffamily\footnotesize,
  b/.style={blk,minimum width=11mm,minimum height=10mm,font=\sffamily\footnotesize,inner sep=2pt},
  bb/.style={b,fill=btint}, bg/.style={b,fill=gtint,draw=golddk}, bn/.style={b,fill=navy,text=white},
  dv/.style={blk,minimum width=8mm,minimum height=6mm,inner sep=1pt,font=\sffamily\footnotesize},
  rate/.style={font=\sffamily\scriptsize,text=golddk,align=center},
  bus/.style={font=\sffamily\scriptsize,text=muted}]
  \draw[region,fill=gtint!35] (-0.35,-3.05) rectangle (4.75,1.3);
  \draw[region,fill=bgpale!80] (4.9,-3.05) rectangle (15.95,1.3);
  \node[lblb,anchor=north west] at (-0.3,1.25) {CMOS, 1.2\,V};
  \node[lblb,anchor=north west] at (4.95,1.25) {HBT CML, 2.5\,V analog supply};
  \node[bg] (s1) at (1.35,0) {16:8};
  \node[bg] (s2) at (3.45,0) {8:4};
  \node[bb] (cv) at (5.65,0) {CMOS\\$\to$CML};
  \node[bb] (s3) at (7.75,0) {4:2};
  \node[bb] (s4) at (9.85,0) {2:1\\$\times$3};
  \node[bb] (pd) at (11.95,0) {pre-\\drivers};
  \node[bn,minimum width=14mm] (dr) at (14.0,0) {3-tap\\driver};
  \node[draw=navy,fill=navy,minimum size=2mm,inner sep=0pt] (pad) at (15.2,0) {};
  \draw[sig,line width=1.4pt] (-0.25,0) -- (s1.west);
  \foreach \a/\b/\n/\r in {s1/s2/8/{625\\Mb/s}, s2/cv/4/{1.25\\Gb/s}, cv/s3/4/{1.25\\Gb/s}, s3/s4/2/{2.5\\Gb/s}, s4/pd/3/{5\\Gb/s}, pd/dr/3/{5\\Gb/s}}{
    \draw[sig] (\a.east) -- (\b.west);
    \path (\a.east) -- (\b.west) coordinate[pos=0.5] (m);
    \draw[navy] ($(m)+(-0.08,-0.12)$) -- ($(m)+(0.08,0.12)$);
    \node[bus,anchor=south] at ($(m)+(0,0.12)$) {\n};
    \node[rate,anchor=north] at ($(m)+(0,-0.12)$) {\r};
  }
  \draw[navy] (0.18,-0.12) -- (0.34,0.12); \node[bus,anchor=south] at (0.26,0.12) {16};
  \node[rate,anchor=north] at (0.26,-0.12) {312.5\\Mb/s};
  \draw[navy,line width=0.75pt] (dr.east) -- (pad.west);
  \draw[sig] (pad.east) -- (15.85,0);
  \node[rate,anchor=north] at (15.45,-0.15) {5\,Gb/s\\1\,Vppd};
  \node[bn,minimum width=14mm] (pll) at (14.0,-2.2) {PLL\\2.5\,GHz};
  \node[dv] (d1) at (8.8,-2.2) {$\div2$};
  \node[dv] (d2) at (4.25,-2.2) {$\div2$};
  \node[dv] (d3) at (2.4,-2.2) {$\div2$};
  \draw[clk] (pll.west) -- (d1.east);
  \draw[clk] (d1.west) -- (d2.east);
  \draw[clk] (d2.west) -- (d3.east);
  \draw[clk] (d3.west) -- (-0.25,-2.2);
  \foreach \x/\t in {9.85/s4, 7.75/s3, 5.65/cv, 3.45/s2, 1.35/s1}{
    \node[dot,fill=gold] at (\x,-2.2) {};
    \draw[clk] (\x,-2.2) -- (\t.south);
  }
  \node[rate,anchor=north] at (11.9,-2.3) {2.5\,GHz};
  \node[rate,anchor=north] at (6.7,-2.3) {1.25\,GHz};
  \node[rate,anchor=north] at (3.25,-2.55) {625\,MHz};
  \node[rate,anchor=north,align=left] at (0.75,-2.3) {312.5\,MHz\\to core};
\end{tikzpicture}
\caption{\textbf{The \OS{} transmitter is a tree of 2:1 stages, each clocked at half its output rate.} Slash numbers give the number of parallel lanes; gold text gives the rate of each lane or clock. CMOS does 16:4. HBT CML does the 4:2 and the final half-rate 2:1, which is built three times to feed the three FFE taps.}
\label{fig:ch9-arch}
\end{figure}

%% =====================================================================
\section{The serializer}\label{sec:serializer}

\subsection{Clocks, not gates, set the speed}

A 130\,nm CMOS gate switches in tens of picoseconds, yet a 5\,GHz CMOS serializer does not work, because the clock has to be \emph{distributed}. A clock passing through a chain of inverters keeps its shape only if each half-period lets every inverter swing fully. Below about 8 fan-out-of-4 delays per period, the pulses shrink at each stage and vanish~\cite{palermo720ckt}. The \textbf{FO4 delay} $t_\mathrm{FO4}$, the delay of an inverter driving four copies of itself, is the natural unit of CMOS speed.

\begin{keyeq}[The clock-distribution limit]
\begin{equation}
  T_\mathrm{clk}\ \ge\ 8\,t_\mathrm{FO4}\qquad\Longrightarrow\qquad \fb \le \frac{n}{8\,t_\mathrm{FO4}},
  \label{eq:ch9-fo4}
\end{equation}
where $n$ is the number of bits per clock period: 1 at full rate, 2 at half rate, 4 at quarter rate.
\end{keyeq}

\begin{example}{Can CMOS clock the last stage?}\label{we:ch9-fo4}
\textbf{Given:} $t_\mathrm{FO4}\approx\qty{42}{\pico\second}$ for 1.2\,V \sgx{} CMOS at the typical corner and about \qty{55}{\pico\second} at the slow corner. These are 130\,nm estimates: measure yours in \tool{ngspice} with a chain of inverters, each loaded by four copies of itself.

\textbf{Full rate} needs $T_\mathrm{clk}=\qty{200}{\pico\second}=4.8\,t_\mathrm{FO4}$, and \cref{eq:ch9-fo4} allows only \qty{3.0}{\giga\bit\per\second}. \textbf{Half rate} needs $T_\mathrm{clk}=\qty{400}{\pico\second}$: $9.5\,t_\mathrm{FO4}$ typical but $7.3\,t_\mathrm{FO4}$ slow, a limit of 5.95 and \qty{4.55}{\giga\bit\per\second}. CMOS fails at the slow corner (\cref{fig:ch9-fo4}).

\textbf{Conclusion:} the final 2:1 stage and its clock must be CML, whose speed is set by $\RL\CL$ rather than FO4. Stages at 625\,MHz and below ($T_\mathrm{clk}\ge29\,t_\mathrm{FO4}$) stay in CMOS.
\end{example}

\begin{figure}[tbp]
\centering
\includegraphics{ch09_fo4}
\caption{\textbf{The 8-FO4 rule rules out full-rate CMOS and leaves half-rate CMOS without margin.} Maximum bit rate from \cref{eq:ch9-fo4} at the typical (42\,ps) and slow (55\,ps) FO4 delays. Only bars above the dashed line can serialize 5\,Gb/s. Quarter and eighth rate pass on paper but need multi-phase clocks and high fan-in output muxes.}
\label{fig:ch9-fo4}
\end{figure}

\subsection{The half-rate CML 2:1 multiplexer}

A \textbf{half-rate} multiplexer has no full-rate flip-flop at its output. The clock \emph{itself} is the select: the even bit goes out while CK is high and the odd bit while CK is low~\cite{palermo720ckt}. \Cref{fig:ch9-mux}(a) shows the HBT version, a stacked (series-gated) CML gate. The clock pair $Q_5/Q_6$ steers the tail current $I_T$ into one of the two data pairs, and the selected pair steers it into one of the shared loads $\RL$. The deselected pair carries no current. The data bases sit at CML levels and the clock bases one $V_\mathrm{BE}$ lower, through emitter followers, so every device keeps its $V_\mathrm{CE}$ near 0.8\,V on the 2.5\,V supply.

\begin{figure}[tbp]
\centering
\begin{tikzpicture}[american, scale=0.64, transform shape, font=\small]
  \node[lblb,anchor=west,font=\sffamily\normalsize\bfseries] at (-2.2,4.0) {(a)};
  \draw (-0.3,4.0) -- (9.1,4.0); \draw (4.4,4.0) node[vcc]{$V_\mathrm{CC}=\qty{2.5}{\volt}$};
  \draw (0,0) node[npn](Q1){};
  \draw (2.4,0) node[npn,xscale=-1](Q2){};
  \draw (6.4,0) node[npn](Q3){};
  \draw (8.8,0) node[npn,xscale=-1](Q4){};
  \draw (1.2,-2.3) node[npn](Q5){};
  \draw (7.6,-2.3) node[npn,xscale=-1](Q6){};
  \foreach \q/\t/\dx in {Q1/Q_1/0.5,Q2/Q_2/-0.5,Q3/Q_3/0.5,Q4/Q_4/-0.5,Q5/Q_5/0.5,Q6/Q_6/-0.5}{\node[text=navy] at ($(\q.center)+(\dx,0.05)$) {$\t$};}
  \draw (0,2.0) to[R, l=$\RL$] (0,4.0);
  \draw (8.8,1.4) to[R, l_=$\RL$] (8.8,4.0);
  \draw (Q1.C) -- (0,2.0) to[short,*-] (6.4,2.0) -- (Q3.C);
  \draw (Q2.C) -- (2.4,1.4) to[short,*-] (6.25,1.4) arc[start angle=180,end angle=0,radius=0.15] -- (8.8,1.4) node[circ]{} -- (Q4.C);
  \draw (0,2.0) to[short,-o] (-1.0,2.0) node[left]{$\mathrm{out}^-$};
  \draw (8.8,1.4) to[short,-o] (9.8,1.4) node[right]{$\mathrm{out}^+$};
  \draw (Q1.B) to[short,-o] ++(-0.45,0) node[left]{$D_\mathrm{e}^+$};
  \draw (Q2.B) to[short,-o] ++(0.45,0) node[above]{$D_\mathrm{e}^-$};
  \draw (Q3.B) to[short,-o] ++(-0.45,0) node[above]{$D_\mathrm{o}^+$};
  \draw (Q4.B) to[short,-o] ++(0.45,0) node[right]{$D_\mathrm{o}^-$};
  \draw (Q1.E) -- (0,-1.15) -- (2.4,-1.15) -- (Q2.E);
  \draw (1.2,-1.15) node[circ]{} -- (Q5.C);
  \draw (Q3.E) -- (6.4,-1.15) -- (8.8,-1.15) -- (Q4.E);
  \draw (7.6,-1.15) node[circ]{} -- (Q6.C);
  \draw (Q5.B) to[short,-o] ++(-0.45,0) node[left]{$\mathrm{CK}^+$};
  \draw (Q6.B) to[short,-o] ++(0.45,0) node[right]{$\mathrm{CK}^-$};
  \draw (Q5.E) -- (1.2,-3.45) -- (7.6,-3.45) -- (Q6.E);
  \draw (4.4,-3.45) node[circ]{} to[I, l=$I_T$] (4.4,-5.0) node[ground]{};
\end{tikzpicture}

\vspace{2mm}
\begin{tikzpicture}[x=4.2mm,y=4.0mm,font=\sffamily\footnotesize,timing/slope=0.12,
  timing/xunit=4.2mm,timing/yunit=4.0mm,timing/d/text/.style={font=\sffamily\scriptsize}]
  \node[lblb,font=\sffamily\small\bfseries] at (-11.5,0.6) {(b)};
  \foreach \s in {0,4,8,12}{
    \fill[gtint] (\s,-2.75) rectangle (\s+2,-1.45);
    \fill[btint] (\s+2,-4.35) rectangle (\s+4,-3.05);
  }
  \foreach \s in {0,2,...,16}{\draw[gridc] (\s,0.25) -- (\s,-6.0);}
  \timing[draw=navy,line width=0.7pt] at (0,-1.0) {2H2L2H2L2H2L2H2L};
  \timing[draw=golddk,line width=0.7pt] at (0,-2.6) {3D{$e_0$}4D{$e_1$}4D{$e_2$}4D{$e_3$}1D{}};
  \timing[draw=bblue,line width=0.7pt] at (0,-4.2) {1D{}4D{$o_0$}4D{$o_1$}4D{$o_2$}3D{$o_3$}};
  \timing[draw=navy,line width=0.9pt] at (0,-5.8) {2D{$e_0$}2D{$o_0$}2D{$e_1$}2D{$o_1$}2D{$e_2$}2D{$o_2$}2D{$e_3$}2D{$o_3$}};
  \node[anchor=east,text=navy] at (-0.4,-0.5) {CK (2.5\,GHz)};
  \node[anchor=east,text=golddk] at (-0.4,-2.1) {$D_\mathrm{e}$ (2.5\,Gb/s)};
  \node[anchor=east,text=bblue] at (-0.4,-3.7) {$D_\mathrm{o}$ (after latch)};
  \node[anchor=east,text=navy] at (-0.4,-5.3) {OUT (5\,Gb/s)};
  \draw[<->,muted] (4,0.55) -- node[above,lbl]{1\,UI} (6,0.55);
  \draw[<->,muted] (0,-6.4) -- node[below,lbl]{$T_\mathrm{clk}=2\,\UI=\qty{400}{\pico\second}$} (4,-6.4);
  \node[lbl,anchor=west,align=left] at (16.4,-2.1) {shaded: selected\\(must be stable)};
  \node[lbl,anchor=west,align=left] at (16.4,-3.7) {changes mid\\CK-high};
\end{tikzpicture}
\caption{\textbf{A half-rate CML multiplexer: the clock picks which data pair gets the current.} (a) When $\mathrm{CK}^+$ is high, $Q_5$ conducts and the even pair $Q_1/Q_2$ drives the loads; when it is low, the odd pair $Q_3/Q_4$ does. The hop marks a crossing without connection. (b) Each input changes only while it is deselected. The even data is selected in the gold windows and changes half a UI after CK falls; an extra latch delays the odd data by half a period so it changes while CK is high. The output UIs follow the clock halves exactly.}
\label{fig:ch9-mux}
\end{figure}

\begin{keyeq}[CML gate design]
\begin{equation}
  \Delta V = I_T\RL,\qquad \tau=\RL\CL,\qquad P = V_\mathrm{CC}I_T=\frac{V_\mathrm{CC}\,\Delta V\,\CL}{\tau}.
  \label{eq:ch9-cml}
\end{equation}
Choose the swing $\Delta V$ first, then $\RL$ from the speed target $\tau$, then $I_T$.
\end{keyeq}

\paragraph{Glitch-free timing and retiming} The select must never switch while a selected input changes. In \cref{fig:ch9-mux}(b) each input changes in the middle of the half where it is deselected: the even data directly, the odd data after one extra latch. Data from the 4:2 stage, which runs on the divided 1.25\,GHz clock, is retimed by latches on the 2.5\,GHz clock, and transitions placed mid-way through the deselected half leave $\pm\qty{100}{\pico\second}$ for divider delay, wiring, setup and hold (\qty{800}{\pico\second} at the CMOS-to-CML boundary). A skew beyond the window does not degrade the eye gracefully: it shifts the bits by one position.

\paragraph{Duty cycle becomes UI width} A clock with duty cycle $D$ produces even UIs of $2D\Tb$ and odd UIs of $2(1-D)\Tb$. A 45\,\% duty cycle gives alternating 180 and 220\,ps bits, \qty{20}{\pico\second} of deterministic jitter before the channel acts (\cref{ch:jitter}). Restore the clock to 50\,\% with a CML buffer chain or a duty-cycle corrector right before the mux.

\begin{insight}
Build the last two serializer stages and the 2.5\,GHz clock buffers in HBT CML, and everything from 625\,MHz down in 1.2\,V CMOS. Half rate removes the 5\,GHz clock but turns duty cycle into UI width: budget for a duty-cycle corrector, and simulate every clock-domain boundary at all corners.
\end{insight}

%% =====================================================================
\section{Termination}\label{sec:termination}

A wave that meets an impedance step reflects (\cref{sec:reflections}). A link can absorb it at the source, at the load or at both ends (\cref{fig:ch9-term}).

\begin{figure}[tbp]
\centering
\begin{tikzpicture}[american, scale=0.72, transform shape, font=\small]
  \tikzset{res/.style={anchor=north,font=\sffamily\small,text=navy,align=center}}
  % (a) series
  \node[lblb,anchor=west] at (-0.6,1.1) {(a) series};
  \draw (0,-1.0) node[ground]{} to[V, l=$V_s$] (0,0) to[R, l={$R_S=\Zo$}] (1.8,0) to[TL, l=$\Zo$] (3.9,0) to[short,-o] (4.4,0);
  \node[lbl,anchor=west] at (4.45,0) {high-Z};
  \node[res] at (2.4,-1.4) {RX sees $V_s$\\echo absorbed at TX};
  % (b) parallel
  \begin{scope}[xshift=6.6cm]
  \node[lblb,anchor=west] at (-0.6,1.1) {(b) parallel};
  \draw (0,-1.0) node[ground]{} to[I, l=$I$] (0,0) -- (0.6,0) to[TL, l=$\Zo$] (2.7,0) -- (3.3,0) to[R, l={$R_T=\Zo$}] (3.3,-1.0) node[ground]{};
  \node[res] at (1.8,-1.4) {RX sees $I\Zo$\\echo returns ($\Gamma_S=+1$)};
  \end{scope}
  % (c) double
  \begin{scope}[xshift=12.6cm]
  \node[lblb,anchor=west] at (-0.6,1.1) {(c) double};
  \draw (0,-1.0) node[ground]{} to[I, l=$I$] (0,0) -- (1.0,0) to[R, l={$R_S$}, *-] (1.0,-1.0) node[ground]{};
  \draw (1.0,0) -- (1.5,0) to[TL, l=$\Zo$] (3.5,0) -- (4.0,0) to[R, l={$R_T$}] (4.0,-1.0) node[ground]{};
  \node[res] at (2.0,-1.4) {RX sees $I\Zo/2$\\echoes absorbed at both ends};
  \end{scope}
\end{tikzpicture}
\caption{\textbf{Three ways to terminate a line.} (a) Series (source) termination launches $V_s/2$, which the open end doubles; it needs a high-impedance RX, which pad capacitance prevents at 5\,Gb/s. (b) Parallel (load) termination: an echo from an imperfect RX returns to the high-impedance source, reflects again and arrives a round trip later as ISI. (c) Double termination ($R_S=R_T=\Zo$) costs half the swing but absorbs echoes at both ends. It is the \OS{} CML driver; a voltage-mode driver with $R_T$ is also doubly terminated.}
\label{fig:ch9-term}
\end{figure}

\paragraph{Why double termination wins} Nothing is perfectly matched: pad capacitance, resistor tolerance and the package always reflect a few percent of the wave at the RX. With $R_S=\Zo$ that echo dies at the TX ($\Gamma_S=0$) instead of returning as a delayed copy of the bit, which would be ISI at a fixed, unknown cursor where no FFE or DFE tap is waiting. The price is swing: half the source current heats the transmitter's own termination.

\paragraph{On-chip resistors in \sgx{}} Place both terminations on the die at the pad. An off-chip resistor leaves bond wire, package and pad as an unterminated stub; on chip, the package becomes part of the line and only the lumped pad capacitance remains (\cref{sec:esd-pad}). Use an unsilicided poly resistor: linear and well matched, but spreading by roughly $\pm20$--30\,\% over process and temperature (check the PDK~\cite{ihppdk}). $\pm20\,\%$ still gives 19\,dB of return loss (\cref{sec:reflections}). A 3-bit trim in 6--8\,\% steps is still worth adding, because a CML driver's swing is proportional to its load resistance. Size each TX resistor for the full 20\,mA.

%% =====================================================================
\section{The CML output driver}\label{sec:cml-driver}

The \OS{} driver (\cref{fig:ch9-cml}) is \cref{fig:ch9-term}(c) in differential form. An HBT pair $Q_1/Q_2$ steers the tail current $I_T$ into one of two loads $\RL=\Zo=\qty{50}{\ohm}$, which are also the TX back-terminations. The RX terminates the line in \qty{100}{\ohm} differential.

\begin{figure}[tbp]
\centering
\begin{tikzpicture}[american, scale=0.7, transform shape, font=\small]
  \draw (-0.4,3.8) -- (3.0,3.8); \draw (1.3,3.8) node[vcc]{$V_\mathrm{CC}=\qty{2.5}{\volt}$};
  \draw (0,0) node[npn](Q1){};
  \draw (2.6,0) node[npn,xscale=-1](Q2){};
  \node[text=navy] at ($(Q1.center)+(0.5,0.05)$) {$Q_1$};
  \node[text=navy] at ($(Q2.center)+(-0.5,0.05)$) {$Q_2$};
  \draw (Q1.C) -- (0,1.5) to[R, l=$\RL$, *-] (0,3.8);
  \draw (Q2.C) -- (2.6,1.0) to[R, l_=$\RL$, *-] (2.6,3.8);
  \draw (0,1.5) -- (2.45,1.5) arc[start angle=180,end angle=0,radius=0.15] -- (4.4,1.5) -- (4.4,2.6) -- (5.2,2.6);
  \draw (2.6,1.0) -- (4.8,1.0) -- (4.8,0.0) -- (5.2,0.0);
  \node[text=navy,anchor=south] at (1.2,1.5) {$\mathrm{out}^-$};
  \node[text=navy,anchor=south] at (3.7,1.0) {$\mathrm{out}^+$};
  \node[draw=navy,fill=navy,minimum size=2mm,inner sep=0pt] at (5.2,2.6) {};
  \node[draw=navy,fill=navy,minimum size=2mm,inner sep=0pt] at (5.2,0.0) {};
  \node[lbl,anchor=south] at (5.2,2.75) {pads};
  \draw (5.3,2.6) to[TL, l=$\Zo$] (8.5,2.6) -- (9.5,2.6);
  \draw (5.3,0.0) to[TL, l_=$\Zo$] (8.5,0.0) -- (9.5,0.0);
  \draw (9.5,2.6) to[R, l_=\qty{50}{\ohm}] (9.5,1.3) to[R, l_=\qty{50}{\ohm}] (9.5,0.0);
  \draw (9.5,1.3) to[short,*-] (10.5,1.3) to[C, l=$C_\mathrm{CM}$] (10.5,0.2) node[ground]{};
  \node[lblb] at (9.5,3.25) {RX};
  \node[lbl] at (6.9,3.4) {channel (REF)};
  \draw (Q1.B) to[short,-o] ++(-0.5,0) node[left]{$D^+$};
  \draw (Q2.B) to[short,-o] ++(0.5,0) node[below]{$D^-$};
  \draw (Q1.E) -- (0,-1.15) -- (2.6,-1.15) -- (Q2.E);
  \draw (1.3,-2.3) node[npn](Q3){};
  \node[text=navy] at ($(Q3.center)+(0.5,0.05)$) {$Q_3$};
  \draw (1.3,-1.15) node[circ]{} -- (Q3.C);
  \draw (Q3.B) to[short,-o] ++(-0.5,0) node[left]{$V_\mathrm{B}$};
  \draw (Q3.E) to[R, l=$R_E$] (1.3,-4.3) node[ground]{};
  \draw[-{Stealth[length=2mm]},c4,line width=0.8pt] (1.75,-1.3) -- (1.75,-1.75) node[right,text=c4,font=\small] {$I_T$};
\end{tikzpicture}
\caption{\textbf{The \OS{} CML driver.} $Q_1/Q_2$ steer the tail current $I_T$, set by $Q_3$, $R_E$ and $V_\mathrm{B}$, into one of the \qty{50}{\ohm} loads $\RL$. The loads double as TX back-terminations, so each output sees $\RL\parallel\Zo=\qty{25}{\ohm}$. The RX terminates in 2$\times$\qty{50}{\ohm} with an AC-grounded center tap. 20\,mA gives 1\,Vppd.}
\label{fig:ch9-cml}
\end{figure}

\paragraph{Swing} When $Q_2$ takes all of $I_T$, the current divides between its own load and the line, so each output sees $\RL\parallel\Zo=\qty{25}{\ohm}$ and swings $I_T\cdot\qty{25}{\ohm}$ single-ended. The differential voltage flips sign with the data, so $V_{d,\mathrm{pp}}=2\cdot I_T\cdot\qty{25}{\ohm}=I_T\Zo$.

\paragraph{Voltage mode} The alternative is a \textbf{voltage-mode} (VM) driver, usually a source-series-terminated (SST) one: each output is an inverter-like switch pair in series with a resistor, $R_\mathrm{on}+R_S=\Zo$~\cite{p4:wong2004}. Its current flows only around the loop of TX resistor, line, \qty{100}{\ohm} RX termination, line and TX resistor: $4\Zo=\qty{200}{\ohm}$. A supply $V_s$ drives $I=V_s/4\Zo$, the RX sees $\pm V_s/2$, and $V_{d,\mathrm{pp}}=V_s$.

\begin{keyeq}[Driver current for a differential swing $V_{d,\mathrm{pp}}$]
\begin{equation}
  I_\mathrm{CML}=\frac{V_{d,\mathrm{pp}}}{\Zo},\qquad
  I_\mathrm{VM,\,SE\ term.}=\frac{V_{d,\mathrm{pp}}}{2\Zo},\qquad
  I_\mathrm{VM,\,diff.\ term.}=\frac{V_{d,\mathrm{pp}}}{4\Zo}.
  \label{eq:ch9-drv}
\end{equation}
CML needs four times the current of VM into a differential termination: half of it flows in its own termination, and it never reverses through the line.
\end{keyeq}

\begin{example}{CML versus SST at \OS{} swings}\label{we:ch9-drv}
\textbf{CML:} 800\,mVppd needs $I_T=0.8/50=\qty{16}{\milli\ampere}$, \qty{40}{\milli\watt} from 2.5\,V. At 1\,Vppd (the model value, $A=\qty{0.5}{\volt}$) it needs \qty{20}{\milli\ampere} and \textbf{\qty{50}{\milli\watt}}. The output common mode sits $I_T\times\qty{25}{\ohm}=\qty{0.5}{\volt}$ below $V_\mathrm{CC}$, and each output moves between 1.75 and 2.25\,V.

\textbf{SST:} 800\,mVppd needs $I=0.8/200=\qty{4}{\milli\ampere}$ from $V_s\approx\qty{0.8}{\volt}$: \qty{3.2}{\milli\watt}, or \qty{4.8}{\milli\watt} through a regulator from 1.2\,V (6\,mW at 1\,Vppd). Keep $R_\mathrm{on}\le\qty{10}{\ohm}$, so a $\pm30\,\%$ spread moves the output impedance by only $\pm\qty{3}{\ohm}$.

\textbf{Verdict:} SST saves about 44\,mW (\cref{fig:ch9-current}), but \OS{} uses CML: a passive-resistor impedance with no calibration loop, constant supply current, FFE taps that leave the impedance untouched, and fast HBTs instead of CMOS pre-drivers passing 200\,ps (4.8\,FO4) pulses. SST is the power upgrade for a second tape-out.
\end{example}

\begin{figure}[tbp]
\centering
\includegraphics{ch09_current}
\caption{\textbf{A CML driver pays four times the current of a voltage-mode driver for the same swing.} (a) Supply current from \cref{eq:ch9-drv}; the band is the \OS{} range of 0.8--1.0\,Vppd. With 6\,dB of de-emphasis an SST driver draws 1.38 times more on average, because its main and post slices fight; CML draws the same current at every FFE setting. (b) Output-stage power: 40--50\,mW for HBT CML from 2.5\,V versus about 6\,mW for SST.}
\label{fig:ch9-current}
\end{figure}

\paragraph{Why HBTs} A CML stage must switch its tail current fully. A bipolar pair splits it as $\Delta I=I_T\tanh(V_\mathrm{id}/2V_T)$, so $V_\mathrm{id}\approx4V_T\approx\qty{100}{\milli\volt}$ steers 96\,\% of it at any current. A MOS pair needs $\sqrt2\,V_\mathrm{OD}$, typically 200--300\,mV. The HBT's $\gm/I_C=1/V_T=\qty{38.7}{\per\volt}$ is also several times a MOSFET's.

\begin{example}{HBT versus MOS in a 2.5\,GHz CML stage}\label{we:ch9-hbt}
\textbf{Given:} a mux or pre-driver stage loaded by $\CL=\qty{40}{\femto\farad}$, with $f_\mathrm{3dB}\ge\qty{7.5}{\giga\hertz}$, so $\tau\le\qty{21}{\pico\second}$. By \cref{eq:ch9-cml}, $\RL=\qty{500}{\ohm}$ gives $\tau=\qty{20}{\pico\second}$; $\Delta V=\qty{300}{\milli\volt}$ then needs $I_T=\qty{0.6}{\milli\ampere}$, or \qty{1.5}{\milli\watt} from 2.5\,V.

\textbf{HBT pair:} at the crossing each device carries 0.3\,mA, $\gm=\qty{11.6}{\milli\siemens}$, and the gain is $\gm\RL=\Delta V/2V_T=\mathbf{5.8}$. Full switching needs about 104\,mV.

\textbf{MOS pair} at the same current with $V_\mathrm{OD}=\qty{150}{\milli\volt}$: $\gm=\qty{4}{\milli\siemens}$, gain $\Delta V/V_\mathrm{OD}=\mathbf{2}$, full switching at 212\,mV. More gain needs wider devices, more $\CL$ and more current. A chain of HBT stages can therefore run at smaller swing, with fewer and smaller stages.
\end{example}

\begin{pitfall}[Headroom and the $\mathrm{BV}_\mathrm{CEO}$ limit]
\sgx{} HBTs break down at about $\mathrm{BV}_\mathrm{CEO}\approx\qty{1.6}{\volt}$ (check the PDK). On 2.5\,V the \emph{off} driver transistor is at risk: its collector sits at the high output level. Stack up \cref{we:ch9-drv} at 1\,Vppd. Outputs: 2.25\,V (off) and 1.75\,V (on). Bases high: 1.9\,V, from a pre-driver whose common-mode resistor pulls its outputs 0.6\,V below $V_\mathrm{CC}$. Common emitter: $1.9-V_\mathrm{BE}\approx\qty{1.0}{\volt}$. So $V_\mathrm{CE}=0.75$\,V on and 1.25\,V off, with 1.0\,V left for the tail. The 0.35\,V margin is eaten by $V_\mathrm{CC}+5\,\%$, overshoot and reflections; an emitter-follower pre-driver with bases near 1.6\,V would leave the off device at 1.5\,V nominal. Check every $V_\mathrm{CE}$ at every corner \emph{with the real RX termination attached}: a far end terminated to ground moves the output common mode.
\end{pitfall}

%% =====================================================================
\section{The FFE inside the driver}\label{sec:ffe-impl}

\Cref{ch:txffe} chose the taps; here we build them. Put one differential pair per tap on the same pair of loads, drive each with its own copy of the data, and set each tail current to the tap weight (\cref{fig:ch9-ffe}). The loads add the currents, which forms the FIR sum of \cref{sec:ffe-taps}. Negative side taps cost nothing: swapping a pair's base inputs inverts its contribution.

\begin{figure}[tbp]
\centering
\begin{tikzpicture}[american, scale=0.66, transform shape, font=\small]
  \draw (-0.6,4.3) -- (11.8,4.3); \draw (5.6,4.3) node[vcc]{$V_\mathrm{CC}=\qty{2.5}{\volt}$};
  \foreach \x/\n in {0/a,4.8/c,9.6/e}{\draw (\x,0) node[npn](Q\n){};}
  \foreach \x/\n in {1.6/b,6.4/d,11.2/f}{\draw (\x,0) node[npn,xscale=-1](Q\n){};}
  \foreach \n/\t/\dx in {a/Q_5/0.5,b/Q_6/-0.5,c/Q_1/0.5,d/Q_2/-0.5,e/Q_3/0.5,f/Q_4/-0.5}{\node[text=navy] at ($(Q\n.center)+(\dx,0.05)$) {$\t$};}
  \draw (0,2.2) to[R, l=$\RL$] (0,4.3);
  \draw (11.2,1.5) to[R, l_=$\RL$] (11.2,4.3);
  \draw (Qa.C) -- (0,2.2) to[short,*-] (9.6,2.2) -- (Qe.C);
  \draw (4.8,2.2) node[circ]{} -- (Qc.C);
  \draw (Qb.C) -- (1.6,1.5) to[short,*-] (4.65,1.5) arc[start angle=180,end angle=0,radius=0.15] -- (9.45,1.5) arc[start angle=180,end angle=0,radius=0.15] -- (11.2,1.5) node[circ]{} -- (Qf.C);
  \draw (6.4,1.5) node[circ]{} -- (Qd.C);
  \draw (0,2.2) to[short,-o] (-1.1,2.2) node[left]{$\mathrm{out}^-$};
  \draw (11.2,1.5) to[short,-o] (12.3,1.5) node[right]{$\mathrm{out}^+$};
  \draw (Qa.B) to[short,-o] ++(-0.45,0) node[left]{$D_{n+1}^-$};
  \draw (Qb.B) to[short,-o] ++(0.45,0) node[below=1pt]{$D_{n+1}^+$};
  \draw (Qc.B) to[short,-o] ++(-0.45,0) node[above=1pt]{$D_{n}^+$};
  \draw (Qd.B) to[short,-o] ++(0.45,0) node[below=1pt]{$D_{n}^-$};
  \draw (Qe.B) to[short,-o] ++(-0.45,0) node[above=1pt]{$D_{n-1}^-$};
  \draw (Qf.B) to[short,-o] ++(0.45,0) node[right]{$D_{n-1}^+$};
  \foreach \l/\r/\c/\lab in {a/b/0.8/{I_{-1}=|c_{-1}|I_T}, c/d/5.6/{I_0=c_0I_T}, e/f/10.4/{I_1=|c_1|I_T}}{
    \draw (Q\l.E) -- ($(Q\l.E)+(0,-0.4)$) -| (Q\r.E);
    \draw (\c,-1.17) node[circ]{} to[I, l={$\lab$}] (\c,-2.9) node[ground]{};
  }
  \node[lblb,font=\sffamily\normalsize\bfseries] at (0.8,-3.75) {pre tap};
  \node[lblb,font=\sffamily\normalsize\bfseries] at (5.6,-3.75) {main tap};
  \node[lblb,font=\sffamily\normalsize\bfseries] at (10.4,-3.75) {post tap};
\end{tikzpicture}
\caption{\textbf{A 3-tap FFE is three CML pairs on one pair of loads.} Each pair is driven by its own copy of the data and steers a tail current proportional to its tap weight. The pre and post pairs have their inputs swapped, which gives $c_{-1},c_1<0$. The tails are current DACs sharing one reference, with $I_{-1}+I_0+I_1=I_T=\qty{20}{\milli\ampere}$, so the supply current never changes. Hops mark crossings without connection.}
\label{fig:ch9-ffe}
\end{figure}

\begin{keyeq}[Current-summing FFE]
\begin{equation}
  v_{d}[n]=\frac{\Zo}{2}\,I_T\sum_{k=-1}^{1}c_k\,a_{n-k},\qquad c_k=\pm\frac{I_k}{I_T},\qquad \sum_k\lvert c_k\rvert=\frac{I_{-1}+I_0+I_1}{I_T}=1 .
  \label{eq:ch9-ffe}
\end{equation}
The peak-swing constraint of \cref{sec:ffe-taps} holds by construction: the taps are slices of one fixed tail current.
\end{keyeq}

\paragraph{Making the delayed copies} No 5\,GHz clock is needed. Call the half-rate streams $E_m$ (bit $2m$) and $O_m$ (bit $2m+1$). The main mux selects $(E_m,O_m)$. The post-tap mux must emit bit $2m-1$ in the even slot and bit $2m$ in the odd slot, so it selects $(O_{m-1},E_m)$; the pre-tap mux selects $(O_m,E_{m+1})$. About four extra latches on the 2.5\,GHz clock keep these inputs glitch-free (\cref{fig:ch9-mux}b).

\begin{example}{Tap currents for the $-6$\,dB preset}\label{we:ch9-ffe}
The de-emphasis of a 2-tap setting is $\mathrm{DE}=20\log_{10}\big[(\lvert c_0\rvert+\lvert c_1\rvert)/(\lvert c_0\rvert-\lvert c_1\rvert)\big]$, the ratio of transition to repeated-bit swing. For 6\,dB that ratio is 2; with $\lvert c_0\rvert+\lvert c_1\rvert=1$, $\lvert c_0\rvert-\lvert c_1\rvert=0.5$, so $c_0=0.75$ and $c_1=-0.25$, the PCIe preset of \cref{sec:ffe-taps}. (The often-quoted $c_1=-1/3$ is the 9.5\,dB setting.)

\textbf{At 1\,Vppd} ($I_T=\qty{20}{\milli\ampere}$): $I_0=\qty{15}{\milli\ampere}$, $I_1=\qty{5}{\milli\ampere}$. Transition bits swing $(15+5)\,\mathrm{mA}\times\Zo=\qty{1.0}{\Vppd}$, repeated bits $(15-5)\times50=\qty{0.5}{\Vppd}$. The supply stays at \qty{50}{\milli\watt} for every preset.

\textbf{Resolution:} build $I_T$ from 32 cells of \qty{0.625}{\milli\ampere} (5 bits). The $-6$\,dB preset is exact (24 and 8 cells). The $-3.5$\,dB preset needs 5.3 cells; five give $c_1=-0.156$ and 3.2\,dB, which is harmless because the eye is flat near its optimum (\cref{ch:txffe}). A pre-tap moves in steps of $1/32=0.031$.
\end{example}

\paragraph{The resolution trade-off} Range and resolution cost capacitance, not current. Sizing each pair for its largest current (20\,mA main, 8\,mA for $\lvert c_1\rvert\le0.4$, 3\,mA for $\lvert c_{-1}\rvert\le0.15$) loads the output with 31\,mA worth of collectors, 1.55 times a plain driver. A \textbf{segmented} driver instead uses 32 identical cells, each with a selector in its pre-driver that assigns it to the pre, main or post stream: the output always sees exactly 20\,mA of devices, and the cost moves into the pre-drivers.

\begin{insight}
Build the FFE from current. The tail current is fixed, so power and output impedance do not depend on the preset, and $\sum\lvert c_k\rvert=1$ holds automatically. Five bits are enough. Spend layout effort on matching the three tap paths: a 10\,ps skew between $D_n$ and $D_{n-1}$ moves the de-emphasis edge by 0.05\,UI.
\end{insight}

%% =====================================================================
\section{Pad, ESD and bandwidth}\label{sec:esd-pad}

At the pad the driver's collector capacitance, the termination's parasitics, the ESD diodes and the bond pad hang on one node (\cref{fig:ch9-tcoil}a). That node sees $\RL\parallel\Zo=\qty{25}{\ohm}$, so the total capacitance $C_p$ makes a pole at
\begin{equation}
  f_\mathrm{3dB}=\frac{1}{2\pi(\Zo/2)C_p}=\frac{1}{\pi\Zo C_p}
  \label{eq:ch9-padpole}
\end{equation}
and spoils the return loss (\cref{sec:reflections}). The \textbf{ESD} protection cannot be left out: dual diodes from the pad to both rails, plus a supply clamp, steer an electrostatic discharge away from the HBT and thin-oxide junctions. Diodes rated for a few kilovolts of human-body-model discharge are large and often dominate $C_p$.

\begin{example}{The \OS{} pad capacitance budget}\label{we:ch9-pad}
\textbf{Illustrative values} (extract the real ones from layout): bond pad \qty{80}{\femto\farad}, ESD diodes \qty{250}{\femto\farad}, driver collectors and termination parasitics \qty{300}{\femto\farad}: $C_p=\qty{0.63}{\pico\farad}$.

\textbf{Return loss at $\fN$:} $x=\omega\Zo C_p=0.495$ and $\lvert S_{11}\rvert=x/\sqrt{4+x^2}=0.24$, or \textbf{$-12.4$\,dB}. It passes a 10\,dB target ($C_p\le\qty{0.85}{\pico\farad}$) but fails 15\,dB ($C_p\le\qty{0.46}{\pico\farad}$). \textbf{Bandwidth:} $f_\mathrm{3dB}=1/(\pi\cdot50\cdot0.63\,\mathrm{pF})=\qty{10.1}{\giga\hertz}$, four times $\fN$.

\textbf{System check:} the REF model assumes \qty{300}{\femto\farad} per pad. With \code{c\_pad=0.63e-12} in \code{channel()} the loss at $\fN$ grows from 17.4 to \qty{17.8}{\deci\bel} and the FFE-only PDA eye drops from 62 to 54\,mV. Always rerun the system model with the extracted pad capacitance.
\end{example}

\begin{figure}[tbp]
\centering
\includegraphics{ch09_pad}
\caption{\textbf{Pad capacitance costs return loss long before it costs bandwidth.} (a) Return loss at 2.5\,GHz versus total pad capacitance. The 0.63\,pF example passes the 10\,dB line (limit 0.85\,pF). A T-coil with \qty{3}{\ohm} of loss per coil holds $\lvert S_{11}\rvert$ near $-25$\,dB. (b) Response at the capacitive node of a doubly terminated pad: the T-coil moves the 0.63\,pF pole from 10.1 to 14.3\,GHz with a maximally flat response.}
\label{fig:ch9-pad}
\end{figure}

\paragraph{The T-coil} When the budget does not close, a \textbf{bridged T-coil} separates the capacitance from the line~\cite{p4:galal2003,p4:kossel2008}. Two coupled inductors connect pad and termination, the capacitive load hangs on their center tap, and a small $C_B$ bridges the ends (\cref{fig:ch9-tcoil}b). With the right values the inductive and capacitive impedances cancel exactly~\cite{p4:razavi2015tcoil}.

\begin{keyeq}[Bridged T-coil, $\zeta=1/\sqrt2$]
\begin{equation}
  L_1=L_2=\frac{R^2C}{2(1+k)},\quad k=\frac{4\zeta^2-1}{4\zeta^2+1}=\frac13,\quad C_B=\frac{C}{16\zeta^2}=\frac{C}{8},\quad Z_\mathrm{in}=R\ \text{at all }f.
  \label{eq:ch9-tcoil}
\end{equation}
The total inductance is $L_1+L_2+2M=R^2C$. From a current source, the capacitor node then has 2.83 times the bandwidth of the plain $RC$. In a doubly terminated pad the plain pole already sees $R/2$, so the bandwidth gain is $\sqrt2$; the bigger gain is in return loss.
\end{keyeq}

\begin{figure}[tbp]
\centering
\begin{tikzpicture}[american, scale=0.78, transform shape, font=\small]
  \node[lblb,anchor=west,font=\sffamily\small\bfseries] at (-1.0,2.4) {(a) everything on one node};
  \draw (-0.6,1.6) -- (3.6,1.6) node[right]{$V_\mathrm{CC}$};
  \draw (-0.6,-1.6) -- (3.6,-1.6);
  \draw (-0.6,0) -- (3.4,0);
  \draw (-0.2,0) to[I, l_=$i_\mathrm{drv}$, *-] (-0.2,-1.6);
  \draw (0.8,0) to[R, l=$\RL$, *-] (0.8,1.6);
  \draw (1.8,0) to[D*, l_=$D_1$, *-] (1.8,1.6);
  \draw (1.8,-1.6) to[D*, l_=$D_2$] (1.8,0);
  \draw (2.8,0) to[C, l=$C_\mathrm{pad}$, *-] (2.8,-1.6);
  \node[draw=navy,fill=navy,minimum size=2mm,inner sep=0pt] at (3.4,0) {};
  \draw (3.5,0) to[TL, l=$\Zo$] (5.6,0);
  \begin{scope}[xshift=8.6cm]
  \node[lblb,anchor=west,font=\sffamily\small\bfseries] at (-1.0,2.4) {(b) bridged T-coil};
  \draw (-0.7,0) node[left]{T} to[short,o-] (0,0) to[L, l=$L_2$, *-*] (2.2,0) to[L, l=$L_1$, -*] (4.4,0);
  \node[text=navy] at (4.55,-0.3) {P};
  \draw (0,0) to[R, l=$\RL$] (0,-1.6) node[ground]{};
  \draw (0,0) -- (0,2.0) to[C, l=$C_B$] (4.4,2.0) -- (4.4,0);
  \node[text=navy] at (2.2,0.35) {X};
  \draw (1.7,0) -- (1.7,-0.35) to[C, l_=$C$] (1.7,-1.6) node[ground]{};
  \draw (2.7,0) -- (2.7,-0.35) to[I, l=$i_\mathrm{drv}$] (2.7,-1.6) node[ground]{};
  \draw[<->,muted,line width=0.6pt] (1.0,0.85) to[bend left=30] (3.4,0.85);
  \node[text=navy] at (2.2,1.45) {$k=\tfrac13$};
  \node[draw=navy,fill=navy,minimum size=2mm,inner sep=0pt] at (4.6,0) {};
  \draw (4.7,0) to[TL, l=$\Zo$] (6.6,0);
  \end{scope}
\end{tikzpicture}
\caption{\textbf{A T-coil hides the pad capacitance from the line.} (a) Without it, driver, termination $\RL$, ESD diodes $D_1/D_2$ and pad share one node, and $C_p=C_\mathrm{drv}+C_\mathrm{ESD}+C_\mathrm{pad}$ sets both the pole of \cref{eq:ch9-padpole} and the return loss. (b) With a bridged T-coil the load $C=C_\mathrm{ESD}+C_\mathrm{drv}$ sits at the center tap X, and $\RL$ (returned to $V_\mathrm{CC}$) at the far end T. The line at P sees $\RL$ at all frequencies.}
\label{fig:ch9-tcoil}
\end{figure}

For $C=\qty{0.63}{\pico\farad}$ and $R=\qty{50}{\ohm}$, \cref{eq:ch9-tcoil} gives $L_1=L_2=\qty{0.59}{\nano\henry}$, $M=\qty{0.20}{\nano\henry}$ and $C_B=\qty{79}{\femto\farad}$ (often the coil's own inter-winding capacitance): better than 25\,dB of return loss at $\fN$ with \qty{3}{\ohm} per coil (\cref{fig:ch9-pad}). In \sgx{} it fits in roughly $\qty{60}{\micro\meter}\times\qty{60}{\micro\meter}$ on the thick top metals; take $k$, $C_B$ and self-resonance from EM simulation (\cref{sec:layout-hs}).

\begin{insight}[Design rule of thumb]
At 5\,Gb/s a T-coil is \emph{optional}. Below about 0.85\,pF the pad already meets 10\,dB of return loss at $\fN$, and its pole at 7.5\,GHz or above costs little ISI. Add a T-coil if ESD and driver push $C_p$ toward 1\,pF, or if the return-loss target is tighter (15\,dB at $\fN$). Either way, give the system model your extracted $C_p$, not the 300\,fF of REF.
\end{insight}

%% =====================================================================
\section{TX design checklist}\label{sec:ch9-checklist}

\Cref{tab:ch9-check} collects the checks and a first-cut power estimate of about 100\,mW, half of it in the driver (\cref{sec:power-budget} adds the RX and clocking).

\begin{table}[tbp]
\centering
\caption{\textbf{TX checklist and first-cut power at 1\,Vppd.} CML gates as in \cref{we:ch9-hbt}, about 1.2\,mA each with emitter followers. Confirm every line in \tool{ngspice}.}
\label{tab:ch9-check}
\small
\begin{tabularx}{\linewidth}{@{}l L r@{}}
\toprule
Block & Check before layout & Power\\
\midrule
3-tap CML driver & $I_T=20$\,mA (16\,mA at 0.8\,Vppd); every $V_\mathrm{CE}<\mathrm{BV}_\mathrm{CEO}$ at $V_\mathrm{CC}+5\,\%$ & 50\,mW\\
Pre-drivers ($\times3$) & full switching ($\ge4V_T$); tap paths matched to a few ps & 10\,mW\\
2:1 muxes ($\times3$), latches ($\times4$) & glitch-free phases; 5-bit tap DACs, $-3.5$/$-6$\,dB presets & 21\,mW\\
Clock buffers, 4:2, $\div2$ & duty cycle $50\pm1\,\%$; clock-to-data windows at all corners & 20\,mW\\
CMOS 16:4 (1.2\,V) & FO4 measured; timing closed at 625\,MHz & 3\,mW\\
Termination and pad & poly $\RL$ for 20\,mA; $C_p<0.85$\,pF or T-coil; RL $\ge10$\,dB at $\fN$ & --\\
\midrule
\textbf{Total TX} & 41\,mA from 2.5\,V, 2.5\,mA from 1.2\,V & \textbf{$\approx$104\,mW}\\
\bottomrule
\end{tabularx}
\end{table}

\begin{takeaways}
\begin{itemize}
  \item A CMOS clock tree needs 8\,FO4 per period. With $t_\mathrm{FO4}\approx42$\,ps (55\,ps slow), full-rate CMOS cannot reach 5\,Gb/s and half-rate CMOS fails at the slow corner, so the final 2:1 mux and its 2.5\,GHz clock are HBT CML.
  \item A half-rate mux uses the clock as the select. Each input must change only while deselected, and the clock duty cycle becomes the UI width.
  \item Double termination absorbs echoes at both ends for half the swing. A CML driver needs $I_T=V_{d,\mathrm{pp}}/\Zo$: 20\,mA and 50\,mW from 2.5\,V at 1\,Vppd, four times a voltage-mode driver.
  \item HBTs switch fully with about 100\,mV: a 2.5\,GHz CML stage has a gain of 5.8 against 2 for MOS at the same current.
  \item The off driver transistor sees the highest $V_\mathrm{CE}$; keep it below $\mathrm{BV}_\mathrm{CEO}\approx1.6$\,V at every corner.
  \item The FFE is current-steering pairs on shared loads; $-6$\,dB is $c_0=0.75$, $c_1=-0.25$ (15 and 5\,mA). Five bits are enough.
  \item Keep $C_p$ below 0.85\,pF for 10\,dB return loss at $\fN$; a T-coil ($k=1/3$, $C_B=C/8$) absorbs larger loads.
\end{itemize}
\end{takeaways}

\tikzsetfigurename{10_rx-}
\chapter{Receiver Front-End and Samplers}\label{ch:rx}
\chapterlead{The receiver must turn a 100\,mV, 200\,ps-wide eye into clean logic levels, $10^{12}$ times in a row without a mistake. This chapter builds the circuits that do it.}

\begin{objectives}
\begin{itemize}
  \item Draw the \OS{} receiver from pad to 16-bit word, and set its termination, common mode and coupling.
  \item List what a sampler must deliver (sensitivity, offset, aperture, metastability, kick-back, noise) and derive the regeneration time constant of the StrongARM and the HBT CML latch.
  \item Compute the metastability error probability and explain why edge samplers need extra latch stages.
  \item Estimate sampler offset from Pelgrom's law and size a 5-bit calibration DAC.
  \item Build a circuit-level noise and sensitivity budget that supports the system budget of \cref{sec:noisebudget}.
  \item Deserialize half-rate decisions to 16 bits at 312.5\,MHz and cross into the core clock domain.
\end{itemize}
\end{objectives}

%% =====================================================================
\section{The receiver front-end}\label{sec:rx-frontend}

\Cref{fig:ch10-arch} shows the receive path. The \qty{100}{\ohm} termination absorbs the wave at the pad. The CTLE (\cref{sec:ctle-hbt}) boosts the high frequencies, and the DFE summer subtracts the post-cursor echoes of earlier decisions (\cref{sec:dfe-timing}). Four samplers then decide. At half rate, $D_0$ and $D_1$ sample the data at \ang{0} and \ang{180} of the 2.5\,GHz clock, one bit each per clock period. $E_0$ and $E_1$ sample the transitions between them at \ang{90} and \ang{270}. The deserializer turns these four 2.5\,Gb/s streams into 16 data bits and 16 edge bits at \qty{312.5}{\mega\hertz}. The CDR logic uses both to steer the phase interpolator (\cref{ch:cdr}), which shifts the four clock phases until $E_0$ and $E_1$ sit on the zero crossings.

\begin{figure}[tbp]
\centering
\begin{tikzpicture}[x=1cm,y=1cm,font=\sffamily\footnotesize,
  b/.style={blk,minimum width=13mm,minimum height=10mm,font=\sffamily\footnotesize,inner sep=2pt},
  bb/.style={b,fill=btint}, bg/.style={b,fill=gtint,draw=golddk}, bn/.style={b,fill=navy,text=white},
  sm/.style={blk,fill=btint,minimum width=16mm,minimum height=7mm,font=\sffamily\footnotesize,inner sep=1pt},
  ch/.style={font=\sffamily\scriptsize,text=muted}]
  \node[draw=navy,fill=navy,minimum size=2mm,inner sep=0pt] (pad) at (0,0) {};
  \draw[sig] (-0.9,0) node[left,text=navy,align=right]{from\\channel} -- (pad.west);
  \node[bb] (term) at (1.35,0) {100\,$\Omega$\\+ ESD};
  \node[bb] (ctle) at (3.3,0) {CTLE\\0--12\,dB};
  \node[sumnode] (sum) at (5.0,0) {$\Sigma$};
  \draw[navy,line width=0.75pt] (pad.east) -- (term.west);
  \draw[sig] (term) -- (ctle);
  \draw[sig] (ctle) -- (sum);
  \node[ch,anchor=north] at (3.3,-0.6) {Ch.~7};
  \node[ch,anchor=north,align=center] at (5.0,-0.4) {DFE\\summer};
  % samplers
  \draw[region,fill=gtint!25] (6.15,-2.55) rectangle (8.65,2.55);
  \foreach \y/\n/\t in {1.95/D0/{$D_0$ \ang{0}}, 0.65/E0/{$E_0$ \ang{90}}, -0.65/D1/{$D_1$ \ang{180}}, -1.95/E1/{$E_1$ \ang{270}}}{
    \node[sm] (\n) at (7.4,\y) {\t};
    \draw[sig] (6.0,\y) -- (\n.west);
  }
  \draw[navy,line width=0.75pt] (sum.east) -- (6.0,0);
  \draw[navy,line width=0.75pt] (6.0,-1.95) -- (6.0,1.95);
  \node[dot] at (6.0,0) {};
  \node[lblb,anchor=south] at (7.4,2.6) {samplers};
  % deserializer
  \node[blk,fill=gtint,draw=golddk,minimum width=13mm,minimum height=50mm] (des) at (10.0,0) {1:16\\deser.};
  \foreach \y/\n in {1.95/D0, 0.65/E0, -0.65/D1, -1.95/E1}{\draw[sig] (\n.east) -- (des.west |- 0,\y);}
  \node[ch,anchor=south] at (9.0,1.95) {2.5\,Gb/s};
  \draw[sig,line width=1.4pt] (des.east |- 0,1.4) -- (13.4,1.4) node[above left,text=navy,align=right]{16 data bits\\312.5\,MHz};
  \node[bg] (cdr) at (12.4,-1.0) {CDR logic\\BB-PD + filter};
  \draw[sig,line width=1.4pt] (des.east |- 0,-1.0) -- (cdr.west);
  \node[ch,anchor=south] at (11.0,-1.0) {data + edge};
  % clocking
  \node[bn] (pll) at (5.0,-3.6) {PLL\\I/Q 2.5\,GHz};
  \node[bb] (pi) at (7.4,-3.6) {phase\\interp.};
  \draw[clk] (pll) -- (pi);
  \draw[sigg] (cdr.south) |- node[pos=0.75,below,text=golddk,font=\sffamily\scriptsize]{PI code} (pi.east);
  \draw[clk] (pi.north) -- (7.4,-2.55) node[pos=0.5,right,text=golddk,font=\sffamily\scriptsize]{4 phases};
  \node[ch,anchor=north] at (10.2,-3.95) {Ch.~12--13};
  % DFE feedback
  \draw[sigg] (8.4,1.95) node[dot,fill=golddk]{} |- (5.0,3.15) -- (sum.north);
  \node[font=\sffamily\scriptsize,text=golddk,anchor=south] at (6.6,3.15) {decisions $\hat a_{n-1},\hat a_{n-2}$ $\times$ taps $w_1,w_2$ (Ch.~8)};
\end{tikzpicture}
\caption{\textbf{The \OS{} receiver.} After termination, CTLE and DFE summer, four half-rate samplers decide on the four phases of the 2.5\,GHz clock: data at \ang{0}/\ang{180}, edges at \ang{90}/\ang{270}. The decisions feed back to the summer (gold) and, deserialized to 312.5\,MHz, drive the CDR logic that steers the phase interpolator. The DFE uses separate even and odd summers in practice (\cref{sec:dfe-timing}).}
\label{fig:ch10-arch}
\end{figure}

\paragraph{Termination and common mode} The RX termination is two \qty{50}{\ohm} poly resistors in series across the inputs, with their center tap AC-grounded by a capacitor $C_\mathrm{CM}$ (\cref{fig:ch10-term}). Differential signals see \qty{100}{\ohm}. Common-mode noise sees a low impedance and is absorbed too. The termination sits right at the pad, inside the ESD diodes, and shares the pad capacitance budget of \cref{sec:esd-pad}. The CTLE input pair needs its own common mode, and the input stage must stay linear for the largest signal: long runs on a short channel arrive at nearly the full $\pm\qty{0.5}{\volt}$. A switchable input attenuator, or the CTLE's own gain control, handles short channels (\cref{sec:ctle-hbt}).

\paragraph{DC or AC coupling} With \textbf{DC coupling} the RX inherits the TX common mode: about 2.0\,V from the CML driver of \cref{sec:cml-driver}, which suits an HBT CTLE on the same 2.5\,V supply. \textbf{AC coupling} with series capacitors lets each side set its own common mode, at the price of a high-pass corner. A run of $N$ identical bits then droops by $1-e^{-N\Tb/\tau}$ (\textbf{baseline wander}). Off-chip \qty{100}{\nano\farad} capacitors in the \qty{100}{\ohm} loop give $\tau=\qty{10}{\micro\second}$, and PRBS31's 31-bit run droops by only 0.06\,\%. On chip, \qty{2}{\pico\farad} into a \qty{10}{\kilo\ohm} bias resistor gives $\tau=\qty{20}{\nano\second}$ and a 27\,\% droop, which needs DC-balanced (scrambled or 8b/10b) data. \OS{} is DC-coupled for near-end loopback and has board footprints for off-chip capacitors.

\begin{figure}[tbp]
\centering
\begin{tikzpicture}[american, scale=0.72, transform shape, font=\small]
  \node[draw=navy,fill=navy,minimum size=2mm,inner sep=0pt] at (0,1.6) {};
  \node[draw=navy,fill=navy,minimum size=2mm,inner sep=0pt] at (0,-1.6) {};
  \node[left=3pt,text=navy] at (0,1.6) {$\mathrm{in}^+$};
  \node[left=3pt,text=navy] at (0,-1.6) {$\mathrm{in}^-$};
  \draw (0.1,1.6) -- (4.0,1.6); \draw (0.1,-1.6) -- (4.0,-1.6);
  \node[blk,minimum width=9mm,minimum height=6mm,font=\sffamily\footnotesize] (e1) at (1.0,2.75) {ESD};
  \node[blk,minimum width=9mm,minimum height=6mm,font=\sffamily\footnotesize] (e2) at (1.0,-2.75) {ESD};
  \draw (1.0,1.6) node[circ]{} -- (e1.south); \draw (1.0,-1.6) node[circ]{} -- (e2.north);
  \draw (2.4,1.6) to[R, l_=\qty{50}{\ohm}, *-*] (2.4,0) to[R, l_=\qty{50}{\ohm}, -*] (2.4,-1.6);
  \draw (2.4,0) -- (3.3,0) to[C, l=$C_\mathrm{CM}$] (3.3,-1.1) node[ground]{};
  \node[lbl,anchor=north west,align=left] at (2.55,1.45) {termination\\at the pad};
  % optional AC coupling
  \draw[region] (4.0,-2.3) rectangle (7.7,2.3);
  \node[lbl,anchor=north] at (5.85,2.25) {optional on-chip AC coupling};
  \draw (4.0,1.6) to[C, l=$C_\mathrm{AC}$] (5.6,1.6) -- (6.6,1.6);
  \draw (4.0,-1.6) to[C, l_=$C_\mathrm{AC}$] (5.6,-1.6) -- (6.6,-1.6);
  \draw (6.2,1.6) to[R, l=$R_B$, *-*] (6.2,0) to[R, l=$R_B$, -*] (6.2,-1.6);
  \draw (6.2,0) -- (7.1,0) node[right,text=navy]{$V_\mathrm{CM}$};
  \draw (6.6,1.6) to[short,-o] (8.4,1.6) node[right]{to CTLE$^+$};
  \draw (6.6,-1.6) to[short,-o] (8.4,-1.6) node[right]{to CTLE$^-$};
\end{tikzpicture}
\caption{\textbf{RX input network.} Two \qty{50}{\ohm} resistors terminate the differential signal at the pad; the AC-grounded center tap also absorbs common-mode noise. An optional AC-coupling stage ($C_\mathrm{AC}$, $R_B$) lets the RX set its own common mode $V_\mathrm{CM}$, at the cost of baseline wander for long runs.}
\label{fig:ch10-term}
\end{figure}

%% =====================================================================
\section{Samplers}\label{sec:sampler}

A \textbf{sampler} (clocked comparator, or slicer) looks at its input at one instant and amplifies the sign of the difference to a full logic level. It is a positive-feedback latch, and everything we want from it follows from that. For \OS{} it must deliver:
\begin{itemize}
  \item \textbf{Sensitivity:} a correct decision at $10^{-12}$ for inputs down to about \qty{20}{\milli\volt}pp, the value budgeted in \cref{sec:noisebudget};
  \item \textbf{Offset} below about \qty{5}{\milli\volt} after calibration (\cref{sec:offset});
  \item \textbf{Aperture} well below 1\,UI, so that it samples a point rather than averaging the eye, and a sampling bandwidth well above $\fN$;
  \item \textbf{Regeneration} fast enough that even tiny inputs resolve in time (\cref{sec:metastability}), and a clock-to-output delay that fits the DFE loop of \cref{sec:dfe-timing};
  \item \textbf{Low kick-back:} the large internal swings must not couple back into the CTLE and disturb the other three samplers;
  \item \textbf{Low noise:} at least the $kT/C$ noise of its internal nodes, $\sqrt{kT/C}=\qty{0.45}{\milli\volt}$ at \qty{20}{\femto\farad}.
\end{itemize}

\subsection{The StrongARM latch}

The \textbf{StrongARM} latch (\cref{fig:ch10-sa-sch}) is the standard CMOS comparator~\cite{p4:kobayashi1993,p4:razavi2015sa}. It draws no static current and produces rail-to-rail outputs. Its operation has four phases:
\begin{enumerate}
  \item \textbf{Reset} (CLK low): $S_1$--$S_4$ precharge out$^\pm$, X and X$'$ to $V_\mathrm{DD}$, and $M_7$ is off.
  \item \textbf{Integration} (CLK rises): $M_7$ turns on. The input pair $M_1/M_2$ discharges X and X$'$, then the output nodes, with currents $I_D\pm\gm\,v_\mathrm{in}/2$. The common mode falls and a differential voltage grows linearly.
  \item \textbf{Regeneration:} when the outputs have fallen by $\lvert V_\mathrm{TP}\rvert$, the PMOS $M_5/M_6$ turn on. The two cross-coupled inverters ($M_3/M_5$, $M_4/M_6$) now amplify the difference exponentially. The input pair is in triode and no longer sees the input.
  \item \textbf{Decision:} one output returns to $V_\mathrm{DD}$ and the other falls to ground. An SR latch holds the result while the next reset begins~\cite{p4:nikolic2000}.
\end{enumerate}

\begin{figure}[tbp]
\centering
\begin{tikzpicture}[american, scale=0.68, transform shape, font=\small]
  \draw (-4.6,4.2) -- (9.0,4.2); \draw (2.2,4.2) node[vcc]{$V_\mathrm{DD}=\qty{1.2}{\volt}$};
  \draw (0,3.1) node[pmos,xscale=-1](M5){};
  \draw (4.4,3.1) node[pmos](M6){};
  \draw (0,1.2) node[nmos,xscale=-1](M3){};
  \draw (4.4,1.2) node[nmos](M4){};
  \draw (0,-0.7) node[nmos](M1){};
  \draw (4.4,-0.7) node[nmos,xscale=-1](M2){};
  \draw (2.2,-2.6) node[nmos](M7){};
  \draw (-1.8,3.1) node[pmos](S1){};
  \draw (6.2,3.1) node[pmos,xscale=-1](S2){};
  \draw (-3.6,1.0) node[pmos](S3){};
  \draw (8.0,1.0) node[pmos,xscale=-1](S4){};
  \foreach \m/\t/\dx in {M5/M_5/-0.55,M6/M_6/0.55,M3/M_3/-0.55,M4/M_4/0.55,M1/M_1/0.55,M2/M_2/-0.55,M7/M_7/0.55,S1/S_1/0.5,S2/S_2/-0.5,S3/S_3/0.5,S4/S_4/-0.5}{\node[text=navy] at ($(\m.center)+(\dx,0.1)$) {$\t$};}
  \draw (M5.S) -- (0,4.2); \draw (M6.S) -- (4.4,4.2); \draw (S1.S) -- (-1.8,4.2); \draw (S2.S) -- (6.2,4.2);
  \draw (S3.S) -- (-3.6,4.2); \draw (S4.S) -- (8.0,4.2);
  \draw (M5.D) -- (M3.D); \draw (M6.D) -- (M4.D);
  \draw (M3.S) -- (M1.D); \draw (M4.S) -- (M2.D);
  \draw (0,2.15) node[circ]{}; \draw (4.4,2.15) node[circ]{};
  \draw (S1.D) |- (0,2.15); \draw (S2.D) |- (4.4,2.15);
  \draw (0,0.27) node[circ]{} node[above left]{X}; \draw (4.4,0.27) node[circ]{} node[above right]{X$'$};
  \draw (S3.D) |- (0,0.27); \draw (S4.D) |- (4.4,0.27);
  \draw (M5.G) -- (4.4,2.15); \draw (M6.G) -- (0,2.15);
  \draw (M3.G) -- (4.4,2.15); \draw (M4.G) -- (0,2.15);
  \draw (-1.8,2.15) to[short,*-o] (-2.6,2.15) node[below]{$\mathrm{out}^-$};
  \draw (6.2,2.15) to[short,*-o] (7.0,2.15) node[below]{$\mathrm{out}^+$};
  \draw (M1.S) -- (0,-1.7) -- (4.4,-1.7) -- (M2.S);
  \draw (2.2,-1.7) node[circ]{} -- (M7.D);
  \draw (M7.S) -- (2.2,-3.6) node[ground]{};
  \draw (M1.G) to[short,-o] ++(-0.4,0) node[left]{$V_\mathrm{in}^+$};
  \draw (M2.G) to[short,-o] ++(0.4,0) node[right]{$V_\mathrm{in}^-$};
  \draw (M7.G) -- ++(-0.3,0) node[left]{CLK};
  \foreach \s in {S1,S3}{\draw (\s.G) -- ++(-0.25,0) node[left,font=\footnotesize]{CLK};}
  \foreach \s in {S2,S4}{\draw (\s.G) -- ++(0.25,0) node[right,font=\footnotesize]{CLK};}
\end{tikzpicture}
\caption{\textbf{The StrongARM latch.} The input pair $M_1/M_2$ integrates the input onto X/X$'$ and the outputs. The cross-coupled inverters $M_3/M_5$ and $M_4/M_6$ (crossing wires) regenerate it, and $S_1$--$S_4$ precharge every internal node while CLK is low. The outputs drive an SR latch (not shown). No current flows once a decision is made.}
\label{fig:ch10-sa-sch}
\end{figure}

\begin{derivation}[Regeneration time constant]
In regeneration each inverter, biased near its switching point, is a transconductance $\gm$ (NMOS plus PMOS) driving its own output capacitance $C$ from the other output: $C\,\dot v_1=-\gm v_2$ and $C\,\dot v_2=-\gm v_1$. Subtracting, the difference $\Delta V=v_1-v_2$ obeys $C\,\dot{\Delta V}=\gm\Delta V$, a growing exponential. The integration phase sets its starting value: the outputs fall at $I_D/C$, so regeneration begins after $t_1=C\lvert V_\mathrm{TP}\rvert/I_D$, and by then the input pair has built $\Delta V(t_1)=g_\mathrm{m,in}v_\mathrm{in}t_1/C=A_\mathrm{int}\,v_\mathrm{in}$.
\end{derivation}

\begin{keyeq}[StrongARM regeneration and decision time]
\begin{equation}
  \Delta V(t)=A_\mathrm{int}\,v_\mathrm{in}\,e^{(t-t_1)/\tau},\qquad \tau=\frac{C}{\gm},\qquad
  t_\mathrm{dec}=t_1+\tau\ln\frac{V_L}{A_\mathrm{int}\,v_\mathrm{in}} .
  \label{eq:ch10-sa}
\end{equation}
Each factor of 10 less input adds $\tau\ln10=2.3\tau$ to the decision time. $V_L$ is the swing the next stage needs.
\end{keyeq}

\Cref{fig:ch10-sa} shows a behavioral model with the numbers of \cref{we:ch10-meta}: $C=\qty{20}{\femto\farad}$, $I_D=\qty{200}{\micro\ampere}$ and $\lvert V_\mathrm{TP}\rvert=0.4$\,V give $t_1=\qty{40}{\pico\second}$ and $A_\mathrm{int}=3$; the cross-coupled pair gives $\tau=20\,\mathrm{fF}/1.5\,\mathrm{mS}\approx\qty{13}{\pico\second}$. A 1\,mV input decides in 109\,ps, a 10\,\textmu V input in 169\,ps.

\begin{figure}[tbp]
\centering
\includegraphics{ch10_sa}
\caption{\textbf{A latch amplifies exponentially: halving the input only adds a fixed delay.} Behavioral StrongARM model ($t_1=40$\,ps, $A_\mathrm{int}=3$, $\tau=13$\,ps). (a) Output nodes for 10\,mV and 0.1\,mV inputs: both fall together during integration, then split. (b) Output difference on a log scale for inputs from 100\,mV to 10\,\textmu V. The straight lines have slope $1/\tau$, and each decade of input costs 30\,ps. Inputs below about 4\,\textmu V miss the 180\,ps deadline.}
\label{fig:ch10-sa}
\end{figure}

\subsection{The HBT CML latch}

The CML latch (\cref{fig:ch10-cml}) is the bipolar alternative~\cite{palermo720ckt}. While CK is high, the track pair $Q_1/Q_2$ is an ordinary CML buffer and the outputs follow the input with time constant $\RL C$. When CK falls, the tail current switches to the cross-coupled pair $Q_3/Q_4$, which holds and amplifies whatever difference the outputs carried at that instant. Two latches on opposite clock phases form a master--slave flip-flop.

\begin{figure}[tbp]
\centering
\begin{tikzpicture}[american, scale=0.66, transform shape, font=\small]
  \node[lblb,font=\sffamily\normalsize\bfseries] at (-1.6,4.0) {(a)};
  \draw (-0.3,4.0) -- (7.7,4.0); \draw (3.7,4.0) node[vcc]{$V_\mathrm{CC}=\qty{2.5}{\volt}$};
  \draw (0,0) node[npn](Q1){};
  \draw (2.0,0) node[npn,xscale=-1](Q2){};
  \draw (4.4,0) node[npn,xscale=-1](Q3){};
  \draw (7.4,0) node[npn](Q4){};
  \draw (1.0,-2.3) node[npn](Q5){};
  \draw (5.9,-2.3) node[npn,xscale=-1](Q6){};
  \foreach \q/\t/\dx in {Q1/Q_1/0.5,Q2/Q_2/-0.5,Q3/Q_3/-0.5,Q4/Q_4/0.5,Q5/Q_5/0.5,Q6/Q_6/-0.5}{\node[text=navy] at ($(\q.center)+(\dx,0.05)$) {$\t$};}
  \draw (0,2.0) to[R, l=$\RL$] (0,4.0);
  \draw (7.4,1.4) to[R, l_=$\RL$] (7.4,4.0);
  \draw (Q1.C) -- (0,2.0) to[short,*-] (4.4,2.0) -- (Q3.C);
  \draw (Q2.C) -- (2.0,1.4) to[short,*-] (4.25,1.4) arc[start angle=180,end angle=0,radius=0.15] -- (7.4,1.4) node[circ]{} -- (Q4.C);
  \draw (0,2.0) to[short,-o] (-1.0,2.0) node[left]{$\mathrm{out}^-$};
  \draw (7.4,1.4) to[short,-o] (8.4,1.4) node[right]{$\mathrm{out}^+$};
  \draw (Q3.B) -- (7.4,1.05) node[circ]{};
  \draw (Q4.B) -- (4.4,1.05) node[circ]{};
  \draw (Q1.B) to[short,-o] ++(-0.45,0) node[left]{$D^+$};
  \draw (Q2.B) to[short,-o] ++(0.45,0) node[below]{$D^-$};
  \draw (Q1.E) -- (0,-1.15) -- (2.0,-1.15) -- (Q2.E);
  \draw (1.0,-1.15) node[circ]{} -- (Q5.C);
  \draw (Q3.E) -- (4.4,-1.15) -- (7.4,-1.15) -- (Q4.E);
  \draw (5.9,-1.15) node[circ]{} -- (Q6.C);
  \draw (Q5.B) to[short,-o] ++(-0.45,0) node[left]{$\mathrm{CK}^+$};
  \draw (Q6.B) to[short,-o] ++(0.45,0) node[right]{$\mathrm{CK}^-$};
  \draw (Q5.E) -- (1.0,-3.45) -- (5.9,-3.45) -- (Q6.E);
  \draw (3.45,-3.45) node[circ]{} to[I, l=$I_T$] (3.45,-5.0) node[ground]{};
  \node[lblb] at (1.0,-0.5) {track};
  \node[lblb] at (5.9,-0.6) {latch};
  \begin{scope}[xshift=10.8cm,yshift=-0.5cm]
    \node[lblb,font=\sffamily\normalsize\bfseries,anchor=west] at (-0.6,2.4) {(b) master--slave};
    \node[blkb,minimum width=14mm,minimum height=12mm] (L1) at (0.8,0.6) {latch\\(master)};
    \node[blkb,minimum width=14mm,minimum height=12mm] (L2) at (3.6,0.6) {latch\\(slave)};
    \draw[sig] (-0.6,0.6) node[left,text=navy]{$D$} -- (L1.west);
    \draw[sig] (L1.east) -- (L2.west);
    \draw[sig] (L2.east) -- (5.2,0.6) node[right,text=navy]{$Q$};
    \draw[clk] (0.8,-1.0) node[below,text=golddk]{$\overline{\mathrm{CK}}$} -- (L1.south);
    \draw[clk] (3.6,-1.0) node[below,text=golddk]{CK} -- (L2.south);
  \end{scope}
\end{tikzpicture}
\caption{\textbf{The HBT CML latch.} (a) With CK high the track pair $Q_1/Q_2$ drives the loads; with CK low the tail current moves to the cross-coupled pair $Q_3/Q_4$ (crossing wires), which regenerates the stored difference. The swing is $I_T\RL$ and the supply current never changes. (b) A master latch on $\overline{\mathrm{CK}}$ followed by a slave on CK forms a flip-flop that samples on the rising edge.}
\label{fig:ch10-cml}
\end{figure}

\begin{keyeq}[CML latch regeneration]
\begin{equation}
  C\,\dot{\Delta V}=\Big(\gm-\frac{1}{\RL}\Big)\Delta V\qquad\Longrightarrow\qquad \tau=\frac{\RL C}{\gm\RL-1},\qquad \gm=\frac{I_T/2}{V_T}.
  \label{eq:ch10-cml}
\end{equation}
The cross-coupled pair adds $\gm\Delta V$ of current and the loads remove $\Delta V/\RL$; regeneration needs a loop gain $\gm\RL>1$.
\end{keyeq}

\paragraph{Why it is fast in \sgx{}} With $I_T=\qty{1}{\milli\ampere}$, each side carries 0.5\,mA and $\gm=\qty{19.3}{\milli\siemens}$; $\RL=\qty{400}{\ohm}$ gives a \qty{400}{\milli\volt} swing and $\gm\RL=7.7$. With $C=\qty{30}{\femto\farad}$, \cref{eq:ch10-cml} gives $\tau=\qty{1.8}{\pico\second}$, and 3--\qty{5}{\pico\second} once base resistance and $C_\mu$ are included. That is three to seven times faster than the StrongARM, because the HBT's $\gm/I=1/V_T$ and its 300\,GHz $f_T$ give large $\gm$ with little capacitance. The latch also kicks back less, since its current never changes, and its sampling instant does not move with the input common mode. Its cost is static power: \qty{2.5}{\milli\watt} per latch from 2.5\,V.

\paragraph{The \OS{} choice} All four samplers are HBT CML master--slave flip-flops. The data samplers must, because the direct-feedback DFE needs a clock-to-output delay of 25--35\,ps (\cref{sec:dfe-timing}); a 1.2\,V StrongARM with its SR latch needs 100\,ps or more for small inputs. The edge samplers use the same circuit so that data and edge samplers have the same aperture delay, as the next paragraph shows. A StrongARM remains the right choice where a decision can wait, such as the DFE's error sampler (\cref{sec:adaptation}) or slow monitors.

\paragraph{Aperture} A real sampler does not take an instantaneous sample. Its output is a weighted average $\int v_\mathrm{in}(t)h(t)\,dt$ with $\int h=1$, where $h(t)$ is the \textbf{sampling function}~\cite{p4:kim2009}. The \textbf{aperture} $w_{80}$ is the time that holds the central 80\,\% of the weight. Its Fourier transform is the sampler's own bandwidth. \Cref{fig:ch10-aperture} compares the two latches. The StrongARM integrates during the 40\,ps after the clock edge; the CML latch remembers the input with its track time constant $\RL C=\qty{12}{\pico\second}$ \emph{before} the edge. Both apertures are about 25--30\,ps (0.12--0.14\,UI), and both bandwidths, 12--15\,GHz, cost less than 0.2\,dB at $\fN$. But their centroids differ by 23\,ps. A StrongARM edge sampler next to CML data samplers would make the CDR lock 0.12\,UI away from the true eye center. Keep $w_{80}\le0.2$\,UI with sharp clock edges, and use one sampler type for data and edge.

\begin{figure}[tbp]
\centering
\includegraphics{ch10_aperture}
\caption{\textbf{A sampler averages its input over its aperture.} (a) Sampling functions $h(t)$ of a StrongARM (weight after the clock edge, during integration) and an HBT CML latch (weight before the edge, through the track time constant), with their cumulative weights. Both have $w_{80}\approx$ 25--30\,ps, but their centroids are 23\,ps apart. (b) The same functions as filters: both pass $\fN$ with less than 0.2\,dB of loss.}
\label{fig:ch10-aperture}
\end{figure}

%% =====================================================================
\section{Metastability}\label{sec:metastability}

If the input is close enough to zero, the latch has not reached a valid logic level $V_L$ when its time $T$ runs out. It is then \textbf{metastable}, and downstream logic may read it either way, or both ways at once. By \cref{eq:ch10-sa} the latch fails only if its starting value is below $V_Le^{-T/\tau}$, that is, if $\lvert v_\mathrm{in}\rvert<V_Le^{-T/\tau}/A$, where $A$ is the gain in front of the regenerative node.

\begin{keyeq}[Metastability error probability]
\begin{equation}
  P_\mathrm{meta}=\frac{2V_L}{A\,\Delta V}\,e^{-T/\tau},\qquad\text{and for cascaded stages}\quad P_\mathrm{meta}\propto e^{-(T_1/\tau_1+T_2/\tau_2+\dots)} .
  \label{eq:ch10-meta}
\end{equation}
The input is assumed spread uniformly over a range $\Delta V$; the fraction that lands in the failing window of width $2V_Le^{-T/\tau}/A$ is $P_\mathrm{meta}$. Each latch stage behind the first multiplies the probability by its own $e^{-T/\tau}$.
\end{keyeq}

\begin{example}{StrongARM metastability, data versus edge samplers}\label{we:ch10-meta}
\textbf{Timing:} the half-rate clock period is 400\,ps, so evaluation lasts 200\,ps. Integration takes $t_1=\qty{40}{\pico\second}$ and the SR latch needs about 20\,ps, which leaves $T=\qty{140}{\pico\second}$ of regeneration with $\tau=\qty{13}{\pico\second}$: a gain of $e^{10.8}=4.8\times10^4$.

\textbf{Data sampler:} reaching $V_L=\qty{0.6}{\volt}$ needs $0.6/4.8\times10^4=\qty{12}{\micro\volt}$ at the latch, or \qty{4}{\micro\volt} at the input with $A_\mathrm{int}=3$. An equalized eye never puts the data sample that close to zero, so data samplers are safe.

\textbf{Edge sampler:} it samples the zero crossing \emph{on purpose}, so its input is often tiny. Take the input spread over $\pm\qty{100}{\milli\volt}$ ($\Delta V=\qty{0.2}{\volt}$) and $A=3$: $P_\mathrm{meta}=(2\cdot0.6/0.6)\,e^{-10.8}=\mathbf{4.2\times10^{-5}}$ per sample, or $10^5$ undefined decisions per second at 2.5\,GS/s. That is far too many for a digital CDR loop.

\textbf{Fix:} two more latch stages, each with 200\,ps and $\tau\approx\qty{20}{\pico\second}$, multiply by $e^{-20}$: $P_\mathrm{meta}=\mathbf{8.7\times10^{-14}}$, below the BER target. A metastable edge decision then costs one bang-bang vote, which the loop filter averages away.

\textbf{HBT CML latch:} with $\tau\approx\qty{4}{\pico\second}$ and the same 140\,ps, $P_\mathrm{meta}\approx4\times10^{-16}$ in the master latch alone (with $V_L=\qty{0.2}{\volt}$).
\end{example}

\begin{figure}[tbp]
\centering
\includegraphics{ch10_meta}
\caption{\textbf{Metastability falls exponentially with regeneration time, and every extra latch stage adds its own exponent.} \Cref{eq:ch10-meta} with $\Delta V=0.2$\,V and $A=3$. A lone StrongARM leaves $4.2\times10^{-5}$ at the 140\,ps of \cref{we:ch10-meta}; two extra CMOS latches bring it below $10^{-12}$. The HBT CML latch crosses $10^{-12}$ after about 110\,ps.}
\label{fig:ch10-meta}
\end{figure}

\begin{insight}[Edge samplers need extra stages]
A data sampler sees an open eye; an edge sampler is aimed at the zero crossing, so it is metastable thousands of times more often. Always put at least two more latch stages between an edge sampler and the CDR logic, and count their regeneration time. In \OS{} the CML slave latch and the first CMOS flip-flops of the deserializer provide them.
\end{insight}

%% =====================================================================
\section{Offset and calibration}\label{sec:offset}

Two transistors drawn identically are not identical. Random dopant and edge fluctuations give the threshold voltages of a matched pair a random difference whose standard deviation falls with gate area~\cite{p4:pelgrom1989}. A sampler with mismatched input devices therefore has an input-referred \textbf{offset} $V_\mathrm{os}$: it decides as if its threshold were $V_\mathrm{os}$ instead of zero.

\begin{keyeq}[Pelgrom mismatch and calibration sizing]
\begin{equation}
  \sigma_{V_\mathrm{os}}\approx\frac{A_{VT}}{\sqrt{WL}},\qquad
  \text{range}=\pm3\sigma_{V_\mathrm{os}},\qquad
  \mathrm{LSB}=\frac{6\sigma_{V_\mathrm{os}}}{2^n},\qquad \lvert V_\mathrm{os,res}\rvert\le\frac{\mathrm{LSB}}{2}.
  \label{eq:ch10-pelgrom}
\end{equation}
Halving $\sigma$ needs four times the area, and so four times the input capacitance. A DAC is cheaper.
\end{keyeq}

\begin{example}{Sizing the offset DAC}\label{we:ch10-offset}
\textbf{Given:} an input pair of $W/L=\qty{4}{\micro\meter}/\qty{0.13}{\micro\meter}$ ($WL=\qty{0.52}{\micro\meter\squared}$) and $A_{VT}\approx\qty{3.5}{\milli\volt\micro\meter}$. That $A_{VT}$ is an assumption for 130\,nm; extract the real value by Monte Carlo if the \sgx{} models include mismatch.

\textbf{Offset:} $\sigma=3.5/\sqrt{0.52}=\qty{4.9}{\milli\volt}$, so $\pm3\sigma=\pm\qty{14.6}{\milli\volt}$. Untrimmed, only 70\,\% of samplers meet $\lvert V_\mathrm{os}\rvert<\qty{5}{\milli\volt}$, and all four in one receiver do so on only 24\,\% of chips (\cref{fig:ch10-offset}).

\textbf{DAC:} a residual of $\pm\qty{1}{\milli\volt}$ needs $\mathrm{LSB}\le\qty{2}{\milli\volt}$, so $n\ge\log_2(29.2/2)=3.9$. Choose \textbf{5 bits}: LSB = \qty{0.91}{\milli\volt}, residual $\pm\qty{0.46}{\milli\volt}$. After trim, 99.8\,\% of samplers are within 1\,mV.

\textbf{The no-calibration alternative} needs $\sigma=\qty{1}{\milli\volt}$, or $WL=\qty{12.2}{\micro\meter\squared}$: 24 times the area, with the capacitance, kick-back and slower $\tau$ that come with it.
\end{example}

\begin{figure}[tbp]
\centering
\includegraphics{ch10_offset}
\caption{\textbf{Mismatch spreads sampler offsets over tens of millivolts; a 5-bit DAC pulls them inside half a millivolt.} Monte Carlo of 20\,000 samplers with $\sigma=4.9$\,mV (\cref{we:ch10-offset}). (a) Distribution before and after trimming with a $\pm3\sigma$, 5-bit DAC. (b) Fraction of samplers within a given offset. Untrimmed, 70\,\% of single samplers but only 24\,\% of four-sampler receivers meet the 5\,mV budget of \cref{sec:noisebudget}.}
\label{fig:ch10-offset}
\end{figure}

\paragraph{Calibrating the \OS{} receiver} The HBT pairs of the CML samplers match better than MOS pairs, but the CTLE, the summer and the load resistors add offsets of their own. \OS{} therefore trims the whole path, one DAC per sampler, by injecting a differential current into the DFE summer load (\cref{fig:ch10-cal}). The total current stays constant, so the common mode does not move. Referred to the input with unity front-end gain, one LSB of 0.91\,mV across the \qty{250}{\ohm} summer load is \qty{3.6}{\micro\ampere}, and the full $\pm16$\,LSB range is $\pm\qty{58}{\micro\ampere}$. The procedure is a short digital loop: short the inputs to the common mode, count ones and zeros for a few thousand cycles, step the code toward 50/50, repeat until it toggles by one LSB, then freeze it. Capacitive trim on the StrongARM's internal nodes is the alternative for CMOS samplers; it costs aperture time.

\begin{figure}[tbp]
\centering
\begin{tikzpicture}[x=1cm,y=1cm,font=\sffamily\footnotesize,
  b/.style={blk,minimum height=10mm,font=\sffamily\footnotesize,inner sep=2pt}]
  \node[b,fill=btint,minimum width=22mm] (fe) at (2.3,0) {CTLE +\\DFE summer};
  \node[b,fill=btint,minimum width=16mm] (smp) at (5.2,0) {sampler};
  \node[b,fill=gtint,draw=golddk,minimum width=30mm] (cnt) at (8.6,0) {count 1s vs 0s\\over $N$ cycles; code $\pm1$};
  \node[b,fill=gtint,draw=golddk,minimum width=24mm] (dac) at (5.2,-1.75) {5-bit offset DAC\\$I_\mathrm{os}=(\mathrm{code}-16)\,I_\mathrm{LSB}$};
  \draw[navy,line width=0.75pt] (-0.4,0.25) node[left,text=navy]{$\mathrm{in}^\pm$} -- (0.2,0.25);
  \draw[navy,line width=0.75pt] (0.2,0.25) -- (0.62,0.5);
  \node[dot] at (0.2,0.25) {};
  \node[dot] (sw) at (0.7,0.25) {};
  \draw[navy,line width=0.75pt] (0.7,0.25) -- (fe.west |- 0,0.25);
  \draw[navy,line width=0.75pt] (0.7,-0.25) -- (fe.west |- 0,-0.25);
  \draw[navy,line width=0.75pt] (0.7,0.25) -- (0.7,-0.25) node[dot]{};
  \node[lbl,anchor=north] at (0.45,-0.35) {$S_\mathrm{cal}$: short to $V_\mathrm{CM}$};
  \draw[sig] (fe) -- (smp);
  \draw[sig] (smp) -- (cnt);
  \draw[sigg] (cnt.south) |- (dac.east);
  \draw[sigg] (dac.west) -| node[pos=0.75,left,text=golddk,font=\sffamily\scriptsize]{current into the summer load} (fe.south);
  \draw[sig] (cnt.east) -- (11.6,0) node[right,text=navy,align=left]{freeze code\\when it toggles};
\end{tikzpicture}
\caption{\textbf{Offset calibration is a slow digital loop.} With the inputs shorted, a sampler with zero offset outputs ones and zeros equally often. The counter steps a 5-bit current DAC in the summer load until the code toggles by one LSB (0.91\,mV, or \qty{3.6}{\micro\ampere} into \qty{250}{\ohm}), then freezes it. Running it at power-up and after large temperature changes keeps the residual offset inside the budget.}
\label{fig:ch10-cal}
\end{figure}

\begin{insight}
Do not buy offset with area. Use minimum-capacitance input devices for speed, calibrate the whole path with a 5-bit DAC sized for $\pm3\sigma$, and leave the 5\,mV of the system budget for drift between calibrations. Simulate offset with Monte Carlo mismatch, never with ideal devices.
\end{insight}

%% =====================================================================
\section{Noise and the sensitivity budget}\label{sec:ch10-noise}

\Cref{sec:noisebudget} budgeted the receiver with three numbers: $\sign=\qty{1.5}{\milli\volt}$ rms input-referred noise, $V_\mathrm{os}=\qty{5}{\milli\volt}$ after calibration, and a sampler sensitivity $V_\mathrm{sens}=\qty{20}{\milli\volt}$. Here we check them against the circuits. The sensitivity of a sampler collects its own Gaussian noise $\sigma_s$ at the BER multiplier, the minimum overdrive $V_\mathrm{min}$ that resolves in time (metastability, hysteresis and kick-back memory), and its residual offset~\cite{palermo720ckt}:

\begin{keyeq}[Sampler sensitivity]
\begin{equation}
  V_\mathrm{sens}\approx 2Q\,\sigma_s+V_\mathrm{min}+2\lvert V_\mathrm{os,res}\rvert,\qquad 2Q=14.07\ \text{at}\ 10^{-12}.
  \label{eq:ch10-sens}
\end{equation}
\end{keyeq}

\begin{example}{Circuit-level noise and sensitivity budget}\label{we:ch10-noise}
All terms are rms and referred to the RX input, conservatively with unity front-end gain.
\begin{center}
\small
\begin{tabular}{@{}lrl@{}}
\toprule
Source & $\sigma$ (mV) & How it was estimated\\
\midrule
Termination and line, \qty{50}{\ohm} diff. & 0.10 & 0.91\,nV/$\sqrt{\mathrm{Hz}}$ over an 11.8\,GHz noise bandwidth\\
CTLE (shot noise, $R_E$, $R_C$) & 0.25 & \cref{sec:ctle-hbt} densities; confirm with \code{.noise}\\
DFE summer loads & 0.32 & $\sqrt{2kT/C}$, $C=\qty{80}{\femto\farad}$\\
Sampler, input-referred & 1.0 & $kT/C$ floor 0.45--0.64\,mV; 0.7--0.85\,mV measured~\cite{p4:kim2009}\\
\midrule
\textbf{Total (RSS)} & \textbf{1.08} & system budget assumes 1.5\\
\bottomrule
\end{tabular}
\end{center}
\textbf{Sensitivity:} with $\sigma_s=\qty{1}{\milli\volt}$, $V_\mathrm{min}\approx\qty{5}{\milli\volt}$ and the 5-bit residual of \qty{0.46}{\milli\volt}, \cref{eq:ch10-sens} gives $14.1+5+0.9=\mathbf{20}$\,mV, the value of \cref{sec:noisebudget}.

\textbf{Margin:} with the circuit noise, the linear budget of \cref{sec:noisebudget} becomes $131-14.07\times1.08-3.5-10-20=\mathbf{82}$\,mV instead of 76\,mV. Since $V_\mathrm{sens}$ already contains the sampler's own noise and residual offset, the system budget counts them twice and is about 15\,mV conservative.
\end{example}

The budget follows Palermo's rule for noise sources~\cite{palermo720ckt}: classify each one as bounded or Gaussian, and as proportional to swing or independent of it. Thermal noise, offset and sensitivity are independent of swing, so they shrink in relative terms as the eye grows. Residual ISI and crosstalk are proportional to swing, and no TX amplitude can beat them. With 82\,mV of margin, the receiver circuits are not the weak link; equalization is.

%% =====================================================================
\section{The deserializer}\label{sec:deserializer}

The first 1:2 of the deserializer is free: a half-rate receiver already delivers even and odd bits on separate samplers. The rest mirrors the transmitter (\cref{fig:ch10-deser}). Each sampler's 2.5\,Gb/s output is converted from CML to CMOS levels and enters a 1:8 tree of CMOS 1:2 stages clocked at 1.25\,GHz, 625\,MHz and 312.5\,MHz. That is 400\,ps per bit at the first CMOS flip-flop, or 9.5\,FO4, which a 1.2\,V CMOS flip-flop handles; a CMOS 1:2 demux is also far less sensitive to duty cycle than the TX mux. The two data lanes interleave into the 16-bit word, and the two edge lanes give 16 edge bits for the CDR logic. The clocks come from the recovered 2.5\,GHz clock through $\div2$ stages, so the whole tree lives in the \textbf{recovered clock domain}.

\begin{figure}[tbp]
\centering
\begin{tikzpicture}[x=1cm,y=1cm,font=\sffamily\footnotesize,
  b/.style={blk,minimum width=11mm,minimum height=8mm,font=\sffamily\footnotesize,inner sep=1pt},
  rate/.style={font=\sffamily\scriptsize,text=golddk,align=center}]
  \draw[region,fill=gtint!35] (2.0,-2.6) rectangle (10.15,1.9);
  \node[lblb,anchor=north west] at (2.05,1.85) {CMOS, 1.2\,V, recovered clock domain};
  \foreach \y/\n/\t in {0.7/a/{$D_0$}, -0.7/b/{$D_1$}}{
    \node[blk,fill=btint,minimum width=11mm,minimum height=8mm] (s\n) at (0.6,\y) {\t};
    \node[b,fill=btint] (c\n) at (2.75,\y) {CML$\to$\\CMOS};
    \node[b,fill=gtint,draw=golddk] (x\n) at (4.55,\y) {1:2};
    \node[b,fill=gtint,draw=golddk] (y\n) at (6.15,\y) {1:2};
    \node[b,fill=gtint,draw=golddk] (z\n) at (7.75,\y) {1:2};
    \draw[sig] (s\n) -- (c\n); \draw[sig] (c\n) -- (x\n); \draw[sig] (x\n) -- (y\n); \draw[sig] (y\n) -- (z\n);
  }
  \node[rate,anchor=south] at (1.65,0.75) {2.5\,Gb/s};
  \node[rate,anchor=south] at (5.35,0.75) {$2\times$1.25};
  \node[rate,anchor=south] at (6.95,0.75) {$4\times$625};
  \node[blk,fill=gtint,draw=golddk,minimum width=13mm,minimum height=22mm] (il) at (9.35,0) {interleave\\16\,b};
  \draw[sig] (za.east) -- (il.west |- 0,0.7); \draw[sig] (zb.east) -- (il.west |- 0,-0.7);
  \node[rate,anchor=south] at (8.6,0.75) {8\,b};
  \node[blk,fill=gtint,draw=golddk,minimum width=15mm,minimum height=12mm] (ff) at (11.6,0) {async\\FIFO};
  \draw[sig,line width=1.3pt] (il.east) -- (ff.west);
  \draw[sig,line width=1.3pt] (ff.east) -- (14.2,0) node[above left,text=navy,align=right]{16 b to core\\312.5\,MHz (PLL)};
  % edges
  \node[lbl,anchor=west,align=left] at (0.1,-1.55) {$E_0$, $E_1$: same 1:8 trees $\to$ 16 edge bits $\to$ CDR logic};
  % clock chain
  \node[blk,fill=navy,text=white,minimum width=13mm,minimum height=7mm] (rc) at (0.6,-3.3) {PI clock};
  \foreach \x/\n in {4.55/k1,6.15/k2,7.75/k3}{\node[blk,minimum width=8mm,minimum height=6mm,inner sep=1pt] (\n) at (\x-0.8,-3.3) {$\div2$};}
  \draw[clk] (rc) -- (k1); \draw[clk] (k1) -- (k2); \draw[clk] (k2) -- (k3); \draw[clk] (k3.east) -- (11.6,-3.3) -- (ff.south);
  \node[rate,anchor=north] at (2.4,-3.45) {2.5\,GHz};
  \node[rate,anchor=north] at (4.55,-3.45) {1.25\,GHz};
  \node[rate,anchor=north] at (6.15,-3.45) {625\,MHz};
  \node[rate,anchor=north] at (9.4,-3.45) {312.5\,MHz (recovered) writes};
  \foreach \x in {4.55,6.15,7.75}{\node[dot,fill=gold] at (\x,-3.3) {}; \draw[clk] (\x,-3.3) -- (\x,-1.1);}
\end{tikzpicture}
\caption{\textbf{The deserializer mirrors the serializer.} The half-rate samplers already provide the first 1:2. Each data lane goes from CML to CMOS and through three 1:2 stages to 8 bits at 312.5\,MHz; the two lanes interleave into the 16-bit word. The edge lanes do the same for the CDR. All clocks are divided from the recovered (PI) clock, and an asynchronous FIFO hands the word to the core clock.}
\label{fig:ch10-deser}
\end{figure}

\paragraph{Crossing into the core clock} The recovered 312.5\,MHz clock and the core's 312.5\,MHz clock, which comes from the local PLL, have the same frequency in near-end loopback, or when both ends share a reference, but an unknown and wandering phase. A small asynchronous FIFO, 4--8 words deep with Gray-coded pointers, absorbs that. If the far end has its own reference ($\pm\qty{100}{\ppm}$ each), the frequencies differ by up to \qty{200}{\ppm}: one word slips every 5000 words, or every \qty{16}{\micro\second}. That needs an elastic buffer with skip and insert symbols, which is a protocol feature. For bring-up, run the PRBS checker directly in the recovered domain, before the FIFO, so a BER measurement needs no clock crossing at all.

%% =====================================================================
\section{RX design checklist}\label{sec:ch10-checklist}

\Cref{tab:ch10-check} collects the RX checks and a first-cut power estimate of about 70\,mW, without the PLL, phase interpolator and CDR logic (\cref{sec:power-budget}).

\begin{table}[tbp]
\centering
\caption{\textbf{RX checklist and first-cut power.} CML currents from the worked examples here and in \cref{ch:ctle,ch:dfe}. Confirm every line in \tool{ngspice} with extracted parasitics.}
\label{tab:ch10-check}
\small
\begin{tabularx}{\linewidth}{@{}l L r@{}}
\toprule
Block & Check before layout & Power\\
\midrule
Termination, ESD, pad & \qty{100}{\ohm} diff., AC-grounded center tap; $C_p<0.85$\,pF; RL $\ge10$\,dB at $\fN$ & --\\
CTLE and buffer & linear range for the long-run level; 0--12\,dB peaking (2\,mA + 3\,mA) & 12.5\,mW\\
DFE summers ($\times2$), taps & settles within the tap-1 budget; offset DAC injection (6\,mA) & 15\,mW\\
4 CML samplers (M--S) & $\tau\le5$\,ps; $t_\mathrm{CK\to Q}\le35$\,ps; same type for data and edge (12\,mA) & 30\,mW\\
I/Q clock buffers & sharp edges ($w_{80}\le0.2$\,UI); phase error $\le\ang{5}$ (5.6\,ps) & 10\,mW\\
Offset DACs ($\times4$) & 5 bits, $\pm3\sigma$ range, LSB 0.91\,mV & 0.6\,mW\\
Deserializer (1.2\,V) & two extra latch stages behind $E_0$/$E_1$; FIFO or in-domain PRBS checker & 5\,mW\\
\midrule
\textbf{Total RX} & 27\,mA from 2.5\,V, 4\,mA from 1.2\,V & \textbf{$\approx$73\,mW}\\
\bottomrule
\end{tabularx}
\end{table}

\begin{takeaways}
\begin{itemize}
  \item The \OS{} receiver is termination, CTLE, DFE summer and four half-rate samplers (data at \ang{0}/\ang{180}, edges at \ang{90}/\ang{270}), followed by a 1:16 deserializer and the CDR logic.
  \item A latch amplifies exponentially, $\Delta V\propto e^{t/\tau}$: $\tau=C/\gm\approx13$\,ps for a 130\,nm StrongARM and $\tau=\RL C/(\gm\RL-1)\approx$ 2--5\,ps for an \sgx{} HBT CML latch.
  \item \OS{} uses HBT CML master--slave samplers everywhere: the DFE loop needs their 25--35\,ps clock-to-output, and one sampler type keeps the data and edge apertures aligned.
  \item $P_\mathrm{meta}=(2V_L/A\Delta V)\,e^{-T/\tau}$. A lone StrongARM edge sampler gives $4\times10^{-5}$; two extra latch stages bring it below $10^{-13}$.
  \item Pelgrom's law gives $\sigma\approx\qty{4.9}{\milli\volt}$ for a small 130\,nm pair; a 5-bit DAC over $\pm3\sigma$ leaves $\pm0.46$\,mV.
  \item The circuit noise is about 1.1\,mV rms and the sensitivity about 20\,mV, which supports the system budget with about 80\,mV of margin.
  \item Half rate gives the first 1:2 for free; CMOS does the rest, and a small FIFO crosses from the recovered to the core clock.
\end{itemize}
\end{takeaways}


\tutpart{Clocking}{Every bit is only as good as the clock edge that samples it. Jitter, clock generation with a phase-locked loop, and clock and data recovery.}
\tikzsetfigurename{11_jitter-}
\chapter{Jitter}\label{ch:jitter}
\chapterlead{A 200\,ps bit leaves no room for sloppy clock edges: this chapter turns ``the clock wobbles'' into numbers you can budget, measure and design against.}

\begin{objectives}
\begin{itemize}
  \item Define time-interval error, period jitter and cycle-to-cycle jitter, and say which one a serial link cares about.
  \item Convert voltage noise into timing jitter through the slew rate.
  \item Separate random from deterministic jitter and combine them with the dual-Dirac model, $\mathrm{TJ} = \mathrm{DJ}_{\delta\delta} + 2Q\,\sigma_\mathrm{RJ}$.
  \item Integrate a phase-noise spectrum $\mathcal{L}(f)$ into rms jitter with the right limits.
  \item Build the \OS{} jitter budget and check it against 0.3\,UI at $10^{-12}$.
\end{itemize}
\end{objectives}

%% =====================================================================
\section{What jitter is}\label{sec:jitter-defs}

An ideal clock has edges at exactly $t = nT$; a real clock has edges at $t_n$, each a little early or late. \textbf{Jitter} is that deviation. Three measures are in common use (\cref{fig:ch11-tie}a):
\begin{equation}
  \mathrm{TIE}_n = t_n - nT, \qquad
  J_{\mathrm{per},n} = T_n - T = \mathrm{TIE}_{n+1} - \mathrm{TIE}_n, \qquad
  J_{\mathrm{cc},n} = T_{n+1} - T_n ,
  \label{eq:ch11-tie}
\end{equation}
where $T_n = t_{n+1} - t_n$ is the $n$-th period. The \textbf{time-interval error} (TIE, also called absolute jitter) is the one a link cares about, because the sampler compares a clock edge with a data edge, not one clock period with the next. \textbf{Period jitter} and \textbf{cycle-to-cycle jitter} matter for digital setup and hold, but they can be tiny while the TIE is large, because small period errors add up.

\begin{figure}[tbp]
\centering
\includegraphics[width=\linewidth]{ch11_tie}
\caption{\textbf{Jitter is measured against an ideal clock, and it accumulates.} (a) TIE$_n$ is the offset of each edge from $nT$; the period $T_n$ differs from $T$ by the period jitter. (b) TIE of the team's 2.5\,GHz LC VCO (73\,fs rms period jitter, \cref{we:ch11-pn}). Free-running, the TIE is a random walk that grows as $\kappa\sqrt{t}$ and passes 10\,ps within 8\,\textmu{}s. Locked in a 3\,MHz PLL it stays near 0.6\,ps rms.}
\label{fig:ch11-tie}
\end{figure}

\Cref{fig:ch11-tie}(b) shows the difference. A free-running oscillator has no memory of where its edges should be, so its TIE random-walks. A phase-locked loop (\cref{ch:pll}) compares its edges with a clean reference and keeps the TIE bounded. The CDR (\cref{ch:cdr}) does the same for the data: it tracks slow jitter and leaves only the fast part to close the eye.

\begin{example}{Three jitter numbers from six edges}\label{we:ch11-defs}
\textbf{Given:} $T = 400$\,ps and edges at 0, 401.0, 799.5, 1201.5, 1600.0 and 2000.8\,ps.
\textbf{TIE:} 0, 1.0, $-0.5$, 1.5, 0, 0.8\,ps: 2.0\,ps peak-to-peak, 0.69\,ps rms about the mean.
\textbf{Period jitter:} 1.0, $-1.5$, 2.0, $-1.5$, 0.8\,ps: 1.42\,ps rms.
\textbf{Cycle-to-cycle:} $-2.5$, 3.5, $-3.5$, 2.3\,ps: 3.0\,ps rms.

\textbf{Lesson:} $J_\mathrm{cc} > J_\mathrm{per} > \mathrm{TIE}$ is typical of a locked clock: a bounded TIE means a long period is usually followed by a short one. Always say which measure a number is.
\end{example}

\paragraph{Voltage noise becomes timing noise}
Near a threshold crossing, a noise voltage moves the crossing by the noise divided by the slope of the edge (\cref{fig:ch11-slope}a):
\begin{keyeq}[Noise-to-jitter conversion]
\begin{equation}
  \sigma_t = \frac{\sigma_v}{\mathrm{SR}}, \qquad \mathrm{SR} = \left.\frac{dv}{dt}\right|_\text{crossing} = 2\pi f_0 A_c \ \text{for a sine of amplitude } A_c .
  \label{eq:ch11-slope}
\end{equation}
\end{keyeq}
A 2.5\,GHz CML clock of 0.4\,V$_\mathrm{ppd}$ ($A_c = 0.2$\,V) slews at \qty{3.14}{\milli\volt\per\pico\second}, so 1\,mV rms of noise costs 0.32\,ps rms. Attenuated to 0.1\,V$_\mathrm{ppd}$ by a lossy route, the same clock turns the same noise into 1.27\,ps (\cref{fig:ch11-slope}b). The relation also works in reverse: \cref{sec:eye} showed that sampling a few picoseconds off the eye centre costs millivolts.

\begin{figure}[tbp]
\centering
\includegraphics[width=\linewidth]{ch11_slope}
\caption{\textbf{Timing jitter is voltage noise divided by the slope.} (a) Noisy edges (noise exaggerated to 30\,mV rms) cross the threshold at spread-out times, with spread $\sigma_v/\mathrm{SR}$. (b) $\sigma_t$ against slew rate. Attenuating a 2.5\,GHz clock from 0.4 to 0.1\,V$_\mathrm{ppd}$ makes it four times more jitter-sensitive.}
\label{fig:ch11-slope}
\end{figure}

\begin{insight}[Design rule of thumb]
Keep clock swings large and edges fast wherever a clock is buffered or sampled: halving the slew rate doubles the jitter from the same noise. This is why \OS{} restores its 2.5\,GHz clocks with CML limiting buffers (\cref{sec:quadrature}) instead of letting them decay along long wires.
\end{insight}

%% =====================================================================
\section{Random and deterministic jitter}\label{sec:ch11-rjdj}

Jitter sources fall into two families with different statistics (\cref{fig:ch11-tree}).
\textbf{Random jitter} (RJ) comes from thermal, shot and flicker noise. It is Gaussian and \emph{unbounded}, and is described by its rms $\sigma_\mathrm{RJ}$; a peak-to-peak RJ is meaningless without a BER.
\textbf{Deterministic jitter} (DJ) is \emph{bounded}. Its parts are \textbf{data-dependent jitter} (DDJ), mainly ISI from the channel plus \textbf{duty-cycle distortion} (DCD); \textbf{periodic jitter} (PJ) from supply tones, reference spurs or spread-spectrum clocking, whose arcsine PDF piles up at its peaks; and \textbf{bounded uncorrelated jitter} (BUJ), mostly crosstalk.

\begin{figure}[tbp]
\centering
\begin{tikzpicture}[font=\sffamily\small,
  nd/.style={blk,minimum width=17mm,minimum height=6.5mm,inner sep=2pt},
  ds/.style={lbl,align=center,text width=30mm},
  ed/.style={navy,line width=0.7pt},
  gauss/.pic={\draw[c1,line width=0.9pt] plot[domain=-0.55:0.55,samples=41,smooth] (\x,{0.38*exp(-(\x/0.17)^2)}); \draw[muted,line width=0.4pt] (-0.6,0)--(0.6,0);},
  dirac2/.pic={\draw[muted,line width=0.4pt] (-0.6,0)--(0.6,0); \foreach \x in {-0.3,0.3}{\draw[golddk,line width=1pt,-{Stealth[length=1.4mm]}] (\x,0)--(\x,0.38);}},
  isi/.pic={\draw[muted,line width=0.4pt] (-0.6,0)--(0.6,0); \foreach \x/\h in {-0.45/0.16,-0.27/0.29,-0.09/0.38,0.09/0.38,0.27/0.29,0.45/0.16}{\draw[golddk,line width=1pt,-{Stealth[length=1.2mm]}] (\x,0)--(\x,\h);}},
  arcs/.pic={\draw[muted,line width=0.4pt] (-0.6,0)--(0.6,0); \draw[golddk,line width=0.9pt] plot[domain=-0.47:0.47,samples=41,smooth] (\x,{0.065/sqrt(1-(\x/0.5)^2)}); \draw[golddk,line width=0.9pt] (-0.47,0.19)--(-0.5,0.38) (0.47,0.19)--(0.5,0.38);},
  buj/.pic={\draw[muted,line width=0.4pt] (-0.6,0)--(0.6,0); \draw[golddk,line width=0.9pt] (-0.5,0)--(-0.35,0.28)--(0.35,0.28)--(0.5,0);}]
  \node[blkn,minimum width=26mm,minimum height=7mm] (tj) at (0,0) {Total jitter (TJ)};
  \node[nd,fill=btint] (rj) at (-5.3,-1.15) {Random (RJ)};
  \node[nd,fill=gtint,draw=golddk,minimum width=26mm] (dj) at (2.0,-1.15) {Deterministic (DJ)};
  \draw[ed] (tj.south) -- ++(0,-0.25) -| (rj.north);
  \draw[ed] (tj.south) -- ++(0,-0.25) -| (dj.north);
  \pic at (-5.3,-2.05) {gauss};
  \node[ds] at (-5.3,-2.65) {Gaussian, unbounded;\\thermal, shot, flicker;\\rms $\sigma_\mathrm{RJ}$, adds in RSS};
  \node[nd] (ddj) at (-0.9,-2.4) {DDJ};
  \node[nd] (pj) at (2.4,-2.4) {PJ};
  \node[nd] (bu) at (5.4,-2.4) {BUJ};
  \foreach \n in {ddj,pj,bu}{\draw[ed] (dj.south) -- ++(0,-0.3) -| (\n.north);}
  \node[nd,minimum width=12mm] (isi) at (-2.1,-3.55) {ISI};
  \node[nd,minimum width=12mm] (dcd) at (0.3,-3.55) {DCD};
  \draw[ed] (ddj.south) -- ++(0,-0.3) -| (isi.north);
  \draw[ed] (ddj.south) -- ++(0,-0.3) -| (dcd.north);
  \pic at (-2.1,-4.5) {isi};   \node[ds,text width=24mm] at (-2.1,-5.0) {channel loss,\\reflections};
  \pic at (0.3,-4.5) {dirac2}; \node[ds,text width=24mm] at (0.3,-5.0) {duty cycle,\\rise/fall mismatch};
  \pic at (2.4,-3.3) {arcs};   \node[ds] at (2.4,-3.85) {sinusoidal: supply tones,\\reference spurs, SSC};
  \pic at (5.4,-3.3) {buj};    \node[ds] at (5.4,-3.85) {uncorrelated with\\the data: crosstalk};
  \node[lbl,text=golddk,align=left,anchor=west] at (3.45,-1.15) {bounded, DJ$_\mathrm{pp}$; adds linearly};
\end{tikzpicture}
\caption{\textbf{Jitter splits into an unbounded random part and bounded deterministic parts.} Each leaf shows the shape of its probability density. RJ terms are independent Gaussians and add in root-sum-square; DJ terms are added linearly, which is conservative.}
\label{fig:ch11-tree}
\figsrc{After Palermo~\cite{palermo720clk}.}
\end{figure}

When independent sources act on the same edge, the PDF of the total is the \emph{convolution} of their PDFs, the rule \cref{sec:stateye} used for ISI. In \cref{fig:ch11-hist} a 2\,ps RJ convolved with 20\,ps$_\mathrm{pp}$ of DJ gives a total whose centre looks like the DJ and whose tails are Gaussian.

\begin{figure}[tbp]
\centering
\includegraphics[width=\linewidth]{ch11_hist}
\caption{\textbf{Total jitter is the convolution of RJ and DJ, and its tails are Gaussian} (log scale). (a) RJ never ends. (b) DJ, here DCD of $\pm4$\,ps convolved with ISI from three earlier bits, stops at $\pm10$\,ps. (c) The convolution. The dashed dual-Dirac model is fitted to the tails between BER $10^{-4}$ and $10^{-12}$; it ignores the centre and returns DJ$_{\delta\delta}$ = 15.6\,ps, less than the true 20\,ps peak-to-peak.}
\label{fig:ch11-hist}
\end{figure}

%% =====================================================================
\section{The dual-Dirac model and total jitter}\label{sec:dualdirac}

At $10^{-12}$ only the far tails of the jitter PDF matter. The \textbf{dual-Dirac model}~\cite{p5:agilentdualdirac} keeps the tails and simplifies the rest: all DJ becomes two impulses at $\pm\mathrm{DJ}_{\delta\delta}/2$, each carrying a Gaussian of rms $\sigma_\mathrm{RJ}$ (\cref{fig:ch11-bathtub}a).

\begin{derivation}[From two Gaussians to total jitter]
Measure the sampling instant $x$ from the left crossing of the eye. Far out, only the Gaussian centred at $+\mathrm{DJ}_{\delta\delta}/2$ can put an edge later than $x$, with probability $\Qf\big((x - \mathrm{DJ}_{\delta\delta}/2)/\sigma_\mathrm{RJ}\big)$. Adding the mirror image from the right crossing gives the \textbf{bathtub curve}
\[
  \BER(x) \approx \Qf\!\left(\frac{x - \mathrm{DJ}_{\delta\delta}/2}{\sigma_\mathrm{RJ}}\right) + \Qf\!\left(\frac{\UI - x - \mathrm{DJ}_{\delta\delta}/2}{\sigma_\mathrm{RJ}}\right).
\]
The target BER is reached where $x - \mathrm{DJ}_{\delta\delta}/2 = Q_\BER\,\sigma_\mathrm{RJ}$ on each side, with $\Qf(Q_\BER) = \BER$. The eye that remains is $\UI - \mathrm{DJ}_{\delta\delta} - 2Q_\BER\,\sigma_\mathrm{RJ}$; the rest is \emph{total jitter}.
\end{derivation}

\begin{keyeq}[Total jitter at a BER (dual-Dirac)]
\begin{equation}
  \mathrm{TJ}(\BER) = \mathrm{DJ}_{\delta\delta} + 2\,Q_\BER\,\sigma_\mathrm{RJ},
  \qquad \text{eye width} = 1\,\UI - \mathrm{TJ}, \qquad 2Q_{10^{-12}} = 14.07 .
  \label{eq:tj}
\end{equation}
\end{keyeq}
$Q_\BER$ is the Gaussian quantile of \cref{eq:ber-q}: 6.00 at $10^{-9}$, 7.03 at $10^{-12}$, 7.94 at $10^{-15}$. The derivation drops two factors of $\tfrac12$ (each impulse carries half the edges, and about half the bits have no transition); keeping them gives $Q = 6.84$. The industry keeps 7.03, which is 3\% conservative, and so do we.

\begin{figure}[tbp]
\centering
\includegraphics[width=\linewidth]{ch11_bathtub}
\caption{\textbf{The bathtub: BER against sampling instant.} (a) Dual-Dirac PDFs at both crossings of the eye, DJ$_{\delta\delta}$ = 40\,ps and $\sigma_\mathrm{RJ}$ = 3\,ps. (b) Bathtubs for DJ$_{\delta\delta}$ = 40\,ps. At $10^{-12}$ the eye is 0.71\,UI wide with the \OS{} budget ($\sigma = 1.32$\,ps), 0.59\,UI with 3\,ps and 0.45\,UI with 5\,ps.}
\label{fig:ch11-bathtub}
\end{figure}

\Cref{fig:ch11-bathtub}(b) completes the teaser of \cref{sec:ber}. Note the factor 14.07 in \cref{eq:tj}: a picosecond of RJ costs 14\,ps of eye at $10^{-12}$, while a picosecond of DJ costs one.

\begin{pitfall}[DJ$_{\delta\delta}$ is not DJ$_\mathrm{pp}$]
DJ$_{\delta\delta}$ is a model parameter fitted to the tails. It is smaller than the true peak-to-peak DJ unless the DJ is two-valued, such as pure DCD (15.6 versus 20\,ps in \cref{fig:ch11-hist}). Use DJ$_{\delta\delta}$ in \cref{eq:tj}, and fit the tails at deep BER (Q of 4 or more): a fit at $10^{-3}$ overestimates $\sigma_\mathrm{RJ}$. A direct TJ($10^{-12}$) measurement needs about $10^{13}$ bits, half an hour per point at 5\,Gb/s.
\end{pitfall}

%% =====================================================================
\section{From phase noise to jitter}\label{sec:phasenoise}

Oscillators are specified in the frequency domain. The \textbf{single-sideband phase noise} $\mathcal{L}(f)$ is the noise power in 1\,Hz at offset $f$ from the carrier, relative to the carrier, in dBc/Hz. For small noise it is half the one-sided phase spectrum $S_\phi(f)$, so integrating it gives the phase variance, and dividing by $2\pi f_0$ turns phase into time:
\begin{keyeq}[Phase noise to rms jitter]
\begin{equation}
  \sigma_t = \frac{\sigma_\phi}{2\pi f_0} = \frac{1}{2\pi f_0}\sqrt{2\!\int_{f_1}^{f_2}\!\mathcal{L}(f)\,df}, \qquad \mathcal{L} \text{ in linear units } \big(10^{\mathcal{L}_\mathrm{dB}/10}\big).
  \label{eq:ch11-pn2j}
\end{equation}
\end{keyeq}
The limits matter. The upper limit $f_2$ is about $f_0/2$. The lower limit $f_1$ is set by what the receiver tracks: jitter far below the CDR bandwidth moves the clock and the data together and does not close the eye. Integrating from 10\,kHz gives the clock's absolute jitter; integrating from the CDR bandwidth gives the jitter that reaches the sampler (\cref{fig:ch11-pn}).

\begin{figure}[tbp]
\centering
\includegraphics[width=\linewidth]{ch11_pn}
\caption{\textbf{Integrating phase noise gives jitter, and the lower limit decides the answer.} (a) The free-running VCO ($-100.8$\,dBc/Hz at 1\,MHz), the PLL's in-band floor, and the output of the 3\,MHz PLL of \cref{sec:pll-loop}: the output follows the floor inside the loop bandwidth and the VCO outside it. (b) rms jitter from $f_1$ to 1.25\,GHz. The PLL output has 1.01\,ps from 10\,kHz, but only 0.61\,ps above a 4\,MHz CDR bandwidth. The free-running VCO diverges as $f_1 \to 0$, which is the random walk of \cref{fig:ch11-tie}(b).}
\label{fig:ch11-pn}
\end{figure}

\paragraph{A useful shortcut} In the $1/f^2$ region of an LC oscillator, $\mathcal{L}(f) = \mathcal{L}(\Delta f)(\Delta f/f)^2$, and the free-running TIE grows as $\kappa\sqrt{t}$ with $\kappa = (\Delta f/f_0)\sqrt{\mathcal{L}(\Delta f)}$. If a loop tracks this noise below a corner $f_c$ (a first-order high-pass), \cref{eq:ch11-pn2j} integrates in closed form to
\begin{equation}
  \sigma_t \approx \kappa/\sqrt{4\pi f_c} ,
  \label{eq:ch11-kappa}
\end{equation}
so jitter falls only as $1/\sqrt{f_c}$: quadrupling the bandwidth halves the VCO's share.

\begin{example}{Phase noise of the team PLL as jitter}\label{we:ch11-pn}
\textbf{Given:} $\mathcal{L} = -100.8$\,dBc/Hz at 1\,MHz~\cite{smacd2026pll}, applied at $f_0 = 2.5$\,GHz.
\textbf{Assumptions} (they change the answer): the value describes the \emph{free-running} VCO with a $1/f^2$ slope and a 100\,kHz $1/f^3$ corner; the in-band floor from reference and charge pump is $-110$\,dBc/Hz; the PLL is the 3\,MHz design of \cref{sec:pll-loop}.

\textbf{Free-running:} $\mathcal{L} = 8.3\times10^{-11}$, so $\kappa = 4\times10^{-4}\sqrt{8.3\times10^{-11}} = 3.65\times10^{-9}\,\sqrt{\mathrm{s}}$. Period jitter $\kappa\sqrt{400\,\mathrm{ps}} = 73$\,fs; TIE after 1\,\textmu{}s $= 3.65$\,ps.
\textbf{Locked, shortcut:} $3.65\times10^{-9}/\sqrt{4\pi\times3\,\mathrm{MHz}} = 0.59$\,ps from the VCO.
\textbf{Locked, full model} (\cref{fig:ch11-pn}): \textbf{1.01\,ps} rms from 10\,kHz, \textbf{0.61\,ps} above 4\,MHz.

\textbf{Check the assumption:} \cref{sec:lcvco} shows that a $Q = 16.8$ tank should be about 20\,dB quieter than $-100.8$\,dBc/Hz, so the number may be a closed-loop value. If so, every VCO term here is pessimistic. Confirm with the team before freezing the PLL bandwidth.
\end{example}

%% =====================================================================
\section{The Open-SERDES jitter budget}\label{sec:jitter-budget}

A \textbf{jitter budget} divides the UI among the blocks that move edges, so each designer gets a number to meet. Add DJ$_{\delta\delta}$ linearly, add RJ in RSS, and combine with \cref{eq:tj}. The \OS{} target is \textbf{TJ $\le$ 0.3\,UI = 60\,ps at $10^{-12}$}, measured at the sampler; the remaining 0.7\,UI covers the sampler aperture and setup time (\cref{sec:sampler}), static phase offsets and margin. Only jitter the CDR does \emph{not} track is counted: the clock RJ rows use the phase noise above the CDR bandwidth (0.61\,ps from \cref{we:ch11-pn}, rounded up to cover the buffers), and the CDR row is the CDR's own tracking error (\cref{tab:ch11-budget}, \cref{fig:ch11-budget}).

\begin{table}[htbp]
\centering
\caption{The \OS{} jitter budget at $10^{-12}$. DJ terms are dual-Dirac values; RJ terms are rms.}
\label{tab:ch11-budget}
\small
\begin{tabularx}{\linewidth}{L r r l}
\toprule
Contributor & DJ$_{\delta\delta}$ (ps) & RJ (ps rms) & Owner \\
\midrule
TX: half-rate mux DCD, driver data-dependent jitter & 8 & -- & \cref{ch:tx} \\
Channel: residual ISI after equalization (DDJ) & 6 & -- & \cref{ch:dfe} \\
RX clock path: I/Q error, duty cycle, PI nonlinearity & 6 & -- & \cref{sec:quadrature} \\
CDR: PI quantization and bang-bang dither & 20 & -- & \cref{sec:picdr} \\
TX clock: PLL plus buffers, not tracked by the CDR & -- & 1.0 & \cref{ch:pll} \\
RX clock: PLL, PI and buffers, not tracked & -- & 0.7 & \cref{ch:pll} \\
Samplers: aperture and noise & -- & 0.5 & \cref{sec:sampler} \\
\midrule
Total & \textbf{40} & \textbf{1.32} (RSS) & \\
\bottomrule
\end{tabularx}
\end{table}

\begin{figure}[htbp]
\centering
\includegraphics[width=\linewidth]{ch11_budget}
\caption{\textbf{The jitter budget fills 29\% of the UI.} DJ terms (gold) stack linearly; the RJ terms are combined in RSS and multiplied by $2Q = 14.07$ (blue). The total, 58.6\,ps, sits just inside the 0.3\,UI line and leaves a 0.71\,UI eye.}
\label{fig:ch11-budget}
\end{figure}

\begin{example}{Does the budget close?}\label{we:ch11-budget}
\textbf{DJ:} $8 + 6 + 6 + 20 = 40$\,ps. \textbf{RJ:} $\sigma = \sqrt{1.0^2 + 0.7^2 + 0.5^2} = 1.32$\,ps, so $2Q\sigma = 14.07\times1.32 = 18.6$\,ps.
\textbf{TJ:} $40 + 18.6 = \mathbf{58.6}$\,ps $=$ \textbf{0.293\,UI}, inside 0.3\,UI with 1.4\,ps to spare; the eye at $10^{-12}$ is 141\,ps wide.

\textbf{Turn it around:} with DJ at 40\,ps, RJ may be at most $(60-40)/14.07 = 1.42$\,ps rms; with RJ fixed, DJ may grow by only 1.4\,ps. Tight budgets are normal.
\textbf{Consistency:} \cref{sec:stateye} used 2\,ps RJ and 20\,ps DJ for the sampling clock; the CDR row here is that 20\,ps, and 2\,ps conservatively covers the RX share of the RJ.
\end{example}

\begin{insight}
The largest line is the CDR's own dither, and the most expensive per picosecond is RJ, because of the factor 14. So: design the CDR loop for low dither (\cref{sec:picdr}); keep the PLL's untracked jitter below 1\,ps rms (\cref{ch:pll}); and keep duty-cycle and quadrature errors small, because they become DJ at the half-rate samplers (\cref{sec:quadrature}).
\end{insight}

%% =====================================================================
\section{Jitter transfer, generation and tolerance}\label{sec:ch11-jtran}

Three more terms describe how a block treats jitter versus jitter frequency. \textbf{Jitter transfer} is the ratio of output to input jitter: a PLL is low-pass to its reference, a CDR low-pass to the data. \textbf{Jitter generation} is the jitter a block adds with a clean input. \textbf{Jitter tolerance} (JTOL) is the largest sinusoidal jitter a receiver can take at a given frequency and still meet the BER. Because the CDR tracks slow jitter, a receiver tolerates many UI at low frequency but only its eye margin at high frequency; \cref{sec:jtol} checks the \OS{} CDR against a mask.

\begin{tryit}[Lab 6: jitter and the bathtub]
\Cref{sec:lab6} adds RJ and DJ to the REF link and plots the bathtub. Reproduce \cref{fig:ch11-bathtub}(b), then change $\sigma_\mathrm{RJ}$ by 0.5\,ps and watch the eye width move by 7\,ps.
\end{tryit}

\begin{takeaways}
\begin{itemize}
  \item A link cares about time-interval error, and only the part the CDR does not track.
  \item Voltage noise becomes jitter through the slope, $\sigma_t = \sigma_v/\mathrm{SR}$: keep clock swings large and edges fast.
  \item RJ is Gaussian and unbounded (rms, adds in RSS); DJ is bounded (DDJ, DCD, PJ, BUJ; adds linearly). Their PDFs convolve.
  \item Dual-Dirac: $\mathrm{TJ} = \mathrm{DJ}_{\delta\delta} + 2Q\sigma_\mathrm{RJ}$ with $2Q = 14.07$ at $10^{-12}$; one picosecond of RJ costs 14\,ps of eye.
  \item $\sigma_t = \sqrt{2\int\mathcal{L}\,df}/(2\pi f_0)$, integrated from the CDR bandwidth. The team PLL model gives 1.01\,ps from 10\,kHz and 0.61\,ps above 4\,MHz.
  \item The \OS{} budget (40\,ps DJ, 1.32\,ps RJ) gives TJ = 58.6\,ps = 0.293\,UI against 0.3\,UI; the CDR's 20\,ps is the largest line.
\end{itemize}
\end{takeaways}

\tikzsetfigurename{12_pll-}
\chapter[Clock Generation: The Phase-Locked Loop]{Clock Generation:\\ The Phase-Locked Loop}\label{ch:pll}
\chapterlead{A crystal gives a clean 50\,MHz; the link needs a clean 2.5\,GHz in four phases. The charge-pump PLL with an LC VCO bridges the two.}

\begin{objectives}
\begin{itemize}
  \item Explain how a phase-frequency detector and charge pump turn a timing error into charge, and why the reset delay removes the dead zone.
  \item Derive the type-II loop's $\omega_n$ and $\zeta$, and design the \OS{} loop filter for a given bandwidth and phase margin.
  \item Use noise transfer functions to choose the loop bandwidth that minimizes output jitter.
  \item Size an LC tank, state the oscillation condition, and use Leeson's model to sanity-check a phase-noise number.
  \item Generate and distribute quadrature (I/Q) clocks, and budget their phase and duty-cycle errors.
\end{itemize}
\end{objectives}

%% =====================================================================
\section{The charge-pump PLL}\label{sec:cppll}

A free-running VCO drifts: its TIE random-walks (\cref{fig:ch11-tie}b) and its frequency wanders with temperature and supply. A \textbf{phase-locked loop} (PLL) fixes both by comparing the divided VCO output with a crystal reference and steering the VCO until the two are aligned. Gardner's analysis of the charge-pump PLL~\cite{p5:gardner1980} underlies this chapter; Razavi~\cite{razaviPLL}, Gardner's book~\cite{gardner2005} and ECEN~720~\cite{palermo720clk} treat it in more depth. Locked, the output frequency is exactly $N f_\mathrm{ref}$. \OS{} uses an integer-N charge-pump PLL (\cref{fig:ch12-arch}) with $f_\mathrm{ref} = 50$\,MHz and $N = 50$, reusing the team's LC-VCO PLL design~\cite{smacd2026pll} retuned from 2.4 to 2.5\,GHz.

\begin{figure}[tbp]
\centering
\begin{tikzpicture}[font=\sffamily\small,scale=0.95,transform shape]
  \node[blk,minimum width=14mm] (xo) at (0,0) {XO\\[-2pt]{\scriptsize 50\,MHz}};
  \node[blkb] (pfd) at (2.6,0) {PFD};
  \node[blkb,minimum width=13mm] (cp) at (4.9,0) {CP\\[-2pt]{\scriptsize $\pm\Icp$}};
  \node[blkb,minimum width=19mm] (vco) at (10.2,0) {LC VCO\\[-2pt]{\scriptsize $\Kvco$, 2.5\,GHz}};
  \node[blkg,minimum width=14mm] (div) at (6.4,-2.6) {$\div N$\\[-2pt]{\scriptsize $N = 50$}};
  \draw[sig] (xo) -- node[above,lbl]{$f_\mathrm{ref}$} (pfd);
  \draw[sig] (pfd.15) -- node[above,lbl]{UP} (cp.165);
  \draw[sig] (pfd.-15) -- node[below,lbl]{DN} (cp.195);
  \draw[navy,line width=0.75pt] (cp.east) -- (6.6,0) node[dot]{} -- (7.9,0) node[dot]{} -- (8.6,0);
  \draw[sig] (8.6,0) -- node[above,lbl]{$V_\mathrm{ctrl}$} (vco.west);
  \draw[color=navy] (6.6,0) to[C, l_=$C_2$] (6.6,-1.5) node[ground]{};
  \draw[color=navy] (7.9,0) to[R, l=$R$] (7.9,-0.85) to[C, l=$C_1$] (7.9,-1.7) node[ground]{};
  \node[lbl,text=bblue] at (7.25,0.45) {loop filter $Z(s)$};
  \draw[sig] (vco.east) -- (12.3,0) node[right,align=left,lbl,text=navy]{2.5\,GHz to I/Q\\generation, TX};
  \draw[sig] (11.7,0) node[dot]{} |- (div.east);
  \draw[sig] (div.west) -| node[pos=0.75,left,lbl]{$f_\mathrm{div}$} (pfd.south);
\end{tikzpicture}
\caption{\textbf{The \OS{} integer-N charge-pump PLL.} The phase-frequency detector (PFD) compares the 50\,MHz reference with the VCO divided by 50; the charge pump (CP) turns the timing difference into charge; the loop filter turns charge into the control voltage. At lock $f_\mathrm{out} = N f_\mathrm{ref} = 2.5$\,GHz.}
\label{fig:ch12-arch}
\end{figure}

\paragraph{Phase-frequency detector and charge pump}
The \textbf{PFD} is two flip-flops with their D inputs tied high (\cref{fig:ch12-pfd}). A rising reference edge sets UP, a rising divider edge sets DN, and when both are high an AND gate resets them after a short delay $T_\mathrm{rst}$. The width of UP minus the width of DN equals the timing difference $\Delta t$ between the edges. If the frequencies differ, one output dominates for many cycles, so the PFD also pulls the frequency in: it is a phase \emph{and} frequency detector. The \textbf{charge pump} switches a current $\Icp$ into the loop filter during UP and out during DN. The average current is $\Icp\,\Delta t/T_\mathrm{ref} = \Icp\,\Delta\phi/2\pi$, so the PFD-plus-CP gain is $\Icp/2\pi$ amperes per radian of reference phase.

\begin{figure}[tbp]
\centering
\begin{tikzpicture}[font=\sffamily\scriptsize,scale=0.85,transform shape,
   ff/.style={draw=navy,line width=0.7pt,fill=tile,minimum width=11mm,minimum height=11mm}]
  \node[ff] (f1) at (0,1.2) {}; \node[ff] (f2) at (0,-1.2) {};
  \foreach \f/\r in {f1/south,f2/north}{
    \node[anchor=west,text=navy] at ([yshift=2.5mm]\f.west) {D};
    \node[anchor=east,text=navy] at ([yshift=2.5mm]\f.east) {Q};
    \node[anchor=\r,text=navy,inner sep=1.5pt] at (\f.\r) {R};
    \draw[navy,line width=0.6pt] ([yshift=-4.4mm]\f.west) -- ++(0.15,0.11) -- ++(-0.15,0.11);
    \draw[navy,line width=0.7pt] ([yshift=2.5mm]\f.west) -- ++(-0.4,0) node[left]{1};}
  \draw[sig] (-1.4,0.95) node[left,text=navy]{REF} -- ([yshift=-2.5mm]f1.west);
  \draw[sig] (-1.4,-1.45) node[left,text=navy]{DIV} -- ([yshift=-2.5mm]f2.west);
  \node[and port,scale=0.7,color=navy] (and) at (2.3,0.45) {};
  \draw[navy,line width=0.7pt] (1.05,1.45) node[circ]{} |- (and.in 1);
  \draw[navy,line width=0.7pt] (1.3,-0.95) node[circ]{} |- (and.in 2);
  \node[blkg,minimum width=9mm,minimum height=5mm,font=\sffamily\scriptsize] (td) at (3.55,0.45) {$T_\mathrm{rst}$};
  \draw[navy,line width=0.7pt] (and.out) -- (td.west);
  \draw[sig] (td.east) -- (4.25,0.45) |- (0.3,0) -- (0,0) node[circ]{};
  \draw[sig] (0,0) -- (f1.south); \draw[sig] (0,0) -- (f2.north);
  \draw[sig] ([yshift=2.5mm]f1.east) -- (5.3,1.45) node[above,pos=0.85,text=navy]{UP};
  \draw[sig] ([yshift=2.5mm]f2.east) -- (5.3,-0.95) node[below,pos=0.85,text=navy]{DN};
  \draw[color=navy] (6.3,2.7) node[vcc]{} to[I, l=$\Icp$] (6.3,1.65) to[nos] (6.3,0.45) to[nos] (6.3,-0.75) to[I, l=$\Icp$] (6.3,-1.8) node[ground]{};
  \draw[dashed,gold,line width=0.6pt] (5.3,1.45) -- (6.05,1.05); \draw[dashed,gold,line width=0.6pt] (5.3,-0.95) -- (6.05,-0.15);
  \draw[sig] (6.3,0.45) node[circ]{} -- (7.5,0.45) node[right,text=navy]{to $Z(s)$};
  \node[text=navy] at (7.0,0.75) {$I_\mathrm{out}$};
\end{tikzpicture}\par\medskip
\begin{tikzpicture}[font=\sffamily\scriptsize,
   tt/.style={timing/xunit=0.245cm,timing/yunit=0.32cm,timing/slope=0.08,line width=0.75pt,draw=navy}]
  \foreach \y/\lab/\s in {0/REF/{2L8H8L8H8L8H6L},-0.62/DIV/{5L8H8L8H8L8H3L},-1.24/UP/{2L4H12L4H12L4H10L},-1.86/DN/{5L1H15L1H15L1H10L}}{
    \node[anchor=east,text=navy] at (-0.1,\y+0.16) {\lab};
    \timing[tt] at (0,\y) {\s};}
  \node[anchor=east,text=navy] at (-0.1,-2.6) {$I_\mathrm{out}$};
  \draw[navy,line width=0.75pt] (0,-2.75) -- (0.49,-2.75) -- (0.49,-2.45) -- (1.225,-2.45) -- (1.225,-2.75) -- (4.41,-2.75) -- (4.41,-2.45) -- (5.145,-2.45) -- (5.145,-2.75) -- (8.33,-2.75) -- (8.33,-2.45) -- (9.065,-2.45) -- (9.065,-2.75) -- (11.76,-2.75);
  \node[anchor=east,text=navy] at (-0.1,-3.3) {$V_\mathrm{ctrl}$};
  \draw[c2,line width=0.75pt] (0,-3.55) -- (0.49,-3.55) -- (0.49,-3.3) -- (1.225,-3.2) -- (1.225,-3.45) -- (4.41,-3.45) -- (4.41,-3.2) -- (5.145,-3.1) -- (5.145,-3.35) -- (8.33,-3.35) -- (8.33,-3.1) -- (9.065,-3.0) -- (9.065,-3.25) -- (11.76,-3.25);
  \draw[c4,<->,line width=0.6pt] (0.49,0.5) -- node[above]{$\Delta t$} (1.225,0.5);
  \draw[muted,dotted] (0.49,0.45) -- (0.49,-3.6); \draw[muted,dotted] (1.225,0.45) -- (1.225,-3.6);
  \draw[golddk,->,line width=0.6pt] (2.9,-2.25) node[right]{both on for $T_\mathrm{rst}$: currents cancel, no dead zone} -- (1.45,-1.75);
  \node[lbl,anchor=west] at (5.4,-3.65) {jumps by $\Icp R$ while pumping; net step $\Icp\Delta t/C_1$};
\end{tikzpicture}
\caption{\textbf{The PFD measures the edge offset as a pulse width; the charge pump turns it into charge.} Top: two reset flip-flops, an AND gate with reset delay $T_\mathrm{rst}$, and the up/down current sources. Bottom: REF leads DIV by $\Delta t$, so UP is wider than DN by $\Delta t$; both are on together for $T_\mathrm{rst}$ and cancel. Each cycle the control voltage jumps by $\Icp R$ while current flows and keeps a net step of $\Icp\Delta t/C_1$.}
\label{fig:ch12-pfd}
\end{figure}

\paragraph{Dead zone and mismatch}
Without the reset delay, a tiny $\Delta t$ would give UP and DN pulses too short to switch the charge pump on, so near lock the loop would get no correction: a \textbf{dead zone} in which the VCO drifts. With $T_\mathrm{rst}$ of 100--200\,ps both pulses always last at least $T_\mathrm{rst}$, and the net charge stays proportional to $\Delta t$ down to zero. The price is mismatch: if the currents differ by $\Delta I$, the loop locks at a static offset $\Delta t = (\Delta I/\Icp)\,T_\mathrm{rst}$ (7.5\,ps for 5\% and 150\,ps), and the current pulses every 20\,ns create a \textbf{reference spur}, a periodic jitter that $C_2$ and the loop must filter.

%% =====================================================================
\section{Loop dynamics}\label{sec:pll-loop}

Near lock every block is linear in \emph{phase}, which gives the model of \cref{fig:ch12-linmodel}. The PFD and CP give $\Icp/2\pi$ (A/rad); the loop filter impedance $Z(s)$ turns current into $V_\mathrm{ctrl}$; the VCO integrates frequency into phase, $\Kvco/s$ (rad/s/V divided by $s$); and the divider scales phase by $1/N$. The same figure marks where each noise source enters.

\begin{figure}[tbp]
\centering
\begin{tikzpicture}[font=\sffamily\small,
  b/.style={blkb,minimum width=13mm,minimum height=8mm},
  sm/.style={sumnode,minimum size=5mm,font=\scriptsize},
  nz/.style={font=\sffamily\scriptsize,text=c4,align=center},
  tshape/.style={line width=0.7pt,draw=golddk}]
  \node[text=navy] (in) at (0,0) {$\phi_\mathrm{ref}$};
  \node[sm] (s0) at (1.15,0) {$+$};
  \node[sm] (s1) at (2.3,0) {$\Sigma$};
  \node[b] (kp) at (3.75,0) {$\dfrac{\Icp}{2\pi}$};
  \node[sm] (s2) at (5.15,0) {$+$};
  \node[b] (zf) at (6.45,0) {$Z(s)$};
  \node[sm] (s3) at (7.75,0) {$+$};
  \node[b] (kv) at (9.1,0) {$\dfrac{\Kvco}{s}$};
  \node[sm] (s4) at (10.5,0) {$+$};
  \node[text=navy] (out) at (12.0,0) {$\phi_\mathrm{out}$};
  \node[blkg,minimum width=11mm,minimum height=7mm] (dv) at (6.45,-1.35) {$1/N$};
  \draw[sig] (in) -- (s0); \draw[sig] (s0) -- (s1); \draw[sig] (s1) -- (kp); \draw[sig] (kp) -- (s2);
  \draw[sig] (s2) -- (zf); \draw[sig] (zf) -- (s3); \draw[sig] (s3) -- (kv); \draw[sig] (kv) -- (s4); \draw[sig] (s4) -- (out);
  \draw[sig] (11.25,0) node[dot]{} |- (dv.east); \draw[sig] (dv.west) -| node[pos=0.95,left,font=\scriptsize]{$-$} (s1.south);
  \foreach \s/\t/\k in {s0/{$\phi_{n,\mathrm{ref}}$}/LP,s2/{$i_{n,\mathrm{CP}}$}/LP,s3/{$v_{n,R}$}/BP,s4/{$\phi_{n,\mathrm{VCO}}$}/HP}{
    \node[nz] (n\s) at ([yshift=12mm]\s) {\t};
    \draw[sig,draw=c4] (n\s) -- (\s);}
  % tiny transfer-shape icons
  \begin{scope}[shift={(1.55,1.05)},x=0.45cm,y=0.3cm]
    \draw[tshape] (0,0.8) -- (0.6,0.8) -- (1,0); \node[font=\sffamily\tiny,text=golddk] at (0.5,-0.4) {LP $\times N$};\end{scope}
  \begin{scope}[shift={(5.55,1.05)},x=0.45cm,y=0.3cm]
    \draw[tshape] (0,0.8) -- (0.6,0.8) -- (1,0); \node[font=\sffamily\tiny,text=golddk] at (0.5,-0.4) {LP};\end{scope}
  \begin{scope}[shift={(8.15,1.05)},x=0.45cm,y=0.3cm]
    \draw[tshape] (0,0) -- (0.35,0.8) -- (0.65,0.8) -- (1,0); \node[font=\sffamily\tiny,text=golddk] at (0.5,-0.4) {BP};\end{scope}
  \begin{scope}[shift={(10.9,1.05)},x=0.45cm,y=0.3cm]
    \draw[tshape] (0,0) -- (0.4,0.8) -- (1,0.8); \node[font=\sffamily\tiny,text=golddk] at (0.5,-0.4) {HP};\end{scope}
\end{tikzpicture}
\caption{\textbf{Linear phase model of the CP-PLL, with its noise sources.} Each source reaches the output through a different filter: reference and charge-pump noise are low-pass filtered (and multiplied by $N$), loop-filter resistor noise is band-pass, and VCO noise is high-pass. The loop bandwidth sets the crossover between ``trust the reference'' and ``trust the VCO''.}
\label{fig:ch12-linmodel}
\end{figure}

\paragraph{Type-II second-order loop} Ignore $C_2$ for now, so $Z(s) = R + 1/(sC_1)$. The open-loop gain is
\[
  LG(s) = \frac{\Icp}{2\pi}\,Z(s)\,\frac{\Kvco}{s}\,\frac1N = \frac{\Icp\Kvco}{2\pi N C_1}\cdot\frac{1 + sRC_1}{s^2} .
\]
It has two integrators (the VCO and $C_1$), so it is a \textbf{type-II} loop: it settles to zero phase error even after a frequency step. The closed-loop transfer $\phi_\mathrm{out}/\phi_\mathrm{ref} = N\,LG/(1+LG)$ becomes, after multiplying through by $s^2$,
\begin{keyeq}[CP-PLL natural frequency and damping]
\begin{equation}
  \frac{\phi_\mathrm{out}}{\phi_\mathrm{ref}} = N\,\frac{2\zeta\omega_n s + \omega_n^2}{s^2 + 2\zeta\omega_n s + \omega_n^2},
  \qquad
  \omega_n = \sqrt{\frac{\Icp\Kvco}{2\pi N C_1}}, \qquad
  \zeta = \frac{R}{2}\sqrt{\frac{\Icp\Kvco C_1}{2\pi N}} .
  \label{eq:ch12-wnz}
\end{equation}
\end{keyeq}
Matching coefficients gives $\omega_n^2 = \Icp\Kvco/(2\pi N C_1)$ and $2\zeta\omega_n = \Icp\Kvco R/(2\pi N)$. Raising $R$ raises the damping without changing $\omega_n$; raising $C_1$ lowers $\omega_n$ and raises $\zeta$. The zero in the numerator causes some \textbf{jitter peaking} at every damping (\cref{fig:ch12-resp}, \cref{tab:ch12-zeta}): the loop amplifies reference jitter near $\omega_n$. The $-3$\,dB bandwidth is $\omega_{-3\mathrm{dB}} = \omega_n\big[1 + 2\zeta^2 + \sqrt{(1+2\zeta^2)^2 + 1}\,\big]^{1/2}$, which tends to $2\zeta\omega_n$ for large $\zeta$.

\begin{table}[htbp]
\centering
\caption{Bandwidth and peaking of the type-II loop of \cref{eq:ch12-wnz} (computed).}
\label{tab:ch12-zeta}
\small
\begin{tabular}{lcccccc}
\toprule
$\zeta$ & 0.5 & 0.707 & 1 & 2 & 4 & 4.5 \\
\midrule
$\omega_{-3\mathrm{dB}}/\omega_n$ & 1.82 & 2.06 & 2.48 & 4.25 & 8.13 & 9.11 \\
Peaking (dB) & 3.3 & 2.1 & 1.25 & 0.40 & 0.12 & 0.09 \\
\bottomrule
\end{tabular}
\end{table}

\begin{figure}[tbp]
\centering
\includegraphics[width=\linewidth]{ch12_resp}
\caption{\textbf{Damping trades peaking and overshoot for bandwidth.} (a) Closed-loop $|H|/N$: low damping gives a sharp, peaky response; high damping a wide, flat one. (b) Output phase after a reference phase step: $\zeta = 0.5$ overshoots by 30\%, $\zeta = 1$ by 13\%. Transmit PLLs use $\zeta \approx 0.7$--1; CDRs, which must not amplify jitter, use $\zeta \gg 1$ (\cref{ch:cdr}).}
\label{fig:ch12-resp}
\end{figure}

\paragraph{The third pole and phase margin} The capacitor $C_2$ in parallel smooths the voltage step $\Icp R$ seen in \cref{fig:ch12-pfd}, which would otherwise modulate the VCO every reference cycle. It adds a pole at $\omega_{p3} = 1/(R C_s)$ with $C_s = C_1C_2/(C_1+C_2)$, and the loop becomes third order. The usual design places the unity-gain frequency $\omega_u$ at the geometric mean of the zero $\omega_z = 1/(RC_1)$ and $\omega_{p3}$, where the phase margin $\Phi_m$ is largest. With $K_C = C_1/C_2$:
\begin{equation}
\begin{gathered}
  K_C = 2\tan^2\Phi_m + 2\tan\Phi_m\sqrt{\tan^2\Phi_m + 1}, \qquad
  \omega_z = \frac{\omega_u}{\sqrt{1+K_C}}, \qquad
  \omega_{p3} = \omega_u\sqrt{1+K_C}, \\
  C_2 = \frac{\Icp\Kvco}{2\pi N\omega_u^2}\sqrt{\frac{\omega_u^2+\omega_z^2}{\omega_u^2+\omega_{p3}^2}}, \qquad C_1 = K_C C_2, \qquad R = \frac{1}{\omega_z C_1} .
\end{gathered}
  \label{eq:ch12-design}
\end{equation}
For $\Phi_m = 60^\circ$, $K_C = 12.9$. Keep $f_u \le f_\mathrm{ref}/10$, or the sampled nature of the PFD makes the continuous-time model wrong.

\begin{example}{Loop filter for the Open-SERDES PLL}\label{we:ch12-lf}
\textbf{Given:} $f_\mathrm{ref} = 50$\,MHz, $N = 50$, $\Kvco = 2\pi\times150$\,MHz/V, $\Icp = 200$\,\textmu{}A, $f_u = 1$\,MHz, $\Phi_m = 60^\circ$.

\textbf{Pole and zero:} $K_C = 12.93$, so $f_z = 1\,\mathrm{MHz}/\sqrt{13.93} = 268$\,kHz and $f_{p3} = 3.73$\,MHz.
\textbf{Components} from \cref{eq:ch12-design}: $C_2 = \mathbf{4.07}$\,pF, $C_1 = K_C C_2 = \mathbf{52.6}$\,pF, $R = 1/(\omega_z C_1) = \mathbf{11.3}$\,k\Ohm.
\textbf{Check} with \cref{eq:ch12-wnz}: $f_n = 537$\,kHz and $\zeta = 1.00$. The full third-order loop, simulated in Python (\cref{fig:ch12-loop}), crosses unity gain at 1.00\,MHz with 60.0$^\circ$ margin and has $f_{-3\mathrm{dB}} = 1.56$\,MHz with 1.7\,dB peaking.
\textbf{Practical:} $f_\mathrm{ref}/f_u = 50$, comfortably above 10. $C_1$ is 0.035\,mm$^2$ of MIM at about 1.5\,fF/\textmu{}m$^2$ (check the PDK value), the largest item in the PLL layout. $R$ adds $\sqrt{4kTR} = 13.7$\,nV/$\sqrt{\mathrm{Hz}}$ of band-pass noise.
\end{example}

\begin{figure}[tbp]
\centering
\includegraphics[width=\linewidth]{ch12_loop}
\caption{\textbf{The designed loop has 60$^\circ$ of phase margin at 1\,MHz.} (a) Open-loop gain and phase of \cref{we:ch12-lf}: the phase is closest to $-180^\circ$ furthest from it at $f_u$, midway between $f_z$ and $f_{p3}$ on a log scale. (b) The third-order closed loop (navy) against the second-order fit with $\zeta = 1$: they agree up to the bandwidth, and $C_2$ adds extra roll-off above it.}
\label{fig:ch12-loop}
\end{figure}

%% =====================================================================
\section{PLL noise and the optimum bandwidth}\label{sec:ch12-noise}

From \cref{fig:ch12-linmodel}, each noise source sees its own transfer function:
\begin{equation}
  \mathrm{NTF}_\mathrm{ref} = N\frac{LG}{1+LG}, \quad
  \mathrm{NTF}_\mathrm{CP} = \frac{2\pi}{\Icp}\,N\frac{LG}{1+LG}, \quad
  \mathrm{NTF}_\mathrm{VCO} = \frac{1}{1+LG}, \quad
  \mathrm{NTF}_R = \frac{\Kvco/s}{1+LG} .
\end{equation}
Reference and charge-pump noise pass at low frequency and are multiplied by $N$ ($+34$\,dB for $N = 50$); VCO noise is suppressed at low frequency and passes above the bandwidth (\cref{fig:ch12-noise}a). Widening the loop therefore lets in more reference noise and keeps out more VCO noise. The best bandwidth is roughly where the free-running VCO noise crosses the in-band floor.

\begin{figure}[tbp]
\centering
\includegraphics[width=\linewidth]{ch12_noise}
\caption{\textbf{The loop bandwidth trades reference noise against VCO noise.} (a) Noise transfer for $f_u$ = 1 and 3\,MHz: the reference and CP path is low-pass, the VCO path high-pass. (b) Integrated output jitter (10\,kHz to 1.25\,GHz) against $f_u$, for the VCO of \cref{we:ch11-pn} and a $-110$\,dBc/Hz in-band floor. The total is flat near its 3\,MHz minimum of 1.01\,ps. If the VCO is 22\,dB quieter (\cref{sec:lcvco}), the optimum moves below 0.3\,MHz (dotted).}
\label{fig:ch12-noise}
\end{figure}

\begin{example}{Choosing the PLL bandwidth}\label{we:ch12-bw}
\textbf{Given:} the VCO noise of \cref{we:ch11-pn} ($-100.8$\,dBc/Hz at 1\,MHz), and an in-band floor of $-110$\,dBc/Hz at the output. That floor corresponds to a normalized PLL figure of merit of $-110 - 10\log_{10}(50\,\mathrm{MHz}) - 20\log_{10}50 = -221$\,dBc/Hz, typical of a good CP-PLL.

\textbf{Sweep} (60$^\circ$ loops of \cref{eq:ch12-design}, \cref{fig:ch12-noise}b): $f_u$ = 0.5, 1, 2, 3, 6\,MHz gives 1.91, 1.34, 1.06, \textbf{1.01}, 1.14\,ps rms. At 3\,MHz the VCO and in-band terms are equal (0.71 and 0.72\,ps).

\textbf{Redesign at 3\,MHz:} capacitors scale as $\Icp/f_u^2$ and $R$ as $f_u/\Icp$. Tripling $\Icp$ to 600\,\textmu{}A keeps $R = 11.3$\,k\Ohm{} and gives $C_1 = 17.5$\,pF, $C_2 = 1.36$\,pF (a third of the area), $f_{-3\mathrm{dB}} = 4.7$\,MHz and $f_\mathrm{ref}/f_u = 17$. The resistor noise adds only 0.11\,ps in RSS. (Keeping 200\,\textmu{}A would need $R = 34$\,k\Ohm{} and a 0.45\,pF $C_2$ that parasitics would swamp.)
\end{example}

\begin{insight}
Choose $f_u \approx 3$\,MHz with $\Icp = 600$\,\textmu{}A. The minimum is shallow (within 15\% from 2 to 6\,MHz), so bandwidth can be traded for area or settling time. But it depends on the VCO's free-running noise, which is uncertain (\cref{sec:lcvco}): pin that down before freezing $f_u$.
\end{insight}

%% =====================================================================
\section{The LC VCO}\label{sec:lcvco}

An LC tank stores energy and loses a little of it each cycle in the tank resistance. The \textbf{cross-coupled pair} of \cref{fig:ch12-vco} replaces that loss: looking into the two drains, the pair presents a negative resistance $-2/\gm$. The tank's loss is a parallel resistance $R_p = Q\omega_0 L$ (differential), so
\begin{keyeq}[LC oscillator: frequency and start-up]
\begin{equation}
  f_0 = \frac{1}{2\pi\sqrt{L\,C_\mathrm{tot}}}, \qquad
  R_p = Q\,\omega_0 L, \qquad
  \text{oscillation if}\ \ \frac{\gm R_p}{2} \ge 1 \ \ (\text{design for 2--3}).
  \label{eq:ch12-lc}
\end{equation}
\end{keyeq}
The amplitude grows until the pair limits, at about $(4/\pi)\,I_\mathrm{tail}R_p/2$ (differential). $C_\mathrm{tot}$ sums the varactors (fine tuning, which sets $\Kvco$), a switched-capacitor bank (coarse bands for process corners) and all parasitics.

\begin{figure}[tbp]
\centering
\begin{tikzpicture}[font=\sffamily\footnotesize,scale=0.9,transform shape]
  \draw[color=navy] (0,4.3) to[L, l=$L/2$] (2.2,4.3) to[L, l=$L/2$] (4.4,4.3);
  \draw[color=navy] (2.2,4.3) -- (2.2,4.9) node[vcc]{$V_\mathrm{DD}$};
  \draw[color=navy] (0,4.3) -- (0,0.9); \draw[color=navy] (4.4,4.3) -- (4.4,0.9);
  \draw[color=navy] (0,3.0) node[circ]{} to[vC, l=$C_v$] (2.2,3.0) to[vC, l=$C_v$] (4.4,3.0) node[circ]{};
  \draw[color=navy] (2.2,3.0) -- (2.2,2.55) node[right,text=navy]{$V_\mathrm{ctrl}$};
  \draw[color=navy] (0,1.85) node[circ]{} to[C, l=$C_s$] (1.4,1.85) to[nos, l_=$b_k$] (3.0,1.85) to[C, l=$C_s$] (4.4,1.85) node[circ]{};
  \node[nmos,xscale=-1,anchor=D] (m1) at (0,0.9) {}; \node[nmos,anchor=D] (m2) at (4.4,0.9) {};
  \node[text=navy,anchor=east] at ([xshift=-1mm]m1.center) {M1}; \node[text=navy,anchor=west] at ([xshift=1mm]m2.center) {M2};
  \draw[color=navy] (m1.G) -- ++(0.45,0) -- (3.0,1.3) -- (4.4,1.3) node[circ]{};
  \draw[color=navy] (m2.G) -- ++(-0.45,0) -- (1.4,1.3) -- (0,1.3) node[circ]{};
  \draw[color=navy] (m1.S) |- (2.2,-1.0) -| (m2.S);
  \draw[color=navy] (2.2,-1.0) node[circ]{} to[I, l=$I_\mathrm{tail}$] (2.2,-2.4) node[ground]{};
  \draw[color=navy] (0,3.65) node[circ]{} -- (-0.6,3.65) node[left,text=navy]{$V_\mathrm{out}^{+}$};
  \draw[color=navy] (4.4,3.65) node[circ]{} -- (5.0,3.65) node[right,text=navy]{$V_\mathrm{out}^{-}$};
  \draw[region] (-0.3,1.55) rectangle (4.7,4.65);
  \node[lbl,anchor=west,align=left] at (5.0,2.4) {tank: $L$ = 4\,nH, $Q$ = 16.8,\\loss $R_p = Q\omega_0L$ = 1.06\,k\Ohm};
  \node[lbl,anchor=west,align=left] at (5.0,0.0) {cross-coupled pair:\\$R_\mathrm{in} = -2/\gm$};
\end{tikzpicture}
\caption{\textbf{Cross-coupled LC VCO.} A center-tapped inductor and the capacitors set $f_0$; varactors $C_v$ give fine tuning through $V_\mathrm{ctrl}$, and switched capacitors $C_s$ (bits $b_k$) select coarse bands. The cross-coupled pair cancels the tank loss when $\gm R_p/2 \ge 1$, and the tail current sets the amplitude. Drawn with MOSFETs; npn13G2 HBTs work the same way, with capacitive cross-coupling to keep $V_\mathrm{CE}$ below BV$_\mathrm{CEO}$.}
\label{fig:ch12-vco}
\end{figure}

\begin{example}{The 2.5\,GHz tank}\label{we:ch12-tank}
\textbf{Given:} the team's 4\,nH differential inductor with $Q = 16.8$~\cite{smacd2026pll}, $f_0 = 2.5$\,GHz.
\textbf{Capacitance:} $C_\mathrm{tot} = 1/(\omega_0^2L) = 1/\big((2\pi\cdot2.5\times10^9)^2\cdot4\times10^{-9}\big) = \mathbf{1.01}$\,pF, including all parasitics.
\textbf{Loss:} $R_p = 16.8\times2\pi\cdot2.5\times10^9\times4\times10^{-9} = \mathbf{1.06}$\,k\Ohm.
\textbf{Start-up:} $\gm \ge 2/R_p = 1.89$\,mS; with a 3$\times$ margin, 5.7\,mS per device (about 0.15\,mA of collector current for an HBT, more for a MOSFET).
\textbf{Tuning:} $\pm10\%$ in frequency needs $C_\mathrm{max}/C_\mathrm{min} = (1.1/0.9)^2 = 1.49$ across varactors and bank together, with the fixed parasitics counted in both.
\end{example}

\paragraph{Phase noise: Leeson's model}
Leeson's empirical model~\cite{p5:leeson1966} captures the three regions of an oscillator's phase noise (\cref{fig:ch12-leeson}a):
\begin{equation}
  \mathcal{L}(\Delta f) = 10\log_{10}\!\left[\frac{2FkT}{P_\mathrm{sig}}\left(1 + \Big(\frac{f_0}{2Q\,\Delta f}\Big)^2\right)\left(1 + \frac{f_{1/f^3}}{\Delta f}\right)\right] ,
  \label{eq:ch12-leeson}
\end{equation}
where $P_\mathrm{sig} = V_\mathrm{pk}^2/(2R_p)$ is the power in the tank and $F$ an excess-noise factor of the active devices. Thermal noise near the carrier is shaped by the tank into a $1/\Delta f^2$ slope, flicker noise adds $1/\Delta f^3$ close in, and far out the noise is flat. For a fixed inductor, $\mathcal{L} \propto 1/(Q\,V_\mathrm{pk}^2)$: a higher $Q$ and a larger swing both help (\cref{fig:ch12-leeson}b).

\begin{figure}[tbp]
\centering
\includegraphics[width=\linewidth]{ch12_leeson}
\caption{\textbf{Leeson's model says the team's tank should be about 21\,dB quieter than the quoted number.} (a) Leeson's model for the 2.5\,GHz tank ($Q = 16.8$, 0.6\,V peak, $F = 2$): $-122$\,dBc/Hz at 1\,MHz. Matching $-100.8$\,dBc/Hz would need 22\,dB of excess noise (dashed). (b) $\mathcal{L}$(1\,MHz) against tank amplitude: even 50\,mV of swing would beat the quoted value.}
\label{fig:ch12-leeson}
\end{figure}

\begin{pitfall}[Is $-100.8$\,dBc/Hz the VCO, or the PLL?]
With $Q = 16.8$, $R_p = 1.06$\,k\Ohm, 0.6\,V of swing and $F = 2$, \cref{eq:ch12-leeson} predicts $-122$\,dBc/Hz at 1\,MHz. The common oscillator figure of merit,
\[\mathrm{FoM} = -\mathcal{L}(\Delta f) + 20\log_{10}(f_0/\Delta f) - 10\log_{10}(P/1\,\mathrm{mW}),\]
would be only 162\,dB for $-100.8$\,dBc/Hz at 5\,mW of VCO power, against 180--190\,dB for typical LC VCOs. The quoted value is probably the \emph{closed-loop} PLL output near the loop bandwidth, or the VCO has excess noise (tail noise, low swing). This book uses it as free-running VCO noise, the pessimistic reading; resolving it is a Phase~0 task, because it moves the optimum of \cref{we:ch12-bw}.
\end{pitfall}

\begin{insight}[Retuning the team's 2.4\,GHz design to 2.5\,GHz]
Keep the 4\,nH inductor: its $Q$ and layout are proven. The tank capacitance must fall from 1.10 to 1.01\,pF, by 86\,fF or 8\%: remove one switched-capacitor unit or shift the band code, and re-check that $V_\mathrm{ctrl}$ sits mid-range at all corners. $R_p$ rises by 4\% with frequency, so the start-up margin improves slightly. Replace the fractional-N divider with an integer $\div50$ and redesign the loop filter for $f_u \approx 3$\,MHz (\cref{we:ch12-bw}). Target $\Kvco \approx 2\pi\times150$\,MHz/V, so that supply and control noise stay small, and cover PVT with a 3--5 bit capacitor bank.
\end{insight}

%% =====================================================================
\section{Quadrature clocks and distribution}\label{sec:quadrature}

The half-rate receiver samples data on the 0$^\circ$ and 180$^\circ$ edges of the 2.5\,GHz clock and the data transitions on the 90$^\circ$ and 270$^\circ$ edges, with four samplers D0, E0, D1, E1 (\cref{sec:sampler}). It needs four phases, I, Q, $\bar{\mathrm{I}}$ and $\bar{\mathrm{Q}}$, spaced by exactly 100\,ps. \OS{} derives them from the 2.5\,GHz VCO with an \textbf{RC polyphase filter} (PPF); the alternative is a 5\,GHz VCO followed by a $\div2$ (\cref{fig:ch12-quad}).

\begin{figure}[tbp]
\centering
\begin{subfigure}[b]{0.62\linewidth}\centering
\begin{tikzpicture}[font=\sffamily\scriptsize,scale=0.72,transform shape]
  \foreach \i/\y in {0/0,1/-1.25,2/-2.5,3/-3.75}{
    \draw[color=navy] (0.6,\y) to[R] (3.0,\y) -- (3.8,\y) to[R] (6.2,\y) -- (6.8,\y);}
  \foreach \y/\yn in {0/-1.25,-1.25/-2.5,-2.5/-3.75}{
    \draw[color=navy] (1.0,\y) node[circ]{} to[C] (2.6,\yn) node[circ]{};
    \draw[color=navy] (4.2,\y) node[circ]{} to[C] (5.8,\yn) node[circ]{};}
  \draw[color=navy] (1.0,-3.75) node[circ]{} to[C] (2.6,-5.0) -- (3.0,-5.0) node[right,text=muted]{to 0$^\circ$};
  \draw[color=navy] (4.2,-3.75) node[circ]{} to[C] (5.8,-5.0) -- (6.2,-5.0) node[right,text=muted]{to 0$^\circ$};
  \foreach \t/\y in {{$V_\mathrm{in}^{+}$}/0,{$V_\mathrm{CM}$}/-1.25,{$V_\mathrm{in}^{-}$}/-2.5,{$V_\mathrm{CM}$}/-3.75}{\node[left,text=navy] at (0.6,\y) {\t};}
  \foreach \t/\y in {I/0,Q/-1.25,{$\bar{\mathrm{I}}$}/-2.5,{$\bar{\mathrm{Q}}$}/-3.75}{\node[right,text=navy] at (6.8,\y) {\t};}
  \node[lbl] at (1.8,0.7) {stage 1: $R_1, C_1$, 2.0\,GHz}; \node[lbl] at (5.0,0.7) {stage 2: $R_2, C_2$, 3.1\,GHz};
\end{tikzpicture}
\caption{Two-stage RC polyphase filter}
\end{subfigure}\hfill
\begin{subfigure}[b]{0.36\linewidth}\centering
\begin{tikzpicture}[font=\sffamily\scriptsize]
  \node[blkb,minimum width=12mm,minimum height=10mm,font=\sffamily\scriptsize] (l1) at (0,0) {D-latch\\CK};
  \node[blkb,minimum width=12mm,minimum height=10mm,font=\sffamily\scriptsize] (l2) at (2.2,0) {D-latch\\$\overline{\mathrm{CK}}$};
  \draw[sig] (l1.east) -- (l2.west);
  \draw[sig] (l2.north) -- ++(0,0.45) -| node[pos=0.25,above]{$\bar{\mathrm{Q}}$ (inverted)} ([xshift=-4mm]l1.west) -- (l1.west);
  \draw[sig] (1.1,0) node[dot]{} -- (1.1,-0.85) -- (3.4,-0.85) node[right]{I};
  \draw[sig] (l2.east) -- (3.4,0) node[right]{Q};
  \draw[clk] (-0.9,-1.45) node[left,text=golddk]{5\,GHz} -- (2.2,-1.45);
  \draw[clk] (0,-1.45) -- (l1.south); \draw[clk] (2.2,-1.45) -- (l2.south);
\end{tikzpicture}
\caption{Alternative: $\div2$ from a 5\,GHz VCO}
\end{subfigure}
\caption{\textbf{Two ways to make quadrature.} (a) Each PPF stage has an $R$ from each input phase and a $C$ from the previous phase (the last capacitor wraps around to the 0$^\circ$ output). Driven differentially, one stage gives outputs exactly 90$^\circ$ apart, with equal amplitude only at $f = 1/(2\pi RC)$; a second stage with a staggered pole keeps them equal across process. (b) A master-slave $\div2$ gives I and Q at 2.5\,GHz from both edges of a 5\,GHz clock; its accuracy depends on the 5\,GHz duty cycle, and it needs a VCO at twice the frequency.}
\label{fig:ch12-quad}
\end{figure}

\paragraph{How accurate is the polyphase filter?}
For a single stage driven differentially, the outputs are $V/(1 + j\omega RC)$ and $j\omega RC\,V/(1 + j\omega RC)$: always 90$^\circ$ apart, with amplitude ratio $\omega RC$. A $\pm20\%$ process shift of $RC$ therefore leaves 1.6--1.9\,dB of amplitude imbalance, while two stages with poles at 2.0 and 3.1\,GHz keep it below 0.07\,dB (\cref{fig:ch12-poly}a). Phase errors come from random mismatch instead: with 1\% mismatch on every $R$ and $C$, Monte Carlo gives 0.57$^\circ$ rms for one stage and 0.59$^\circ$ for two (\cref{fig:ch12-poly}b). The last stage sets the phase error, so it needs large, well-matched resistors and MIM capacitors in a common-centroid layout. The PPF is passive and lossy, $-5.2$\,dB for two unloaded stages, so CML buffers must restore the swing.

\begin{figure}[tbp]
\centering
\includegraphics[width=\linewidth]{ch12_poly}
\caption{\textbf{Two PPF stages fix the amplitude; matching sets the phase.} (a) I/Q amplitude imbalance against frequency. One stage is balanced only at its pole; two staggered stages stay within 0.07\,dB over the whole $\pm20\%$ $RC$ range (shaded) at 2.5\,GHz. (b) Quadrature phase error over 3000 Monte Carlo runs with 1\% $R$ and $C$ mismatch: about 0.6$^\circ$ rms (0.66\,ps) either way.}
\label{fig:ch12-poly}
\end{figure}

\paragraph{What the errors do to the half-rate samplers}
At 2.5\,GHz one degree is 1.11\,ps. The CDR locks the edge samplers (Q, $\bar{\mathrm{Q}}$) onto the data transitions, so an I/Q error $\delta$ moves the data samplers $\delta\times1.11$\,ps away from the eye centre: a static offset that eats timing margin. A duty-cycle error is worse, because the D1 and E1 samplers fire on the falling edges: a 1\% duty error shifts them by 4\,ps relative to D0 and E0, a two-valued DCD seen by every other bit.

\paragraph{Distribution} CML limiting buffers after the PPF restore the swing, remove the residual imbalance and keep the slew rate high (\cref{eq:ch11-slope}). Give them 2--3$\times f_\mathrm{clk}$ of bandwidth (shunt peaking or tuned loads): a buffer near $f_\mathrm{clk}$ amplifies jitter and duty-cycle error. A duty-cycle correction (DCC) loop that servoes the clock's average to zero holds the duty cycle within about $\pm0.5\%$.

\begin{example}{The RX clock-path line of the jitter budget}\label{we:ch12-iq}
\textbf{I/Q error:} two-stage PPF, 1\% mismatch: $3\sigma = 3\times0.59^\circ = 1.8^\circ \to$ \textbf{2.0\,ps}.
\textbf{Duty cycle:} $\pm0.5\%$ after DCC: $0.5\% \times 400$\,ps $\to$ \textbf{2.0\,ps}.
\textbf{PI nonlinearity:} $\pm0.6$\,ps with corrected weights (\cref{sec:picdr}) $\to$ \textbf{1.2\,ps}.
\textbf{Total} $5.2$\,ps, inside the 6\,ps allocated in \cref{tab:ch11-budget}, with 0.8\,ps of margin. Without DCC, a 52\% duty cycle alone would cost 8\,ps.
\end{example}

\begin{specnote}[Open-SERDES clock generation targets]
Integer-N CP-PLL, $f_\mathrm{ref} = 50$\,MHz, $N = 50$, $\Kvco \approx 2\pi\times150$\,MHz/V, $\Icp = 600$\,\textmu{}A, $f_u \approx 3$\,MHz ($f_{-3\mathrm{dB}} \approx 4.7$\,MHz), 60$^\circ$ phase margin. Jitter $\le 1.0$\,ps rms from 10\,kHz and $\le 0.6$\,ps above 4\,MHz. Quadrature from a two-stage PPF and CML buffers: I/Q error $\le 2^\circ$ ($3\sigma$), duty cycle $50 \pm 0.5\%$. Power estimate: PLL about 13\,mW (the team's design draws 12.73\,mW); PPF buffers and distribution 10--15\,mW (four CML stages of about 1.5\,mA at 2.5\,V).
\end{specnote}

\begin{takeaways}
\begin{itemize}
  \item The PFD turns the edge offset into an UP/DN width difference; the charge pump gives $\Icp/2\pi$ A/rad. The reset delay removes the dead zone; current mismatch causes a static offset and reference spurs.
  \item Type-II loop: $\omega_n = \sqrt{\Icp\Kvco/(2\pi NC_1)}$, $\zeta = (R/2)\sqrt{\Icp\Kvco C_1/(2\pi N)}$; peaking is 2.1\,dB at $\zeta = 0.707$ and 0.4\,dB at $\zeta = 2$.
  \item For 60$^\circ$ margin at 1\,MHz: $C_2$ = 4.07\,pF, $C_1$ = 52.6\,pF, $R$ = 11.3\,k\Ohm{} ($f_{-3\mathrm{dB}}$ = 1.56\,MHz); at the 3\,MHz optimum, $\Icp$ = 600\,\textmu{}A, $C_1$ = 17.5\,pF.
  \item Reference noise is low-pass and VCO noise high-pass; the bandwidth that equalizes them, about 3\,MHz here, minimizes jitter at 1.0\,ps rms.
  \item The 4\,nH, $Q = 16.8$ tank needs 1.01\,pF at 2.5\,GHz, has $R_p$ = 1.06\,k\Ohm{}, and oscillates if $\gm R_p/2 \ge 1$. Leeson suggests the quoted $-100.8$\,dBc/Hz is closed-loop: check it.
  \item A two-stage PPF plus CML buffers gives I/Q within about 2$^\circ$; every degree is 1.11\,ps and every 1\% of duty cycle is 4\,ps at the half-rate samplers.
\end{itemize}
\end{takeaways}

\tikzsetfigurename{13_cdr-}
\chapter{Clock and Data Recovery}\label{ch:cdr}
\chapterlead{The receiver has no clock wire: it must find the centre of every bit from the data itself, and keep finding it while the two crystals drift apart.}

\begin{objectives}
\begin{itemize}
  \item Explain why a receiver needs a CDR and compare PLL-based, PI-based and oversampling architectures.
  \item Analyze the Hogge (linear) and Alexander (bang-bang) phase detectors, including the half-rate version used by \OS{}.
  \item Model a linear CDR and a bang-bang CDR: gain $K_\mathrm{BB}$, slewing bandwidth and dither.
  \item Read a jitter-tolerance mask and size the loop to pass it.
  \item Design the \OS{} phase-interpolator CDR: PI linearity, step size, frequency-offset range and a behavioural simulation.
\end{itemize}
\end{objectives}

%% =====================================================================
\section{Why the receiver needs a CDR}\label{sec:ch13-why}

The TX and RX of a serial link run from different crystals, each accurate to about $\pm100$\,ppm. Their bit rates can differ by 200\,ppm, so the data phase drifts past a fixed RX clock by $2\times10^{-4}$\,UI every UI: one whole UI every 5000\,bits, or every microsecond at 5\,Gb/s. A fixed-phase sampler would walk out of the eye within half a microsecond, so the receiver must extract phase and frequency from the data transitions: \textbf{clock and data recovery} (CDR). Line coding or scrambling guarantees the transitions, and the CDR tracks slow jitter as a bonus (\cref{sec:ch11-jtran}).

\Cref{fig:ch13-arch} shows the three common architectures. A \textbf{PLL-based CDR} has a VCO in the loop: lowest dither, but a VCO per lane and a frequency detector to pull in. A \textbf{PI-based dual-loop CDR} shares one clean PLL and rotates its phase per lane with a digitally controlled phase interpolator: synthesizable, free of false lock and robust over PVT, at the cost of quantization and latency. An \textbf{oversampling} CDR samples three times per UI and picks the best sample, with no loop but three times the samplers. \OS{} uses the PI dual loop because it reuses the team's PLL (\cref{ch:pll}).

\begin{figure}[tbp]
\centering
\begin{subfigure}[b]{0.33\linewidth}\centering
\begin{tikzpicture}[font=\sffamily\scriptsize,b/.style={blkb,minimum width=8mm,minimum height=6mm,font=\sffamily\scriptsize,inner sep=1.5pt}]
  \node[text=navy] (d) at (-0.15,0) {D$_\mathrm{in}$};
  \node[b] (pd) at (0.95,0) {PD}; \node[b] (lf) at (2.1,0) {LF}; \node[b] (vc) at (3.25,0) {VCO};
  \node[b,fill=gtint,draw=golddk] (fd) at (2.1,-1.05) {FD};
  \draw[sig] (d) -- (pd); \draw[sig] (pd) -- (lf); \draw[sig] (lf) -- (vc);
  \draw[sig] (vc.east) -- (4.25,0) node[above,text=navy]{CK};
  \draw[sig] (3.95,0) node[dot]{} |- (0.95,0.7) -- (pd.north);
  \draw[sig] (3.95,0) |- (fd.east); \draw[sig] (fd.north) -- (lf.south);
  \draw[sig] (0.6,-1.05) node[left,text=navy]{$f_\mathrm{ref}$} -- (fd.west);
\end{tikzpicture}
\caption{PLL-based}
\end{subfigure}\hfill
\begin{subfigure}[b]{0.34\linewidth}\centering
\begin{tikzpicture}[font=\sffamily\scriptsize,b/.style={blkb,minimum width=8mm,minimum height=6mm,font=\sffamily\scriptsize,inner sep=1.5pt}]
  \node[text=navy] (d) at (-0.15,0) {D$_\mathrm{in}$};
  \node[b] (sa) at (0.95,0) {S}; \node[blkg,minimum width=8mm,minimum height=6mm,font=\sffamily\scriptsize,inner sep=1.5pt] (bb) at (2.15,0) {BBPD}; \node[blkg,minimum width=8mm,minimum height=6mm,font=\sffamily\scriptsize,inner sep=1.5pt] (dl) at (3.45,0) {DLF};
  \node[b] (pl) at (0.95,1.15) {PLL}; \node[b] (pi) at (2.15,1.15) {PI};
  \draw[sig] (d) -- (sa); \draw[sig] (sa) -- (bb); \draw[sig] (bb) -- (dl);
  \draw[sig] (pl) -- node[above]{I/Q} (pi); \draw[sig] (dl.north) |- node[pos=0.3,right]{code} (pi.east);
  \draw[clk] (pi.south) -- ++(0,-0.25) -| (sa.north);
  \draw[sig] (-0.15,1.15) node[left,text=navy]{$f_\mathrm{ref}$} -- (pl);
\end{tikzpicture}
\caption{PI-based dual loop (\OS)}
\end{subfigure}\hfill
\begin{subfigure}[b]{0.3\linewidth}\centering
\begin{tikzpicture}[font=\sffamily\scriptsize,b/.style={blkb,minimum width=8mm,minimum height=6mm,font=\sffamily\scriptsize,inner sep=1.5pt}]
  \node[text=navy] (d) at (-0.15,0) {D$_\mathrm{in}$};
  \node[b] (sa) at (1.0,0) {3$\times$S}; \node[blkg,minimum width=8mm,minimum height=6mm,font=\sffamily\scriptsize,inner sep=1.5pt] (pk) at (2.35,0) {pick};
  \node[b] (pl) at (1.0,1.15) {PLL};
  \draw[sig] (d) -- (sa); \draw[sig] (sa) -- (pk); \draw[sig] (pk) -- (3.35,0) node[above,text=navy]{data};
  \draw[clk] (pl) -- node[right]{3 phases/UI} (sa);
\end{tikzpicture}
\caption{Oversampling}
\end{subfigure}
\caption{\textbf{Three ways to recover the clock.} (a) A VCO inside the loop, with a frequency detector (FD) to pull in. (b) A global PLL plus a per-lane phase interpolator steered by a bang-bang phase detector and a digital loop filter (DLF). (c) Blind oversampling: three samples per UI and a decision on which to keep.}
\label{fig:ch13-arch}
\end{figure}

%% =====================================================================
\section{Phase detectors}\label{sec:pd}

A PLL's PFD does not work on data: random data skips edges, and a missing edge looks like a frequency error. A CDR phase detector must compare edges only where the data has them, and output \emph{nothing} during runs of identical bits (\textbf{tri-state}). Two classic detectors do this.

\paragraph{The Hogge linear detector}
Hogge's detector~\cite{p5:hogge1985} uses two flip-flops and two XORs (\cref{fig:ch13-hogge}). FF1 retimes the data (B), FF2 delays B by half a cycle (A). $X = D_\mathrm{in}\oplus B$ lasts from each data transition to the next rising clock edge, $\Tb/2 + \Delta t$ when the clock is $\Delta t$ late; $Y = B\oplus A$ always lasts $\Tb/2$. X pumps up and Y pumps down, so the net charge per transition is proportional to $\Delta t$: a linear detector of gain $TD\cdot I_P/2\pi$, where the \textbf{transition density} TD is about 0.5 for random data. It tri-states and retimes the data. Its weakness is speed: at 5\,Gb/s X and Y are 100\,ps pulses from gates with 20--30\,ps edges, and the clock-to-Q delay leaves a static error unless matched; at 10\,Gb/s the pulses shrink to 50\,ps, close to the practical limit.

\begin{figure}[tbp]
\centering
\begin{tikzpicture}[font=\sffamily\scriptsize,
   ff/.style={draw=navy,line width=0.7pt,fill=tile,minimum width=10mm,minimum height=10mm}]
  \node[ff] (f1) at (0,0) {FF1}; \node[ff] (f2) at (2.1,0) {FF2};
  \draw[sig] (-1.6,0.2) node[left,text=navy]{D$_\mathrm{in}$} -- ([yshift=2mm]f1.west);
  \draw[sig] ([yshift=2mm]f1.east) -- node[below]{B} ([yshift=2mm]f2.west);
  \draw[navy,line width=0.7pt] ([yshift=2mm]f2.east) -- (3.0,0.2) node[above,text=navy]{A};
  \draw[clk] (-1.6,-0.75) node[left,text=golddk]{CK} -- (1.35,-0.75);
  \draw[clk] (-0.9,-0.75) |- ([yshift=-2.5mm]f1.west);
  \draw[clk] (1.35,-0.75) |- ([yshift=-2.5mm]f2.west); \node[golddk] at (1.15,-0.4) {$\overline{\mathrm{CK}}$};
  \node[xor port,scale=0.7,color=navy] (x1) at (2.4,1.45) {};
  \node[xor port,scale=0.7,color=navy] (x2) at (4.0,0.7) {};
  \draw[navy,line width=0.7pt] (-1.15,0.2) node[circ]{} |- (x1.in 1);
  \draw[navy,line width=0.7pt] (0.85,0.2) node[circ]{} |- (x1.in 2);
  \draw[navy,line width=0.7pt] (1.15,0.2) node[circ]{} |- (x2.in 1);
  \draw[navy,line width=0.7pt] (3.0,0.2) |- (x2.in 2);
  \draw[sig] (x1.out) -- (5.4,1.45) node[right,text=navy]{X $\to$ up ($+I_P$)};
  \draw[sig] (x2.out) -- (5.4,0.7) node[right,text=navy]{Y $\to$ down ($-I_P$)};
  \draw[sig] (3.0,0.2) -- (3.0,-0.35) -- (5.4,-0.35) node[right,text=navy]{retimed data};
\end{tikzpicture}\par\smallskip
\begin{tikzpicture}[font=\sffamily\scriptsize,tt/.style={timing/xunit=0.2cm,timing/yunit=0.3cm,timing/slope=0.08,line width=0.7pt,draw=navy}]
  \foreach \y/\lab/\s in {0/D$_\mathrm{in}$/{8H8L16H16L10H},-0.55/CK/{H4L4H4L4H4L4H4L4H4L4H4L4H4L4HL},-1.1/B/{13H8L16H16L5H},-1.65/A/{17H8L16H16L1H},-2.2/X/{8L5H3L5H11L5H11L5H5L},-2.75/Y/{13L4H4L4H12L4H12L4H1L}}{
    \node[anchor=east,text=navy] at (-0.1,\y+0.15) {\lab};
    \timing[tt] at (0,\y) {\s};}
  \foreach \x in {1.6,3.2,6.4,9.6}{\draw[muted,dotted] (\x,0.4) -- (\x,-2.8);}
  \draw[c4,<->,line width=0.6pt] (1.6,-2.05) -- (2.6,-2.05); \node[c4,anchor=east] at (1.55,-2.05) {$\Tb/2+\Delta t$};
  \draw[c3,<->,line width=0.6pt] (2.6,-2.6) -- (3.4,-2.6); \node[c3,anchor=west] at (3.45,-2.6) {$\Tb/2$};
  \node[lbl,anchor=west] at (11.8,-0.5) {clock late by $\Delta t = \Tb/8$:};
  \node[lbl,anchor=west] at (11.8,-0.9) {X is wider than Y,};
  \node[lbl,anchor=west] at (11.8,-1.3) {net charge is up};
\end{tikzpicture}
\caption{\textbf{The Hogge detector measures phase as a pulse-width difference.} FF1 retimes the data (B), FF2 delays it by half a period (A). X spans from each data transition to the next clock edge, $\Tb/2 + \Delta t$; Y is always $\Tb/2$. With no transition, both stay low: the detector tri-states.}
\label{fig:ch13-hogge}
\figsrc{After Razavi~\cite{razaviOptical}.}
\end{figure}

\paragraph{The Alexander bang-bang detector}
The Alexander detector~\cite{p5:alexander1975} (\cref{fig:ch13-alex}) takes two samples per bit: a data sample at the bit centre and an \textbf{edge sample} half a bit earlier, at the expected transition. Three consecutive samples, S1 (previous data), S2 (edge) and S3 (current data), tell everything. If S1 = S3 there was no transition and the detector says nothing. Otherwise S2 equals whichever data bit it was closer to: if S2 = S1 the edge sample was taken \emph{before} the transition and the clock is early; if S2 = S3 it was taken after, and the clock is late. With $X = S1\oplus S2$ and $Y = S2\oplus S3$:
\begin{center}\small
\begin{tabular}{cccl}
\toprule
$(X, Y)$ & (0,0) & (0,1) & (1,0) \\
\midrule
Meaning & no transition & clock early: delay it & clock late: advance it \\
\bottomrule
\end{tabular}
\end{center}
The output is only the \emph{sign} of the error, hence ``bang-bang'' (BB). It reuses the data samplers, retimes, tri-states and runs at the samplers' full speed: the most common CDR detector.

\begin{figure}[tbp]
\centering
\begin{subfigure}[b]{0.5\linewidth}\centering
\begin{tikzpicture}[font=\sffamily\scriptsize,
   ff/.style={draw=navy,line width=0.7pt,fill=tile,minimum width=9mm,minimum height=8mm}]
  \node[ff] (fd) at (0,0.9) {FF$_D$}; \node[ff] (fe) at (0,-0.9) {FF$_E$}; \node[ff] (fp) at (1.7,0.9) {FF$_P$};
  \draw[sig] (-1.5,0) node[left,text=navy]{D$_\mathrm{in}$} -- (-0.9,0) node[circ]{} |- (fd.west);
  \draw[sig] (-0.9,0) |- (fe.west);
  \draw[sig] (fd.east) -- (fp.west);
  \node[xor port,scale=0.65,color=navy] (xl) at (3.6,1.25) {};
  \node[xor port,scale=0.65,color=navy] (xe) at (3.6,-0.35) {};
  \draw[navy,line width=0.7pt] (fp.east) -- ++(0.25,0) |- (xl.in 1);
  \draw[navy,line width=0.7pt] (fe.east) -- (2.6,-0.9) |- (xe.in 2);
  \draw[navy,line width=0.7pt] (2.6,-0.9) node[circ]{} -- (2.6,0.5) |- (xl.in 2);
  \draw[navy,line width=0.7pt] (0.75,0.9) node[circ]{} -- (0.75,0.15) -| (2.35,-0.2) -- (xe.in 1);
  \draw[sig] (xl.out) -- ++(0.35,0) node[right,text=navy]{late};
  \draw[sig] (xe.out) -- ++(0.35,0) node[right,text=navy]{early};
  \node[lbl] at (2.55,1.25) {S1}; \node[lbl] at (2.85,-1.1) {S2}; \node[lbl] at (0.95,0.35) {S3};
  \node[golddk] at (0,1.55) {CK}; \node[golddk] at (0,-1.55) {$\overline{\mathrm{CK}}$}; \node[golddk] at (1.7,1.55) {CK};
\end{tikzpicture}
\caption{Three flip-flops and two XORs}
\end{subfigure}\hfill
\begin{subfigure}[b]{0.48\linewidth}\centering
\begin{tikzpicture}[font=\sffamily\scriptsize,x=1cm,y=1cm]
  \foreach \yo/\ttl/\tx/\res/\col in {0/{clock early}/0.70/{S2 = S1: (0,1), early}/c1,-1.9/{clock late}/1.30/{S2 = S3: (1,0), late}/c4}{
    \begin{scope}[yshift=\yo cm]
      \draw[navy,line width=0.9pt] (-0.3,0) -- (0.75,0) -- (1.25,1) -- (2.6,1);
      \draw[navy,line width=0.9pt,densely dotted] (-0.3,1) -- (0.75,1) -- (1.25,0) -- (2.6,0);
      \draw[muted,line width=0.4pt] (-0.3,0.5) -- (2.6,0.5);
      \foreach \x/\l in {0/S1,\tx/S2,2.0/S3}{\draw[gold,line width=0.9pt,-{Stealth[length=1.6mm]}] (\x,1.35) -- (\x,1.05); \node[golddk] at (\x,1.5) {\l};}
      \node[anchor=west,text=\col] at (2.7,0.75) {\ttl};
      \node[anchor=west,text=\col] at (2.7,0.3) {\res};
    \end{scope}}
\end{tikzpicture}
\caption{Early and late, for a 0$\to$1 (solid) or 1$\to$0 transition}
\end{subfigure}
\caption{\textbf{The Alexander detector compares the edge sample with its neighbours.} (a) FF$_D$ samples the data (S3), FF$_E$ samples the edge half a period earlier (S2), FF$_P$ holds the previous bit (S1); in practice a fourth flip-flop retimes S2 to CK. (b) If the edge sample lands before the transition it equals S1 (clock early); after it, S3 (clock late).}
\label{fig:ch13-alex}
\end{figure}

\paragraph{Half-rate in \OS{}} At half rate the four samplers D0, E0, D1, E1 fire on I, Q, $\bar{\mathrm{I}}$, $\bar{\mathrm{Q}}$ of the 2.5\,GHz clock (\cref{sec:quadrature,sec:sampler}). Each clock cycle yields two early/late decisions, from (D0, E0, D1) and (D1, E1, D0$'$), so every transition is used. The cost is that I/Q and duty-cycle errors now appear as phase offsets between the two halves (\cref{we:ch12-iq}).

\paragraph{Mueller-M\"uller} This \textbf{baud-rate} detector~\cite{p5:mueller1976} uses one data sample and one error sample per UI and no edge samplers. Its output $z_k = e_k d_{k-1} - e_{k-1}d_k$ averages to $h_1 - h_{-1}$, so it locks where the first pre- and post-cursors are equal, a channel-dependent point. It suits ADC-based receivers; \OS{} has edge samplers and uses Alexander.

%% =====================================================================
\section{The linear CDR model and jitter transfer}\label{sec:ch13-linear}

With a linear detector, a CDR is a PLL with $N = 1$ and a detector gain $TD\cdot I_P/2\pi$. With proportional resistor $R_p$ and integrating capacitor $C_p$, \cref{eq:ch12-wnz} gives
\begin{keyeq}[Linear CDR]
\begin{equation}
\begin{gathered}
  H(s) = \frac{\phi_\mathrm{out}}{\phi_\mathrm{in}} = \frac{2\zeta\omega_n s + \omega_n^2}{s^2 + 2\zeta\omega_n s + \omega_n^2}, \\
  \omega_n = \sqrt{\frac{\Kvco I_P}{2\pi C_p}}, \qquad
  \zeta = \frac{R_p}{2}\sqrt{\frac{\Kvco I_P C_p}{2\pi}}, \qquad
  \omega_{-3\mathrm{dB}} \approx 2\zeta\omega_n = \frac{\Kvco I_P R_p}{2\pi} .
\end{gathered}
  \label{eq:ch13-lin}
\end{equation}
\end{keyeq}
with TD absorbed into $I_P$. $|H|$ is the \textbf{jitter transfer}: the jitter the recovered clock follows. $|1 - H|$ is what the clock does \emph{not} follow, and therefore what reaches the sampler as eye closure (\cref{fig:ch13-jtran}a). Wireline CDRs use $\zeta \gg 1$, typically 4--5, because any peaking would amplify jitter in a chain of repeaters; with $\zeta = 4.5$ the peaking is 0.09\,dB (\cref{tab:ch12-zeta}).

\begin{figure}[tbp]
\centering
\includegraphics[width=\linewidth]{ch13_jtran}
\caption{\textbf{A CDR is a low-pass filter for jitter, and a bang-bang CDR's corner moves with the jitter amplitude.} (a) Linear CDR with $f_{-3\mathrm{dB}}$ = 4\,MHz and $\zeta$ = 4.5: the clock tracks slow jitter ($|H|$), and the error $|1-H|$ rises at 40 then 20\,dB/decade. (b) Simulated jitter transfer of the analog bang-bang CDR of \cref{we:ch13-bb} ($R$ = 200\,\Ohm, $S = 3\times10^6$\,UI/s) for four sinusoidal-jitter amplitudes. The corner halves each time the amplitude doubles; the triangles mark the slewing prediction $\pi S/(2A_\mathrm{pk})$, within 7--22\% of simulation.}
\label{fig:ch13-jtran}
\end{figure}

%% =====================================================================
\section{Bang-bang CDR dynamics}\label{sec:bbcdr}

A BB detector has infinite gain at zero error, but jitter rescues the linear model: with Gaussian edge jitter of rms $\sigma_j$, the \emph{average} output is the sign function smoothed by the jitter PDF:
\[
  \overline{\mathrm{PD}}(\Delta\phi) = TD\cdot E\big[\sgn(\Delta\phi + j)\big] = TD\cdot\operatorname{erf}\!\Big(\frac{\Delta\phi}{\sqrt2\,\sigma_j}\Big),
\]
whose slope at zero, $TD\cdot 2/(\sqrt{2\pi}\sigma_j)$, is the linearized gain (\cref{fig:ch13-bbchar}a):
\begin{keyeq}[Bang-bang detector gain, linearized by jitter]
\begin{equation}
  K_\mathrm{BB} = TD\,\sqrt{\frac{2}{\pi}}\,\frac{1}{\sigma_j} \quad \text{(per UI of phase error)} .
  \label{eq:ch13-kbb}
\end{equation}
\end{keyeq}
The loop gain, and so the bandwidth, is inversely proportional to the input jitter: a clean input makes a fast loop, a noisy input a slow one. Always quote a BB-CDR's bandwidth together with the jitter it was measured with.

\begin{figure}[tbp]
\centering
\includegraphics[width=\linewidth]{ch13_bbchar}
\caption{\textbf{Jitter linearizes the bang-bang detector; large jitter makes the loop slew.} (a) Average detector output against phase error: the ideal sign function smoothed by Gaussian jitter, with slope $\sqrt{2/\pi}/\sigma_j$ at zero, against an ideal linear detector. (b) Analog BB-CDR (proportional path only) with 0.4\,UI$_\mathrm{pp}$ of 5\,MHz input jitter. The input slope, $\omega A_\mathrm{pk} = 6.3\times10^6$\,UI/s, exceeds the loop's slew rate $S = 3\times10^6$\,UI/s, so the recovered phase becomes a triangle that lags the input.}
\label{fig:ch13-bbchar}
\end{figure}

\paragraph{Slewing} For large jitter the detector saturates and the recovered phase moves at its maximum rate, the \textbf{slew rate} $S$. In a VCO-based loop each decision shifts the VCO frequency by $\Delta f_P = \Kvco I_P R_p$ for one bit, moving the phase by $\Delta\phi_P = \Delta f_P\,\Tb$ clock cycles, and decisions arrive at rate $TD/\Tb$; so $S = TD\,\Delta\phi_P/\Tb$. A sinusoid $A_\mathrm{pk}\sin\omega t$ is tracked as long as its peak slope $\omega A_\mathrm{pk}$ is below $S$. Beyond that the output is a triangle of peak $S\,T/4$ (\cref{fig:ch13-bbchar}b), so its gain is $\pi S/(2A_\mathrm{pk}\omega)$ and
\begin{keyeq}[Slew-limited BB-CDR bandwidth]
\begin{equation}
  S = TD\,\frac{\Delta\phi_P}{\Tb}, \qquad
  \omega_{-3\mathrm{dB}} = \frac{\pi S}{2A_\mathrm{pk}}, \qquad
  \text{dither}_\mathrm{pp} \approx 2(D+1)\,\Delta\phi_P .
  \label{eq:ch13-slew}
\end{equation}
\end{keyeq}
The last term is the \textbf{dither}: with no jitter at all, a BB loop hunts around the lock point, and a loop latency of $D$ extra decision periods lets it overshoot by $D$ more steps each way. Large $S$ helps jitter tolerance but costs dither; the integral path must be weak compared with the proportional path, $K_P/K_I \ge 16$ (for an analog loop, $R_pC_p/\Tb \gg 1$), or the loop rings.

\begin{table}[htbp]
\centering
\caption{The slewing model against a measured 90\,nm, 10\,Gb/s BB-CDR: $R_p$ = 10\,\Ohm, $C_p$ = 100\,nF, $\Kvco$ = 1\,GHz/V, $I_P$ = 300\,\textmu{}A, so $\Delta f_P$ = 3\,MHz and, at full rate with $TD$ = 0.5, $S = 1.5\times10^6$\,UI/s.}
\label{tab:ch13-slew}
\small
\begin{tabular}{lccc|lccc}
\toprule
Jitter (UI$_\mathrm{pp}$) & 0.125 & 0.25 & 0.5 & $R_p$ (\Ohm) / $C_p$ (nF) & 5 / 100 & 10 / 100 & 20 / 100 \\
Predicted $f_{-3\mathrm{dB}}$ (MHz) & 6.0 & 3.0 & 1.5 & Measured (MHz) & 1.0 & 1.5 & 2.7 \\
Measured (MHz) & 4.1 & 2.8 & 1.45 & $C_p$ = 10 / 32 / 100\,nF & \multicolumn{3}{c}{2.8 / 2.55 / 2.6} \\
\bottomrule
\end{tabular}
\end{table}

\Cref{tab:ch13-slew} checks \cref{eq:ch13-slew} against measurements. The prediction matches at 0.25 and 0.5\,UI$_\mathrm{pp}$; at 0.125\,UI the loop is no longer slewing (the input slope is below $S$) and the measured corner saturates. The bandwidth scales with $R_p$, the proportional path, and hardly depends on $C_p$: a BB-CDR is designed through its proportional step.

\begin{example}{An analog BB-CDR that fails, and the fix}\label{we:ch13-bb}
\textbf{Given:} a half-rate VCO-based BB-CDR at 5\,Gb/s, $\Kvco$ = 150\,MHz/V, $I_P$ = 100\,\textmu{}A, $R_p$ = 100\,\Ohm, $C_p$ = 1\,nF, $TD$ = 0.5, and a stressed input that leaves a timing margin TM = 0.4\,UI.

\textbf{Step:} $\Delta f_P$ = 1.5\,MHz; one cycle of the 2.5\,GHz clock is 2\,UI, so $\Delta\phi_P = 1.5\times10^6\times200\,\mathrm{ps}\times2 = 6\times10^{-4}$\,UI per bit and $S = 0.5\times6\times10^{-4}/200\,\mathrm{ps} = 1.5\times10^6$\,UI/s.
\textbf{JTOL} (\cref{sec:jtol}): $\mathrm{TM} + 2S/\omega = 0.4 + 0.48 = 0.88$\,UI$_\mathrm{pp}$ at 1\,MHz, below the 1\,UI mask point: \textbf{fails}. Simulation gives 0.95\,UI at 1\,MHz and 0.35\,UI at 2\,MHz (mask 0.5).
\textbf{Fix:} $R_p$ = 200\,\Ohm{} doubles $S$ to $3\times10^6$\,UI/s: $0.4 + 0.95 = 1.35$\,UI at 1\,MHz (simulated 1.36) and 0.61\,UI at 2\,MHz: \textbf{passes}.
\textbf{Costs:} dither with a 3-bit latency is $2\times4\times1.2\times10^{-3} = 0.01$\,UI, negligible. $R_pC_p/\Tb = 1000$, so the loop is well damped. The proportional path alone pulls in only $|\Delta f/f| < TD\,\Delta\phi_P \approx 600$\,ppm, so a frequency detector or reference loop must bring the VCO that close first.
\end{example}

%% =====================================================================
\section{Jitter tolerance}\label{sec:jtol}

\textbf{Jitter tolerance} (JTOL) is the largest sinusoidal jitter (SJ), versus frequency, that the receiver accepts while meeting the BER, measured with a stressed input that already carries the specified RJ and DJ. Let TM be the timing margin left by that stress. A linear CDR leaves the error $|1-H|$ times the input, so
\begin{equation}
  \mathrm{JTOL}(\omega) = \frac{\mathrm{TM}}{|1 - H(j\omega)|}, \qquad
  \mathrm{JTOL}_\mathrm{BB}(\omega) \approx \mathrm{TM} + \frac{2S}{\omega} .
  \label{eq:ch13-jtol}
\end{equation}
Far below the bandwidth the CDR tracks almost everything and JTOL is many UI; far above, JTOL falls to TM. A BB-CDR in its slew-limited region follows $2S/\omega$, a $-20$\,dB/decade line. Standards define a \textbf{JTOL mask} that the curve must stay above. We use a USB~3 Gen~1-style mask: 2\,UI$_\mathrm{pp}$ at 0.5\,MHz falling at $-20$\,dB/decade to 0.2\,UI at 4.9\,MHz, flat to 50\,MHz (check the exact points against the standard you target). Passing the $-20$\,dB/decade part needs $S \gtrsim \pi f(A_\mathrm{mask} - \mathrm{TM})$, about $2.5\times10^6$\,UI/s at 0.5\,MHz.

\begin{figure}[tbp]
\centering
\includegraphics[width=\linewidth]{ch13_jtol}
\caption{\textbf{Jitter tolerance against the mask.} Navy: the mask. Dotted: a linear CDR with 4\,MHz bandwidth. Red and teal: the analog BB-CDR of \cref{we:ch13-bb}; dots are simulations, dashed lines \cref{eq:ch13-jtol}. With $R_p$ = 100\,\Ohm{} the curve dips under the mask at 1--2\,MHz; 200\,\Ohm{} passes. Blue: the simulated \OS{} PI-CDR, which passes with a thin margin at 1--2\,MHz. At high frequency every BB curve flattens near 0.3\,UI, TM minus the loop's own dither.}
\label{fig:ch13-jtol}
\end{figure}

%% =====================================================================
\section{The Open-SERDES PI-based CDR}\label{sec:picdr}

\Cref{fig:ch13-picdr} shows the \OS{} CDR. The PLL supplies I and Q at 2.5\,GHz. Two \textbf{phase interpolators} (PIs) rotate them: the data PI (code $k$) clocks D0 and D1, the edge PI (code $k + 16$, a quarter period later) clocks E0 and E1. A 1:8 demultiplexer feeds the digital loop at 625\,MHz ($T_\mathrm{upd}$ = 1.6\,ns, 8\,UI), where up to eight Alexander early/late decisions per update are reduced by a majority vote to $v \in \{-1, 0, +1\}$. The loop filter is proportional plus integral:
\begin{equation}
  f \leftarrow f + K_I\,v, \qquad \mathrm{acc} \leftarrow \mathrm{acc} + K_P\,v + f, \qquad k = \mathrm{round}(\mathrm{acc}) \bmod 64 .
\end{equation}
The \textbf{frequency register} $f$ learns the TX-RX frequency offset, and the \textbf{phase accumulator} has fractional bits below the PI's 6.25\,ps LSB, so the proportional step $K_P$ can be a fraction of a PI step. The baseline uses $K_P = 1/8$\,LSB (0.78\,ps), $K_P/K_I = 128$ and a latency of two updates ($D = 1$).

\begin{figure}[tbp]
\centering
\begin{tikzpicture}[font=\sffamily\scriptsize,scale=0.95,transform shape,
   b/.style={blkb,minimum width=15mm,minimum height=8mm,font=\sffamily\scriptsize,inner sep=2pt},
   g/.style={blkg,minimum width=14mm,minimum height=8mm,font=\sffamily\scriptsize,inner sep=2pt}]
  \node[b] (pll) at (0.6,2.1) {PLL\\2.5\,GHz};
  \node[b] (iq) at (2.7,2.1) {PPF + CML\\I/Q};
  \node[b,minimum width=19mm] (pi) at (5.1,2.1) {PI$_D$: code $k$\\PI$_E$: $k+16$};
  \node[b,minimum width=17mm] (sa) at (1.6,0) {4 samplers\\D0 E0 D1 E1};
  \node[g] (de) at (3.9,0) {1:8\\demux};
  \node[g] (bb) at (6.0,0) {Alexander\\8 early/late};
  \node[g,minimum width=11mm] (vo) at (8.0,0) {vote\\$v$};
  \node[g,minimum width=10mm] (kp) at (9.8,0.6) {$K_P$};
  \node[g,minimum width=10mm] (ki) at (9.8,-0.6) {$K_I$};
  \node[g,minimum width=15mm] (fr) at (11.6,-0.6) {freq. reg.\\$f \mathrel{+}= K_I v$};
  \node[sumnode] (sm) at (13.2,0.6) {$+$};
  \node[g,minimum width=17mm] (ac) at (13.2,2.1) {phase acc.\\6 bits + fraction};
  \draw[sig] (-0.6,0) node[left,text=navy,align=right]{from\\CTLE/DFE} -- (sa);
  \draw[sig] (pll) -- (iq); \draw[sig] (iq) -- (pi);
  \draw[clk] (pi.south) -- (5.1,1.05) -| node[pos=0.25,above,text=golddk]{four phases} (sa.north);
  \draw[sig] (sa) -- (de); \draw[sig] (de) -- (bb); \draw[sig] (bb) -- (vo);
  \draw[sig] (vo.east) -- (8.85,0) node[dot]{} |- (kp.west); \draw[sig] (8.85,0) |- (ki.west);
  \draw[sig] (ki) -- (fr); \draw[sig] (kp) -- (sm); \draw[sig] (fr.east) -| (sm.south);
  \draw[sig] (sm) -- (ac); \draw[sig] (ac) -- node[above]{code $k$ (6 bits)} (pi.east);
  \draw[sig] (de.south) -- ++(0,-0.55) node[below,text=navy]{data to 1:16 deserializer};
  \draw[region] (3.05,-1.25) rectangle (14.2,1.35);
  \node[lbl,anchor=south east] at (14.15,-1.22) {digital, 625\,MHz, 1.2\,V CMOS};
  \node[lbl,anchor=south west] at (-1.4,2.75) {analog, 2.5\,GHz};
\end{tikzpicture}
\caption{\textbf{The \OS{} PI-based dual-loop CDR.} The PLL and polyphase filter make I/Q; two PIs rotate them under a 6-bit code; four half-rate samplers feed an Alexander detector through a 1:8 demux; a majority vote drives a proportional path ($K_P$) and an integral path whose register holds the frequency offset. The phase accumulator keeps fractional bits below the PI step.}
\label{fig:ch13-picdr}
\end{figure}

\paragraph{The phase interpolator} A PI forms $(1-\alpha)\,\mathrm{I} + \alpha\,\mathrm{Q}$ with two current-steered weights. For sinusoidal I and Q the output phase is $\theta = \arctan[\alpha/(1-\alpha)]$, not $90^\circ\alpha$: linear weights trace the chord between the two phasors, not the arc (\cref{fig:ch13-pi}a). With 16 codes per quadrant, the integral nonlinearity (INL) reaches $\pm4.1^\circ$ ($\pm4.5$\,ps) at $\alpha \approx 0.25$ and 0.75, and the amplitude dips to 0.71 at $\alpha = 0.5$. Weights taken from a sine/cosine table and quantized to 6 bits cut the INL to $\pm0.55^\circ$ (0.6\,ps) (\cref{fig:ch13-pi}b); a limiting buffer after the PI removes the amplitude dip. The law assumes near-sinusoidal inputs, with edges about as slow as the 100\,ps phase spacing, which the PPF and CML buffers of \cref{sec:quadrature} provide; square-wave inputs make the PI jump between phases.

\begin{figure}[tbp]
\centering
\begin{minipage}[b]{0.42\linewidth}\centering
\begin{tikzpicture}[font=\sffamily\scriptsize,scale=1.0]
  \draw[pale,line width=0.6pt] (2,0) arc (0:90:2);
  \draw[->,navy,line width=1pt] (0,0) -- (2,0) node[right]{I};
  \draw[->,navy,line width=1pt] (0,0) -- (0,2) node[above]{Q};
  \draw[c4,dashed,line width=0.8pt] (2,0) -- (0,2);
  \draw[->,c4,line width=1pt] (0,0) -- (1.5,0.5) node[right]{$\alpha = \frac14$};
  \draw[->,muted,densely dotted,line width=0.8pt] (0,0) -- (1.848,0.765);
  \draw[c4,line width=0.6pt] (0.9,0) arc (0:18.4:0.9); \node[c4] at (1.08,0.13) {$\theta$};
  \draw[golddk,line width=0.6pt] (1.25,0.416) arc (18.4:22.5:1.31);
  \node[golddk,anchor=west] at (1.85,1.0) {INL $-4.1^\circ$};
  \draw[golddk,line width=0.4pt] (1.86,0.95) -- (1.33,0.52);
  \node[dot,fill=c4] at (1,1) {}; \node[c4,anchor=south west] at (1,1) {$\alpha=\frac12$: $|v| = 0.71$};
  \node[muted,anchor=north] at (1,-0.15) {$(1-\alpha)\mathrm{I} + \alpha\mathrm{Q}$ moves on the chord};
\end{tikzpicture}
\subcaption{Phasors}
\end{minipage}\hfill
\begin{minipage}[b]{0.56\linewidth}\centering
\includegraphics{ch13_inl}
\subcaption{INL over the 64 codes}
\end{minipage}
\caption{\textbf{Linear weights make the PI nonlinear.} (a) Interpolating between I and Q with linear weights moves along the chord (dashed), so at $\alpha = 1/4$ the phase is 18.4$^\circ$ instead of 22.5$^\circ$ and the amplitude drops. (b) INL for linear weights ($\pm4.1^\circ$ = $\pm4.5$\,ps, repeating every quadrant) and for sine/cosine weights quantized to 6 bits ($\pm0.55^\circ$).}
\label{fig:ch13-pi}
\end{figure}

\paragraph{Frequency-offset range} A $\pm300$\,ppm offset (two $\pm100$\,ppm crystals plus margin) drifts the data by $300\times10^{-6}\times1.6\,\mathrm{ns} = 0.48$\,ps per update, 0.077 PI step; the integral register holds this rate and the PI rotates continuously, wrapping every 833 updates (1.33\,\textmu{}s). The proportional path alone could follow $0.78\,\mathrm{ps}/1.6\,\mathrm{ns} = 488$\,ppm, so $\pm300$\,ppm would use 61\% of its slew: hence the integral path. The PI's ceiling, one step per update, is $6.25\,\mathrm{ps}/1.6\,\mathrm{ns} = 3906$\,ppm. The rotating code sweeps the INL pattern at four times the 750\,kHz rotation rate, a 3\,MHz periodic jitter. Spread-spectrum clocking (5000\,ppm) is out of range.

\begin{example}{PI step size: 64 or 128 steps?}\label{we:ch13-step}
\textbf{Question:} does doubling the PI resolution reduce the CDR's own jitter? \textbf{Method:} the behavioural model of \code{ch13\_sim.py}: +200\,ppm offset, input RJ 1.32\,ps and $\pm10$\,ps DDJ, $K_P$ chosen so that every design has the same slew $S = 2.44\times10^6$\,UI/s (the JTOL requirement), $D$ = 1. ``Wander'' is the recovered-phase error; the peak-to-peak value is taken at the $10^{-5}$ quantiles of $1.2\times10^5$ updates.
\begin{center}\small
\begin{tabular}{llcccc}
\toprule
PI steps & loop clock & $K_P$ & wander rms & wander pp & JTOL vs mask \\
\midrule
64 & 312.5\,MHz & 1 LSB (integer) & 10.7\,ps & 58.1\,ps & 2.6$\times$ \\
64 & 312.5\,MHz & 1/4 LSB & 3.3\,ps & 24.3\,ps & pass \\
128 & 312.5\,MHz & 1/2 LSB & 2.9\,ps & 23.0\,ps & pass \\
\textbf{64} & \textbf{625\,MHz} & \textbf{1/8 LSB} & \textbf{2.5\,ps} & \textbf{19.2\,ps} & pass \\
128 & 625\,MHz & 1/4 LSB & 1.9\,ps & 15.6\,ps & pass \\
\bottomrule
\end{tabular}
\end{center}
\textbf{Answer:} at a fixed loop clock, 128 steps saves only 1.3\,ps. What matters is the proportional step in picoseconds (fixed by JTOL) and the latency in nanoseconds, since dither $\approx 2S\,T_\mathrm{lat}$; an integer-step loop is far too noisy. A 625\,MHz loop halves the latency and brings 64 steps inside the 20\,ps allocation of \cref{tab:ch11-budget}. Keep 64 steps; use 128 only if the 625\,MHz loop fails timing.
\end{example}

\begin{figure}[tbp]
\centering
\includegraphics[width=\linewidth]{ch13_sim}
\caption{\textbf{The \OS{} PI-CDR locks to a +200\,ppm input in 0.2\,\textmu{}s and then tracks it.} Behavioural simulation of the baseline (64 steps, 625\,MHz loop, $K_P$ = 1/8\,LSB, $K_P/K_I$ = 128). (a) The PI code follows the drifting data phase and wraps every 2.0\,\textmu{}s, one full rotation of 2\,UI. (b) The integral register learns the offset and settles at +200\,ppm. (c) Phase error: 74\,ps at start, locked after 0.19\,\textmu{}s, then 2.6\,ps rms tracking error.}
\label{fig:ch13-sim}
\end{figure}

\paragraph{Behavioural check} In \cref{fig:ch13-sim} the baseline starts 0.37\,UI off with +200\,ppm offset: it locks in 0.19\,\textmu{}s, the integral register settles at 200\,ppm within about 1\,\textmu{}s, and the tracking error is 2.6\,ps rms. The same model gives the JTOL curve in \cref{fig:ch13-jtol}: 3.0\,UI$_\mathrm{pp}$ at 0.5\,MHz, 1.08 at 1\,MHz, 0.51 at 2\,MHz (mask 0.5: marginal, within the sweep's 10\% resolution) and 0.32 at high frequency. Its jitter transfer has a $-3$\,dB corner near 9\,MHz for 0.2\,UI$_\mathrm{pp}$ input and 16\,MHz for 0.05\,UI$_\mathrm{pp}$, as \cref{eq:ch13-slew} predicts for a BB loop; \cref{ch:jitter} conservatively assumed 4\,MHz when counting the PLL's untracked jitter.

\begin{insight}
Size a BB-CDR in this order: $S$ from the JTOL mask; $K_P\cdot\mathrm{LSB} = S\,T_\mathrm{upd}$; the shortest latency the logic allows (dither $\propto S\,T_\mathrm{lat}$); then $K_I \le K_P/16$ and an integral register wide enough for the ppm range.
\end{insight}

\begin{tryit}[Lab 7: CDR]
\Cref{sec:lab7} builds a fuller version of \code{ch13\_sim.py}. Reproduce \cref{fig:ch13-sim}, then sweep $K_P$ and the latency $D$ and plot wander against JTOL margin: you will find the trade-off of \cref{we:ch13-step}.
\end{tryit}

\begin{specnote}[Open-SERDES CDR targets]
Half-rate Alexander BBPD; two 64-step PIs (6.25\,ps) with sine/cosine weights (INL $\le 1^\circ$); P+I loop at 625\,MHz, $K_P$ = 1/8\,LSB, $K_P/K_I$ = 128, latency $\le$ 2 updates; $\pm300$\,ppm; lock $<$ 1\,\textmu{}s; wander $\le$ 20\,ps pp; JTOL above the mask with TM = 0.4\,UI. Power estimate: PIs and buffers 6--10\,mW, digital loop about 2\,mW.
\end{specnote}

%% =====================================================================
\section{Exercises}\label{sec:ch13-ex}

\begin{enumerate}[label=\textbf{13.\arabic*},leftmargin=2.4em]
  \item TX and RX crystals are at $+100$ and $-80$\,ppm. How long does the data take to drift one UI at 5\,Gb/s, and how many PI codes per microsecond must the CDR rotate?
  \item In the Hogge detector at 5\,Gb/s, the clock is 20\,ps early. Give the widths of X and Y, and the average current for $I_P$ = 100\,\textmu{}A and $TD$ = 0.5.
  \item A BB-CDR sees $\sigma_j$ = 0.03\,UI. Compute $K_\mathrm{BB}$ for $TD$ = 0.5. By what factor does the small-signal bandwidth change if the channel doubles the jitter?
  \item For a 625\,MHz loop, find the smallest $K_P$ (in PI LSBs) that passes the mask point 2\,UI$_\mathrm{pp}$ at 0.5\,MHz with TM = 0.4\,UI. What dither would $D = 3$ give?
\end{enumerate}

\begin{takeaways}
\begin{itemize}
  \item 200\,ppm drifts one UI per microsecond, so the receiver recovers phase and frequency from the data; \OS{} uses a PI dual loop that reuses the PLL.
  \item Hogge is linear but needs 100\,ps pulses; Alexander reuses the samplers: (0,1) early, (1,0) late, (0,0) no transition.
  \item Linear CDR: a PLL with $N = 1$, $\omega_{-3\mathrm{dB}} \approx \Kvco I_P R_p/2\pi$; only $|1-H|$ of the jitter reaches the sampler.
  \item Bang-bang: $K_\mathrm{BB} = TD\sqrt{2/\pi}/\sigma_j$; slewing gives $\omega_{-3\mathrm{dB}} = \pi S/(2A_\mathrm{pk})$; JTOL $\approx$ TM + $2S/\omega$.
  \item The \OS{} PI-CDR locks to +200\,ppm in 0.2\,\textmu{}s, wanders 2.6\,ps rms and passes the mask.
\end{itemize}
\end{takeaways}


\tutpart{Putting It Together}{The complete Open-SERDES architecture and its budgets, the open-source design and verification flow, and hands-on labs that let you reproduce every key plot in this tutorial.}
\tikzsetfigurename{14_arch-}
% Part V packs many tables and figures: allow fuller float areas (affects only Part V and the appendices).
\setcounter{topnumber}{3}\setcounter{bottomnumber}{2}\setcounter{totalnumber}{5}
\renewcommand{\topfraction}{0.9}\renewcommand{\bottomfraction}{0.6}
\renewcommand{\textfraction}{0.08}\renewcommand{\floatpagefraction}{0.8}

\chapter{The Open-SERDES Architecture}\label{ch:arch}

\chapterlead{Every decision of the last thirteen chapters on a few pages: the blocks, the clocks, the number each block must meet, and what it all costs in power.}

\begin{objectives}
\begin{itemize}
  \item Read the complete \OS{} block diagram and follow a bit from the 16-bit TX word to the 16-bit RX word.
  \item Trace every clock in the chip from the 50\,MHz reference to the last divider, with its frequency and phase.
  \item Use the block specification table to find the target, and the chapter that derives it, for any block you design.
  \item Build a first-cut power budget from tail currents and convert it to energy per bit.
  \item Defend the ten main architecture decisions and name the risks that remain open.
\end{itemize}
\end{objectives}

%% =====================================================================
\section{The complete transceiver}\label{sec:arch-block}

\Cref{fig:arch} puts the whole chip on one page. It is a single-lane transceiver: one transmitter, one receiver and one clock generator on the same \sgx{} die, connected to each other through the off-chip REF channel (\cref{sec:refchannel}) or through an on-chip loopback.

\paragraph{Transmit path} A 16-bit word arrives every 3.2\,ns (312.5\,MHz), from the user or from the on-chip PRBS7/31 generator (\cref{sec:prbs}). CMOS on 1.2\,V performs the first 16:4 serialization; HBT CML on the 2.5\,V analog supply performs the 4:2 and the final half-rate 2:1, built three times so the driver receives $D_{n+1}$, $D_n$ and $D_{n-1}$ for its three FFE taps (\cref{sec:serializer,sec:ffe-impl}). The CML driver launches 1\,Vppd into \qty{50}{\ohm} per side (\cref{sec:cml-driver}).

\paragraph{Receive path} Behind the ESD diodes and the \qty{100}{\ohm} termination (\cref{sec:esd-pad}), a loopback multiplexer selects the pad or the near-end loopback. The CTLE adds 0--12\,dB of peaking (\cref{ch:ctle}), the DFE summer subtracts $w_1$ and $w_2$ times the previous decisions (\cref{sec:dfe-timing}), and four samplers decide: data D0, D1 and edges E0, E1 (\cref{sec:sampler}). The deserializer delivers 16 data and 16 edge bits at 312.5\,MHz (\cref{sec:deserializer}); the PRBS checker counts errors.

\paragraph{Clocking and control} An integer-N PLL with an LC VCO multiplies the 50\,MHz reference by $N=50$ to 2.5\,GHz (\cref{ch:pll}). The TX gets it through buffers and $\div2$ dividers; the RX gets I and Q from a polyphase filter (\cref{sec:quadrature}), rotated by two phase interpolators (PIs) under the digital CDR loop filter (\cref{sec:picdr}). An SPI port writes every setting.

\begin{figure}[!tb]
\centering
% ch14_arch.tex -- fig:arch : the complete Open-SERDES transceiver (Chapter 14)
\begin{tikzpicture}[x=1cm,y=1cm,font=\sffamily\scriptsize,
  a/.style={blkb,minimum height=9mm,minimum width=13mm,font=\sffamily\scriptsize,inner sep=1.5pt},
  d/.style={blkg,minimum height=9mm,minimum width=13mm,font=\sffamily\scriptsize,inner sep=1.5pt},
  rt/.style={font=\sffamily\scriptsize,text=muted,align=center},
  rate/.style={font=\sffamily\scriptsize\bfseries,text=slate,align=center},
  ck/.style={font=\sffamily\scriptsize,text=golddk,align=center},
  cfg/.style={-{Stealth[length=1.6mm]},draw=muted,line width=0.6pt,dotted},
  lb/.style={-{Stealth[length=2.2mm]},draw=c5,line width=0.9pt},
  pad/.style={draw=navy,fill=navy,minimum size=2.2mm,inner sep=0pt},
  bandlbl/.style={font=\sffamily\scriptsize\bfseries,text=slate,anchor=north west}]
% ------------------------------------------------------------ die and bands
\draw[navy,line width=1.1pt,rounded corners=5pt,fill=offwhite] (0,0.05) rectangle (13.0,9.1);
\fill[gtint!35] (0.08,7.22) rectangle (12.92,9.02);
\fill[bgpale]   (0.08,2.82) rectangle (12.92,7.14);
\fill[btint!40] (0.08,0.13) rectangle (12.92,2.74);
\node[bandlbl] at (0.12,9.0) {TRANSMITTER};
\node[bandlbl,anchor=south west] at (0.12,2.9) {CLOCKING AND CDR};
\node[bandlbl] at (0.12,2.72) {RECEIVER};
\node[font=\sffamily\scriptsize\bfseries,text=navy,anchor=south east] at (13.0,9.12) {\OS{} die, IHP \sgx{}};
% ------------------------------------------------------------ TX row (y = 8.0)
\node[d,minimum width=18mm] (prbs) at (1.45,8.0) {PRBS7/31 gen.\\+ data mux};
\node[d,minimum width=13mm] (s164) at (3.95,8.0) {16:4\\CMOS};
\node[a,minimum width=12mm] (s42)  at (5.95,8.0) {4:2\\CML};
\node[a,minimum width=16mm] (s21)  at (8.0,8.0) {2:1 CML $\times$3\\half-rate};
\node[a,minimum width=19mm,line width=1pt] (drv) at (10.55,8.0) {\textbf{3-tap FFE}\\CML driver};
\node[pad] (txpad) at (13.0,8.0) {};
\draw[sig,line width=1.1pt] (-0.8,8.0) -- (prbs.west);
\node[rt,anchor=south] at (-0.4,8.05) {16\,b\\in};
\draw[sig] (prbs) -- node[rt,above=0pt]{16} (s164);
\draw[sig] (s164) -- node[rt,above=0pt]{4} (s42);
\draw[sig] (s42) -- node[rt,above=0pt]{2} (s21);
\draw[sig] (s21) -- node[rt,above=0pt]{3} (drv);
\foreach \x/\r in {2.82/{312.5\,Mb/s}, 4.97/{1.25\,Gb/s}, 6.95/{2.5\,Gb/s}, 9.2/{5\,Gb/s}}{\node[rate] at (\x,8.72) {\r};}
\draw[navy,line width=1.1pt] (drv.east) -- (txpad.west); \coordinate (lbtap) at (12.25,8.0);
\node[rt,anchor=south] at (12.6,8.05) {TX$\pm$};
% ------------------------------------------------------------ channel (off-chip)
\node[blk,minimum width=16mm,minimum height=46mm,fill=white,font=\sffamily\scriptsize,align=center] (chan) at (14.15,4.8)
  {\textbf{Channel}\\[3pt]REF: 1\,m FR-4\\+ 300\,fF pads\\[3pt]$-17.4$\,dB\\at 2.5\,GHz\\[3pt]\textcolor{muted}{off-chip:}\\\textcolor{muted}{SMA, PCB}\\\textcolor{muted}{or cable}};
\draw[sig,line width=1.1pt] (txpad.east) -| (chan.north);
\draw[sig,line width=1.1pt] (chan.south) |- (13.0,1.6);
% ------------------------------------------------------------ RX row (y = 1.6), signal flows right to left
\node[pad] (rxpad) at (13.0,1.6) {};
\node[a,minimum width=14mm] (term) at (11.8,1.6) {\qty{100}{\ohm} term.\\+ ESD};
\node[a,minimum width=8mm,minimum height=7mm] (lbm) at (10.5,1.6) {LB\\mux};
\node[a,minimum width=13mm] (ctle) at (9.3,1.6) {CTLE\\0--12\,dB};
\node[sumnode,minimum size=6.5mm] (sum) at (8.05,1.6) {$\Sigma$};
\node[a,minimum width=16.5mm] (samp) at (6.55,1.6) {4 samplers\\D0 E0 D1 E1};
\node[d,minimum width=13mm] (des) at (4.45,1.6) {1:16\\deser.};
\node[d,minimum width=14mm] (chk) at (2.3,1.6) {PRBS\\checker};
\node[rt,anchor=south] at (12.76,1.63) {RX$\pm$};
\draw[navy,line width=1.1pt] (rxpad.west) -- (term.east);
\draw[sig] (term) -- (lbm);
\draw[sig] (lbm) -- (ctle);
\draw[sig] (ctle) -- (sum);
\draw[sig] (sum) -- (samp);
\node[rt,anchor=south,text=bblue] at (8.05,1.95) {DFE};
\draw[sig] (samp) -- node[rt,above=0pt]{4} (des);
\node[rate,anchor=north] at (5.42,1.1) {2.5\,Gb/s};
\draw[sig] (des) -- node[rt,above=0pt]{32} (chk);
\node[rate,anchor=north] at (3.37,1.1) {312.5\,M};
\draw[sig,line width=1.1pt] (chk.west) -- (-0.8,1.6);
\node[rt,anchor=south] at (-0.4,1.65) {16\,b\\out};
% DFE feedback: decisions -> taps w1, w2 -> summer
\draw[sigb] ([xshift=4mm]samp.south) -- ++(0,-0.52) -| (sum.south);
\node[font=\sffamily\scriptsize,text=bblue,anchor=north] at (7.6,0.6) {taps $w_1$, $w_2$};
% ------------------------------------------------------------ clocking row (y = 5.1)
\node[a,minimum width=19mm] (pll) at (1.6,5.1) {\textbf{PLL}, $N=50$\\LC VCO 2.5\,GHz};
\node[a,minimum width=17mm] (ppf) at (4.25,5.1) {I/Q polyphase\\+ CML buffers};
\node[a,minimum width=21mm] (pi)  at (6.7,5.1) {2 phase interp.\\64 steps (6.25\,ps)};
\node[a,minimum width=55mm,minimum height=5.5mm] (txck) at (5.9,6.55) {TX clock buffers and $\div2$, $\div4$, $\div8$ dividers};
\node[d,minimum width=18mm,minimum height=7.5mm] (lf) at (4.45,3.6) {CDR loop filter\\P + I, 625\,MHz};
\node[d,minimum width=19mm] (spi) at (10.4,5.1) {SPI +\\config registers};
% refclk
\draw[clk] (-0.8,5.1) -- (pll.west);
\node[pad] at (0,5.1) {};
\node[ck,anchor=south] at (-0.4,5.15) {REF\\50\,M};
% PLL -> TX clock tree -> serializer
\draw[clk] (pll.north) |- node[ck,pos=0.3,left=0pt]{2.5\,GHz} (txck.west);
\draw[clk] (3.35,6.55 |- txck.north) -- (3.35,7.05) -| (prbs.south);
\node[ck,anchor=east] at (1.4,7.3) {312.5\,MHz};
\draw[clk] (txck.north -| s164.south) -- node[ck,right=0pt,pos=0.45]{625\,M} (s164.south);
\draw[clk] (txck.north -| s42.south) -- node[ck,right=0pt,pos=0.45]{1.25\,G} (s42.south);
\draw[clk] (txck.north -| s21.south) -- node[ck,right=0pt,pos=0.45]{2.5\,G} (s21.south);
% PLL -> PPF -> PI -> samplers
\draw[clk] (pll.east) -- (ppf.west);
\draw[clk] (ppf.east) -- node[ck,above=0pt]{I, Q} (pi.west);
\draw[clk] (7.15,5.1 |- pi.south) -- coordinate[pos=0.78] (rck) node[ck,pos=0.3,right=0pt]{\ang{0}/\ang{90}/\ang{180}/\ang{270}} (7.15,1.6 |- samp.north);
% recovered clock to the deserializer
\fill[gold] (rck) circle (1.3pt);
\draw[clk] (rck) -| node[ck,pos=0.25,above=0pt]{recovered $\div2^n$} (4.85,1.6 |- des.north);
% CDR: deserializer -> loop filter -> PI code
\draw[sigg] (4.05,1.6 |- des.north) -- node[rt,left=0pt]{edge + data} (4.05,3.6 |- lf.south);
\draw[sigg] (lf.east) -| node[rt,pos=0.25,above=0pt]{code $k$} (6.2,5.1 |- pi.south);
% configuration (dotted)
\draw[cfg] (spi.north) -- node[rt,right=0pt,pos=0.5]{FFE taps} (spi.north |- drv.south);
\draw[cfg] (9.7,5.1 |- spi.south) -- node[rt,left=0pt,pos=0.55]{peaking} (9.7,1.6 |- ctle.north);
\draw[cfg] (spi.west) -- node[rt,above=0pt]{override} (pi.east);
\node[rt,anchor=north west,align=left,text width=19mm] at (10.0,4.45) {also DFE and offset DACs, PLL band, PRBS mode, loopback};
% near-end loopback
\fill[c5] (lbtap) circle (1.5pt);
\draw[lb] (lbtap) -- (12.25,2.4) -| (lbm.north);
\node[font=\sffamily\scriptsize,text=c5,rotate=90,anchor=south] at (12.25,6.1) {near-end loopback};
% ------------------------------------------------------------ legend
\begin{scope}[shift={(0.3,-0.45)}]
  \node[blkb,minimum width=6mm,minimum height=3.2mm,inner sep=0pt] at (0,0) {};
  \node[rt,anchor=west,text=navy] at (0.4,0) {analog / HBT CML (2.5\,V)};
  \node[blkg,minimum width=6mm,minimum height=3.2mm,inner sep=0pt] at (3.85,0) {};
  \node[rt,anchor=west,text=navy] at (4.25,0) {digital CMOS (1.2\,V)};
  \draw[clk] (7.25,0) -- (8.05,0); \node[rt,anchor=west,text=navy] at (8.1,0) {clock};
  \draw[lb] (9.05,0) -- (9.85,0); \node[rt,anchor=west,text=navy] at (9.9,0) {loopback};
  \draw[cfg] (11.25,0) -- (12.05,0); \node[rt,anchor=west,text=navy] at (12.1,0) {configuration};
\end{scope}
\end{tikzpicture}

\caption{\textbf{The complete \OS{} transceiver: a CMOS core at 312.5\,MHz, wrapped in HBT CML wherever a signal or clock runs at 1.25\,GHz or faster.} Read the TX band left to right and the RX band right to left; bold numbers above the TX row are lane rates, and slash numbers count parallel lanes. Notice three things. The 2.5\,GHz PLL clock is used twice: divided down for the TX (fixed phase), and split into I/Q and rotated by the PIs for the RX (recovered phase). The CDR is a digital loop: edge and data bits leave the deserializer, the loop filter computes the PI code, and the samplers move. The near-end loopback (violet) takes the driver output straight to the RX input, so the whole link can be tested without a board channel.}
\label{fig:arch}
\end{figure}

\subsection{The clocking plan}\label{sec:ch14-clock}

\Cref{fig:ch14-clock} follows each clock from the reference to the last flip-flop. All of the frequencies come from one relation between the bit rate, the half-rate clock and the word rate.

\begin{keyeq}[The \OS{} clock frequencies]
\begin{equation}
  f_\mathrm{clk} = \frac{\fb}{2} = N f_\mathrm{ref} = \qty{2.5}{\giga\hertz},\qquad
  f_\mathrm{word} = \frac{\fb}{16} = \qty{312.5}{\mega\hertz},\qquad
  \Delta t_\mathrm{PI} = \frac{1}{64\,f_\mathrm{clk}} = \qty{6.25}{\pico\second}
  \label{eq:ch14-clk}
\end{equation}
\end{keyeq}

The TX needs one clock phase: the final 2:1 mux uses both edges of the 2.5\,GHz clock as its select (\cref{sec:serializer}). The RX needs four phases half a UI apart: data samplers on the edges of I (eye centers), edge samplers on the edges of Q (crossings), all moved together by the CDR in 6.25\,ps steps. The RX dividers run from the \emph{recovered} clock, so the deserializer and the CDR logic sit in a clock domain that is not phase-locked to the TX word clock; a small FIFO crosses back to the core (\cref{sec:deserializer}). In near-end loopback both ends share the PLL, so the frequency offset is zero. Through a board to a second chip, the CDR must absorb up to $\pm300$\,ppm (\cref{sec:picdr}).

\begin{figure}[!tbp]
\centering
% ch14_clock.tex -- fig:ch14-clock : Open-SERDES clocking plan (frequencies and phases)
\begin{tikzpicture}[x=0.94cm,y=1cm,font=\sffamily\scriptsize,
  a/.style={blkb,minimum height=7.5mm,minimum width=12mm,font=\sffamily\scriptsize,inner sep=1.5pt},
  dv/.style={blkb,minimum height=6mm,minimum width=8mm,font=\sffamily\scriptsize,inner sep=1pt},
  fq/.style={font=\sffamily\scriptsize\bfseries,text=golddk,align=center},
  ds/.style={font=\sffamily\scriptsize,text=navy,align=center},
  tp/.style={circle,fill=gold,inner sep=0pt,minimum size=3pt}]
% ---------------- PLL core
\node[fq,anchor=east] (ref) at (0.35,0) {REFCLK\\50\,MHz\\(20\,ns)};
\node[a,minimum width=16mm] (pfd) at (1.4,0) {PFD + CP\\+ loop filter};
\node[a,minimum width=14mm] (vco) at (3.8,0) {LC VCO\\4\,nH, $Q$ = 16.8};
\node[dv] (div50) at (2.6,-1.05) {$\div50$};
\draw[clk] (ref.east) -- (pfd.west);
\draw[sig,draw=navy] (pfd.east) -- node[ds,above=0pt]{$V_\mathrm{ctrl}$} (vco.west);
\coordinate (v) at (5.0,0);
\draw[clk] (vco.east) -- (v);
\fill[gold] (v) circle (1.6pt);
\draw[clk] (v) |- (div50.east);
\draw[clk] (div50.west) -| (pfd.south);
\node[fq,anchor=west,align=left] at (5.1,0) {2.5\,GHz\\\textcolor{muted}{400\,ps}};
% ---------------- TX branch (y = 1.45)
\node[a] (tbuf) at (6.2,1.45) {CML\\buffer};
\node[dv] (t1) at (8.55,1.45) {$\div2$};
\node[dv] (t2) at (10.85,1.45) {$\div2$};
\node[dv] (t3) at (13.15,1.45) {$\div2$};
\draw[clk] (v) |- (tbuf.west);
\draw[clk] (tbuf) -- coordinate[pos=0.3] (tp0) node[fq,pos=0.66,above=1pt]{2.5\,GHz} (t1);
\draw[clk] (t1) -- coordinate[pos=0.3] (tp1) node[fq,pos=0.64,above=1pt]{1.25\,GHz} (t2);
\draw[clk] (t2) -- coordinate[pos=0.3] (tp2) node[fq,pos=0.64,above=1pt]{625\,MHz} (t3);
\draw[clk] (t3) -- coordinate[pos=0.3] (tp3) node[fq,pos=0.68,above=1pt]{312.5\,MHz} (15.45,1.45);
\foreach \p/\t in {tp0/{final 2:1 select\\(both edges: 5\,Gb/s)}, tp1/{4:2 CML,\\CMOS$\to$CML}, tp2/{CMOS 8:4}, tp3/{CMOS 16:8,\\TX word clock}}{
  \node[tp] at (\p) {}; \draw[clk,-] (\p) -- ++(0,-0.35); \node[ds,anchor=north] at ($(\p)+(0,-0.35)$) {\t};}
\node[lblb,anchor=west] at (5.55,2.3) {TX: one fixed phase, divided down};
% ---------------- RX branch (y = -1.75)
\node[a,minimum width=15mm] (ppf) at (6.1,-1.75) {2-stage PPF\\+ CML buffers};
\node[a,minimum width=16mm] (pi) at (8.6,-1.75) {PI$_D$: code $k$\\PI$_E$: $k+16$};
\node[dv] (r1) at (10.75,-1.75) {$\div2$};
\node[dv] (r2) at (12.4,-1.75) {$\div2$};
\node[dv] (r3) at (14.05,-1.75) {$\div2$};
\draw[clk] (v) |- (ppf.west);
\draw[clk] (ppf) -- coordinate (mq) (pi);
\node[ds,text=golddk,anchor=north] at ($(mq)+(0,-0.45)$) {I, Q, $\bar{\mathrm{I}}$, $\bar{\mathrm{Q}}$};
\draw[clk] (pi) -- coordinate[pos=0.25] (rp0) (r1);
\draw[clk] (r1) -- coordinate[pos=0.25] (rp1) (r2);
\draw[clk] (r2) -- coordinate[pos=0.25] (rp2) (r3);
\draw[clk] (r3) -- coordinate[pos=0.25] (rp3) (15.45,-1.75);
\foreach \p/\t/\dx in {rp0/{\textcolor{golddk}{\bfseries 2.5\,GHz}\\samplers D0 D1\\and E0 E1}/-0.1, rp1/{\textcolor{golddk}{\bfseries 1.25\,GHz}\\CMOS 1:2}/0, rp2/{\textcolor{golddk}{\bfseries 625\,MHz}\\CDR loop\\(8\,UI)}/0, rp3/{\textcolor{golddk}{\bfseries 312.5\,MHz}\\1:16 + FIFO}/0.25}{
  \node[tp] at (\p) {}; \draw[clk,-] (\p) -- ++(0,-0.35); \node[ds,anchor=north] at ($(\p)+(\dx,-0.35)$) {\t};}
\node[lblb,anchor=west] at (5.4,-0.85) {RX: four phases, rotated by the CDR};
% ---------------- timing of the four sampling phases (lower panel)
\begin{scope}[shift={(4.9,-5.05)},x=0.0185cm]
  \foreach \c/\n in {0/{bit $n$ (even)},200/{bit $n{+}1$ (odd)},400/{bit $n{+}2$}}{
    \draw[navy,line width=0.8pt,fill=btint] (\c-100,0.55) -- (\c-75,0.95) -- (\c+75,0.95) -- (\c+100,0.55) -- (\c+75,0.15) -- (\c-75,0.15) -- cycle;
    \node[ds] at (\c,0.55) {\n};}
  \draw[navy,line width=0.8pt] (-150,0.95) -- (-125,0.95) -- (-100,0.55) -- (-125,0.15) -- (-150,0.15);
  \node[ds,anchor=east] at (-160,0.55) {data at sampler};
  \draw[golddk,line width=0.9pt] (-150,-0.55) -- (0,-0.55) -- (0,-0.05) -- (200,-0.05) -- (200,-0.55) -- (400,-0.55) -- (400,-0.05) -- (550,-0.05);
  \node[ds,anchor=east] at (-160,-0.3) {clock I (\ang{0})};
  \draw[golddk,line width=0.9pt,dashed] (-150,-1.35) -- (100,-1.35) -- (100,-0.85) -- (300,-0.85) -- (300,-1.35) -- (550,-1.35);
  \node[ds,anchor=east] at (-160,-1.1) {clock Q (\ang{90})};
  \foreach \t/\s in {0/D0,100/E0,200/D1,300/E1,400/D0}{
    \draw[c4,densely dotted,line width=0.7pt] (\t,1.1) -- (\t,-1.55);
    \node[font=\sffamily\scriptsize\bfseries,text=c4,anchor=north] at (\t,-1.55) {\s};}
  \draw[<->,muted,line width=0.5pt] (0,1.25) -- node[ds,above=0pt]{1\,UI = 200\,ps} (200,1.25);
  \draw[<->,muted,line width=0.5pt] (0,-2.05) -- node[ds,below=0pt]{$T_\mathrm{clk}$ = 400\,ps = 2\,UI = 64 PI steps of 6.25\,ps} (400,-2.05);
\end{scope}
\end{tikzpicture}

\caption{\textbf{One 2.5\,GHz PLL serves the whole chip; the TX divides it, the RX rotates it.} Top: the frequency at every node. The TX chain divides by 2, 4 and 8 down to the 312.5\,MHz word clock. The RX chain first makes four phases with the polyphase filter, rotates them with the two PIs, and divides the recovered clock for the deserializer and for the CDR loop, which updates every 8\,UI at 625\,MHz. Bottom: what the four phases do. D0 and D1 sample the eye centers on the two edges of I; E0 and E1 sample the crossings on the edges of Q. One 400\,ps clock period carries two bits and contains 64 PI steps.}
\label{fig:ch14-clock}
\end{figure}

\begin{insight}
Three boundaries organize the chip, and most integration bugs live on them: the \textbf{speed boundary} between CMOS and CML (1.25\,Gb/s in the TX, 2.5\,Gb/s in the RX), which needs level converters and half a slow-clock cycle of timing margin; the \textbf{supply boundary} between 1.2 and 2.5\,V, where every $V_\mathrm{CE}$ must stay below $\mathrm{BV}_\mathrm{CEO}$; and the \textbf{clock-domain boundary} between the TX clock and the recovered RX clock. Draw all three on the floorplan first.
\end{insight}

%% =====================================================================
\section{Block specifications}\label{sec:ch14-specs}

\Cref{tab:specs} is the contract between the system model and the circuit designers: one row per block, the numbers that block must meet, and the chapter where each number is derived. Values marked $^\dagger$ are first estimates from hand analysis or behavioral models. They must be confirmed by simulation in \sgx{}, and some of them will move.

\begin{table}[!tbp]
\centering
\caption{\textbf{\OS{} block specifications.} $^\dagger$\,First estimate, not yet confirmed by circuit simulation in \sgx{}.}
\label{tab:specs}
\footnotesize
\renewcommand{\arraystretch}{1.12}
\begin{tabularx}{\linewidth}{@{}>{\raggedright\arraybackslash}p{19mm} >{\raggedright\arraybackslash}p{27mm} L >{\raggedright\arraybackslash}p{19mm}@{}}
\toprule
\textbf{Block} & \textbf{Key specification} & \textbf{Target value} & \textbf{Derived in}\\
\midrule
Link & rate, BER, channel & 5\,Gb/s NRZ; BER $\le10^{-12}$ over REF ($-17.4$\,dB at 2.5\,GHz); margin up to 20\,dB & \cref{sec:os-targets,sec:refchannel}\\
Serializer & ratio, clocks, jitter & 16:1 from 16\,b at 312.5\,MHz; CMOS 16:4, CML 4:2, half-rate CML 2:1; DCD + DDJ $\le8$\,ps$^\dagger$ & \cref{sec:serializer}\\
TX driver & topology, swing & HBT CML, \qty{50}{\ohm} back-termination per side; 0.8--1.0\,Vppd ($I_T$ = 16--20\,mA); every $V_\mathrm{CE}<\mathrm{BV}_\mathrm{CEO}\approx1.6$\,V at all corners & \cref{sec:cml-driver}\\
TX FFE & taps, range, resolution & $c_{-1}, c_0, c_1$ with $\sum\lvert c_k\rvert=1$; $c_1$ to $-0.25$ (6\,dB), $c_{-1}$ to $-0.1$; 5-bit current cells (steps of 1/32) & \cref{sec:ffe-taps,sec:ffe-impl}\\
Termination, pad & impedance, $C_p$, return loss & \qty{100}{\ohm} diff., $\pm10$\,\% after 3--4-bit trim; $C_p\le0.85$\,pF (0.63\,pF estimated$^\dagger$); RL $\ge10$\,dB at $\fN$ & \cref{sec:termination,sec:esd-pad}\\
CTLE & peaking, poles, noise & 0--12\,dB programmable; $f_{p1}\approx2.5$\,GHz, $f_{p2}\ge7.5$\,GHz; noise $\le0.25$\,mV rms$^\dagger$ & \cref{ch:ctle}\\
DFE & taps, loop timing & half-rate, 2 taps, direct feedback; $w_1$ up to $0.5\,p_0$, $w_2$ to $\pm0.15\,p_0$; tap-1 loop $\approx$140\,ps$^\dagger$ $<$ 1\,UI; loop unrolling as fallback & \cref{sec:dfe-timing}\\
Samplers & count, sensitivity, offset & 4 HBT CML master--slave (D0 E0 D1 E1); $V_\mathrm{sens}\approx20$\,mV; offset $\sigma\approx4.9$\,mV raw$^\dagger$, $\pm0.46$\,mV after 5-bit trim; $\tau\le5$\,ps; aperture $\le0.2$\,UI & \cref{sec:sampler,sec:offset}\\
Deserializer & ratio, output & 1:16 CMOS; 16 data + 16 edge bits at 312.5\,MHz; two extra latches on the edge path; FIFO to the core & \cref{sec:deserializer,sec:metastability}\\
PLL & architecture, jitter & integer-N CP-PLL, $f_\mathrm{ref}$ = 50\,MHz, $N$ = 50; LC VCO 2.5\,GHz (4\,nH, $Q$ = 16.8); $\Kvco\approx2\pi\times150$\,MHz/V; $f_u\approx3$\,MHz; RJ $\le1.0$\,ps rms$^\dagger$ & \cref{ch:pll}\\
Quadrature & accuracy & 2-stage polyphase + CML buffers; I/Q error $\le\ang{2}$ (3$\sigma$); duty cycle $50\pm0.5$\,\% with DCC & \cref{sec:quadrature}\\
PI & resolution, linearity & 2 PIs, 64 steps per 400\,ps (6.25\,ps = 0.031\,UI); INL $\le\ang{1}$ & \cref{sec:picdr}\\
CDR & loop, tracking & half-rate Alexander BBPD; P+I at 625\,MHz, $K_P$ = 1/8 LSB, $K_P/K_I$ = 128; $\pm300$\,ppm; lock $<1$\,\textmu{}s; wander $\le20$\,ps pp$^\dagger$; JTOL above the USB-style mask & \cref{sec:bbcdr,sec:jtol}\\
Timing budget & total jitter & TJ $\le0.3$\,UI (60\,ps) at $10^{-12}$; estimated 58.6\,ps$^\dagger$ & \cref{sec:jitter-budget}\\
Voltage budget & margin at the sampler & $\ge50$\,mV after noise, offset and sensitivity; 76\,mV on REF with 2-tap de-emphasis$^\dagger$ & \cref{sec:noisebudget}\\
Test & on-chip & PRBS7/31 generator and checker; near-end loopback; SPI registers & \cref{sec:labtest}\\
Power & total, per bit & $\approx$216\,mW, 43\,pJ/b$^\dagger$ & \cref{sec:power-budget}\\
\bottomrule
\end{tabularx}
\end{table}

\paragraph{The link budget in one paragraph} Both budgets close (\cref{fig:ch14-budget}). In voltage, REF with only the 2-tap TX de-emphasis leaves a 131\,mV statistical eye at $10^{-12}$; noise ($2Q\sign=21$\,mV), jitter (3.5\,mV), offset (10\,mV) and sensitivity (20\,mV) leave a \textbf{76\,mV} margin by linear addition, 89\,mV statistically (\cref{sec:noisebudget}), and the full partition ($-6$\,dB preset + 2-tap DFE, 172\,mV noise-free PRBS7 eye, \cref{sec:eq-partition}) only adds to it. In time, 40\,ps of DJ plus $2Q\times1.32$\,ps of RJ give TJ = \textbf{58.6\,ps = 0.293\,UI}, just inside 0.3\,UI (\cref{sec:jitter-budget}). Timing is the tighter budget, and its largest line is the CDR's own dither.

\begin{figure}[!tbp]
\centering
\includegraphics[width=\linewidth]{ch14_budget}
\caption{\textbf{Both budgets close, but timing is the tight one.} (a) The 131\,mV statistical eye of REF with 2-tap de-emphasis, minus the terms of \cref{sec:noisebudget}: 76\,mV of margin remains, more than half the eye. (b) One 200\,ps UI minus the deterministic jitter (linear sum) and $2Q$ times the random jitter (RSS): 141\,ps remains, only 1.4\,ps more than the 0.7\,UI target allows.}
\label{fig:ch14-budget}
\end{figure}

%% =====================================================================
\section{Power budget}\label{sec:power-budget}

A CML gate draws its tail current whatever the data, so a CML-heavy chip has a power budget you can write down before any layout: list the tails, multiply by the supply. CMOS logic burns $\alpha C V^2 f$ instead.

\begin{keyeq}[CML power and energy per bit]
\begin{equation}
  P_\mathrm{CML} = V_\mathrm{CC}\sum_i I_{T,i},\qquad P_\mathrm{CMOS}=\alpha\,C_\mathrm{sw}V_\mathrm{DD}^2f,\qquad E_b = \frac{P_\mathrm{total}}{\fb}
  \label{eq:ch14-eb}
\end{equation}
\end{keyeq}

\Cref{tab:ch14-power} collects the estimates of \cref{ch:tx,ch:rx,ch:pll,ch:cdr} and \cref{fig:ch14-power} plots them. Assumptions: 1\,Vppd swing (20\,mA driver tail); CML gates as in \cref{we:ch9-hbt}, about 1.2\,mA each with emitter followers; CML on 2.5\,V, CMOS on 1.2\,V; the team's measured PLL power of 12.73\,mW~\cite{smacd2026pll}. The polyphase buffers, the PIs and the sampler clock buffers are separate circuits and are counted separately.

\begin{table}[!tbp]
\centering
\caption{\textbf{First-cut \OS{} power budget.} Currents from the checklists and target boxes of \cref{ch:tx,ch:rx,ch:pll,ch:cdr}; every line is an estimate to be replaced by simulation.}
\label{tab:ch14-power}
\footnotesize
\begin{tabularx}{\linewidth}{@{}l >{\raggedright\arraybackslash}p{52mm} L r@{}}
\toprule
& \textbf{Blocks} & \textbf{Assumption} & \textbf{mW}\\
\midrule
TX & 3-tap CML driver & $I_T$ = 20\,mA from 2.5\,V (40\,mW at 0.8\,Vppd) & 50.0\\
& CML 4:2, 2:1 ($\times3$), latches, pre-drivers, clock buffers, $\div2$ & about 17 CML gates of 1.2\,mA & 51.0\\
& CMOS 16:4 & 1.2\,V, 625\,MHz & 3.0\\
RX & CTLE + buffer, DFE summers and tap DACs & 5\,mA + 6\,mA & 27.5\\
& 4 CML samplers + their clock buffers & 12\,mA + 4\,mA & 40.0\\
& CMOS 1:16 deserializer, offset DACs & 1.2\,V; 5 bits $\times4$ & 5.6\\
Clocking & PLL (team design) & measured, 130\,nm BiCMOS~\cite{smacd2026pll} & 12.7\\
& polyphase + I/Q buffers, distribution, 2 PIs & 6\,mA + 8\,mW (\cref{sec:picdr}) & 23.0\\
Digital & CDR loop filter, PRBS gen.\ + checker, SPI & CMOS at 625 and 312.5\,MHz & 3.2\\
\midrule
\textbf{Total} & & TX 104, RX 73, clocking 36, digital 3 & \textbf{216}\\
\bottomrule
\end{tabularx}
\end{table}

\begin{figure}[!tbp]
\centering
\includegraphics[width=\linewidth]{ch14_power}
\caption{\textbf{The CML driver alone is almost a quarter of the chip's power; CMOS logic is a rounding error.} (a) The first-cut budget of \cref{tab:ch14-power}, colored by section. The TX takes 48\,\%, mostly static CML current; the four samplers are the largest RX item. (b) Energy per bit against the Purdue OpenSerDes, an all-digital 2\,Gb/s link in SKY130 at about 438\,mW~\cite{openserdes2021}, and against \OS{} with the voltage-mode (SST) driver of \cref{sec:cml-driver}.}
\label{fig:ch14-power}
\end{figure}

\begin{example}{Energy per bit, and where to save}\label{we:ch14-power}
\textbf{Total:} the CML tails add up to 77\,mA from 2.5\,V (192\,mW); the PLL, the CMOS blocks and the DACs add 25\,mW: $P\approx\qty{216}{\milli\watt}$. By \cref{eq:ch14-eb}, $E_b = 216\,\mathrm{mW}/5\,\mathrm{Gb/s} = \mathbf{43.2}$\,pJ/b.

\textbf{Comparison:} the Purdue OpenSerDes spends $438\,\mathrm{mW}/2\,\mathrm{Gb/s} = 219$\,pJ/b~\cite{openserdes2021}. \OS{} runs 2.5 times faster for about half the power, five times less energy per bit, because the HBTs do in a few CML gates what an all-digital 130\,nm CMOS design must do with oversampling and many more gates.

\textbf{Savings, largest first:} (1) the SST driver of \cref{sec:cml-driver} saves about 44\,mW, giving 172\,mW and 34.4\,pJ/b; (2) a swing of 0.8\,Vppd saves 10\,mW; (3) each CML gate removed from the clock paths saves about 3\,mW. Lowering the CML supply would save the most of all, but every $V_\mathrm{CE}$ must still fit below $\mathrm{BV}_\mathrm{CEO}$ with the swing on top (\cref{sec:cml-driver}), so this waits for silicon data.
\end{example}


%% =====================================================================
\section{Design decisions}\label{sec:ch14-decisions}

\Cref{tab:decisions} lists the ten choices that shape \OS{}. Each row says what was chosen, what was rejected and why, and where the argument is made in full. Use it in design reviews: if someone proposes to change a row, the chapter in the last column says which numbers must be recomputed.

\begin{table}[!tbp]
\centering
\caption{\textbf{The ten main \OS{} architecture decisions.}}
\label{tab:decisions}
\footnotesize
\renewcommand{\arraystretch}{1.12}
\begin{tabularx}{\linewidth}{@{}>{\raggedright\arraybackslash}p{17mm} >{\raggedright\arraybackslash}p{22mm} >{\raggedright\arraybackslash}p{21mm} L >{\raggedright\arraybackslash}p{16mm}@{}}
\toprule
\textbf{Decision} & \textbf{Chosen} & \textbf{Alternatives} & \textbf{Why} & \textbf{Ch.}\\
\midrule
Modulation & NRZ & PAM4 & on REF NRZ wins by 2.6\,dB; PAM4 pays only above about 24\,dB and needs 3 slicers and linear circuits & \cref{sec:pam4}\\
Clock rate & half rate, 2.5\,GHz & full rate, quarter rate & full rate needs 5\,GHz clocks; quarter rate needs 8 calibrated phases & \cref{sec:serializer,sec:quadrature}\\
Fast logic & SiGe HBT CML & CMOS, MOS CML & full switching at 100\,mV, gain 5.8 vs.\ 2; CMOS misses the 8\,FO4 rule at 2.5\,GHz, slow corner & \cref{sec:serializer}\\
TX driver & CML & SST (voltage mode) & resistor-set impedance, constant current, taps leave the impedance alone; SST ($-44$\,mW) is the upgrade & \cref{sec:cml-driver}\\
TX FFE & 3 taps & 2 taps, 4+ taps & post-tap does most of the work; pre-tap for $p_{-1}$ on longer channels (on REF, 2 taps win: Lab~3) & \cref{ch:txffe,sec:lab3}\\
RX equalizer & CTLE 0--12\,dB + 2-tap DFE & CTLE only, more DFE taps, RX FFE & DFE removes $p_1$, $p_2$ without noise; CTLE adds gain, adaptivity and reach above 22\,dB & \cref{sec:eq-partition}\\
VCO & LC, 4\,nH, $Q$ = 16.8 & ring VCO & 15--30\,dB lower phase noise; reuses the team's proven design & \cref{sec:lcvco}\\
CDR & PI dual loop, digital BB & VCO-based BB, linear (Hogge) & reuses the clean PLL; synthesizable; PVT-independent gain; no false lock & \cref{sec:picdr}\\
Quadrature & 2-stage polyphase + CML & 5\,GHz VCO + $\div2$, quadrature VCO & keeps the 2.5\,GHz VCO; $\div2$ needs a 5\,GHz VCO with an accurate duty cycle & \cref{sec:quadrature}\\
On-chip test & PRBS7/31 + checker, loopback & external BERT only & scope and PC suffice; $10^{-12}$ in 10 minutes; loopback isolates the chip & \cref{sec:labtest}\\
\bottomrule
\end{tabularx}
\end{table}

%% =====================================================================
\section{Risks and open questions}\label{sec:ch14-risks}

The architecture closes on paper. It can still fail in these places, roughly in order of the cost of a late fix:
\begin{itemize}
  \item \textbf{ESD and pad capacitance.} ESD diodes can take most of the 0.85\,pF budget alone; get the \sgx{} ESD cell capacitance early, and keep the T-coil of \cref{sec:esd-pad} as the fallback.
  \item \textbf{$\mathrm{BV}_\mathrm{CEO}$ headroom.} The off driver transistor sits at 1.25--1.5\,V against about 1.6\,V; simulate every $V_\mathrm{CE}$ at every corner with the real termination.
  \item \textbf{PLL phase noise.} The quoted $-100.8$\,dBc/Hz at 1\,MHz is far worse than Leeson predicts (\cref{sec:lcvco}); the 1\,ps RJ line depends on what it really is.
  \item \textbf{Inductor EM modeling.} The 4\,nH tank and any T-coil must be EM-simulated on TopMetal2 (\cref{sec:layout-hs}) and matched to the team's measured $Q$.
  \item \textbf{Timing loops.} DFE tap 1 uses about 140 of its 200\,ps and the CDR dither 19 of its 20\,ps; both need post-layout confirmation.
  \item \textbf{Area, pads and the channel.} The PLL is 0.347\,mm$^2$~\cite{smacd2026pll}, the MIM capacitors add tenths of a mm$^2$, and the high-speed pads need SMA-grade access; REF must be replaced by measured S-parameters (\cref{sec:ch16-next}).
\end{itemize}

\begin{takeaways}
\begin{itemize}
  \item \OS{} is a single-lane transceiver: CMOS at 312.5\,MHz in the core, HBT CML at 1.25\,GHz and above, joined by level converters, two supplies and two clock domains.
  \item One 2.5\,GHz LC PLL feeds both sides: the TX divides it; the RX makes four phases and rotates them in 6.25\,ps steps under digital CDR control.
  \item The block specification table is the contract: every block has a number and a chapter; $^\dagger$ marks estimates awaiting simulation.
  \item The voltage budget leaves about 76\,mV; the timing budget leaves a 141\,ps eye with 1.4\,ps to spare.
  \item First-cut power is about 216\,mW (43\,pJ/b), five times less energy per bit than OpenSerDes, dominated by static CML current; the SST driver is the first saving.
\end{itemize}
\end{takeaways}

\tikzsetfigurename{15_flow-}
\chapter{Design and Verification Flow with Open Tools}\label{ch:flow}

\chapterlead{From a Python model to a GDS file and a test bench, with nothing but open-source tools: which tool does what, how the model and the circuits keep each other honest, and how to prove the chip works.}

\begin{objectives}
\begin{itemize}
  \item Name the open tool used at each stage of the \OS{} flow and say what it checks.
  \item Close the loop between the serdespy model and the transistor-level circuits by back-annotating simulated responses.
  \item Choose a testbench, a measurement and a pass/fail criterion for every block, including corners and Monte Carlo.
  \item Apply the layout rules that matter at 5\,Gb/s: symmetry, current density, metal width, inductors, extraction and decoupling.
  \item Plan the lab bring-up: test board, loopback, on-chip PRBS checking and the measurements that compare silicon with the model.
\end{itemize}
\end{objectives}

The flow in \cref{fig:ch15-flow} runs left to right, but it is not a pipeline. Every stage feeds numbers back to the system model, and the chip is ready for tape-out only when the model, rerun with those numbers, still shows margin at every corner.

\begin{figure}[!tbp]
\centering
\includegraphics[width=\linewidth]{ch15_flow_infographic}
\caption{\textbf{The \OS{} design flow, end to end.} Specifications and the serdespy system model come first; schematics in \tool{Xschem} with \tool{ngspice}, layout in \tool{KLayout}, extraction and EM simulation, and sign-off (DRC, LVS, \tool{CACE} datasheets) follow; then tape-out and lab test. The arrows back to the model are the important ones: each stage returns measured numbers that replace assumptions.}
\label{fig:ch15-flow}
\end{figure}

%% =====================================================================
\section{The open toolchain}\label{sec:toolchain}

Every tool in this chapter ships in one Docker image, \textbf{IIC-OSIC-TOOLS}~\cite{iicosic}, together with the IHP \sgx{} PDK~\cite{ihppdk}; pinning one image tag gives the whole team identical tools. The \OS{} repository builds on the JKU \textbf{analog mixed-signal chip template}~\cite{jkutemplate} (it needs image 2026.06 or later). Its recursive Makefiles make one \code{make all} run every testbench including PVT and Monte Carlo, run DRC, LVS and extraction on the analog macros, harden the digital logic and assemble the chip digital-on-top with seal ring, pad frame and fill, run a full-chip LVS without black boxes, and draw the bonding diagram. GitHub Actions rerun it on every commit, so a broken block is noticed the day it breaks.

\Cref{fig:ch15-tools} maps the tools onto the flow, and \cref{tab:ch15-tools} says what each one is for.

\begin{figure}[!tbp]
\centering
% ch15_tools.tex -- fig:ch15-tools : the open tool map of the Open-SERDES flow
\providecommand{\chFifteenT}[1]{{\color{slate}\bfseries #1}}
\begin{tikzpicture}[x=1cm,y=1cm,font=\sffamily\scriptsize,
  st/.style={blk,minimum width=19mm,minimum height=15mm,font=\sffamily\scriptsize,inner sep=2pt,align=center,text width=18mm},
  sa/.style={st,fill=btint},
  sd/.style={st,fill=gtint,draw=golddk},
  sn/.style={st,fill=navy,text=white},
  hd/.style={font=\sffamily\scriptsize\bfseries},
  fb/.style={-{Stealth[length=2mm]},draw=c4,line width=0.7pt,dashed}]
\draw[region,fill=bgpale!70] (-1.2,-2.6) rectangle (14.95,3.62);
\node[font=\sffamily\scriptsize\bfseries,text=slate,anchor=north west] at (-1.15,3.57) {IIC-OSIC-TOOLS container  +  JKU \texttt{ihp-sg13g2-ams-chip-template} (\texttt{make all})};
% system model (spans both lanes)
\node[sn,minimum height=46mm,minimum width=18mm,text width=16mm] (sys) at (-0.2,0) {\textbf{Specs and system model}\\[4pt]\textcolor{goldlt}{serdespy}\\\textcolor{goldlt}{PyBERT}\\\textcolor{goldlt}{scikit-rf}\\[4pt]\textcolor{pale}{eye, BER, budgets}};
% analog lane
\node[hd,text=bblue,anchor=south west] at (1.15,2.13) {ANALOG / CUSTOM: driver, CML, CTLE, samplers, PLL, PI};
\node[sa] (a1) at (2.25,1.3) {\textbf{Schematic + sim}\\[2pt]\chFifteenT{Xschem}\\\chFifteenT{ngspice} \chFifteenT{Xyce}\\\chFifteenT{VACASK}};
\node[sa] (a2) at (4.55,1.3) {\textbf{Characterize}\\[2pt]\chFifteenT{CACE}\\PVT corners,\\Monte Carlo};
\node[sa] (a3) at (6.85,1.3) {\textbf{Layout}\\[2pt]\chFifteenT{KLayout}\\\chFifteenT{Magic}\\PCells};
\node[sa] (a4) at (9.15,1.3) {\textbf{Extract + EM}\\[2pt]\chFifteenT{KLayout-PEX}\\\chFifteenT{Magic}\\\chFifteenT{openEMS} \chFifteenT{Palace}};
% digital lane
\node[hd,text=golddk,anchor=north west] at (1.15,-2.1) {DIGITAL / SEMI-CUSTOM: PRBS, CDR logic, CMOS 16:4 and 1:16, SPI};
\node[sd] (d1) at (2.25,-1.3) {\textbf{RTL + verify}\\[2pt]\chFifteenT{Verilator}\\\chFifteenT{Icarus}\\\chFifteenT{cocotb}};
\node[sd] (d2) at (4.55,-1.3) {\textbf{Synthesis}\\[2pt]\chFifteenT{Yosys}\\(\sgx{} std.\\cells)};
\node[sd] (d3) at (6.85,-1.3) {\textbf{Place, route}\\[2pt]\chFifteenT{LibreLane}\\\chFifteenT{OpenROAD}};
\node[sd] (d4) at (9.15,-1.3) {\textbf{Timing + RC}\\[2pt]\chFifteenT{OpenSTA}\\\chFifteenT{OpenRCX}\\(in LibreLane)};
% merge: top level and sign-off
\node[st,minimum height=46mm,text width=18mm] (top) at (11.65,0) {\textbf{Top-level}\\[2pt]digital-on-top\\assembly\\(\chFifteenT{LibreLane}),\\pad ring, seal,\\fill\\[3pt]full-chip\\simulation};
\node[st,minimum height=46mm,minimum width=18mm,text width=16mm,fill=gtint!60] (so) at (13.95,0) {\textbf{Sign-off}\\[2pt]DRC: \chFifteenT{KLayout},\\\chFifteenT{Magic}\\LVS: \chFifteenT{KLayout},\\\chFifteenT{Netgen}\\[3pt]$\to$ GDS,\\bonding\\diagram};
% arrows
\foreach \u/\v in {a1/a2,a2/a3,a3/a4,d1/d2,d2/d3,d3/d4}{\draw[sig] (\u) -- (\v);}
\draw[sig] (sys.east |- a1) -- (a1); \draw[sig] (sys.east |- d1) -- (d1);
\draw[sig] (a4.east) -- (a4.east -| top.west); \draw[sig] (d4.east) -- (d4.east -| top.west);
\draw[sig] (top) -- (so);
% back-annotation loop
\draw[fb] (a4.north) -- ++(0,0.95) -| node[pos=0.25,below=0pt,text=c4,font=\sffamily\scriptsize]{back-annotate: $H(f)$, TX FIR, $\sigma$, $V_\mathrm{os}$ (\cref{sec:model-loop})} (sys.north);
\draw[fb] (a2.south) -- ++(0,-0.25) -| node[pos=0.25,below=0pt,text=c4,font=\sffamily\scriptsize]{iterate} (a1.south);
\end{tikzpicture}

\caption{\textbf{Two lanes, one container.} The analog lane (HBT CML driver, CTLE, samplers, PLL, PIs) is hand-designed and verified block by block; the digital lane (PRBS, CDR loop filter, CMOS serializer and deserializer stages, SPI) is synthesized by \tool{LibreLane}. Both merge in a digital-on-top assembly and one sign-off. The red dashed arrows are the feedback paths: corners and Monte Carlo send the designer back to the schematic, and extracted responses go back to the system model.}
\label{fig:ch15-tools}
\end{figure}

\begin{table}[!tbp]
\centering
\caption{\textbf{Open tools in the \OS{} flow.} All are in the IIC-OSIC-TOOLS image except serdespy, PyBERT and scikit-rf, which install with \code{pip}.}
\label{tab:ch15-tools}
\footnotesize
\renewcommand{\arraystretch}{1.08}
\begin{tabularx}{\linewidth}{@{}>{\raggedright\arraybackslash}p{40mm} L l@{}}
\toprule
\textbf{Tool} & \textbf{Purpose in \OS{}} & \textbf{Stage}\\
\midrule
\tool{serdespy}~\cite{serdespy}, \tool{PyBERT}~\cite{pybert} & link models, eyes, BER (\cref{ch:labs}); PyBERT adds a GUI and IBIS-AMI & system model\\
\tool{scikit-rf}~\cite{p2:skrf2022} & Touchstone files, mixed-mode S-parameters, de-embedding & model, EM, lab\\
\tool{Xschem}~\cite{p6:xschem} & schematics, symbols, testbench netlists for every simulator & schematic\\
\tool{ngspice}~\cite{p6:ngspice} & \code{.op}, \code{.ac}, \code{.noise}, \code{.tran}, transient noise, Verilog-A via OSDI & simulation\\
\tool{Xyce}~\cite{p6:xyce}, \tool{VACASK}~\cite{p6:vacask} & parallel SPICE for large post-layout runs; Verilog-A simulator with PSS and HB & simulation\\
\tool{CACE}~\cite{p6:cace} & YAML datasheet: runs testbenches over corners, checks min/typ/max limits & characterization\\
\tool{KLayout}~\cite{p6:klayout}, \tool{Magic}~\cite{p6:magic} & layout editing; IHP DRC and LVS decks; RC extraction (\tool{Magic}) & layout, sign-off\\
\tool{Netgen}~\cite{p6:netgen}, \tool{KLayout-PEX}~\cite{p6:klayoutpex} & LVS comparison; parasitic extraction for \sgx{} & sign-off, extraction\\
\tool{openEMS}~\cite{p6:openems}, \tool{Palace}~\cite{p6:palace} & FDTD and FEM field solvers for inductors, pads and package & EM\\
\tool{Verilator}~\cite{p6:verilator}, \tool{Icarus}~\cite{p6:icarus}, \tool{cocotb}~\cite{p6:cocotb} & Verilog simulation and lint; Python testbenches for PRBS, CDR, SPI & RTL\\
\tool{Yosys}~\cite{p6:yosys}, \tool{LibreLane}~\cite{p6:librelane} & synthesis to \sgx{} cells; RTL-to-GDS flow (successor of OpenLane~2~\cite{p6:openlane2}) & digital\\
\tool{OpenROAD}~\cite{p6:openroad}, \tool{OpenSTA}~\cite{p6:opensta} & placement, clock tree, routing; static timing (inside \tool{LibreLane}) & digital\\
\bottomrule
\end{tabularx}
\end{table}

%% =====================================================================
\section{The model--circuit loop}\label{sec:model-loop}

The labs of \cref{ch:labs} describe each block by an ideal transfer function. The circuits will not match those ideals, and the only way to know whether the difference matters is to put the circuit's real response back into the model (\cref{fig:ch15-loop}). Because the link is linear up to the sampler, each simulated block becomes one more factor in the cascade:

\begin{keyeq}[Back-annotating the link model]
\begin{equation}
  H_\mathrm{link}(f) = H_\mathrm{TX}(f)\,\Hch(f)\,H_\mathrm{CTLE,ckt}(f),\qquad
  p_\mathrm{TX}(t) = s(t) - s(t-\Tb),\qquad
  \sign^2 = \sum_i \sigma_{n,i}^2 .
  \label{eq:ch15-annot}
\end{equation}
\end{keyeq}

Here $s(t)$ is the simulated step response of the driver into its load, so $p_\mathrm{TX}(t)$ is its real 1-UI pulse, with finite edges and any ringing; sampling it at $\Tb$ spacing gives the effective FIR taps. The CTLE enters through its \code{.ac} response, and each noise source through its rms value from \code{.noise} or Monte Carlo.

\begin{figure}[!tbp]
\centering
% ch15_loop.tex -- fig:ch15-loop : the model <-> circuit loop
\begin{tikzpicture}[x=1cm,y=1cm,font=\sffamily\scriptsize,
  bx/.style={blk,minimum width=52mm,minimum height=15mm,font=\sffamily\scriptsize,align=left,text width=49mm,inner sep=3pt},
  al/.style={font=\sffamily\scriptsize\itshape,text=golddk,align=center}]
\node[bx,fill=navy,text=white] (m) at (0,2.4) {\textbf{1. System model} (\textcolor{goldlt}{serdespy})\\REF channel, TX FIR, CTLE, DFE,\\noise and jitter $\to$ eye, BER, margins};
\node[bx,fill=gtint,draw=golddk] (s) at (7.6,2.4) {\textbf{2. Block specs} (\cref{tab:specs})\\swing, FFE range, peaking, $\sign$,\\$V_\mathrm{os}$, $V_\mathrm{sens}$, RJ, DJ, power};
\node[bx,fill=btint] (c) at (7.6,0) {\textbf{3. Circuit} (\tool{Xschem} + \tool{ngspice})\\\texttt{.ac}, \texttt{.tran}, \texttt{.noise}, transient noise,\\PVT corners, Monte Carlo, extracted RC};
\node[bx,fill=white] (b) at (0,0) {\textbf{4. Back-annotate the model}\\CTLE \texttt{.ac} $\to$ measured $\Hctle(f)$\\driver step $\to$ TX edge and FIR taps\\sampler MC $\to$ $\sign$, $V_\mathrm{os}$; PLL $\to$ RJ};
\draw[sig] (m) -- node[al,above=1pt]{derive targets} (s);
\draw[sig] (s) -- node[al,right=1pt]{design} (c);
\draw[sig] (c) -- node[al,below=1pt]{simulate, extract} (b);
\draw[sig] (b) -- node[al,left=1pt]{re-run eye, BER\\(\cref{lst:ch15-ac})} (m);
\node[font=\sffamily\scriptsize\bfseries,text=c4,align=center] at (3.8,1.2) {repeat at every corner\\until the margin is $>0$};
\end{tikzpicture}

\caption{\textbf{The model sets the specs; the circuits return their real responses; the model decides.} Each pass replaces one ideal block with a simulated one. The loop ends when the margin is positive at every corner with every block back-annotated.}
\label{fig:ch15-loop}
\end{figure}

\Cref{lst:ch15-ac} loads an \tool{ngspice} AC sweep into the Lab~4 model. In \tool{ngspice}, \code{wrdata ctle\_ac.txt v(outp,outn)} after an \code{ac dec 100 1meg 100g} writes frequency, real and imaginary parts in three columns. The sweep is resampled onto the model's 10\,MHz grid, the end values are held outside it, and the DC point is forced real, as \code{sdp.freq2impulse} requires.

\begin{lstlisting}[caption={Back-annotating a circuit CTLE (tested with a synthetic \code{wrdata} file).},label={lst:ch15-ac},captionpos=b]
from common import *                        # F grid, channel(), impulse(), ...
f_sp, re, im = np.loadtxt("ctle_ac.txt", usecols=(0, 1, 2), unpack=True)
H_sp = re + 1j * im                         # wrdata ctle_ac.txt v(outp,outn)
mag = np.interp(F, f_sp, np.abs(H_sp))      # resample onto the model grid;
phs = np.interp(F, f_sp, np.unwrap(np.angle(H_sp)))  # holds end values
H_ckt = mag * np.exp(1j * phs)
H_ckt[0] = mag[0]                           # DC must be real
h_link, t = impulse(channel() * H_ckt)      # REF x circuit CTLE
bits = np.tile(sdp.prbs7(1), 40)
y = through(tx_wave(bits), h_link)
report_eye("REF + ngspice CTLE", y, bits)
\end{lstlisting}

\begin{example}{What a slow second pole costs}\label{we:ch15-pole}
\textbf{Check:} feed \cref{lst:ch15-ac} the ideal \code{ctle(8)} response sampled from 1\,MHz to 100\,GHz. It returns 60.0\,mV and 0.56\,UI, exactly the Lab~5 pre-DFE eye, so the resampling loses nothing.

\textbf{A realistic circuit:} suppose the extracted CTLE has its output pole at $f_{p2}=6$\,GHz instead of the 10\,GHz of the model, because the four samplers load it more than planned. The 12\,dB eye drops from 104.0 to 98.0\,mV and the 14\,dB eye from 113.3 to 107.4\,mV, about 6\,\%. At 5\,GHz the loss is 9\,\%. That is affordable against 76\,mV of margin, but it is exactly the kind of number that only the loop reveals.
\end{example}

%% =====================================================================
\section{Testbenches}\label{sec:testbenches}

Every specification in \cref{tab:specs} needs a testbench that measures it and a pass/fail limit. \Cref{tab:ch15-tb} lists the key ones and \cref{fig:ch15-tb} sketches two of them.

\begin{table}[!tbp]
\centering
\caption{\textbf{Key \OS{} testbenches.} Limits from \cref{tab:specs}; every analog row runs at all PVT corners, and the mismatch-sensitive rows also run Monte Carlo.}
\label{tab:ch15-tb}
\footnotesize
\renewcommand{\arraystretch}{1.1}
\begin{tabularx}{\linewidth}{@{}>{\raggedright\arraybackslash}p{19mm} >{\raggedright\arraybackslash}p{36mm} L >{\raggedright\arraybackslash}p{30mm}@{}}
\toprule
\textbf{Block} & \textbf{Testbench} & \textbf{Measure} & \textbf{Pass}\\
\midrule
TX serializer + driver & \code{.tran}, PRBS7 into \qty{100}{\ohm} + $C_p$; then through the channel model & swing, edges, DCD, eye, every $V_\mathrm{CE}$ & 0.8--1.0\,Vppd; DCD + DDJ $\le8$\,ps; $V_\mathrm{CE}<\mathrm{BV}_\mathrm{CEO}$\\
Termination, pad & \code{.ac} or \code{.sp} with ESD and pad parasitics & $S_{11}$, $C_p$ & RL $\ge10$\,dB at 2.5\,GHz\\
CTLE & \code{.ac} at every code; \code{.noise}; large-signal \code{.tran} & peaking, $f_z$, $f_{p2}$, noise, linearity & 0--12\,dB; $f_{p2}\ge7.5$\,GHz; $\le0.25$\,mV\\
DFE summer & \code{.tran} with tap 1 toggling every UI & settling inside the tap-1 loop & $t_\mathrm{CK\to Q}+t_\mathrm{settle}+t_\mathrm{setup}\le$ 1\,UI\\
Sampler & DC-step sweep counting $P(1)$, Monte Carlo; tiny-input sweep; transient noise & $V_\mathrm{os}$, $\sigma$, $\tau$ & $\sigma\le1$\,mV, $\tau\le5$\,ps, trim range $\ge\pm3\sigma$\\
PLL & lock transient; VCO phase noise (transient noise, or PSS/pnoise in \tool{VACASK}); Python loop model & lock time, $\mathcal{L}(f)$, RJ & RJ $\le1$\,ps rms\\
PPF + PI & \code{.tran} at all 64 codes & phase vs.\ code, amplitude & I/Q $\le\ang{2}$, INL $\le\ang{1}$\\
CDR logic & \tool{cocotb} with a Python phase model driving the early/late inputs & lock, ppm range, dither & $\pm300$\,ppm; $\le20$\,ps pp\\
PRBS, SPI & \tool{cocotb} on \tool{Verilator} or \tool{Icarus}; error injection & error count, resync & exact counts\\
\bottomrule
\end{tabularx}
\end{table}

\begin{figure}[!tbp]
\centering
% ch15_tb.tex -- fig:ch15-tb : two key testbenches
\begin{tikzpicture}[x=1cm,y=1cm,font=\sffamily\scriptsize,
  b/.style={blk,minimum height=11mm,minimum width=18mm,font=\sffamily\scriptsize,inner sep=2pt,align=center},
  ba/.style={b,fill=btint}, bd/.style={b,fill=gtint,draw=golddk}, bw/.style={b,fill=white},
  pl/.style={font=\sffamily\footnotesize\bfseries,text=golddk,anchor=west}]
% ---------------- (a) TX eye testbench
\node[pl] at (-0.4,1.25) {(a) TX eye through the channel};
\node[bd] (src) at (0.75,0) {PRBS7 bits\\PWL file or\\XSPICE source};
\node[ba,minimum width=22mm,line width=1pt] (dut) at (3.4,0) {\textbf{DUT:} 2:1 muxes\\+ 3-tap driver\\(extracted)};
\node[ba] (pad) at (5.95,0) {pad + ESD\\$C_p\approx$ 0.6\,pF};
\node[bw,minimum width=24mm] (ch) at (8.6,0) {channel model\\RLGC ladder / LTRA,\\or $h(t)$ as FIR};
\node[ba] (rx) at (11.2,0) {\qty{100}{\ohm} load\\+ 300\,fF};
\node[bd] (py) at (13.6,0) {\texttt{wrdata} $\to$\\Python eye,\\\texttt{eye\_scan}};
\foreach \u/\v in {src/dut,dut/pad,pad/ch,ch/rx,rx/py}{\draw[sig] (\u) -- (\v);}
\node[b,minimum width=16mm,minimum height=7mm] (ck) at (3.4,-1.35) {ideal 2.5\,GHz\\clock (or PLL)};
\draw[clk] (ck) -- (dut);
\node[font=\sffamily\scriptsize,text=muted,anchor=north,align=center] at (8.6,-0.75) {pass: eye $\ge$ model eye $-$10\%\\at all corners};
% ---------------- (b) sampler offset / sensitivity / metastability
\node[pl] at (-0.4,-2.45) {(b) Sampler offset, noise and metastability};
\node[bd] (vin) at (0.75,-3.7) {$V_\mathrm{in}$ steps\\$-5\ldots+5$\,mV\\around $V_\mathrm{CM}$};
\node[ba,minimum width=24mm,line width=1pt] (sa) at (3.6,-3.7) {\textbf{DUT:} CML sampler\\+ offset DAC\\Monte Carlo mismatch};
\node[bd] (cnt) at (6.65,-3.7) {count 1s over\\$N$ clocks $\to$\\$P(1)$ vs $V_\mathrm{in}$};
\node[bw,minimum width=30mm] (fit) at (9.75,-3.7) {fit $P(1)=\Phi\!\left(\frac{V_\mathrm{in}-V_\mathrm{os}}{\sigma}\right)$\\$\to$ $V_\mathrm{os}$ histogram, $\sigma$};
\node[bw,minimum width=24mm] (met) at (13.2,-3.7) {tiny $V_\mathrm{in}$ sweep:\\$t_\mathrm{CK\to Q}$ vs $\ln V_\mathrm{in}$\\slope $\to$ $\tau$};
\node[b,minimum width=16mm,minimum height=7mm] (ck2) at (3.6,-5.05) {2.5\,GHz clock\\(transient noise on)};
\draw[clk] (ck2) -- (sa);
\foreach \u/\v in {vin/sa,sa/cnt,cnt/fit}{\draw[sig] (\u) -- (\v);}
\draw[sig] (4.6,-3.7 |- sa.south) -- (4.6,-4.7) -| (met.south);
\node[font=\sffamily\scriptsize,text=muted,anchor=north] at (9.0,-4.75) {pass: $\sigma\le1$\,mV, $\lvert V_\mathrm{os}\rvert\le0.5$\,mV after trim, $\tau\le5$\,ps};
\end{tikzpicture}

\caption{\textbf{Two testbenches that decide whether the link works.} (a) The TX eye: PRBS bits drive the extracted serializer and driver into the pad and a channel model; the waveform goes back to Python for the eye measurement. (b) The sampler: a slow input step sweep under Monte Carlo gives $P(1)$ curves whose center is the offset and whose width is the noise, and a sweep of tiny inputs gives the regeneration time constant $\tau$ from the slope of $t_\mathrm{CK\to Q}$ against $\ln V_\mathrm{in}$ (\cref{sec:metastability}).}
\label{fig:ch15-tb}
\end{figure}

\paragraph{The channel in SPICE} The REF channel has skin-effect loss that grows as $\sqrt f$ and dielectric loss that grows as $f$ (\cref{sec:loss}). Neither open simulator models that directly in a transient. \tool{ngspice} offers the lossless line (\code{T}), the SPICE3 lossy line LTRA (\code{O}), and the single and coupled lossy lines TXL and CPL; all of them take constant per-unit-length $R$, $L$, $G$, $C$. It has no HSPICE-style W-element with frequency-dependent tables and no element that reads a Touchstone file; its \code{.sp} analysis \emph{computes} S-parameters of a circuit, it does not import them. \tool{Xyce} has TRA, LTRA and a lumped line (YTRANSLINE), and a linear YLIN device that reads Touchstone data, documented mainly for harmonic-balance and other frequency-domain analyses. So there are three practical options: (i) a long ladder of RLGC sections, (ii) a behavioral source that convolves with $h(t)$, or (iii) splitting the problem at the channel.

\begin{keyeq}[Sections needed for a lumped line model]
\begin{equation}
  \Delta\ell\le\frac{\lambda_\mathrm{min}}{10}=\frac{v}{10\,f_\mathrm{max}}
  \quad\Rightarrow\quad N\ge\frac{10\,\ell\,f_\mathrm{max}}{v}
  \label{eq:ch15-seg}
\end{equation}
\end{keyeq}

For REF, $\ell=1$\,m, $v=c/\sqrt{4}=1.5\times10^8$\,m/s and $f_\mathrm{max}=10$\,GHz give $N\ge667$ sections, each needing extra R--L branches to mimic the skin effect. That is slow and hard to fit. Option (iii) is exact for a matched link and is what \OS{} uses: simulate the TX into a \qty{50}{\ohm}-per-side load in SPICE, export the differential waveform, apply the channel in Python with \code{through(x, h)}, and feed the result back into the RX testbench as a PWL source behind \qty{50}{\ohm} per side.

\paragraph{Corners, Monte Carlo and CACE} The \sgx{} model libraries have typical, best-case and worst-case sections for MOS devices, HBTs and resistors; combine them with supply $\pm5$\,\% and 0--85\,$^\circ$C. Mismatch-sensitive specs (sampler offset, I/Q phase, mirror ratios) need Monte Carlo, and the number of runs sets how well you know the answer.

\begin{keyeq}[How many Monte Carlo runs]
\begin{equation}
  \frac{\Delta\hat\sigma}{\sigma}\approx\frac{1}{\sqrt{2N}}
  \label{eq:ch15-mc}
\end{equation}
\end{keyeq}

Fifty runs know $\sigma$ to about 10\,\%, 200 runs to 5\,\% (for the sampler, $4.9\pm0.25$\,mV). \tool{CACE} turns all of this into a datasheet: a YAML file per block lists each parameter with its minimum, typical and maximum limits, the testbench that measures it and the conditions to sweep; \tool{CACE} runs every combination and reports pass or fail, so \cref{tab:specs} becomes an executable checklist.

%% =====================================================================
\section{Layout for high speed}\label{sec:layout-hs}

At 5\,Gb/s a layout can lose the margin that the schematic won. Six habits keep it.

\paragraph{Symmetry} Draw every differential circuit as two mirror-image halves and route each pair as a pair. Place HBT pairs and mirror devices common-centroid (ABBA) with dummies at the ends. A 10\,ps skew between FFE tap paths moves the de-emphasis edge by 0.05\,UI (\cref{sec:ffe-impl}); 1\,\% of clock duty-cycle error costs 4\,ps (\cref{sec:quadrature}).

\paragraph{Current density} An HBT reaches its peak $f_T$ at one collector current density. Size each emitter so that the switching transistors sit near that point at their tail current, from the npn13G2 $f_T$-versus-$J_C$ curves of the PDK; past the peak, $f_T$ collapses.

\paragraph{Metal width} Every DC path needs $w\ge I/J_\mathrm{max}$, with the electromigration limit $J_\mathrm{max}$ of that layer and temperature from the \sgx{} rules. Route the 20\,mA driver tail, its loads and the 2.5\,V straps on TopMetal1 and TopMetal2; a few ohms of strap under 77\,mA of CML current shift every common-mode level in the chip.

\paragraph{Inductors and EM} Put the 4\,nH VCO inductor (and any T-coil) on TopMetal2, away from other metal. IHP's \code{gds2openEMS} and \code{gds2palace} scripts take the GDS and a stackup file, run \tool{openEMS} (FDTD) or \tool{Palace} (FEM), and return S-parameters for a $\pi$-model or a wideband fit~\cite{ihppdk}. Match $L$ and $Q$ to the team's measured inductor before trusting the flow.

\paragraph{Extraction and isolation} Extract every analog block with \tool{KLayout-PEX} or \tool{Magic} and resimulate; clock-to-data coupling and the wiring capacitance on the CTLE output are the usual surprises. Surround the CML blocks, the VCO and the samplers with substrate guard rings, away from the switching CMOS logic.

\paragraph{Decoupling} CML draws constant current; CMOS draws its charge in spikes at the clock edges, and that charge must come from on-chip capacitors.

\begin{example}{Decoupling the deserializer}\label{we:ch15-decap}
\textbf{Switched capacitance:} the 1:16 deserializer burns 5\,mW at 1.2\,V and 312.5\,MHz, so $\alpha C_\mathrm{sw}=P/(V_\mathrm{DD}^2f)=5\,\mathrm{mW}/(1.44\times312.5\,\mathrm{MHz})=\qty{11}{\pico\farad}$ is switched every cycle.

\textbf{Worst case:} if all of it switches on one edge, before the package can refill the rail, the droop is $\Delta V/V\approx C_\mathrm{sw}/(C_\mathrm{sw}+C_\mathrm{dec})$. For 5\,\% droop, $C_\mathrm{dec}\ge19\,C_\mathrm{sw}=\qty{210}{\pico\farad}$: about 0.14\,mm$^2$ of MIM capacitor at roughly 1.5\,fF/\textmu{}m$^2$ (check the PDK value). Inside the synthesized digital macros, \tool{LibreLane} adds decap cells~\cite{jkutemplate}; budget the rest as MIM next to the macro, and count it in the die area.
\end{example}

\paragraph{The digital flow} The PRBS logic, the CDR loop filter, the SPI registers and the slow CMOS serializer and deserializer stages are Verilog, verified with \tool{cocotb} and hardened by \tool{LibreLane}. Synthesize what runs at 625\,MHz or slower; place the 1.25\,GHz deserializer stage and the CMOS-to-CML converters by hand from standard cells.

\paragraph{Sign-off} DRC with the IHP \tool{KLayout} decks (\tool{Magic} as a second opinion), LVS with \tool{KLayout} or \tool{Netgen}, then antenna and density checks after fill, and the template's full-chip LVS without black boxes~\cite{jkutemplate}: run it whenever a macro changes.

\subsection{Repository structure}\label{sec:ch15-repo}

The \OS{} repository keeps the recursive layout of the JKU template: every analog block is a macro with its own schematics, testbenches, layout, verification and Makefile, so a block can be built, simulated and verified alone. The listing shows these folders for \code{tx\_serializer}; every other macro has the same ones. The digital half (PRBS, the CMOS serializer and deserializer stages, the CDR loop filter and SPI) is not a macro: it lives in \code{rtl/} and is hardened with the chip, digital-on-top. \OS{} adds four folders of its own: \code{spec/}, the one specification file that the link model and every testbench read; \code{model/}; \code{channels/}; and \code{lab/}, for the test board and the measured data of \cref{sec:labtest}. The others (\code{flow/}, \code{layout/}, \code{netlist/}, \code{ip/}, \code{packaging/}, \code{release/} and \code{doc/}) come with the template.

\begin{lstlisting}[style=shell,language={},basicstyle=\ttfamily\scriptsize]
open-serdes/                  # from ihp-sg13g2-ams-chip-template
|-- Makefile                  # make model && make all
|-- spec/                     # one spec file: model and testbenches
|-- model/                    # link model, stat. eye, GUI, labs (Ch. 16)
|-- channels/                 # S-parameters: REF, REF-10, 802.3, boards
|-- macros/                   # analog macros, one per block
|   |-- tx_serializer/        # CML 4:2 and 2:1, CMOS-to-CML
|   |   |-- schematic/  testbenches/  layout/  netlist/
|   |   |-- verification/     # drc/  lvs/  cace/ (limits from spec/)
|   |   |-- scripts/
|   |   |-- final/            # GDS, LEF, lib for the chip assembly
|   |   `-- Makefile          # make all: build, sim, DRC, LVS
|   |-- tx_ffe_driver/        # 3-tap CML driver, termination, ESD, pad
|   |-- rx_ctle/              # termination, ESD, loopback mux, CTLE
|   |-- rx_sampler_dfe/       # 2-tap DFE summer, four samplers
|   |-- pll_lcvco/            # integer-N PLL, LC VCO, dividers
|   |-- cdr_pi/               # I/Q generation, two phase interpolators
|   `-- cml_lib/              # shared HBT CML cells
|-- rtl/                      # PRBS, CMOS SER/DES, CDR loop filter, SPI
|-- flow/                     # LibreLane: rtl/ hardened digital-on-top
|-- verification/             # link regression, eye checks, full-chip DRC/LVS
|-- layout/  netlist/  ip/  packaging/  release/
|-- lab/                      # test board, bring-up scripts, data
`-- doc/                      # documentation (from the template)
\end{lstlisting}

The template's \code{make macro FROM=inverter NAME=rx\_ctle} creates a new macro from its example. The macro's CACE datasheet reads every limit from \code{spec/}, so a number changes in one place only.

%% =====================================================================
\section{Testing the chip in the lab}\label{sec:labtest}

\Cref{fig:ch15-labsetup} shows the bring-up bench. The test board carries the chip, SMA connectors for TX$\pm$, RX$\pm$, the reference clock and a divided clock-monitor output, local regulators and decoupling, and an SPI header. Keep the \qty{50}{\ohm} traces from the pads to the SMAs short and identical within each pair, and de-embed them with \tool{scikit-rf} from a measured test structure.

\begin{figure}[!tbp]
\centering
% ch15_lab.tex -- fig:ch15-labsetup : the Open-SERDES bring-up bench
\begin{tikzpicture}[x=1cm,y=1cm,font=\sffamily\scriptsize,
  ins/.style={blk,minimum width=22mm,minimum height=10mm,font=\sffamily\scriptsize,align=center,inner sep=2pt,fill=white},
  sma/.style={draw=golddk,fill=gold,circle,minimum size=2.6mm,inner sep=0pt},
  cab/.style={line width=1.1pt,draw=navy},
  alt/.style={line width=1.0pt,draw=c5,densely dashed,-{Stealth[length=2mm]}},
  cabc/.style={line width=1.0pt,draw=gold,dashed},
  lb/.style={font=\sffamily\scriptsize,text=muted,align=center}]
% PCB
\draw[draw=c3,line width=1pt,fill=c3!8,rounded corners=4pt] (0,0) rectangle (6.6,4.2);
\node[font=\sffamily\scriptsize\bfseries,text=c3,anchor=south west] at (0.1,0.04) {test PCB: low-loss laminate, \qty{50}{\ohm} lines};
\node[blk,fill=navy,text=white,minimum width=19mm,minimum height=19mm,font=\sffamily\scriptsize,align=center] (chip) at (3.4,2.2) {\textbf{\OS{}}\\QFN or\\chip-on-board\\[2pt]\textcolor{pale}{PRBS gen.,}\\\textcolor{pale}{checker, LB}};
\node[blkb,minimum width=13mm,minimum height=8mm,font=\sffamily\scriptsize] (ldo) at (1.05,0.95) {LDOs\\+ decaps};
\node[blkg,minimum width=13mm,minimum height=7mm,font=\sffamily\scriptsize] (hdr) at (1.05,3.45) {SPI\\header};
\node[sma] (txp) at (6.6,3.35) {}; \node[sma] (txn) at (6.6,2.85) {};
\node[sma] (rxp) at (6.6,1.55) {}; \node[sma] (rxn) at (6.6,1.05) {};
\node[sma] (ref) at (0,2.2) {}; \node[sma] (mon) at (3.4,4.2) {};
\foreach \s in {txp,txn,rxp,rxn}{\draw[navy,line width=0.8pt] (chip.east |- \s) -- (\s);}
\draw[gold,line width=0.8pt] (ref) -- (chip.west);
\draw[gold,line width=0.8pt] (mon) -- (chip.north);
\draw[navy,line width=0.6pt] (ldo) -- (chip.south west);
\draw[golddk,line width=0.6pt] (hdr) -- (chip.north west);
\node[lb,anchor=south east] at (6.5,3.42) {TX$\pm$};
\node[lb,anchor=north east] at (6.5,0.98) {RX$\pm$};
\node[lb,anchor=west] at (3.5,3.75) {clock monitor};
% channel board between TX and RX
\node[blk,fill=white,minimum width=21mm,minimum height=11mm,font=\sffamily\scriptsize,align=center] (chb) at (8.7,2.2) {channel board\\REF-like 1\,m FR-4\\or matched cables};
\draw[cab] (txp) -- (7.15,3.35) |- ([yshift=2.5mm]chb.west);
\draw[cab,-{Stealth[length=2.2mm]}] ([yshift=-2.5mm]chb.west) -| (7.15,1.55) -- (rxp);
% alternatives
\node[ins] (scope) at (11.45,3.65) {sampling scope:\\TX eye, jitter};
\node[ins] (bert) at (11.45,0.75) {BERT: stressed RX,\\SJ for JTOL};
\draw[alt] (7.15,3.35) |- (scope.west);
\draw[alt] (bert.west) -| (7.15,1.55);
\node[lb,text=c5,anchor=south] at (9.0,3.67) {or: scope on TX};
\node[lb,text=c5,anchor=north] at (9.0,0.73) {or: BERT into RX};
% bench instruments
\node[ins] (pna) at (3.4,5.55) {spectrum / phase-noise\\analyzer: PLL $\mathcal{L}(f)$};
\draw[cabc] (mon) -- (pna);
\node[ins] (pc) at (-2.05,3.65) {PC + USB--SPI:\\Python config,\\error counters};
\draw[cab,draw=golddk] (pc) -- (pc -| hdr.west);
\node[ins] (sg) at (-2.05,2.2) {50\,MHz reference\\(generator or\\crystal module)};
\draw[cabc] (sg) -- (ref);
\node[ins] (psu) at (-2.05,0.75) {supplies 1.2, 2.5,\\3.3\,V + current\\meters (power)};
\draw[cab,draw=c4] (psu) -- (psu -| ldo.west);
\end{tikzpicture}

\caption{\textbf{The chip tests itself; the instruments check that it tells the truth.} The baseline loop needs only a reference generator, supplies, a PC and a channel board: the on-chip PRBS generator drives TX, the channel board brings the signal back to RX, and the on-chip checker counts errors. A sampling scope on TX (eye and jitter), a BERT into RX (stressed eye, sinusoidal jitter for JTOL) and a phase-noise analyzer on the clock monitor are added when available.}
\label{fig:ch15-labsetup}
\end{figure}

\paragraph{Bring-up order} (1) Supplies and currents against \cref{tab:ch14-power}. (2) SPI read-back of every register. (3) PLL lock, seen on the clock monitor. (4) TX eye with PRBS7 on a scope. (5) Near-end loopback with the on-chip checker: this tests the chip without the board. (6) The external loop through the channel board, sweeping FFE presets and CTLE codes. (7) Freeze the CDR, step the PI code through the override register and count errors at each step: an on-chip bathtub, with the offset DACs giving the vertical scan.

\begin{keyeq}[Bits needed to claim a BER]
\begin{equation}
  N_\mathrm{bits}\ge\frac{-\ln(1-\mathrm{CL})}{\BER}\approx\frac{3}{\BER}\quad\text{(zero errors, CL = 95\,\%)}
  \label{eq:ch15-bits}
\end{equation}
\end{keyeq}

Claiming $10^{-12}$ at 95\,\% confidence takes $3\times10^{12}$ error-free bits, or 600\,s at 5\,Gb/s: ten minutes per setting, which is why the checker is on chip.

\begin{table}[!tbp]
\centering
\caption{\textbf{Measurement plan: each silicon measurement and the model result it is compared with.}}
\label{tab:ch15-meas}
\footnotesize
\begin{tabularx}{\linewidth}{@{}>{\raggedright\arraybackslash}p{24mm} L >{\raggedright\arraybackslash}p{40mm}@{}}
\toprule
\textbf{Measurement} & \textbf{Setup} & \textbf{Target or comparison}\\
\midrule
TX eye and jitter & PRBS7, sampling scope at the TX SMAs & 0.8--1.0\,Vppd; edges, DCD vs.\ simulation\\
BER and bathtub & loopback or channel board; on-chip checker; PI-code sweep & BER $\le10^{-12}$ (10\,min); bathtub vs.\ Lab~6\\
Equalizer sweep & FFE presets $\times$ CTLE codes $\times$ DFE on/off & best setting vs.\ \cref{sec:eq-partition}\\
Jitter tolerance & BERT with sinusoidal jitter into RX & above the mask of \cref{sec:jtol}\\
PLL phase noise & clock monitor into a phase-noise analyzer & $\mathcal{L}(1\,\mathrm{MHz})$; RJ $\le1$\,ps rms\\
Return loss & VNA at the RX SMAs, de-embedded & RL $\ge10$\,dB at 2.5\,GHz\\
Power & current per supply rail in each mode & vs.\ \cref{tab:ch14-power}\\
\bottomrule
\end{tabularx}
\end{table}

\Cref{tab:ch15-meas} is the measurement plan. Each row ends in a comparison with a model number, which closes the last loop of the flow: when silicon and model disagree, the model is updated first, and the next design starts from it.

\begin{takeaways}
\begin{itemize}
  \item One container, IIC-OSIC-TOOLS, holds every tool and the \sgx{} PDK; the JKU template turns the whole flow into \code{make all} with regression on every commit.
  \item The system model is not finished when the circuits start: back-annotate the CTLE \code{.ac} response, the driver step and the noise and offset of every block, and rerun the eye at each corner.
  \item Open SPICE engines have only constant-RLGC lossy lines; for REF, simulate the TX into a matched load and apply the channel in Python.
  \item Run every analog testbench at PVT corners; $N$ Monte Carlo runs know $\sigma$ to $1/\sqrt{2N}$; let \tool{CACE} check every limit of the specification table.
  \item Layout for 5\,Gb/s means mirror symmetry, HBTs at their peak-$f_T$ current density, metal sized for the DC current, EM-simulated inductors on TopMetal2, extraction, guard rings and real decoupling area.
  \item The chip tests itself: PRBS generator, checker and loopback verify $10^{-12}$ in ten minutes; instruments confirm the eye, jitter tolerance, phase noise and power against the model.
\end{itemize}
\end{takeaways}

\tikzsetfigurename{16_labs-}
\chapter{Hands-on Labs with serdespy}\label{ch:labs}

\chapterlead{Seven short Python scripts reproduce the key numbers of this tutorial, from the closed eye of the bare channel to a bang-bang CDR locking onto real data.}

\begin{objectives}
\begin{itemize}
  \item Set up serdespy on a current Python and run every lab from \code{tutorial/code/}.
  \item Turn the REF channel into a pulse response, cursors and eye, and see why it must be equalized.
  \item Design a TX FFE, a CTLE and a DFE in the model and measure what each buys.
  \item Count bit errors with jitter and noise, and compare them with the semi-analytic BER and the bathtub.
  \item Build a behavioral half-rate bang-bang PI-CDR and watch it track a frequency offset.
\end{itemize}
\end{objectives}

Every lab has the same shape: a goal, the few lines of code that matter, the figure and numbers you should get, and questions that push past the script. The figures here are drawn by \code{figsrc/ch16\_labN.py} with exactly the computations of the labs, so your printed numbers must match to the last digit; each lab also writes its own PNGs.

%% =====================================================================
\section{Lab 0: set-up}\label{sec:lab0}

\paragraph{Goal} A working serdespy environment, and a map of the labs.

\begin{lstlisting}[style=shell]
python3 -m venv ~/serdes-venv           # tested with Python 3.13
source ~/serdes-venv/bin/activate
pip install serdespy                     # pulls numpy, scipy, matplotlib, scikit-rf
cd tutorial/code
python lab1_channel.py                   # PNGs appear in this folder
\end{lstlisting}

serdespy~1.0 still uses \code{np.complex\_}, which NumPy~2 removed; \code{common.py}, imported by every lab, restores it with \code{np.complex\_ = np.complex128} before serdespy is imported. It also defines the simulation grid (64 samples per UI, $\Delta t=3.125$\,ps, 0--160\,GHz in 10\,MHz steps) and the shared functions: \code{channel()} for REF, \code{ctle()}, \code{impulse()}, \code{pulse()}, \code{tx\_wave()}, \code{through()}, the eye measurement \code{eye\_scan()} (serdespy has none) and the Gaussian BER estimate \code{ber\_q()}. \Cref{tab:ch16-labs} maps the labs to the chapters.

\begin{table}[!tbp]
\centering
\caption{\textbf{The labs.} Run times on a laptop; every script is self-contained apart from \code{common.py}.}
\label{tab:ch16-labs}
\footnotesize
\begin{tabularx}{\linewidth}{@{}c l L l r@{}}
\toprule
\textbf{Lab} & \textbf{Script} & \textbf{What it shows (key number)} & \textbf{Theory} & \textbf{Time}\\
\midrule
1 & \code{lab1\_channel.py} & REF $S_{21}$, $h(t)$, $p(t)$, cursors ($-17.38$\,dB; PDA $-0.161$\,V) & \cref{sec:pulse-response,sec:refchannel} & 3\,s\\
2 & \code{lab2\_prbs\_eye.py} & PRBS7 eye before and after REF ($-77.8$\,mV) & \cref{sec:eye} & 2\,s\\
3 & \code{lab3\_tx\_ffe.py} & zero-forcing TX FFE (2-tap 162.7\,mV, 3-tap 144.8\,mV) & \cref{ch:txffe} & 3\,s\\
4 & \code{lab4\_ctle.py} & CTLE peaking sweep (113.3\,mV at 14\,dB) & \cref{ch:ctle} & 4\,s\\
5 & \code{lab5\_dfe.py} & DFE taps and BER ($3.2\times10^{-14}$ at 10\,mV) & \cref{ch:dfe} & 3\,s\\
6 & \code{lab6\_jitter\_ber.py} & counted vs.\ estimated BER, bathtub (0.45\,UI) & \cref{ch:jitter} & 21\,s\\
7 & \code{lab7\_bbcdr.py} & bang-bang PI-CDR (lock in 0.13\,\textmu{}s, 0 errors) & \cref{ch:cdr} & 7\,s\\
\bottomrule
\end{tabularx}
\end{table}

\begin{pitfall}[serdespy quirks that cost an afternoon each]
\begin{itemize}
  \item \code{sdp.Transmitter(data, levels, f)}: \code{f} is the \emph{Nyquist} frequency (pass 2.5e9 for 5\,Gb/s), whatever the docstring says.
  \item Call \code{Transmitter.FIR()} before \code{oversample()}, with taps \code{[pre, main, post]} or \code{[main, post]}.
  \item \code{sdp.channel\_coefficients()} always plots, at 1200\,dpi by default: pass \code{res=150}.
  \item \code{sdp.forcing\_ffe()} sets the main tap to 1: rescale to $\sum\lvert c_k\rvert=1$ yourself.
  \item \code{Receiver.noise()} discards any FFE or DFE already applied: add noise first, then equalize. Use \code{Receiver.FFE()}, not \code{FFE\_BR()}, which misaligns the data.
  \item serdespy has no CDR and no MLSE, whatever the README lists; Lab 7 builds the CDR.
\end{itemize}
\end{pitfall}

%% =====================================================================
\section{Lab 1: the channel}\label{sec:lab1}

\paragraph{Goal} Build REF from serdespy's ABCD-matrix functions, plot $\lvert S_{21}\rvert$, and turn it into the impulse response, the pulse response and its cursors (\cref{sec:pulse-response,sec:refchannel}).

\begin{lstlisting}
from common import *

H = channel()                                  # complex S21 on the grid F
h, t = impulse(H)                              # h is per-sample: sum(h) = H(0)
p = pulse(h)
n_pre, n_post = 2, 10
cur = sdp.channel_coefficients(p, t, SPS, n_pre, n_post, res=150,
                               title="Pulse response and cursors")
isi = np.sum(np.abs(cur)) - cur[n_pre]
print(f"Peak-distortion eye (1 V swing, {n_pre}+{n_post} cursors) = "
      f"{cur[n_pre] - isi:+.4f} V  (negative = closed)")
\end{lstlisting}

\paragraph{What you should see} \Cref{fig:ch16-lab1}. The loss is $-0.59$\,dB at DC, $-10.40$\,dB at 1.25\,GHz, \textbf{$-17.38$\,dB at 2.5\,GHz} and $-30.07$\,dB at 5\,GHz. The cursors are $p_{-1}=0.0486$, $p_0=0.3396$, $p_1=0.1817$, $p_2=0.0909$, $p_3=0.0537$\,V, with $p_1/p_0=0.535$. The ISI magnitudes add to 0.5008\,V, more than the main cursor, so the peak-distortion eye is \textbf{$-0.161$\,V}: closed.

\begin{figure}[!tbp]
\centering
\includegraphics[width=\linewidth]{ch16_lab1}
\caption{\textbf{Lab 1: a third of the pulse arrives in its own UI, and the tail outweighs it.} (a) $\lvert S_{21}\rvert$ of REF; the dot is the 2.5\,GHz Nyquist loss. (b) The 1-UI, 1\,V pulse response (pale) and its cursors at 1\,UI spacing from the peak: pre-cursor red, main gold, post-cursors blue.}
\label{fig:ch16-lab1}
\end{figure}

\paragraph{Questions and extensions} (1)~Call \code{channel(LOSS\_FAMILY[10])} and \code{channel(LOSS\_FAMILY[25])} (0.572 and 1.441\,m). How do $p_1/p_0$ and the PDA eye change with loss? (2)~Set \code{c\_pad=1e-12}. Where in the tail do the pad reflections appear, and why can a short DFE not cancel them?

%% =====================================================================
\section{Lab 2: a PRBS eye through the channel}\label{sec:lab2}

\paragraph{Goal} Send PRBS7 NRZ through REF and measure the eye at both ends (\cref{sec:prbs,sec:eye}).

\begin{lstlisting}
prbs = sdp.prbs7(1)                       # one period: 127 bits, seed = 1
bits = np.tile(prbs, 40)                  # 5080 bits = 1.016 us of data
tx = sdp.Transmitter(bits, LEVELS, F_NYQ)  # NB: 3rd argument = Nyquist frequency
tx.oversample(SPS)
x = tx.signal_ideal
H = channel()
h, t = impulse(H)
y = through(x, h)
hts_tx, ph_tx, _, _ = report_eye("TX output", x, bits)
hts_rx, ph_rx, S, A = report_eye("Channel output", y, bits)
save_eye(y, ph_rx, "Eye after 1 m channel (17.4 dB @ 2.5 GHz)", "lab2_eye_rx.png")
\end{lstlisting}

\paragraph{What you should see} \Cref{fig:ch16-lab2}. The TX eye is 1000\,mV and 1.00\,UI. After REF the eye is \textbf{$-77.8$\,mV} and 0\,UI wide: closed at every phase. The mean levels are $+169.8$ and $-165.3$\,mV, at a best phase of 6/64 UI, and the channel delay is 6.722\,ns.

\begin{figure}[!tbp]
\centering
\includegraphics[width=\linewidth]{ch16_lab2}
\caption{\textbf{Lab 2: the levels survive, the eye does not.} (a) 40 bits of PRBS7 at the TX (amber) and after REF with the 6.7\,ns delay removed (blue): single bits barely cross zero. (b) Density eye at the channel output, centerd on the best phase: the inner eye is closed by 77.8\,mV.}
\label{fig:ch16-lab2}
\end{figure}

\paragraph{Questions and extensions} (1)~Repeat with \code{sdp.prbs13(1)} (tile it 2 times). \Cref{sec:pda} quotes $-181$\,mV: why is PRBS7 optimistic? (2)~Count the longest run of identical bits in PRBS7 and in PRBS13. Which cursor sum does each run length expose?

%% =====================================================================
\section{Lab 3: TX FFE}\label{sec:lab3}

\paragraph{Goal} Compute zero-forcing 2-tap and 3-tap TX FFEs from the cursors, impose the peak-swing constraint, and measure the eyes (\cref{ch:txffe}).

\begin{lstlisting}
cur = sdp.channel_coefficients(p, t, SPS, 1, 1, res=150)
c3 = sdp.forcing_ffe(1, cur)              # [c-1, c0, c+1], c0 = 1
c3 = c3 / np.sum(np.abs(c3))              # peak-swing constraint sum|c| = 1
c2 = sdp.forcing_ffe(0, cur[1:])          # [c0, c+1]: classic de-emphasis
c2 = c2 / np.sum(np.abs(c2))
# ...
bits = np.tile(sdp.prbs7(1), 40)
for name, c, png in ((None, None, None), ("2-tap", c2, None),
                     ("3-tap", c3, "lab3_eye_ffe3.png")):
    y = through(tx_wave(bits, fir=c), h)
    hts, ph, S, A = report_eye(f"Channel + {name or 'no'} FFE", y, bits)
\end{lstlisting}

\paragraph{What you should see} \Cref{fig:ch16-lab3}. The 3-tap taps are $[-0.0853, 0.5959, -0.3188]$ (DC gain 0.192, 14.34\,dB of Nyquist boost), the 2-tap taps $[0.6515, -0.3485]$ (10.37\,dB). The main cursor drops from 0.3396 to 0.1749\,V, the price of equalizing at the TX. The eyes are \textbf{162.7\,mV, 0.89\,UI with 2 taps} and \textbf{144.8\,mV, 0.83\,UI with 3 taps}.

\begin{figure}[!tbp]
\centering
\includegraphics[width=\linewidth]{ch16_lab3}
\caption{\textbf{Lab 3: zero-forcing flattens the channel to Nyquist, but the pre-tap costs more than it gives.} (a) REF (blue), the 3-tap FFE (amber, 0\,dB at Nyquist and $-14.3$\,dB at DC) and their product (teal), flat to about 2.5\,GHz. (b), (c) PRBS7 eyes with the 2-tap and the 3-tap FFE.}
\label{fig:ch16-lab3}
\end{figure}

\paragraph{Questions and extensions} (1)~Zero-forcing does not maximize the eye. Hold $c_1$ and sweep $c_{-1}$ from 0 to $-0.1$, renormalizing each time: where is the best eye, and why does the pre-tap lose under $\sum\lvert c_k\rvert=1$? (2)~Quantize the taps to multiples of 1/32 (the 5-bit driver of \cref{sec:ffe-impl}). How much eye is lost?

%% =====================================================================
\section{Lab 4: CTLE}\label{sec:lab4}

\paragraph{Goal} Sweep the peaking of a one-zero, two-pole CTLE, find the best eye, and translate it into a degenerated differential pair (\cref{sec:ctle-degen}).

\begin{lstlisting}
FP1, FP2 = 2.5e9, 10e9                    # pole 1 at Nyquist, output pole 10 GHz
H = channel()
bits = np.tile(sdp.prbs7(1), 40)
x = tx_wave(bits)
peaks, eh, ew = np.arange(0, 21, 2), [], []
for pk in peaks:
    h_link, t = impulse(H * ctle(pk, FP1, FP2))   # channel x CTLE
    hts, ph, S, A = eye_scan(through(x, h_link), bits)
    eh.append(hts[ph]); ew.append(np.sum(hts > 0) / SPS)
best = peaks[int(np.argmax(eh))]
# ...
gm, a = 20e-3, 10**(best / 20)            # example: gm = 20 mS per side
RS = 2 * (a - 1) / gm                     # A = 1 + gm*RS/2
CS = 1 / (RS * 2 * np.pi * FP1 / a)       # wz = 1/(RS*CS)
RD, CL = 1 / gm, 1 / (2 * np.pi * FP2 / gm)   # gm*RD = 1, wp2 = 1/(RD*CL)
\end{lstlisting}

\paragraph{What you should see} \Cref{fig:ch16-lab4}. The eye opens at 6\,dB (31.6\,mV), reaches 104.0\,mV at 12\,dB, peaks at \textbf{113.3\,mV, 0.94\,UI at 14\,dB} and falls again beyond. With 14\,dB nominal, the real peak-to-DC gain is only 12.11\,dB because the second pole eats some of it; the CTLE is $-14.00$\,dB at DC and $-3.10$\,dB at Nyquist. The circuit values are $R_S=\qty{401}{\ohm}$, $C_S=0.795$\,pF, $R_D=\qty{50}{\ohm}$ and $\CL=318$\,fF.

\begin{figure}[!tbp]
\centering
\includegraphics[width=\linewidth]{ch16_lab4}
\caption{\textbf{Lab 4: more peaking opens the eye until the CTLE starts making ISI of its own.} (a) CTLE responses for 4, 8 and 14\,dB, all 0\,dB at high frequency. (b) Eye height against peaking; the shaded band is the 0--12\,dB range of the \OS{} CTLE. (c) The eye at the 14\,dB optimum.}
\label{fig:ch16-lab4}
\end{figure}

\paragraph{Questions and extensions} (1)~The optimum lies outside the \OS{} range. What does the DFE of Lab~5 do with the 8\,dB eye that the CTLE alone cannot? (2)~Move \code{FP2} to 6\,GHz and re-sweep; compare with \cref{we:ch15-pole}.

%% =====================================================================
\section{Lab 5: DFE}\label{sec:lab5}

\paragraph{Goal} Take the post-cursors left by a moderate 8\,dB CTLE as DFE taps, measure the eye against the number of taps, and estimate the BER (\cref{ch:dfe}).

\begin{lstlisting}
PEAK = 8                                   # CTLE peaking [dB] (lab 4 model)
h_link, t = impulse(H * ctle(PEAK))        # channel x CTLE
y = through(tx_wave(bits), h_link)         # CTLE output waveform
cur = sdp.channel_coefficients(pulse(h_link), t, SPS, 1, 8, res=150)
h_pre, h0, post = cur[0], cur[1], cur[2:]
for n in (0, 1, 2, 3, 4, 5, 8):
    rx = sdp.Receiver(y, SPS, F_NYQ, LEVELS, shift=True)  # centre eye at 0
    if n:
        rx.nrz_DFE(post[:n])
    report_eye(f"CTLE {PEAK} dB + {n}-tap DFE", rx.signal, bits)
# ... final design: 3 taps, then ber_q(S[ph], A[ph], sigma) for 5..20 mV
\end{lstlisting}

\paragraph{What you should see} \Cref{fig:ch16-lab5}. With the 8\,dB CTLE, $p_{-1}=0.0138$, $p_0=0.1832$ and $p_1$--$p_3$ = 0.0552, 0.0235, 0.0160\,V. The eye grows from 60.0\,mV with no DFE to 110.7, 129.5 and \textbf{143.0\,mV} with 1, 2 and 3 taps, then saturates near 155\,mV. With 3 taps the feedback is $\pm27.6$, $\pm11.8$ and $\pm8.0$\,mV, and the estimated BER is $5.9\times10^{-46}$, \textbf{$3.2\times10^{-14}$}, $6.0\times10^{-8}$ and $1.8\times10^{-5}$ at $\sigma$ = 5, 10, 15 and 20\,mV. A noise-free run through \code{sdp.prbs\_checker} gives 0 errors in 5024 bits. The lab uses 3 taps; the chip has 2, which gives 129.5\,mV here.

\begin{figure}[!tbp]
\centering
\includegraphics[width=\linewidth]{ch16_lab5}
\caption{\textbf{Lab 5: the first DFE tap is worth more than the next five together.} (a) Eye height against the number of DFE taps after an 8\,dB CTLE; teal is the \OS{} choice, gold the lab's. (b), (c) The eye before and after a 3-tap DFE. The steps at $\pm0.5$\,UI in (c) are an artifact of serdespy's per-UI feedback; only the sampling instant is meaningful.}
\label{fig:ch16-lab5}
\end{figure}

\paragraph{Questions and extensions} (1)~Reproduce the \OS{} operating point of \cref{sec:eq-partition}: TX taps $[0.75, -0.25]$ ($-6$\,dB), a 0\,dB CTLE and 2 DFE taps. You should get 171.8\,mV. (2)~The worst-case $Q$ at 10\,mV is $143.0/2/10=7.15$, yet the BER is $3.2\times10^{-14}$, not $Q(7.15)$. Why is \code{ber\_q} more optimistic than the worst case?

%% =====================================================================
\section{Lab 6: jitter, noise and BER}\label{sec:lab6}

\paragraph{Goal} Add 5\,ps rms TX random jitter and RX noise to the Lab~5 link, count real bit errors, and compare them with the estimate and the bathtub (\cref{sec:dualdirac}).

\begin{lstlisting}
np.random.seed(2026)                       # repeatable jitter and noise
PEAK, N_TAPS, RJ_UI = 8, 3, 0.025          # RJ = 2.5 % UI = 5 ps rms
bits = np.tile(sdp.prbs7(1), 800)          # 101 600 bits
y = through(tx_wave(bits, jitter_ui=RJ_UI), h_link)   # TX RJ -> channel -> CTLE
rx = sdp.Receiver(y, SPS, F_NYQ, LEVELS, shift=True)  # sampling phase -> index 0
rx.nrz_DFE(taps)
hts, ph, S, A = report_eye("With 5 ps TX RJ, no noise", rx.signal, bits)
for sigma in (10e-3, 15e-3, 20e-3, 25e-3):
    rx.noise(sigma)                        # resets to pre-DFE signal + AWGN
    rx.nrz_DFE(taps)                       # real decisions: error propagation
    d = sdp.nrz_a2d(rx.signal, SPS, 0.0)[skip:-5]
    n_err = sdp.prbs_checker(7, sdp.prbs7(1), d)[0]
    est = ber_q(S[0], A[0], sigma)
# ... bathtub: ber_q(S[o % SPS], A[o % SPS], sigma) at every phase offset o
\end{lstlisting}

\paragraph{What you should see} \Cref{fig:ch16-lab6}. With jitter but no noise the eye is 115.6\,mV and 0.73\,UI. Counting over 101\,544 bits:
\begin{center}\small
\begin{tabular}{@{}lcccc@{}}
\toprule
$\sigma$ (mV) & 10 & 15 & 20 & 25\\
\midrule
errors counted & 0 & 1 & 6 & 43\\
counted BER & $<3\times10^{-5}$ & $9.9\times10^{-6}$ & $5.9\times10^{-5}$ & $4.2\times10^{-4}$\\
estimated BER & $1.4\times10^{-12}$ & $1.4\times10^{-7}$ & $2.4\times10^{-5}$ & $3.9\times10^{-4}$\\
\bottomrule
\end{tabular}
\end{center}
At 25\,mV count and estimate agree within 10\,\%; the single 15\,mV error is a true rare event (a 63.6\,mV margin hit by $4.25\sigma$ of noise). With $\sigma=5$\,mV the bathtub is open \textbf{0.45\,UI at $10^{-12}$} (0.66\,UI at $10^{-5}$); at 20\,mV it never reaches $10^{-5}$, and the counted points fall on the estimate.

\begin{figure}[!tbp]
\centering
\includegraphics[width=\linewidth]{ch16_lab6}
\caption{\textbf{Lab 6: counted errors confirm the estimate where counting is possible; the estimate takes over below $10^{-5}$.} (a) Bathtubs from \code{ber\_q} for 5 and 20\,mV of noise, with error counts at 20\,mV (dots) at offsets up to $\pm0.375$\,UI. (b) BER at the eye center against noise. The open triangle at 10\,mV is the 95\,\% bound $3/N$ of \cref{eq:ch15-bits} for zero errors in $N$ bits.}
\label{fig:ch16-lab6}
\end{figure}

\paragraph{Questions and extensions} (1)~How many bits would you have to simulate to see one error at $\sigma=10$\,mV? How long would that take? (2)~Raise the RJ to 6\,ps rms (\code{RJ\_UI = 0.03}). How much does the $10^{-12}$ opening at 5\,mV shrink? Compare with the 14$\sigma$ rule of \cref{eq:tj}.

%% =====================================================================
\section{Lab 7: a bang-bang PI-CDR}\label{sec:lab7}

\paragraph{Goal} serdespy has no CDR, so \code{lab7\_bbcdr.py} builds one on the Lab~4 link (REF + 14\,dB CTLE, 5\,ps RJ, 5\,mV noise): a half-rate Alexander detector, a P+I loop filter at 625\,MHz and a 64-step PI, with the RX clock 200\,ppm fast and starting 74\,ps off the eye center, as in \cref{fig:ch13-sim}. It samples the real waveform where the PI says: per UI one data sample D and one edge sample E half a UI earlier, as the four half-rate samplers deliver. A transition with $E$ equal to the previous $D$ means the clock is early (the loop body below is excerpted; the script also counts errors and has a fixed-clock mode).

\begin{lstlisting}
for k in range(n_upd):
    n = k * UPD + np.arange(UPD)                      # RX bit indices of this update
    # ... phi_d = data phase: phi0 + ppm offset + optional sinusoidal jitter
    phi_c = np.round(acc) * LSB                       # PI phase (UI)
    t_d = t0 + SPS * (n - phi_d + phi_c)              # data-sampler instants
    d = sample(y, t_d) + sigma * rng.standard_normal(UPD) > 0
    e = sample(y, t_d - SPS / 2) + sigma * rng.standard_normal(UPD) > 0
    prev = np.concatenate(([d[0] if d_prev is None else d_prev], d[:-1]))
    trans = prev != d                                 # a transition before each bit
    votes = np.where(e == prev, 1, -1) * trans        # +1: clock early -> later
    pipe.append(np.sign(votes.sum()))                 # majority vote of the block
    v = pipe.pop(0)                                   # ... seen d_lat updates later
    if cdr:
        freg += ki * v                                # integral path (frequency)
        acc += kp * v + freg                          # proportional path + phase
\end{lstlisting}

\paragraph{What you should see} \Cref{fig:ch16-lab7}, with the baseline of \cref{sec:picdr} ($K_P$ = 1/8 LSB, $K_P/K_I$ = 128, one update of extra latency). The loop locks in \textbf{0.13\,\textmu{}s} (the phase-only model of \cref{fig:ch13-sim}, with a stricter lock test, reports 0.19\,\textmu{}s); the integral register settles at 199\,ppm against the true 200; the tracking error is \textbf{2.57\,ps rms} and 17.3\,ps peak-to-peak, around a static offset of $-7.4$\,ps from the maximum-height eye center. There are \textbf{0 bit errors} in 59\,680 bits with the CDR, against \textbf{1372} (BER $2.3\times10^{-2}$) with a fixed clock, which slides across an eye edge once per microsecond. Raising $K_P$ to 1/2 and 1\,LSB raises the wander to 5.13 and 9.90\,ps rms (31.1 and 57.5\,ps pp): the trade-off of \cref{we:ch13-step}.

\begin{figure}[!tbp]
\centering
\includegraphics[width=\linewidth]{ch16_lab7}
\caption{\textbf{Lab 7: the PI rotates continuously to follow a 200\,ppm offset, and the errors stop.} (a) PI code (navy) over the code the data phase requires (gold); it wraps every 2.0\,\textmu{}s. (b) Clock-minus-data phase: the 74\,ps initial error is pulled in by 0.13\,\textmu{}s, then the loop dithers around a $-7.4$\,ps lock point. (c) Cumulative errors: fixed clock (red, one burst per UI of drift) and CDR (blue, none).}
\label{fig:ch16-lab7}
\end{figure}

\paragraph{Questions and extensions} (1)~Sweep $K_P$ from 1/16 to 1\,LSB and the latency \code{d\_lat} from 0 to 4. Plot the rms and peak-to-peak wander and compare with dither $\approx2S\,T_\mathrm{lat}$ (\cref{sec:bbcdr}). (2)~Set \code{ki=0} and raise \code{ppm} to 300, 450 and 550. Where does the proportional-only loop lose lock, and does it match the 488\,ppm of \cref{sec:picdr}? Watch the static phase offset grow with the frequency offset before it fails. (3)~Add sinusoidal jitter with \code{sj\_pp} and \code{sj\_f}. Find the largest amplitude at 1 and 4.9\,MHz with zero errors, and plot your JTOL points on \cref{fig:ch13-jtol}. Why does the loop lock 7\,ps before the eye center?

%% =====================================================================
\section{Where to go next}\label{sec:ch16-next}

\paragraph{Build the GUI demo} The proposal lists a GUI front-end for the link model among its deliverables, in the spirit of the Julia SerDes framework (JLSD) it cites. The labs are the engine: put sliders for loss, taps, peaking and noise in front of \code{channel()}, the FFE, \code{ctle()}, the DFE and \code{lab7\_bbcdr.simulate()}, and redraw the eye and bathtub on every change. Cross-check the eyes against PyBERT~\cite{pybert}, a mature open link simulator with IBIS-AMI support.

\paragraph{Use a measured channel} Load the test board's 4-port Touchstone file. \code{sdp.four\_port\_to\_diff} assumes ports 1$\to$2 and 3$\to$4 and returns $V_\mathrm{out}/V_s$, half of $S_{21}$ for a matched link; the data must start near DC (\tool{scikit-rf} can extrapolate).
\begin{lstlisting}
import skrf as rf
ntw = rf.Network("board.s4p")               # measured channel, ports 1->2, 3->4
Hd, fd, hd, td = sdp.four_port_to_diff(ntw, None, 50, 50, t_d=DT)
h = 2 * hd                                  # S21-like: same scale as impulse(channel())
\end{lstlisting}
On a synthetic 4-port built from REF this returns exactly $-17.38$\,dB and $p_0=0.3396$\,V; Labs 2--7 then run unchanged. Next: crosstalk, adaptation with \code{sdp.lms\_equalizer}, and the statistical eye of \cref{sec:stateye}.

\begin{takeaways}
\begin{itemize}
  \item serdespy 1.0 runs on NumPy 2 with one shim; remember its quirks: Nyquist frequency in \code{Transmitter}, FIR before oversampling, \code{res=150}, renormalized taps, noise before equalization.
  \item REF loses 17.38\,dB at Nyquist; its PDA eye is $-0.161$\,V and its PRBS7 eye $-77.8$\,mV.
  \item Under the peak-swing constraint the 2-tap FFE (162.7\,mV) beats the zero-forcing 3-tap (144.8\,mV); a CTLE alone peaks at 113.3\,mV at 14\,dB; 8\,dB plus a 2- or 3-tap DFE gives 129.5 or 143.0\,mV.
  \item Counted errors confirm \code{ber\_q} where counting is possible; below $10^{-5}$ only the estimate speaks.
  \item A bang-bang PI-CDR locks to 200\,ppm in 0.13\,\textmu{}s with 2.6\,ps rms wander and no errors (a fixed clock: 1372).
\end{itemize}
\end{takeaways}


\appendix
\tikzsetfigurename{A_appendices-}
%% Appendices (\appendix is issued in tutorial.tex)
\providecommand{\appGrp}[1]{\rowcolor{btint}\multicolumn{2}{@{}l}{\sffamily\bfseries\color{navy}#1}\\}
\providecommand{\appBox}{$\square$}

%% =====================================================================
\chapter{Formula Sheet}\label{app:formulas}

The equations of this tutorial on two pages, in the notation of \cref{ch:intro}: $\Tb$ = 1\,UI = 200\,ps, $\fN=2.5$\,GHz, symbols $a_k=\pm1$, TX levels $\pm A$ with $A=0.5$\,V, cursors $p_k$ of a 1\,V, 1-UI pulse.

{\small
\renewcommand{\arraystretch}{1.32}
\begin{tabularx}{\linewidth}{@{}>{\centering\arraybackslash}p{0.47\linewidth} L@{}}
\toprule
\textbf{Formula} & \textbf{Meaning}\\
\midrule
\appGrp{Signals and BER (\cref{ch:signals,ch:linkanalysis})}
$\fb=1/\Tb,\quad \fN=1/(2\Tb)$ & bit rate and Nyquist frequency: 5\,Gb/s, 2.5\,GHz\\
$y(t)=A\sum_n a_n\,p(t-n\Tb)+n(t)$ & received signal: superposition of shifted pulse responses\\
$p_k=p(\ts+k\Tb)$ & cursors: $p_0$ main, $k<0$ pre-, $k>0$ post-cursors\\
$\mathrm{EH}=2A\big(p_0-\sum_{k\neq0}\lvert p_k\rvert\big)$ & peak-distortion (worst-case) inner eye height\\
$\BER=\Qf\!\left(\mathrm{EH}/2\sign\right),\ \Qf(x)=\tfrac12\erfc(x/\sqrt2)$ & BER for Gaussian noise of rms $\sign$ at the slicer\\
$\mathrm{EH}\ge2Q\sign=14.07\,\sign$ & eye needed for $10^{-12}$ ($Q=7.03$)\\
$S(f)=A^2\Tb\,\mathrm{sinc}^2(f\Tb)$ & NRZ power spectral density; first null at $\fb$\\
$M=\mathrm{EH}_\mathrm{ISI}-2Q\sign-\Delta V_\mathrm{jit}-2V_\mathrm{os}-V_\mathrm{sens}$ & linear voltage margin at the sampler\\
\appGrp{Channel (\cref{ch:channel})}
$v=c/\sqrt{\varepsilon_r},\quad t_d=\ell/v$ & propagation speed and delay; FR-4: 6.67\,ps/mm\\
$\Gamma=\frac{Z_L-\Zo}{Z_L+\Zo},\quad \mathrm{RL}=-20\log_{10}\lvert S_{11}\rvert$ & reflection coefficient and return loss\\
$\alpha_D\approx\pi f\sqrt{\varepsilon_r}\tan\delta/c$ & dielectric loss (Np/m), grows as $f$; skin loss grows as $\sqrt f$\\
$\lvert S_{11}\rvert=x/\sqrt{4+x^2},\ x=\omega\Zo C_p$ & pad capacitance on a doubly terminated node\\
$f_\mathrm{3dB}=1/(\pi\Zo C_p)$ & pad pole (node sees $\Zo/2$)\\
$\mathrm{IL}(\fN)-\mathrm{IL}(\fN/2)>9.54$\,dB & PAM4 can win over NRZ only above this octave loss\\
\appGrp{Equalization (\cref{ch:txffe,ch:ctle,ch:dfe})}
$y_n=\sum_k c_k x_{n-k},\quad\sum_k\lvert c_k\rvert=1$ & TX FFE under the peak-swing constraint\\
$\mathrm{DE}=20\log_{10}\frac{\lvert c_0\rvert+\lvert c_1\rvert}{\lvert c_0\rvert-\lvert c_1\rvert}$ & de-emphasis of a 2-tap FFE ($c_1=-0.25$: 6\,dB)\\
$\Hctle(s)=\frac{1}{A}\frac{1+s/\omega_z}{(1+s/\omega_{p1})(1+s/\omega_{p2})}$ & CTLE model, peaking $A=\omega_{p1}/\omega_z$\\
$A=1+\gm R_S/2,\ \omega_z=\frac{1}{R_SC_S},\ \omega_{p2}=\frac{1}{R_D\CL}$ & degenerated differential pair\\
$z_n=y_n-\sum_k w_k\hat d_{n-k},\quad w_k=A\,p_k$ & DFE; zero-forcing taps equal the post-cursors\\
$t_\mathrm{CK\to Q}+t_\mathrm{settle}+t_\mathrm{setup}\le\Tb$ & DFE tap-1 timing (direct feedback)\\
$w_k\leftarrow w_k+\mu\,\sgn(e_n)\,\sgn(\hat d_{n-k})$ & sign-sign LMS adaptation\\
\bottomrule
\end{tabularx}
}

{\small
\renewcommand{\arraystretch}{1.32}
\begin{tabularx}{\linewidth}{@{}>{\centering\arraybackslash}p{0.47\linewidth} L@{}}
\toprule
\textbf{Formula} & \textbf{Meaning}\\
\midrule
\appGrp{Circuits (\cref{ch:tx,ch:rx})}
$\Delta V=I_T\RL,\ \tau=\RL\CL,\ P=V_\mathrm{CC}I_T$ & CML gate: swing, time constant, static power\\
$V_{d,\mathrm{pp}}=I_T\Zo$ (CML), $I=V_{d,\mathrm{pp}}/4\Zo$ (SST) & driver current, doubly terminated\\
$\gm=I_C/V_T$; full switching at $\approx4V_T$ & HBT pair: 38.7\,mS/mA, about 100\,mV\\
$\tau=C/\gm$ (StrongARM), $\tau=\frac{\RL C}{\gm\RL-1}$ (CML) & latch regeneration time constant\\
$P_\mathrm{meta}=\frac{2V_L}{A\,\Delta V}\,e^{-T/\tau}$ & metastability probability per decision\\
$\sigma_{\Delta V_t}=A_{VT}/\sqrt{WL}$ & Pelgrom mismatch (offset)\\
$V_\mathrm{sens}\approx2Q\sigma_s+V_\mathrm{min}+2\lvert V_\mathrm{os,res}\rvert$ & sampler sensitivity\\
$\overline{v_n^2}=4kTR\,\Delta f,\quad v_{n,\mathrm{rms}}=\sqrt{kT/C}$ & resistor and sampled ($kT/C$) noise\\
\appGrp{Jitter, PLL and CDR (\cref{ch:jitter,ch:pll,ch:cdr})}
$\mathrm{TJ}=\mathrm{DJ}_{\delta\delta}+2Q\,\sigma_\mathrm{RJ}$ & total jitter at a BER (dual-Dirac)\\
$\sigma_\phi^2=2\!\int\!\mathcal{L}(f)\,df,\quad\sigma_t=\sigma_\phi/(2\pi f_0)$ & phase noise to rms jitter\\
$\omega_n=\sqrt{\frac{\Icp\Kvco}{2\pi NC_1}},\ \zeta=\frac{R}{2}\sqrt{\frac{\Icp\Kvco C_1}{2\pi N}}$ & CP-PLL natural frequency and damping\\
$f_0=\frac{1}{2\pi\sqrt{LC}},\ R_p=Q\omega_0L,\ \gm R_p/2\ge1$ & LC tank and start-up\\
$K_\mathrm{BB}=\mathrm{TD}\sqrt{2/\pi}/\sigj$ & bang-bang gain, linearized by jitter\\
$S=K_P\cdot\mathrm{LSB}/T_\mathrm{upd},\ \mathrm{JTOL}\approx\mathrm{TM}+2S/\omega$ & BB-CDR slew rate and slew-limited JTOL\\
$\Delta t_\mathrm{PI}=T_\mathrm{clk}/N_\mathrm{steps},\ \text{dither}\approx2S\,T_\mathrm{lat}$ & PI step (6.25\,ps) and limit-cycle dither\\
\appGrp{Flow and test (\cref{ch:arch,ch:flow})}
$E_b=P/\fb$ & energy per bit (216\,mW $\to$ 43\,pJ/b)\\
$N_\mathrm{bits}\ge3/\BER$ & error-free bits for a BER claim at 95\,\% confidence\\
$N\ge10\,\ell f_\mathrm{max}/v,\quad\Delta\hat\sigma/\sigma\approx1/\sqrt{2N}$ & lumped-line sections; Monte Carlo accuracy\\
\bottomrule
\end{tabularx}
}

%% =====================================================================
\chapter{Tables and Constants}\label{app:qtable}

\Cref{tab:app-q} converts a target BER into the Gaussian multiplier $Q$, and \cref{tab:app-db} converts loss in decibels into voltage and power ratios. \Cref{tab:app-const} collects the constants used in the examples.

\begin{table}[htbp]
\centering
\begin{minipage}[t]{0.46\linewidth}\centering
\captionof{table}{\textbf{BER against $Q$}, with $\BER=\Qf(Q)$.}
\label{tab:app-q}
\small
\begin{tabular}{@{}lcc@{}}
\toprule
BER & $Q$ & $2Q$ (eye/$\sigma$)\\
\midrule
$10^{-3}$ & 3.090 & 6.18\\
$10^{-4}$ & 3.719 & 7.44\\
$10^{-5}$ & 4.265 & 8.53\\
$10^{-6}$ & 4.753 & 9.51\\
$10^{-7}$ & 5.199 & 10.40\\
$10^{-8}$ & 5.612 & 11.22\\
$10^{-9}$ & 5.998 & 12.00\\
$10^{-10}$ & 6.361 & 12.72\\
$10^{-11}$ & 6.706 & 13.41\\
\rowcolor{gtint}$10^{-12}$ & \textbf{7.034} & \textbf{14.07}\\
$10^{-13}$ & 7.349 & 14.70\\
$10^{-14}$ & 7.651 & 15.30\\
$10^{-15}$ & 7.941 & 15.88\\
\bottomrule
\end{tabular}
\end{minipage}\hfill
\begin{minipage}[t]{0.5\linewidth}\centering
\captionof{table}{\textbf{Decibels to ratios} (loss).}
\label{tab:app-db}
\small
\begin{tabular}{@{}rcc@{}}
\toprule
Loss (dB) & Voltage ratio & Power ratio\\
\midrule
0.5 & 0.944 & 0.891\\
1 & 0.891 & 0.794\\
3 & 0.708 & 0.501\\
6 & 0.501 & 0.251\\
10 & 0.316 & 0.100\\
12 & 0.251 & 0.0631\\
14 & 0.200 & 0.0398\\
\rowcolor{gtint}17.4 (REF) & 0.135 & 0.0182\\
20 & 0.100 & 0.0100\\
30 & 0.0316 & 0.00100\\
40 & 0.0100 & 0.000100\\
\bottomrule
\end{tabular}
\end{minipage}
\end{table}

\begin{table}[htbp]
\centering
\caption{\textbf{Handy constants and \OS{} conversions} (at $T=300$\,K where it matters).}
\label{tab:app-const}
\small
\begin{tabularx}{\linewidth}{@{}l l L@{}}
\toprule
\textbf{Quantity} & \textbf{Value} & \textbf{Use}\\
\midrule
Boltzmann constant $k$ & $1.381\times10^{-23}$\,J/K & noise\\
Elementary charge $q$ & $1.602\times10^{-19}$\,C & shot noise, $V_T$\\
Thermal voltage $kT/q$ & 25.85\,mV & HBT $\gm=I_C/V_T$: 38.7\,mS per mA\\
$kT$ & $4.14\times10^{-21}$\,J = $-173.8$\,dBm/Hz & noise floors\\
$\sqrt{4kTR}$, $R=\qty{50}{\ohm}$ & 0.91\,nV/$\sqrt{\mathrm{Hz}}$ & termination noise\\
$\sqrt{kT/C}$, $C=50$\,fF & 0.29\,mV & sampled noise\\
Vacuum permittivity $\varepsilon_0$ & $8.854\times10^{-12}$\,F/m & capacitance\\
Speed of light $c$ & $2.998\times10^8$\,m/s & line delay\\
FR-4 ($\varepsilon_r=4$): $v=c/2$ & $1.50\times10^8$\,m/s = 6.67\,ps/mm = 169\,ps/in & reflections, flight time\\
1\,UI at 5\,Gb/s & 200\,ps & everything\\
1 degree of the 2.5\,GHz clock & 1.11\,ps & I/Q and PI errors\\
1 PI step & 6.25\,ps = 0.031\,UI & CDR\\
\bottomrule
\end{tabularx}
\end{table}

%% =====================================================================
\chapter{Design-Review Checklists}\label{app:checklists}

Use these lists in every design review. Each item points to the section that explains it; the numbers are the targets of \cref{tab:specs}.

{\small
\setlist[itemize]{label=\appBox,leftmargin=1.3em,itemsep=0.1em,topsep=0.2em}
\noindent\begin{minipage}[t]{0.485\linewidth}
\textbf{\sffamily\color{navy}Transmitter}
\begin{itemize}
  \item FO4 measured in the PDK; CMOS stages close timing at the slow corner (\cref{sec:serializer}).
  \item Half-rate mux: glitch-free phases, duty cycle $50\pm1$\,\%.
  \item Driver $I_T$ for 0.8--1.0\,Vppd; every $V_\mathrm{CE}<\mathrm{BV}_\mathrm{CEO}$ at $V_\mathrm{CC}+5$\,\%, with the real far-end termination.
  \item FFE presets $-3.5$/$-6$\,dB exact within one LSB; tap paths matched to a few ps.
  \item Termination trimmed to $\pm10$\,\%; $C_p\le0.85$\,pF; RL $\ge10$\,dB at $\fN$.
  \item Post-layout eye through REF within 10\,\% of the model.
\end{itemize}
\textbf{\sffamily\color{navy}Channel and board}
\begin{itemize}
  \item Loss at $\fN$ (not $\fb$) known; REF replaced by measured S-parameters.
  \item Via stubs short or back-drilled; first resonance above 5\,GHz.
  \item Pairs length-matched; TX and RX pairs separated (NEXT).
  \item SMA launches and bond wires modeled; reflections in the pulse tail checked.
\end{itemize}
\textbf{\sffamily\color{navy}Receiver}
\begin{itemize}
  \item CTLE: 0--12\,dB at every corner; $f_{p2}$ with the real sampler load.
  \item DFE tap-1 loop $\le1$\,UI post-layout; unrolling ready if not.
  \item Sampler $\sigma\le1$\,mV; offset DAC $\ge\pm3\sigma$, 5 bits; $\tau\le5$\,ps.
  \item Two extra latches behind the edge samplers (metastability).
  \item All four samplers simulated together (kick-back).
  \item Deserializer FIFO and recovered-clock domain crossing checked.
\end{itemize}
\end{minipage}\hfill
\begin{minipage}[t]{0.485\linewidth}
\textbf{\sffamily\color{navy}PLL and clocks}
\begin{itemize}
  \item VCO tunes 2.5\,GHz with $V_\mathrm{ctrl}$ mid-range at all corners; band code range covers PVT.
  \item Loop: 60$^\circ$ phase margin, $f_u\approx3$\,MHz, $f_\mathrm{ref}/f_u\ge10$.
  \item VCO phase noise simulated open-loop; RJ $\le1$\,ps rms integrated.
  \item Loop-filter capacitors MIM, not leaky MOS.
  \item I/Q error $\le\ang{2}$ ($3\sigma$), duty cycle $50\pm0.5$\,\% after DCC.
  \item Clock buffers have 2--3$\times f_\mathrm{clk}$ bandwidth.
\end{itemize}
\textbf{\sffamily\color{navy}CDR}
\begin{itemize}
  \item Phase-detector truth table verified for both half-rate phases.
  \item $K_P$ from the JTOL slew requirement; $K_P/K_I\ge16$ (baseline 128).
  \item Latency counted in updates; wander $\le20$\,ps pp in simulation.
  \item Integral register covers $\pm300$\,ppm; lock $<1$\,\textmu{}s.
  \item PI INL $\le\ang{1}$ at all 64 codes; code wrap glitch-free.
  \item Lab~7 model rerun with the real PI and latency.
\end{itemize}
\textbf{\sffamily\color{navy}Tape-out}
\begin{itemize}
  \item DRC clean (KLayout decks), antenna and density after fill.
  \item LVS clean per macro and full chip without black boxes.
  \item Post-layout corners and Monte Carlo pass in \tool{CACE}.
  \item Model back-annotated with extracted responses; margin $>0$.
  \item ESD on every pad; supply domains, decaps and guard rings reviewed.
  \item Test access: PRBS, checker, loopback, PI override, clock monitor.
  \item Bonding diagram, pinout and test-board plan signed off.
\end{itemize}
\end{minipage}
}

%% =====================================================================
\chapter{Glossary}\label{app:glossary}

{\small
\begin{description}[leftmargin=31mm,labelwidth=29mm,labelsep=2mm,itemsep=0.05em,parsep=0pt,topsep=0pt,font=\sffamily\bfseries\color{navy}]
  \item[Alexander PD] bang-bang phase detector that compares a data, an edge and the next data sample
  \item[Aperture] time window over which a sampler effectively averages its input
  \item[Bathtub curve] BER against sampling position across the UI
  \item[BER] bit-error rate: errors divided by bits; \OS{} targets $10^{-12}$
  \item[BiCMOS] process with both bipolar (HBT) and CMOS transistors, such as \sgx{}
  \item[$\mathrm{BV}_\mathrm{CEO}$] collector-emitter breakdown voltage of an HBT with the base open
  \item[CDR] clock and data recovery: extracts the sampling phase from the data transitions
  \item[CML] current-mode logic: differential pairs steering a constant tail current into resistors
  \item[Common centroid] layout that interleaves matched devices so gradients cancel
  \item[CTLE] continuous-time linear equalizer: an analog high-frequency boost
  \item[Cursor] sample of the pulse response at an integer number of UIs from the main sample
  \item[DCD] duty-cycle distortion: unequal high and low times, a two-valued DJ
  \item[De-emphasis] TX FFE setting that lowers repeated bits instead of raising transitions
  \item[DFE] decision-feedback equalizer: subtracts ISI predicted from past decisions
  \item[DJ, RJ] deterministic (bounded) and random (Gaussian) jitter
  \item[Dual-Dirac model] jitter model with DJ as two impulses convolved with Gaussian RJ
  \item[ESD] electrostatic discharge, and the diodes that protect the pads from it
  \item[Eye diagram] overlay of the received waveform folded at one or two UIs
  \item[FFE] feed-forward equalizer: an FIR filter, here in the TX driver
  \item[HBT] heterojunction bipolar transistor; \sgx{} npn13G2 with $f_T\approx300$\,GHz
  \item[Insertion loss] attenuation through the channel, $-20\log_{10}\lvert S_{21}\rvert$
  \item[ISI] inter-symbol interference: energy of a bit that spills into its neighbors
  \item[JTOL] jitter tolerance: largest sinusoidal jitter, against frequency, that still meets the BER
  \item[Loopback] internal path from the TX output to the RX input for self-test
  \item[LVS, DRC] layout-versus-schematic and design-rule checks before tape-out
  \item[Metastability] a latch input so small that the decision is late or undefined
  \item[Monte Carlo] simulation repeated with random device mismatch
  \item[NRZ, PAM4] two-level and four-level pulse-amplitude modulation
  \item[Nyquist frequency] half the symbol rate: 2.5\,GHz for 5\,Gb/s NRZ
  \item[PDA] peak-distortion analysis: worst-case eye from the sum of absolute cursors
  \item[Phase interpolator] circuit that mixes two clock phases to make any phase between them
  \item[Phase noise] single-sideband noise power near a clock, $\mathcal{L}(f)$, in dBc/Hz
  \item[PLL] phase-locked loop: multiplies a clean reference up to the VCO frequency
  \item[Polyphase filter] RC network that splits a differential clock into I and Q
  \item[PRBS] pseudo-random binary sequence from a linear-feedback shift register
  \item[Pulse response] channel output for one isolated 1-UI pulse
  \item[Return loss] reflection at a port, $-20\log_{10}\lvert S_{11}\rvert$
  \item[S-parameters] frequency-domain description of a network by incident and reflected waves
  \item[Sampler] clocked comparator (latch) that decides each bit
  \item[Serializer] multiplexer tree that turns a parallel word into a serial stream
  \item[SST driver] source-series-terminated (voltage-mode) driver
  \item[Statistical eye] eye at a target BER computed from the ISI and noise distributions
  \item[T-coil] coupled-inductor network that hides pad capacitance from the line
  \item[UI] unit interval: one bit time, 200\,ps
  \item[Zero-forcing] equalizer that forces selected cursors to zero, ignoring noise
\end{description}
}

%% =====================================================================
\chapter{Further Reading}\label{app:reading}

\paragraph{Start here} Palermo's ECEN\,720 course~\cite{palermo720} is the backbone of this tutorial; its lectures on channels and link analysis~\cite{palermo720ch}, equalization~\cite{palermo720eq}, circuits~\cite{palermo720ckt} and clocking~\cite{palermo720clk} go deeper on every chapter. The forthcoming book by Shekhar, Palermo and Hanumolu~\cite{shekhar2026} covers the same ground in textbook form. Benjamin Liu's handwritten notes on equalization~\cite{liuEq} and on PLL-based CDRs~\cite{liuCDR} are the quickest way to see each idea derived by hand.

\paragraph{Signal integrity and systems} Dally and Poulton~\cite{dally1998} is the classic on signaling, noise budgets and timing in digital systems. Hall and Heck~\cite{hall2009} treat transmission lines, losses, S-parameters and channel modeling in the detail you need to replace REF with a real board. Stojanovi\'c \emph{et al.}~\cite{stojanovic2005} show a complete adaptive PAM2/4 transceiver and why the equalization partition matters. The IEEE 802.3 backplane clauses~\cite{ieee8023} define the channels, masks and test methods of real 2.5 and 5\,Gb/s Ethernet links.

\paragraph{Circuits} Razavi's analog text~\cite{razaviAnalog} is the reference for differential pairs, noise and mismatch. His optical-communications book~\cite{razaviOptical} covers CML, limiting amplifiers, CDRs and multiplexers at the transistor level, much as \OS{} builds them.

\paragraph{Clocking} Razavi's PLL book~\cite{razaviPLL} takes the charge-pump PLL from the loop equations to the VCO and divider circuits. Gardner~\cite{gardner2005} is the classic on loop dynamics and the source of most PLL design formulas. The team's 2.4\,GHz LC-VCO PLL paper~\cite{smacd2026pll} is the starting point for the \OS{} PLL.

\paragraph{Open-source design} The IHP \sgx{} open PDK~\cite{ihppdk}, the IIC-OSIC-TOOLS image~\cite{iicosic} and the JKU chip template~\cite{jkutemplate} are everything you need to tape out; read the template tutorial before your first \code{make all}. The Purdue OpenSerDes paper~\cite{openserdes2021} is the closest open-source prior art, an all-digital 2\,Gb/s link in SKY130. For modeling, serdespy~\cite{serdespy} powers the labs of \cref{ch:labs}, and PyBERT~\cite{pybert} is the natural next step, with a GUI and IBIS-AMI support.


\backmatter
\cleardoublepage
\renewcommand{\bibname}{References}
\phantomsection\addcontentsline{toc}{chapter}{References}
\bibliographystyle{unsrtnat}
\bibliography{tutorial/bib/core,tutorial/bib/p1,tutorial/bib/p2,tutorial/bib/p3,tutorial/bib/p4,tutorial/bib/p5,tutorial/bib/p6}
\end{document}
